\documentclass[11pt, a4paper]{book}

% ----- Pacchetti -----
\usepackage[utf8]{inputenc}
\usepackage[T1]{fontenc}
\usepackage{amsmath, amssymb, amsthm}
\usepackage{mathtools}
\usepackage{geometry}
\usepackage{hyperref}
\usepackage{enumitem}
\usepackage{xcolor}
\usepackage{mdframed}
\usepackage{tocloft}
\usepackage{booktabs}
\usepackage{array}

\geometry{top=2.5cm, bottom=2.5cm, left=3cm, right=3cm}

% ----- Colori tematici -----
\definecolor{defcolor}{RGB}{0,70,127}
\definecolor{thmcolor}{RGB}{120,0,0}
\definecolor{excolor}{RGB}{0,100,60}
\definecolor{notecolor}{RGB}{100,60,0}
\definecolor{boxbg}{RGB}{245,248,255}
\definecolor{exbg}{RGB}{245,255,248}

% ----- Ambienti teorema -----
\theoremstyle{definition}
\newtheorem{definizione}{Definizione}[section]
\newtheorem{esempio}[definizione]{Esempio}
\newtheorem{esercizio}[definizione]{Esercizio}

\theoremstyle{plain}
\newtheorem{teorema}[definizione]{Teorema}
\newtheorem{proposizione}[definizione]{Proposizione}
\newtheorem{lemma}[definizione]{Lemma}
\newtheorem{corollario}[definizione]{Corollario}
\newtheorem{assioma}[definizione]{Assioma}

\theoremstyle{remark}
\newtheorem{osservazione}[definizione]{Osservazione}
\newtheorem{nota}[definizione]{Nota}

% ----- Box colorati -----
\newmdenv[
  backgroundcolor=boxbg,
  linecolor=defcolor,
  linewidth=1.5pt,
  roundcorner=4pt,
  skipabove=8pt,
  skipbelow=8pt
]{defbox}

\newmdenv[
  backgroundcolor=exbg,
  linecolor=excolor,
  linewidth=1.5pt,
  roundcorner=4pt,
  skipabove=8pt,
  skipbelow=8pt
]{exbox}

% ----- Comandi matematici -----
\newcommand{\N}{\mathbb{N}}
\newcommand{\Z}{\mathbb{Z}}
\newcommand{\Q}{\mathbb{Q}}
\newcommand{\R}{\mathbb{R}}
\newcommand{\C}{\mathbb{C}}
\newcommand{\K}{\mathbb{K}}
\newcommand{\Fp}{\mathbb{F}_p}
\newcommand{\mcd}{\mathrm{mcd}}
\newcommand{\im}{\mathrm{im}\,}
\newcommand{\Id}{\mathrm{Id}}
\newcommand{\Arg}{\mathrm{Arg}}
\newcommand{\OT}{\mathrm{OT}}
\DeclareMathOperator{\ord}{ord}
\DeclareMathOperator{\car}{char}
\DeclareMathOperator{\rg}{rg}
\DeclareMathOperator{\Span}{Span}
\newcommand{\GL}{\mathrm{GL}}
\newcommand{\End}{\mathrm{End}}
\newcommand{\Ann}{\mathrm{Ann}}
\newcommand{\val}{\mathrm{Val}}
\DeclareMathOperator{\Aut}{Aut}
\DeclareMathOperator{\Tr}{Tr}
\newcommand{\VV}{\mathcal{V}}   % gruppo delle traslazioni V(S)
\newcommand{\RA}{\mathrm{RA}}   % riferimento affine
\DeclareMathOperator{\Isom}{Isom}
\newcommand{\dista}{d}           % distanza

% ----- Titolo -----
\title{
  \Huge\bfseries Algebra Lineare e Geometria\\[1em]
  \large Una presentazione ragionata\\[2em]
  \normalsize Basato sugli appunti di\\
  \large Kieran G.\ O'Grady\\
  \normalsize ``Sapienza'' Università di Roma
}
\author{}
\date{Versione ampliata -- 2024}

% ========================================
\begin{document}
% ========================================

\maketitle
\tableofcontents

% ============================================================
\chapter*{Prefazione}
\addcontentsline{toc}{chapter}{Prefazione}
% ============================================================

Questi appunti nascono come rielaborazione dei materiali del corso di \emph{Algebra Lineare e Geometria} tenuto dal prof.\ Kieran G.\ O'Grady presso la Sapienza Università di Roma. L'obiettivo di questa versione è trasformare una raccolta di definizioni, teoremi e dimostrazioni in un testo più discorsivo, che accompagni il lettore con motivazioni, esempi concreti e commenti sul ``perché'' delle scelte formali.

\medskip
\noindent\textbf{A chi è rivolto questo testo.} Il libro è pensato per studenti del primo anno di Matematica, Fisica o Ingegneria che si avvicinano per la prima volta alla matematica rigorosa. Non si presuppone alcuna conoscenza di algebra astratta, ma si assume familiarità con il calcolo elementare e la geometria delle scuole superiori.

\medskip
\noindent\textbf{Come leggere questo testo.} Il Capitolo~1 è di natura preliminare: contiene gli strumenti logico-algebrici su cui si costruisce tutto il resto. Alcune sezioni (insiemi, applicazioni, induzione) possono essere scorrevoli per chi ha già una solida formazione; altre (anelli, campi, gruppi) introducono strutture fondamentali che verranno usate sistematicamente nei capitoli successivi. Si consiglia di non saltare le dimostrazioni: spesso il vero apprendimento avviene capendo \emph{perché} un risultato è vero, non solo conoscendo il risultato.

% ============================================================
\chapter{Preliminari}
% ============================================================

\section*{Introduzione al capitolo}

Prima di addentrarci nell'algebra lineare propriamente detta, dobbiamo costruire un linguaggio comune. Questo capitolo introduce i concetti fondamentali: insiemi, applicazioni, relazioni, induzione matematica e le strutture algebriche di base (anelli, campi, gruppi).

Potrebbe sembrare che questo materiale sia puramente tecnico e distante dalla geometria. In realtà, queste strutture sono il \emph{vocabolario} senza il quale è impossibile parlare precisamente di spazi vettoriali, trasformazioni lineari e tutto ciò che seguirà. Un campo, per esempio, non è solo un insieme di numeri: è il contesto in cui ``vivono'' i coefficienti dei nostri spazi vettoriali. Comprendere cosa è un campo apre la porta a lavorare su $\R$, $\C$, $\Q$, ma anche su strutture finite come $\mathbb{F}_p$, che giocano un ruolo fondamentale in crittografia e teoria dei codici.

\section{Insiemi}

\subsection{La nozione intuitiva e il principio di estensione}

La nozione di \emph{insieme} è così fondamentale da risultare difficile da definire senza ricorrere a concetti altrettanto primitivi. Intuitivamente, un insieme è una \emph{collezione di oggetti}, detti \emph{elementi} dell'insieme. Per esempio:
\begin{itemize}
  \item l'insieme $I$ degli italiani,
  \item l'insieme $\N$ dei numeri naturali $\{0, 1, 2, 3, \ldots\}$,
  \item l'insieme vuoto $\emptyset$, che non contiene alcun elemento.
\end{itemize}

Un insieme può essere descritto \emph{elencando} i suoi elementi oppure \emph{specificando una proprietà} che li caratterizza. Per esempio, l'insieme dei numeri naturali pari non maggiori di $10$ può essere scritto come
\[
  X \coloneqq \{0, 2, 4, 6, 8, 10\}
  \quad\text{oppure come}\quad
  X \coloneqq \{x \in \N \mid x \text{ pari}, \, x \leq 10\}.
\]
Il simbolo $\coloneqq$ indica una \emph{definizione}: il simbolo a sinistra viene definito come l'espressione a destra. La barra verticale ``$\mid$'' si legge ``tale che''.

\begin{defbox}
\begin{assioma}[Principio di estensione]
Un insieme è caratterizzato dagli elementi che gli appartengono: se $X$ e $Y$ sono insiemi, allora $X = Y$ se e solo se $X$ e $Y$ hanno esattamente gli stessi elementi.
\end{assioma}
\end{defbox}

\begin{osservazione}
Questo assioma chiarisce che un insieme \emph{non} dipende dall'ordine in cui elenchiamo i suoi elementi né da eventuali ripetizioni: $\{1, 2, 3\} = \{3, 1, 2\} = \{1, 1, 2, 3\}$.
\end{osservazione}

\subsection{Insiemi numerici fondamentali}

Gli insiemi numerici seguenti saranno onnipresenti nel testo:
\begin{itemize}
  \item $\N = \{0, 1, 2, \ldots\}$: i numeri \emph{naturali} (con lo zero).
  \item $\Z = \{0, \pm 1, \pm 2, \ldots\}$: i numeri \emph{interi}.
  \item $\Q = \{p/q \mid p \in \Z,\, q \in \Z \setminus \{0\}\}$: i numeri \emph{razionali}.
  \item $\R$: i numeri \emph{reali}, che incluono per esempio $\sqrt{2}$, $\pi$, $e$.
  \item $\C$: i numeri \emph{complessi} (introdotti in Sezione~\ref{sec:complessi}).
\end{itemize}

Abbiamo le inclusioni $\N \subset \Z \subset \Q \subset \R \subset \C$.

\medskip
Dati $a, b \in \R$ con $a \leq b$, definiamo gli \emph{intervalli}:
\[
  [a,b] \coloneqq \{x \in \R \mid a \leq x \leq b\}, \quad
  (a,b) \coloneqq \{x \in \R \mid a < x < b\},
\]
e analogamente gli intervalli semi-aperti $[a,b)$ e $(a,b]$. Per un estremo infinito usiamo $[a,+\infty) \coloneqq \{x \in \R \mid x \geq a\}$ e così via.

Infine, per $a \in \Z$ definiamo l'insieme dei multipli di $a$:
\[
  (a) \coloneqq \{x \in \Z \mid x = na \text{ per qualche } n \in \Z\}.
\]

\subsection{Operazioni tra insiemi}

\begin{defbox}
\begin{definizione}
Siano $X$ e $Y$ insiemi. Definiamo:
\begin{enumerate}
  \item \textbf{Unione}: $X \cup Y \coloneqq \{x \mid x \in X \text{ oppure } x \in Y\}$.
  \item \textbf{Intersezione}: $X \cap Y \coloneqq \{x \mid x \in X \text{ e } x \in Y\}$.
  \item \textbf{Differenza}: $X \setminus Y \coloneqq \{x \in X \mid x \notin Y\}$.
  \item \textbf{Inclusione}: $X \subseteq Y$ se ogni elemento di $X$ è anche elemento di $Y$.
\end{enumerate}
\end{definizione}
\end{defbox}

\begin{exbox}
\begin{esempio}
Consideriamo i multipli di $2$ e di $3$ in $\Z$:
\[
  (2) \cup \{x \in \Z \mid x \text{ dispari}\} = \Z, \quad
  (2) \cap (3) = (6), \quad
  (4) \cap (6) = (12).
\]
La prima uguaglianza dice che ogni intero è pari o dispari. La seconda dice che un intero divisibile sia per $2$ che per $3$ è divisibile per $6$. La terza dice che un multiplo sia di $4$ che di $6$ è un multiplo di $12$ (il minimo comune multiplo).
\end{esempio}
\end{exbox}

\begin{osservazione}
Per il Principio di estensione, $X = Y$ se e solo se $X \subseteq Y$ e $Y \subseteq X$. Questo è il modo standard per dimostrare l'uguaglianza di due insiemi: si dimostrano le due inclusioni separatamente.
\end{osservazione}

Queste operazioni si generalizzano a famiglie arbitrarie di insiemi $\{X_i\}_{i \in I}$:
\[
  \bigcup_{i \in I} X_i \coloneqq \{x \mid x \in X_i \text{ per qualche } i \in I\}, \quad
  \bigcap_{i \in I} X_i \coloneqq \{x \mid x \in X_i \text{ per ogni } i \in I\}.
\]

\begin{defbox}
\begin{definizione}[Prodotto cartesiano]
Siano $X_1, \ldots, X_n$ insiemi. Il \emph{prodotto cartesiano} $X_1 \times \cdots \times X_n$ è l'insieme delle $n$-uple ordinate $(x_1, \ldots, x_n)$ con $x_i \in X_i$. Se $X_1 = \cdots = X_n = X$, scriviamo $X^n$.
\end{definizione}
\end{defbox}

\begin{exbox}
\begin{esempio}
$\R^n$ è l'insieme delle $n$-uple di numeri reali: è lo spazio in cui vivranno i vettori colonna che useremo in tutto il corso. Il piano cartesiano è $\R^2$, lo spazio tridimensionale è $\R^3$.
\end{esempio}
\end{exbox}

% --------------------------------------------------------
\section{Applicazioni}
% --------------------------------------------------------

\subsection{Definizione e prime proprietà}

Le \emph{applicazioni} (o \emph{funzioni}) sono il meccanismo con cui si collegano insiemi diversi. Formalizziamo un'idea intuitiva: una ``regola'' che a ogni elemento del dominio associa un elemento del codominio.

\begin{defbox}
\begin{definizione}
Un'\emph{applicazione} $f \colon X \to Y$ è una legge che associa a ogni $x \in X$ un unico elemento $f(x) \in Y$. L'insieme $X$ è il \emph{dominio} e $Y$ è il \emph{codominio} di $f$.

Due applicazioni $f_1 \colon X_1 \to Y_1$ e $f_2 \colon X_2 \to Y_2$ sono uguali se e solo se $X_1 = X_2$, $Y_1 = Y_2$, e $f_1(x) = f_2(x)$ per ogni $x$.
\end{definizione}
\end{defbox}

\begin{exbox}
\begin{esempio}
Ecco alcune applicazioni importanti:
\begin{enumerate}
  \item $\Id_X \colon X \to X$, $x \mapsto x$: l'applicazione \emph{identità}.
  \item $f \colon \Z \to \Z$, $x \mapsto x - 5$: uno spostamento.
  \item $g \colon \R \times \R \to \R$, $(a,b) \mapsto ab$: la moltiplicazione.
  \item $h \colon \R \to \R$, $x \mapsto x^3 - x$: una funzione polinomiale.
\end{enumerate}
\end{esempio}
\end{exbox}

\begin{defbox}
\begin{definizione}
Un'applicazione $f \colon X \to Y$ è:
\begin{itemize}
  \item \emph{iniettiva} se $f(x_1) = f(x_2)$ implica $x_1 = x_2$ (elementi distinti hanno immagini distinte);
  \item \emph{suriettiva} se $f(X) = Y$ (ogni elemento del codominio è raggiunto);
  \item \emph{biunivoca} (o \emph{biettiva}) se è sia iniettiva che suriettiva.
\end{itemize}
\end{definizione}
\end{defbox}

\begin{exbox}
\begin{esempio}
Consideriamo le applicazioni $f, g, h \colon \R \to \R$ definite da
\[
  f(x) = x^2 + 1, \quad g(x) = x^3, \quad h(x) = x^3 - x.
\]
\begin{itemize}
  \item $f$ non è iniettiva ($f(1) = f(-1) = 2$) né suriettiva (l'immagine è $[1,+\infty)$).
  \item $g$ è biunivoca: $g^{-1}(y) = y^{1/3}$.
  \item $h$ è suriettiva ma non iniettiva: $h(0) = h(1) = h(-1) = 0$ (più avanti si verificherebbe, ma $h(-1)=(-1)^3 - (-1) = 0$, $h(0)=0$, $h(1)=0$).
\end{itemize}
\end{esempio}
\end{exbox}

\begin{osservazione}
Dominio e codominio fanno parte della definizione di applicazione. L'applicazione $f(x) = x^2 + 1$ con codominio $\R$ non è suriettiva, ma se restringiamo il codominio a $[1,+\infty)$ diventa suriettiva. Formalmente, si tratta di applicazioni \emph{diverse}.
\end{osservazione}

\subsection{Composizione e inversa}

\begin{defbox}
\begin{definizione}
Siano $f \colon X \to Y$ e $g \colon Y \to Z$. La \emph{composizione} $g \circ f \colon X \to Z$ è definita da
\[
  (g \circ f)(x) \coloneqq g(f(x)).
\]
\end{definizione}
\end{defbox}

\begin{osservazione}
L'ordine nella composizione è importante: $g \circ f$ si legge ``prima $f$, poi $g$''. In generale $f \circ g \neq g \circ f$.
\end{osservazione}

La composizione è \emph{associativa}: $(h \circ g) \circ f = h \circ (g \circ f)$. Inoltre $f \circ \Id_X = \Id_Y \circ f = f$.

\begin{defbox}
\begin{definizione}
Sia $f \colon X \to Y$ biunivoca. L'\emph{inversa} $f^{-1} \colon Y \to X$ associa a ogni $y \in Y$ l'unico $x \in X$ con $f(x) = y$. Si ha $f \circ f^{-1} = \Id_Y$ e $f^{-1} \circ f = \Id_X$.
\end{definizione}
\end{defbox}

\subsection{Immagine, controimmagine e cardinalità}

Dati $f \colon X \to Y$, $A \subseteq X$, $B \subseteq Y$:
\[
  f(A) \coloneqq \{y \in Y \mid \exists\, x \in A : f(x) = y\} \quad \text{(immagine di } A\text{)},
\]
\[
  f^{-1}(B) \coloneqq \{x \in X \mid f(x) \in B\} \quad \text{(controimmagine di } B\text{)}.
\]
L'\emph{immagine} di $f$ è $\im f \coloneqq f(X)$.

\medskip
Si dice che $X$ ha la stessa \emph{cardinalità} di $Y$ (scritto $X \approx Y$) se esiste una biiezione $f \colon X \to Y$. Per insiemi finiti questo equivale ad avere lo stesso numero di elementi. Un risultato non banale è:

\begin{teorema}[Schröder--Bernstein]
Siano $X$, $Y$ insiemi tali che esiste un'applicazione iniettiva $X \to Y$ e una iniettiva $Y \to X$. Allora $X \approx Y$.
\end{teorema}

% --------------------------------------------------------
\section{Relazioni}
% --------------------------------------------------------

\subsection{Relazioni d'ordine e di equivalenza}

Una \emph{relazione} su un insieme $X$ è semplicemente un sottoinsieme $R \subseteq X \times X$: diciamo che $x_1 R x_2$ se $(x_1, x_2) \in R$.

\begin{defbox}
\begin{definizione}
Una relazione $R$ su $X$ è:
\begin{itemize}
  \item \emph{di ordine} se è riflessiva ($xRx$), antisimmetrica ($xRy$ e $yRx \Rightarrow x=y$) e transitiva ($xRy$ e $yRz \Rightarrow xRz$). È \emph{totale} se per ogni $x, y$ vale almeno una tra $xRy$ e $yRx$.
  \item \emph{di equivalenza} se è riflessiva, \emph{simmetrica} ($xRy \Rightarrow yRx$) e transitiva.
\end{itemize}
\end{definizione}
\end{defbox}

\begin{exbox}
\begin{esempio}[Congruenza modulo $n$]
Fissato $n \in \Z$, diciamo che $x \equiv y \pmod{n}$ se $n \mid (x - y)$, cioè $(x-y) \in (n)$. Questa è una relazione di equivalenza su $\Z$. Per esempio, modulo $10$, due interi sono equivalenti se e solo se hanno la stessa ultima cifra decimale.
\end{esempio}
\end{exbox}

\subsection{Classi di equivalenza e insiemi quoziente}

Data una relazione di equivalenza $\sim$ su $X$, la \emph{classe di equivalenza} di $x_0$ è
\[
  [x_0] \coloneqq \{x \in X \mid x \sim x_0\}.
\]
L'\emph{insieme quoziente} $X/{\sim}$ è l'insieme di tutte le classi di equivalenza.

\begin{defbox}
\begin{proposizione}
Le classi di equivalenza di $\sim$ formano una \emph{partizione} di $X$: ogni elemento di $X$ appartiene a esattamente una classe.
\end{proposizione}
\end{defbox}

\begin{osservazione}[Fondamentale]
Data un'applicazione $f \colon X \to Y$, si dice che $f$ \emph{discende} al quoziente $X/{\sim}$ se $x_1 \sim x_2 \Rightarrow f(x_1) = f(x_2)$. In tal caso esiste un'unica $\bar{f} \colon X/{\sim} \to Y$ tale che $f = \bar{f} \circ \pi$, dove $\pi \colon X \to X/{\sim}$, $x \mapsto [x]$, è l'applicazione quoziente.
\end{osservazione}

\begin{exbox}
\begin{esempio}
L'insieme quoziente $\Z/(n)$ ha $n$ elementi: $[0], [1], \ldots, [n-1]$. Questa struttura sarà cruciale quando costruiremo i campi finiti $\mathbb{F}_p$.
\end{esempio}
\end{exbox}

% --------------------------------------------------------
\section{Induzione matematica}
% --------------------------------------------------------

L'\emph{induzione matematica} è uno dei principali strumenti di dimostrazione in matematica. L'idea è: se una proprietà vale per $n = 0$ (o $n = 1$) e se vale per $n$ implica che vale per $n+1$, allora vale per tutti i naturali.

\begin{defbox}
\begin{assioma}[di Peano -- Principio di induzione]
Sia $X \subseteq \N$ un insieme tale che $0 \in X$ e, ogni volta che $n \in X$, anche $n+1 \in X$. Allora $X = \N$.
\end{assioma}
\end{defbox}

\begin{exbox}
\begin{esempio}[Somma dei primi $n$ interi]
Dimostriamo per induzione che per ogni $n \geq 1$:
\[
  1 + 2 + \cdots + n = \frac{n(n+1)}{2}. \tag{$\star$}
\]
\textbf{Base}: per $n = 1$, il membro sinistro è $1$ e il destro è $\frac{1 \cdot 2}{2} = 1$. \checkmark

\textbf{Passo induttivo}: supponiamo che ($\star$) valga per $n$. Allora:
\[
  1 + 2 + \cdots + n + (n+1) = \frac{n(n+1)}{2} + (n+1) = (n+1)\!\left(\frac{n}{2} + 1\right) = \frac{(n+1)(n+2)}{2}.
\]
Questa è esattamente ($\star$) con $n$ sostituito da $n+1$. \checkmark

Per il principio di induzione, ($\star$) vale per ogni $n \geq 1$.
\end{esempio}
\end{exbox}

\begin{osservazione}
Una variante utile: nel passo induttivo si può assumere che la proprietà $P_m$ valga per \emph{tutti} gli $m \leq n-1$ e dedurne la validità per $m = n$. Questo è il cosiddetto \emph{principio di induzione forte}.
\end{osservazione}

% --------------------------------------------------------
\section{L'algoritmo euclideo}
% --------------------------------------------------------

\subsection{Divisione con resto e massimo comun divisore}

La divisione con resto è un fatto elementare ma fondamentale dell'aritmetica degli interi.

\begin{proposizione}[Divisione con resto]
Siano $a, b \in \Z$ con $b \neq 0$. Esistono unici $q \in \Z$ (quoziente) e $r \in \N$ (resto) tali che
\[
  a = b \cdot q + r, \quad 0 \leq r < |b|.
\]
\end{proposizione}

\begin{defbox}
\begin{teorema}[Massimo comun divisore]
Siano $a, b \in \Z$, non entrambi nulli. Esiste un unico (a meno di segno) $m \in \Z$ tale che:
\begin{enumerate}
  \item $m \mid a$ e $m \mid b$;
  \item se $d \mid a$ e $d \mid b$, allora $d \mid m$;
  \item esistono $x, y \in \Z$ con $ax + by = m$ \emph{(identità di Bézout)}.
\end{enumerate}
Il \emph{massimo comun divisore} è $\mcd\{a,b\} \coloneqq \max\{d \in \Z \mid d \mid a,\, d \mid b\}$.
\end{teorema}
\end{defbox}

\begin{osservazione}
L'identità di Bézout è tutt'altro che ovvia: dice che il massimo comun divisore di $a$ e $b$ è una \emph{combinazione lineare intera} di $a$ e $b$. Per esempio $\mcd\{6, 10\} = 2 = (-1) \cdot 6 + 1 \cdot 10$. Questo fatto è la chiave per dimostrare l'invertibilità in $\mathbb{F}_p$.
\end{osservazione}

\subsection{L'algoritmo}

L'algoritmo euclideo calcola $\mcd\{a,b\}$ sfruttando le seguenti tre proprietà:
\begin{enumerate}
  \item[(F1)] $\mcd\{a,b\} = \mcd\{b,a\}$,
  \item[(F2)] se $b = 0$, allora $\mcd\{a,b\} = |a|$,
  \item[(F3)] se $a = b \cdot q + r$ (divisione con resto), allora $\mcd\{a,b\} = \mcd\{b,r\}$.
\end{enumerate}

\begin{exbox}
\begin{esempio}
Calcoliamo $\mcd\{861, 308\}$:
\[
  861 = 308 \cdot 2 + 245, \quad
  308 = 245 \cdot 1 + 63, \quad
  245 = 63 \cdot 3 + 56, \quad
  63 = 56 \cdot 1 + 7, \quad
  56 = 7 \cdot 8 + 0.
\]
Quindi $\mcd\{861, 308\} = 7$. La sequenza dei resti $245, 63, 56, 7, 0$ decresce strettamente, garantendo la terminazione.
\end{esempio}
\end{exbox}

% --------------------------------------------------------
\section{Anelli e campi}
\label{sec:anelliCampi}
% --------------------------------------------------------

\subsection{Motivazione}

Fino ad ora abbiamo lavorato con strutture concrete ($\Z$, $\Q$, $\R$). Ma in algebra moderna vogliamo trattare in modo \emph{uniforme} sistemi numerici diversi, identificando le proprietà che \emph{davvero} contano. Un \emph{anello} cattura l'essenziale della somma e della moltiplicazione; un \emph{campo} aggiunge la divisione. Questa astrazione ci permetterà, nei capitoli successivi, di definire spazi vettoriali su $\R$, su $\C$, ma anche su $\mathbb{F}_p$ -- tutti con la stessa teoria.

\subsection{Anelli}

\begin{defbox}
\begin{definizione}[Anello]
Un \emph{anello} è un insieme $A$ con due operazioni $+$ (somma) e $\cdot$ (moltiplicazione) tali che:
\begin{enumerate}
  \item $(A, +)$ è un \emph{gruppo abeliano}: esiste $0 \in A$ neutro per la somma, la somma è commutativa e associativa, ogni elemento ha un opposto.
  \item La moltiplicazione è \emph{associativa}: $(a \cdot b) \cdot c = a \cdot (b \cdot c)$.
  \item Vale la \emph{distributività}: $a \cdot (b + c) = a \cdot b + a \cdot c$.
\end{enumerate}
L'anello è \emph{commutativo} se $a \cdot b = b \cdot a$ per ogni $a, b$. Ha un'\emph{unità} se esiste $1 \neq 0$ con $1 \cdot a = a \cdot 1 = a$ per ogni $a$.

D'ora in poi, con ``anello'' intendiamo un anello commutativo con unità, salvo avviso contrario.
\end{definizione}
\end{defbox}

\begin{exbox}
\begin{esempio}
\begin{itemize}
  \item $\Z$, $\Q$, $\R$, $\C$ con le operazioni usuali sono anelli commutativi con unità.
  \item $\N$ \emph{non} è un anello (mancano gli opposti additivi: non esiste $n \in \N$ con $1 + n = 0$).
  \item L'insieme dei numeri interi pari $2\Z$ è un anello (commutativo) ma \emph{senza} unità ($1 \notin 2\Z$).
  \item L'anello $\Z/n\Z$ (costruito in Sezione~\ref{sec:Zmodulo}) è un anello commutativo con unità, finito.
\end{itemize}
\end{esempio}
\end{exbox}

\begin{proposizione}
Sia $A$ un anello. Allora $0$ e $1$ sono unici, e per ogni $a \in A$ vale $0 \cdot a = 0$.
\end{proposizione}

\begin{proof}
Per l'unicità di $0$: se $0'$ è un altro neutro per la somma, allora $0 = 0 + 0' = 0'$. Analogamente per $1$.

Per $0 \cdot a = 0$: si ha $0 \cdot a = (0 + 0) \cdot a = 0 \cdot a + 0 \cdot a$. Aggiungendo l'opposto di $0 \cdot a$ a entrambi i membri, si ottiene $0 = 0 \cdot a$.
\end{proof}

\subsection{Campi}

\begin{defbox}
\begin{definizione}[Campo]
Un anello $K$ è un \emph{campo} se ogni elemento $z \neq 0$ ha un inverso moltiplicativo: esiste $w \in K$ con $w \cdot z = 1$. Denotiamo tale inverso $z^{-1}$.
\end{definizione}
\end{defbox}

In un campo si può dividere per qualsiasi elemento non nullo. I campi più comuni sono $\Q$, $\R$, $\C$. Ma $\Z$ \emph{non} è un campo: $2$ non ha inverso moltiplicativo in $\Z$ (non esiste $n \in \Z$ con $2n = 1$).

\begin{proposizione}
In un campo $K$: l'inverso di ogni $z \neq 0$ è unico; vale la legge di cancellazione ($zw = zw' \Rightarrow w = w'$ se $z \neq 0$); se $zw = 0$ allora $z = 0$ oppure $w = 0$.
\end{proposizione}

\subsection{Esempi: sottocampi e estensioni algebriche}

\begin{defbox}
\begin{definizione}
Sia $A$ un anello e $B \subseteq A$. $B$ è un \emph{sottoanello} se $0, 1 \in B$ e $b_1 - b_2, b_1 b_2 \in B$ per ogni $b_1, b_2 \in B$. Se $A$ è un campo, $B$ è un \emph{sottocampo} se inoltre $b^{-1} \in B$ per ogni $0 \neq b \in B$.
\end{definizione}
\end{defbox}

\begin{defbox}
\begin{definizione}
Siano $K \subseteq F$ campi e $\alpha \in F$. Definiamo
\[
  K[\alpha] \coloneqq \{c_0 + c_1 \alpha + \cdots + c_n \alpha^n \mid c_i \in K\}.
\]
Diciamo che $\alpha$ è \emph{algebrico su $K$} se soddisfa un'equazione $c_0 \alpha^m + c_1 \alpha^{m-1} + \cdots + c_m = 0$ con $c_i \in K$ e $c_0 \neq 0$; altrimenti $\alpha$ è \emph{trascendente su $K$}.
\end{definizione}
\end{defbox}

\begin{exbox}
\begin{esempio}
\textbf{($\Q[\sqrt{2}]$ è un campo.)}
Poiché $(\sqrt{2})^2 = 2 \in \Q$, si ha $\Q[\sqrt{2}] = \{a + b\sqrt{2} \mid a, b \in \Q\}$. Dobbiamo verificare che ogni elemento non nullo ha inverso. Se $a + b\sqrt{2} \neq 0$, allora $a^2 - 2b^2 \neq 0$ (perché altrimenti $\sqrt{2} = a/b$ sarebbe razionale), e
\[
  (a + b\sqrt{2})^{-1} = \frac{a}{a^2 - 2b^2} - \frac{b}{a^2 - 2b^2}\sqrt{2} \in \Q[\sqrt{2}].
\]
\end{esempio}
\end{exbox}

\begin{teorema}
Siano $K \subseteq F$ campi e $\alpha \in F$. Il sottoanello $K[\alpha]$ è un campo se e solo se $\alpha$ è algebrico su $K$.
\end{teorema}

Questo teorema spiega perché $\Q[\pi]$ non è un campo (Lindemann, 1882, ha dimostrato che $\pi$ è trascendente su $\Q$), mentre $\Q[\sqrt{d}]$ per $d \in \Q$ positivo non un quadrato perfetto razionale è sempre un campo.

\subsection{Campi finiti: $\mathbb{F}_p$}
\label{sec:Zmodulo}

Fino ad ora i campi che abbiamo incontrato ($\Q$, $\R$, $\C$) erano infiniti. Esistono però campi con un numero finito di elementi, con applicazioni fondamentali in crittografia, teoria dei codici e combinatoria.

\begin{osservazione}
Il Lemma~\ref{lem:congruenza} qui sotto garantisce che le operazioni di $\Z$ siano compatibili con la congruenza, rendendo $\Z/n\Z$ un anello.
\end{osservazione}

\begin{lemma}
\label{lem:congruenza}
Se $a \equiv a' \pmod{n}$ e $b \equiv b' \pmod{n}$, allora $a + b \equiv a' + b' \pmod{n}$ e $ab \equiv a'b' \pmod{n}$.
\end{lemma}

\begin{proof}
Esistono $s, t \in \Z$ con $a' = a + sn$ e $b' = b + tn$. Allora $a' + b' = (a+b) + (s+t)n$ e $a'b' = ab + (at + sb + stn)n$, che sono congruenti ad $a+b$ e $ab$ rispettivamente.
\end{proof}

Quindi su $\Z/(n)$ si definiscono $[a] + [b] \coloneqq [a+b]$ e $[a] \cdot [b] \coloneqq [ab]$. L'anello risultante è un campo se e solo se $n$ è primo.

\begin{proposizione}
Se $p$ è primo, allora $\mathbb{F}_p \coloneqq \Z/(p)$ è un campo.
\end{proposizione}

\begin{proof}
Sia $[a] \neq [0]$, cioè $p \nmid a$. Poiché $p$ è primo, $\mcd\{a, p\} = 1$. Per l'identità di Bézout esistono $x, y \in \Z$ con $ax + py = 1$. Allora $[a] \cdot [x] = [ax] = [1 - py] = [1]$, quindi $[x]$ è l'inverso di $[a]$.
\end{proof}

\begin{exbox}
\begin{esempio}
In $\mathbb{F}_5 = \{[0],[1],[2],[3],[4]\}$:
\begin{center}
\begin{tabular}{c|ccccc}
  $\cdot$ & $[1]$ & $[2]$ & $[3]$ & $[4]$ \\
  \hline
  $[2]$   & $[2]$ & $[4]$ & $[1]$ & $[3]$ \\
  $[3]$   & $[3]$ & $[1]$ & $[4]$ & $[2]$ \\
\end{tabular}
\end{center}
In particolare $[2]^{-1} = [3]$ perché $[2] \cdot [3] = [6] = [1]$.
\end{esempio}
\end{exbox}

\subsubsection{La caratteristica di un campo}

\begin{defbox}
\begin{definizione}
La \emph{caratteristica} di un campo $K$ è $0$ se $n \cdot 1 \neq 0$ per ogni $n \in \Z \setminus \{0\}$, altrimenti è il minimo $n > 0$ con $n \cdot 1 = 0$.
\end{definizione}
\end{defbox}

\begin{osservazione}
La caratteristica di un campo, se non zero, è sempre un numero primo. Infatti se $\mathrm{char}\, K = n = ab$ con $1 < a, b < n$, allora $0 = n \cdot 1 = (a \cdot 1)(b \cdot 1)$, e per la legge del prodotto nullo uno dei fattori sarebbe zero, contraddicendo la minimalità di $n$.

Per esempio: $\mathrm{char}\, \Q = \mathrm{char}\, \R = \mathrm{char}\, \C = 0$; $\mathrm{char}\, \mathbb{F}_p = p$.
\end{osservazione}

\subsection{Omomorfismi e isomorfismi di anelli}

\begin{defbox}
\begin{definizione}
Un'\emph{applicazione} $f \colon A \to B$ tra anelli è un \emph{omomorfismo} se $f(1) = 1$ e per ogni $a_1, a_2 \in A$:
\[
  f(a_1 + a_2) = f(a_1) + f(a_2), \quad f(a_1 \cdot a_2) = f(a_1) \cdot f(a_2).
\]
Un omomorfismo biunivoco è un \emph{isomorfismo}. Un isomorfismo $f \colon A \to A$ è un \emph{automorfismo}.
\end{definizione}
\end{defbox}

\begin{exbox}
\begin{esempio}
\begin{enumerate}
  \item La proiezione $\Z \to \Z/(n)$, $a \mapsto [a]$, è un omomorfismo suriettivo (non iniettivo se $n > 1$).
  \item L'applicazione $f \colon \Q[\sqrt{2}] \to \Q[\sqrt{2}]$ con $f(a + b\sqrt{2}) = a - b\sqrt{2}$ è un automorfismo (la \emph{coniugazione} di $\Q[\sqrt{2}]$).
\end{enumerate}
\end{esempio}
\end{exbox}

\begin{lemma}
Un omomorfismo di campi $f \colon K \to F$ è sempre iniettivo.
\end{lemma}

\begin{proof}
Se $f$ non fosse iniettivo, esisterebbero $x \neq y$ con $f(x) = f(y)$, cioè $f(x - y) = 0$ con $x - y \neq 0$. Ma allora $1 = f(1) = f((x-y)(x-y)^{-1}) = f(x-y) \cdot f((x-y)^{-1}) = 0$, contraddizione.
\end{proof}

% --------------------------------------------------------
\section{Polinomi e funzioni razionali}
% --------------------------------------------------------

\subsection{Polinomi in una indeterminata}

Dato un campo $K$, l'\emph{anello dei polinomi} $K[x]$ è l'insieme delle espressioni formali
\[
  p = a_0 + a_1 x + a_2 x^2 + \cdots + a_d x^d, \quad a_i \in K.
\]
con somma e prodotto definiti dalle regole usuali. Il \emph{grado} di $p$ (se $a_d \neq 0$) è $\deg p = d$; per convenzione $\deg 0 = -\infty$.

\begin{osservazione}
$K[x]$ è un anello ma non un campo: $x$ non ha inverso (non esiste $q \in K[x]$ con $xq = 1$). Tuttavia, come vedremo, condivide molte proprietà con $\Z$: ammette divisione con resto, massimo comun divisore, e una nozione di ``elementi irriducibili'' analoga ai numeri primi.
\end{osservazione}

\subsection{Divisione con resto tra polinomi}

\begin{proposizione}[Divisione con resto]
Siano $a, b \in K[x]$ con $b \neq 0$. Esistono unici $q, r \in K[x]$ con
\[
  a = b \cdot q + r, \quad \deg r < \deg b.
\]
\end{proposizione}

\begin{exbox}
\begin{esempio}
Dividiamo $a = x^3 + 2x + 1$ per $b = x^2 - 1$ in $\Q[x]$:
\[
  x^3 + 2x + 1 = (x^2 - 1) \cdot x + (3x + 1).
\]
Quindi $q = x$ e $r = 3x + 1$. Verifica: $\deg r = 1 < 2 = \deg b$. \checkmark
\end{esempio}
\end{exbox}

\subsection{Radici e fattorizzazione}

\begin{defbox}
\begin{definizione}
Un $\alpha \in K$ è una \emph{radice} (o \emph{zero}) di $p \in K[x]$ se $p(\alpha) = 0$. La \emph{molteplicità} di $\alpha$ come radice di $p$ è il massimo $n \in \N$ tale che $(x - \alpha)^n \mid p$.
\end{definizione}
\end{defbox}

\begin{lemma}[Ruffini]
$\alpha \in K$ è radice di $p \in K[x]$ se e solo se $(x - \alpha)$ divide $p$.
\end{lemma}

\begin{proof}
Se $(x-\alpha) \mid p$, allora $p = (x-\alpha)q$ e $p(\alpha) = 0$. Viceversa, scriviamo $p(x) = \sum_{i=0}^n c_i (x-\alpha)^i$ (sostituendo $x = (x-\alpha)+\alpha$ e espandendo); allora $p(\alpha) = c_0 = 0$ implica che $(x-\alpha)$ divide $p$.
\end{proof}

\begin{proposizione}
Un polinomio non nullo $p \in K[x]$ di grado $n$ ha al più $n$ radici (contate con molteplicità).
\end{proposizione}

\begin{osservazione}
Questo risultato è fondamentale: due polinomi di grado $\leq d$ che coincidono su più di $d$ punti devono essere identici. Se $K$ è infinito, due funzioni polinomiali uguali su tutto $K$ hanno gli stessi coefficienti.
\end{osservazione}

\subsection{Campo delle funzioni razionali}

Analogamente alla costruzione di $\Q$ da $\Z$, si costruisce il \emph{campo delle funzioni razionali} $K(x)$: le sue elementi sono le classi di equivalenza di coppie $(p, q)$ con $p, q \in K[x]$, $q \neq 0$, dove $(p, q) \sim (a, b)$ se $pb = aq$. La classe di $(p,q)$ si scrive $p/q$.

% --------------------------------------------------------
\section{Fattoriali e coefficienti binomiali}
% --------------------------------------------------------

\subsection{Fattoriale e combinatorica}

Per $n \in \N$, $n \geq 1$, il \emph{fattoriale} è
\[
  n! \coloneqq n \cdot (n-1) \cdots 2 \cdot 1, \quad 0! \coloneqq 1.
\]

\begin{proposizione}
Il numero di ordinamenti totali di un insieme con $n$ elementi è $n!$.
\end{proposizione}

Il \emph{coefficiente binomiale} è
\[
  \binom{n}{i} \coloneqq \frac{n!}{i!(n-i)!} = \frac{n(n-1)\cdots(n-i+1)}{i!}.
\]

\begin{proposizione}
Il numero di sottoinsiemi di cardinalità $i$ di un insieme con $n$ elementi è $\binom{n}{i}$.
\end{proposizione}

\begin{exbox}
\begin{esempio}
In quanti modi si possono scegliere 2 studenti da un gruppo di 5? La risposta è $\binom{5}{2} = \frac{5!}{2! \cdot 3!} = \frac{5 \cdot 4}{2} = 10$.
\end{esempio}
\end{exbox}

\subsection{Formula del binomio di Newton}

\begin{proposizione}[Newton]
Per ogni anello $A$, ogni $x, y \in A$ e ogni $n \in \N$:
\[
  (x + y)^n = \sum_{i=0}^n \binom{n}{i} x^{n-i} y^i.
\]
\end{proposizione}

\begin{proof}
Per induzione su $n$. Il caso $n=0$ è banale. Per il passo induttivo si moltiplica l'uguaglianza per $n$ per $(x+y)$ e si usa il fatto che $\binom{n}{j} + \binom{n}{j-1} = \binom{n+1}{j}$ (identità di Pascal).
\end{proof}

% --------------------------------------------------------
\section{Numeri complessi}
\label{sec:complessi}
% --------------------------------------------------------

\subsection{Costruzione e operazioni}

Storicamente i numeri complessi nacquero per dare senso alle radici di equazioni come $x^2 + 1 = 0$. La costruzione rigorosa è elegante: si prende $\R^2$ e si definisce una moltiplicazione su misura.

\begin{defbox}
\begin{definizione}
L'insieme dei \emph{numeri complessi} è $\C \coloneqq \R^2$ con operazioni:
\[
  (a_1, b_1) + (a_2, b_2) \coloneqq (a_1 + a_2, b_1 + b_2),
\]
\[
  (a_1, b_1) \cdot (a_2, b_2) \coloneqq (a_1 a_2 - b_1 b_2,\; a_1 b_2 + a_2 b_1).
\]
Con queste operazioni, $\C$ è un campo. Identificando $a \in \R$ con $(a, 0)$ e ponendo $i \coloneqq (0, 1)$, ogni numero complesso si scrive $z = a + bi$ con $a, b \in \R$.
\end{definizione}
\end{defbox}

\begin{osservazione}
Perché questa moltiplicazione? Proprio perché $i \cdot i = (0,1)(0,1) = (-1, 0) = -1$: l'unità immaginaria $i$ è una radice di $x^2 + 1 = 0$.
\end{osservazione}

Se $z = a + bi$, la \emph{parte reale} è $\Re(z) = a$, la \emph{parte immaginaria} è $\Im(z) = b$, il \emph{modulo} è $|z| = \sqrt{a^2 + b^2}$, e il \emph{coniugato} è $\bar{z} = a - bi$.

\begin{exbox}
\begin{esempio}
L'inverso di $z = a + bi \neq 0$ è
\[
  z^{-1} = \frac{\bar{z}}{|z|^2} = \frac{a}{a^2+b^2} - \frac{b}{a^2+b^2}i.
\]
Per esempio $(1+i)^{-1} = \frac{1-i}{2} = \frac{1}{2} - \frac{1}{2}i$.
\end{esempio}
\end{exbox}

\subsection{Forma polare e radici $n$-esime}

Ogni $z \neq 0$ si può scrivere in \emph{forma polare}:
\[
  z = \rho(\cos\theta + i\sin\theta), \quad \rho = |z|,\; \theta = \mathrm{Arg}(z).
\]
Le formule di addizione per seno e coseno danno la regola di moltiplicazione:
\[
  |z_1 z_2| = |z_1| \cdot |z_2|, \quad \mathrm{Arg}(z_1 z_2) = \mathrm{Arg}(z_1) + \mathrm{Arg}(z_2).
\]

\begin{exbox}
\begin{esempio}[Radici $n$-esime dell'unità]
Le $n$ radici $n$-esime di $1$ sono i numeri complessi $\zeta_k = e^{2\pi i k/n} = \cos\frac{2\pi k}{n} + i\sin\frac{2\pi k}{n}$ per $k = 0, 1, \ldots, n-1$. Nel piano complesso formano i vertici di un poligono regolare con $n$ lati inscritto nel cerchio unitario.
\end{esempio}
\end{exbox}

\subsection{Il teorema fondamentale dell'algebra}

\begin{teorema}[Fondamentale dell'Algebra]
Ogni polinomio $p \in \C[z]$ di grado $n \geq 1$ ha almeno una radice in $\C$.
\end{teorema}

Applicando ripetutamente il Lemma di Ruffini, si ottiene che ogni polinomio di grado $n$ su $\C$ si fattorizza completamente in fattori lineari:
\[
  z^n + a_1 z^{n-1} + \cdots + a_n = (z - c_1)(z - c_2) \cdots (z - c_n), \quad c_i \in \C.
\]

\begin{osservazione}
$\C$ è l'unico campo che ``completa'' $\R$ in questo senso: ogni polinomio con coefficienti reali si fattorizza su $\C$. Questa proprietà si chiama \emph{chiusura algebrica}: $\C$ è algebricamente chiuso. $\R$ non lo è ($x^2 + 1$ non ha radici reali).
\end{osservazione}

% --------------------------------------------------------
\section{Gruppi}
% --------------------------------------------------------

\subsection{Definizione e motivazione}

I \emph{gruppi} astraggono la nozione di \emph{simmetria}. Ogni volta che abbiamo un insieme di trasformazioni che si possono comporre, dove ogni trasformazione ha un'inversa e la composizione è associativa, siamo in presenza di un gruppo.

\begin{defbox}
\begin{definizione}[Gruppo]
Un \emph{gruppo} è un insieme $G$ con un'operazione $\cdot \colon G \times G \to G$ tale che:
\begin{enumerate}
  \item Esiste un elemento \emph{neutro} $1 \in G$ con $1 \cdot g = g \cdot 1 = g$ per ogni $g$.
  \item Ogni $g \in G$ ha un \emph{inverso} $g^{-1} \in G$ con $g \cdot g^{-1} = g^{-1} \cdot g = 1$.
  \item L'operazione è \emph{associativa}: $(g \cdot h) \cdot k = g \cdot (h \cdot k)$.
\end{enumerate}
Il gruppo è \emph{abeliano} (o \emph{commutativo}) se $g \cdot h = h \cdot g$ per ogni $g, h$.
\end{definizione}
\end{defbox}

\begin{exbox}
\begin{esempio}
\begin{enumerate}
  \item $(A, +)$ per ogni anello $A$: gruppo abeliano.
  \item $K^\times = K \setminus \{0\}$ per ogni campo $K$: gruppo (abeliano) con la moltiplicazione.
  \item $S(X)$, l'insieme delle biiezioni $X \to X$ con la composizione: gruppo (\emph{gruppo simmetrico} di $X$). Per $|X| \geq 3$ non è abeliano.
  \item $\mathrm{Isom}(E^2)$, le isometrie del piano euclideo con la composizione: gruppo.
\end{enumerate}
\end{esempio}
\end{exbox}

\begin{proposizione}
In un gruppo $G$: l'elemento neutro è unico; l'inverso di ogni $g$ è unico; vale la legge di cancellazione ($gh = gk \Rightarrow h = k$).
\end{proposizione}

\subsection{Sottogruppi}

\begin{defbox}
\begin{definizione}
Un \emph{sottogruppo} di $G$ è un sottoinsieme $H \subseteq G$ non vuoto che contiene $1$ ed è chiuso per l'operazione e per gli inversi.
\end{definizione}
\end{defbox}

\begin{exbox}
\begin{esempio}
$\Z \subset \Q$ è un sottogruppo rispetto alla somma. Per ogni $n$, l'insieme $(n) = n\Z$ dei multipli di $n$ è un sottogruppo di $(\Z, +)$.
\end{esempio}
\end{exbox}

\subsection{Omomorfismi di gruppi e azioni}

\begin{defbox}
\begin{definizione}
Un'applicazione $f \colon G \to H$ tra gruppi è un \emph{omomorfismo} se $f(g_1 \cdot g_2) = f(g_1) \cdot f(g_2)$ per ogni $g_1, g_2$. Un omomorfismo biunivoco è un \emph{isomorfismo}.
\end{definizione}
\end{defbox}

\begin{lemma}
Se $f \colon G \to H$ è un omomorfismo, allora $f(1) = 1$ e $f(g^{-1}) = f(g)^{-1}$.
\end{lemma}

\begin{exbox}
\begin{esempio}
$\log \colon (\R^+, \cdot) \to (\R, +)$ è un isomorfismo di gruppi: la proprietà $\log(ab) = \log a + \log b$ è proprio la condizione di omomorfismo. L'inverso è $\exp \colon (\R, +) \to (\R^+, \cdot)$, isomorfismo grazie a $e^{a+b} = e^a \cdot e^b$.
\end{esempio}
\end{exbox}

\begin{defbox}
\begin{definizione}
Un'\emph{azione} di $G$ su un insieme $X$ è un omomorfismo $\phi \colon G \to S(X)$, equivalentemente un'applicazione $G \times X \to X$, $(g,x) \mapsto gx$, tale che $1 \cdot x = x$ e $(gh)x = g(hx)$.
\end{definizione}
\end{defbox}

\subsection{Il gruppo simmetrico $S_n$ e il prodotto diretto}

\begin{defbox}
\begin{definizione}
Il \emph{gruppo simmetrico} $S_n$ è il gruppo delle biiezioni di $\{1, \ldots, n\}$ in sé (con la composizione). Ha $n!$ elementi.
\end{definizione}
\end{defbox}

\begin{defbox}
\begin{definizione}
Il \emph{prodotto diretto} $G_1 \times G_2$ è il prodotto cartesiano con l'operazione $(g_1, g_2) \cdot (h_1, h_2) \coloneqq (g_1 h_1, g_2 h_2)$.
\end{definizione}
\end{defbox}

\begin{exbox}
\begin{esempio}
$\Z/2\Z \times \Z/2\Z$ ha $4$ elementi come $\Z/4\Z$, ma i due gruppi non sono isomorfi: in $\Z/2\Z \times \Z/2\Z$ ogni elemento ha ordine $\leq 2$ (si verifica che $(x,y)+(x,y) = (0,0)$ per ogni $(x,y)$), mentre in $\Z/4\Z$ l'elemento $[1]$ ha ordine $4$.

Al contrario, $\Z/2\Z \times \Z/3\Z \cong \Z/6\Z$ (esercizio!).
\end{esempio}
\end{exbox}

\subsection{Gruppi finiti e ordine di un elemento}

\begin{proposizione}
Sia $G$ un gruppo finito di ordine $n = |G|$. Allora $g^n = 1$ per ogni $g \in G$.
\end{proposizione}

\begin{proof}
Si considera la moltiplicazione per $g$ sull'insieme $G$: è una biiezione, quindi la funzione $h \mapsto gh$ permuta gli elementi di $G$. Il prodotto $\prod_{h \in G} h$ e il prodotto $\prod_{h \in G} (gh) = g^n \prod_{h \in G} h$ sono uguali a meno del fattore $g^n$, da cui $g^n = 1$.
\end{proof}

\begin{corollario}[Piccolo Teorema di Fermat]
Se $p$ è primo e $p \nmid a$, allora $a^{p-1} \equiv 1 \pmod{p}$.
\end{corollario}

\begin{proof}
$(\mathbb{F}_p^\times, \cdot)$ è un gruppo di ordine $p - 1$. Applicando la proposizione precedente, $[a]^{p-1} = [1]$ in $\mathbb{F}_p$.
\end{proof}

\begin{teorema}
Ogni gruppo di ordine primo $p$ è isomorfo a $\Z/p\Z$.
\end{teorema}

% --------------------------------------------------------
\section*{Esercizi}
\addcontentsline{toc}{section}{Esercizi del Capitolo 1}
% --------------------------------------------------------

\begin{esercizio}
Siano $X_1 = \{0,2,4,6,8\}$, $X_2 = \{1,2,4,5,6\}$, $X_3 = \{0,4,8\}$. Determinate $X_i \cup X_j$ e $X_i \cap X_j$ per ogni $1 \leq i < j \leq 3$.
\end{esercizio}

\begin{esercizio}
Dimostrate le \emph{Formule di De Morgan}: per insiemi $X, Y, Z$,
\[
  X \setminus (Y \cup Z) = (X \setminus Y) \cap (X \setminus Z), \quad
  X \setminus (Y \cap Z) = (X \setminus Y) \cup (X \setminus Z).
\]
\end{esercizio}

\begin{esercizio}
Stabilite se le seguenti applicazioni sono iniettive, suriettive o biunivoche. Per quelle biunivoche, trovate l'inversa.
\begin{enumerate}
  \item $\alpha \colon \Q \to \Q$, $\alpha(x) = x^3$.
  \item $\delta \colon \R \to [2, +\infty)$, $\delta(x) = e^x + e^{-x}$.
  \item $\varepsilon \colon \R \to \R$, $\varepsilon(x) = x^3 + x$.
\end{enumerate}
\end{esercizio}

\begin{esercizio}
Dimostrate per induzione che $1^2 + 2^2 + \cdots + n^2 = \dfrac{n(n+1)(2n+1)}{6}$.
\end{esercizio}

\begin{esercizio}
Calcolate i seguenti elementi del campo $\mathbb{F}_5$:
\[
  3 \cdot 4^{-1}, \quad 3^5 \cdot 2^{-2}, \quad (1 + 2 + \cdots + 9) \cdot 3^{-10}.
\]
\end{esercizio}

\begin{esercizio}
Calcolate $(1 - 3i)(5 + 2i)$, $(1-i)^{-1}$, $(1+i)^{10}$ e le radici quadrate di $2i$.
\end{esercizio}

\begin{esercizio}[Piccolo Teorema di Fermat -- approccio diretto]
Sia $p$ primo e $x \in \mathbb{F}_p^\times$. Dimostrate che $x^{p-1} = 1$ in $\mathbb{F}_p$ seguendo i passi:
\begin{enumerate}[label=(\alph*)]
  \item Dimostrate che tra $x, x^2, \ldots, x^{p-1}$ ce n'è uno uguale a $1$ in $\mathbb{F}_p$.
  \item Sia $m$ il minimo con $x^m = 1$: dimostrate che $m \mid (p-1)$.
  \item Concludete.
\end{enumerate}
\end{esercizio}

\begin{esercizio}
Dimostrate che $\Z/2\Z \times \Z/3\Z \cong \Z/6\Z$ come gruppi.
\end{esercizio}

% ============================================================
\chapter{Spazi vettoriali}
% ============================================================

\section*{Introduzione al capitolo}

Gli spazi vettoriali sono la struttura centrale di tutta l'algebra lineare. L'idea di fondo è semplice: vogliamo catturare in un'unica definizione astratta tutto ciò che è in comune a sistemi apparentemente diversi — le $n$-uple di numeri reali $\R^n$, i polinomi $K[x]$, le funzioni continue, i vettori geometrici nel piano. Tutti questi sistemi ammettono una nozione di \emph{somma} e di \emph{moltiplicazione per uno scalare}, e questi due tipi di operazioni si comportano nello stesso modo in tutti i casi.

Il vantaggio dell'astrazione è duplice. Prima di tutto, ogni teorema che dimostriamo per gli spazi vettoriali in generale vale automaticamente per ognuno di questi sistemi. Secondo, l'astrazione ci obbliga a chiederci quali proprietà siano \emph{davvero} necessarie per dimostrare un certo risultato, rendendo più netto il nostro pensiero.

I concetti chiave di questo capitolo sono:
\begin{itemize}
  \item \emph{sottospazi}: le parti di uno spazio vettoriale che sono a loro volta spazi vettoriali;
  \item \emph{combinazioni lineari} e \emph{sottospazio generato} da un insieme di vettori;
  \item \emph{indipendenza lineare}: quando nessun vettore di un insieme è «ridondante»;
  \item \emph{base} e \emph{dimensione}: la base è il mattone costruttivo di ogni spazio vettoriale, e la dimensione è il numero di mattoni necessari.
\end{itemize}

% --------------------------------------------------------
\section{Definizione e primi esempi}
% --------------------------------------------------------

\subsection{L'archetipo: $K^n$}

Il prototipo di spazio vettoriale è $K^n$, l'insieme delle $n$-uple di elementi di un campo $K$. Su $K^n$ sono definite due operazioni naturali:
\begin{align}
  (x_1,\ldots,x_n) + (y_1,\ldots,y_n) &\coloneqq (x_1+y_1,\ldots,x_n+y_n), \label{eq:sumKn}\\
  \lambda \cdot (x_1,\ldots,x_n) &\coloneqq (\lambda x_1,\ldots,\lambda x_n). \label{eq:scalarKn}
\end{align}
Queste operazioni hanno proprietà che conosciamo bene: la somma è commutativa e associativa, esiste un elemento neutro $(0,\ldots,0)$, ogni vettore ha un opposto, e la moltiplicazione per scalare è distributiva. La definizione astratta di spazio vettoriale generalizza esattamente queste proprietà.

\subsection{Definizione assiomatica}

\begin{defbox}
\begin{definizione}[Spazio vettoriale]
\label{def:spaziovett}
Sia $K$ un campo. Uno \emph{spazio vettoriale su $K$} è un insieme $V$ con due operazioni,
\[
  V \times V \to V,\ (v_1,v_2) \mapsto v_1 + v_2 \quad \text{(somma)},
\]
\[
  K \times V \to V,\ (\lambda,v) \mapsto \lambda v \quad \text{(moltiplicazione per scalare)},
\]
tali che:
\begin{enumerate}
  \item esiste $0 \in V$ con $0 + v = v$ per ogni $v \in V$ \emph{(neutro additivo)};
  \item per ogni $v \in V$ esiste $w \in V$ con $v + w = 0$ \emph{(opposto)};
  \item $(u+v)+w = u+(v+w)$ per ogni $u,v,w \in V$ \emph{(associatività della somma)};
  \item $v + w = w + v$ per ogni $v,w \in V$ \emph{(commutatività della somma)};
  \item $1 \cdot v = v$ per ogni $v \in V$;
  \item $(\lambda+\mu)v = \lambda v + \mu v$ per ogni $v \in V$, $\lambda,\mu \in K$;
  \item $\lambda(v+w) = \lambda v + \lambda w$ per ogni $v,w \in V$, $\lambda \in K$;
  \item $(\lambda\mu)v = \lambda(\mu v)$ per ogni $v \in V$, $\lambda,\mu \in K$.
\end{enumerate}
Gli elementi di $V$ si chiamano \emph{vettori}, quelli di $K$ si chiamano \emph{scalari}.
\end{definizione}
\end{defbox}

\begin{osservazione}
Le prime quattro proprietà dicono esattamente che $(V,+)$ è un gruppo abeliano. Le restanti quattro riguardano l'interazione tra la moltiplicazione per scalare e le strutture di campo di $K$ e di gruppo di $V$.
\end{osservazione}

\begin{osservazione}
Si usa il simbolo $0$ sia per l'elemento neutro del campo $K$ che per il vettore nullo di $V$. Il contesto chiarisce sempre di quale dei due si tratta.
\end{osservazione}

\subsection{Esempi fondamentali}

\begin{exbox}
\begin{esempio}
\label{ex:esempiVett}
\begin{enumerate}
  \item \textbf{$K^n$}: come sopra, con le operazioni~\eqref{eq:sumKn}--\eqref{eq:scalarKn}. È il caso più importante.
  \item \textbf{Vettori geometrici nel piano $\mathcal{V}(E^2)$}: due segmenti orientati $\overrightarrow{AB}$ e $\overrightarrow{CD}$ si dicono \emph{equipollenti} se $AB \parallel CD$ e $AC \parallel BD$. L'insieme delle classi di equipollenza, con la somma definita dalla regola del parallelogramma e la moltiplicazione per scalare definita scalando la lunghezza, è uno spazio vettoriale reale di dimensione $2$.
  \item \textbf{Polinomi $K[x]$}: con la somma e il prodotto per scalare usuali, $K[x]$ è uno spazio vettoriale su $K$ (di dimensione infinita).
  \item \textbf{Funzioni $K^X$}: l'insieme di tutte le funzioni $f \colon X \to K$, con le operazioni $(f+g)(x) \coloneqq f(x)+g(x)$ e $(\lambda f)(x) \coloneqq \lambda f(x)$, è uno spazio vettoriale su $K$.
  \item \textbf{$\C$ come spazio reale}: poiché $\R \subset \C$, i numeri complessi formano uno spazio vettoriale su $\R$ (di dimensione $2$, con base $\{1, i\}$).
\end{enumerate}
\end{esempio}
\end{exbox}

\begin{osservazione}[Su cosa «dimentica» la definizione]
Nella definizione di spazio vettoriale non compaiono né la norma di un vettore né il prodotto scalare tra vettori. Queste strutture aggiuntive (studiate nei Capitoli~6 e~7) sono importanti, ma non sono necessarie per i risultati fondamentali sull'algebra lineare. Concentrarsi sulla struttura minima ci permette di dimostrarli una volta sola, in piena generalità.
\end{osservazione}

\subsection{Prime proprietà}

\begin{proposizione}
Sia $V$ uno spazio vettoriale su $K$. Allora:
\begin{enumerate}
  \item il vettore neutro $0$ è unico;
  \item per ogni $v \in V$, l'opposto è unico; lo denotiamo $-v$;
  \item $0_K \cdot v = 0_V$ per ogni $v \in V$;
  \item $\lambda \cdot 0_V = 0_V$ per ogni $\lambda \in K$;
  \item $(-1) \cdot v = -v$ per ogni $v \in V$.
\end{enumerate}
\end{proposizione}

\begin{proof}
(1) Se $0_1$ e $0_2$ sono due neutri, allora $0_1 = 0_1 + 0_2 = 0_2$.

(2) Se $v + w_1 = 0$ e $v + w_2 = 0$, allora $w_1 = w_1 + 0 = w_1 + (v + w_2) = (w_1 + v) + w_2 = 0 + w_2 = w_2$.

(3) $0_K \cdot v = (0_K + 0_K) \cdot v = 0_K \cdot v + 0_K \cdot v$; aggiungendo $-(0_K \cdot v)$ si ottiene $0_V = 0_K \cdot v$.

(5) $(-1)v + v = (-1)v + 1 \cdot v = (-1+1)v = 0_K \cdot v = 0_V$, quindi $(-1)v = -v$.
\end{proof}

% --------------------------------------------------------
\section{Sottospazi}
% --------------------------------------------------------

\subsection{Definizione e criteri}

Un \emph{sottospazio} è una parte di uno spazio vettoriale che è a sua volta uno spazio vettoriale rispetto alle stesse operazioni. Il seguente criterio è pratico e conciso.

\begin{defbox}
\begin{definizione}[Sottospazio]
Sia $V$ uno spazio vettoriale su $K$. Un sottoinsieme $W \subseteq V$ è un \emph{sottospazio vettoriale} di $V$ se:
\begin{enumerate}
  \item[(a)] $W \neq \emptyset$;
  \item[(b)] per ogni $v_1, v_2 \in W$, si ha $(v_1 + v_2) \in W$;
  \item[(c)] per ogni $v \in W$ e $\lambda \in K$, si ha $\lambda v \in W$.
\end{enumerate}
\end{definizione}
\end{defbox}

\begin{osservazione}
Le condizioni (b) e (c) si possono condensare in una sola: $W$ è un sottospazio se e solo se è non vuoto e $\lambda v_1 + \mu v_2 \in W$ per ogni $v_1, v_2 \in W$ e $\lambda, \mu \in K$. Un sottospazio contiene sempre $0$: basta prendere $\lambda = 0$ nella condizione (c).
\end{osservazione}

\begin{exbox}
\begin{esempio}
\label{ex:sottospazi}
\begin{enumerate}
  \item \textbf{Soluzioni di un'equazione omogenea}: il sottoinsieme $U = \{(x_1,\ldots,x_n) \in K^n \mid a_1 x_1 + \cdots + a_n x_n = 0\}$ è un sottospazio di $K^n$. Al contrario, $W = \{(x_1,\ldots,x_n) \mid a_1 x_1 + \cdots + a_n x_n = 1\}$ \emph{non} è un sottospazio: non contiene il vettore nullo.

  La differenza tra $U$ e $W$ rispecchia la distinzione fondamentale tra sistemi lineari \emph{omogenei} (secondo membro $= 0$) e \emph{non omogenei}: le soluzioni di un sistema omogeneo formano un sottospazio, quelle di un sistema non omogeneo no.

  \item \textbf{Polinomi di grado $\leq d$}: $K[x]_{\leq d} \coloneqq \{a_0 + a_1 x + \cdots + a_d x^d \mid a_i \in K\}$ è un sottospazio di $K[x]$.

  \item \textbf{Sottospazi di $\mathcal{V}(E^2)$}: I soli sottospazi del piano vettoriale sono $\{0\}$, le rette per l'origine, e l'intero piano. In particolare, una retta che non passa per l'origine \emph{non} è un sottospazio.
\end{enumerate}
\end{esempio}
\end{exbox}

\begin{proposizione}
\label{prop:intersez}
L'intersezione $\bigcap_{i \in I} W_i$ di una qualunque famiglia di sottospazi di $V$ è ancora un sottospazio di $V$.
\end{proposizione}

\begin{proof}
Il vettore $0$ appartiene a ogni $W_i$, quindi all'intersezione. Se $v_1, v_2 \in \bigcap W_i$ e $\lambda \in K$, allora $v_1, v_2 \in W_i$ per ogni $i$, e poiché ciascun $W_i$ è un sottospazio, $(v_1+v_2)$ e $\lambda v_1$ appartengono a ogni $W_i$, dunque all'intersezione.
\end{proof}

\begin{exbox}
\begin{esempio}[Sistemi lineari omogenei]
Le soluzioni di un sistema lineare omogeneo
\[
  \sum_{j=1}^n a_{ij} x_j = 0, \quad i = 1, \ldots, m
\]
sono l'intersezione degli $m$ sottospazi (uno per ogni equazione) di $K^n$. Per la Proposizione~\ref{prop:intersez}, l'insieme delle soluzioni è automaticamente un sottospazio di $K^n$.
\end{esempio}
\end{exbox}

% --------------------------------------------------------
\section{Combinazioni lineari e sottospazi generati}
% --------------------------------------------------------

\begin{defbox}
\begin{definizione}[Combinazione lineare]
Sia $V$ uno spazio vettoriale su $K$ e $v_1, \ldots, v_n \in V$. Un vettore $v \in V$ è una \emph{combinazione lineare} di $v_1, \ldots, v_n$ se esistono $\lambda_1, \ldots, \lambda_n \in K$ tali che
\[
  v = \lambda_1 v_1 + \lambda_2 v_2 + \cdots + \lambda_n v_n.
\]
Per convenzione, solo il vettore $0$ è combinazione lineare della lista vuota.
\end{definizione}
\end{defbox}

\begin{exbox}
\begin{esempio}
I vettori standard di $K^n$,
\[
  e_1 = (1,0,\ldots,0),\quad e_2 = (0,1,0,\ldots,0),\quad \ldots,\quad e_n = (0,\ldots,0,1),
\]
generano $K^n$: ogni $(x_1,\ldots,x_n) = x_1 e_1 + \cdots + x_n e_n$ è combinazione lineare di $e_1,\ldots,e_n$.
\end{esempio}
\end{exbox}

\begin{defbox}
\begin{definizione}[Sottospazio generato]
Sia $S \subseteq V$. Il \emph{sottospazio generato da $S$}, denotato $\langle S \rangle$ o $\mathrm{Span}(S)$, è l'insieme di tutte le combinazioni lineari finite di vettori di $S$:
\[
  \langle S \rangle \coloneqq \{\lambda_1 v_1 + \cdots + \lambda_n v_n \mid n \in \N,\, v_i \in S,\, \lambda_i \in K\}.
\]
\end{definizione}
\end{defbox}

\begin{proposizione}
$\langle S \rangle$ è il più piccolo sottospazio di $V$ che contiene $S$: è un sottospazio, contiene $S$, ed è contenuto in ogni sottospazio che contiene $S$.
\end{proposizione}

\begin{defbox}
\begin{definizione}[Somma di sottospazi]
Se $U, W \subseteq V$ sono sottospazi, la loro \emph{somma} è
\[
  U + W \coloneqq \langle U \cup W \rangle = \{u + w \mid u \in U,\, w \in W\}.
\]
\end{definizione}
\end{defbox}

\begin{osservazione}
L'\emph{unione} $U \cup W$ in generale non è un sottospazio. Per esempio in $\R^2$, $U = \langle e_1 \rangle$ e $W = \langle e_2 \rangle$: $e_1 \in U$, $e_2 \in W$, ma $e_1 + e_2 \notin U \cup W$. La somma $U + W$, invece, è l'intero piano $\R^2$.
\end{osservazione}

\begin{exbox}
\begin{esempio}[Spazi finitamente generati e non]
\begin{itemize}
  \item $K^n$ è finitamente generato: lo generano $e_1, \ldots, e_n$.
  \item $K[x]$ \emph{non} è finitamente generato: se $f_1, \ldots, f_m$ fossero generatori, ogni polinomio in $\langle f_1, \ldots, f_m \rangle$ avrebbe grado $\leq \max_j \deg f_j$, ma in $K[x]$ ci sono polinomi di grado arbitrario.
  \item $K[x]_{\leq d}$ è finitamente generato da $\{1, x, x^2, \ldots, x^d\}$.
\end{itemize}
\end{esempio}
\end{exbox}

% --------------------------------------------------------
\section{Dipendenza e indipendenza lineare}
% --------------------------------------------------------

L'indipendenza lineare cattura l'idea di «non ridondanza»: un insieme di vettori è linearmente indipendente se nessuno di essi è esprimibile come combinazione lineare degli altri.

\begin{defbox}
\begin{definizione}[Indipendenza lineare]
Una lista di vettori $v_1, \ldots, v_n \in V$ è \emph{linearmente indipendente} se l'unica relazione lineare
\[
  \lambda_1 v_1 + \lambda_2 v_2 + \cdots + \lambda_n v_n = 0
\]
è quella con $\lambda_1 = \lambda_2 = \cdots = \lambda_n = 0$. È \emph{linearmente dipendente} se esiste una relazione lineare con almeno un $\lambda_i \neq 0$.
\end{definizione}
\end{defbox}

\begin{osservazione}
Attenzione alla logica: «linearmente indipendente» è una proprietà della \emph{lista} $(v_1, \ldots, v_n)$, non di ciascun vettore singolarmente. In particolare:
\begin{itemize}
  \item se $v_i = 0$ per qualche $i$, la lista è dipendente (coefficiente $1$ davanti a $v_i$, zero altrove);
  \item se $v_i = v_j$ per $i \neq j$, la lista è dipendente ($v_i - v_j = 0$).
\end{itemize}
\end{osservazione}

\begin{exbox}
\begin{esempio}
\begin{itemize}
  \item $e_1, \ldots, e_n \in K^n$ sono linearmente indipendenti: se $\sum \lambda_i e_i = 0$, confrontando le componenti si trova $\lambda_i = 0$ per ogni $i$.
  \item $(2,2), (3,3) \in \R^2$ sono linearmente dipendenti: $3(2,2) - 2(3,3) = (0,0)$.
  \item $1, x, x^2, \ldots, x^n, \ldots \in K[x]$ sono linearmente indipendenti: per definizione di uguaglianza tra polinomi.
  \item Due vettori $v_1 = (a,b)$, $v_2 = (c,d) \in K^2$ sono linearmente dipendenti se e solo se $ad - bc = 0$ (il determinante della matrice $2 \times 2$ con righe $v_1$ e $v_2$).
\end{itemize}
\end{esempio}
\end{exbox}

\begin{proposizione}
I vettori $v_1, \ldots, v_n$ sono linearmente dipendenti se e solo se esiste un indice $i$ tale che $v_i \in \langle v_1, \ldots, v_{i-1}, v_{i+1}, \ldots, v_n \rangle$.
\end{proposizione}

\begin{proof}
Se $\sum_{j=1}^n \lambda_j v_j = 0$ con $\lambda_i \neq 0$, si isola $v_i = -\lambda_i^{-1} \sum_{j \neq i} \lambda_j v_j$. Viceversa, se $v_i = \sum_{j \neq i} a_j v_j$, allora $\sum_{j \neq i} a_j v_j - v_i = 0$ è una relazione non banale.
\end{proof}

% --------------------------------------------------------
\section{Basi e dimensione}
% --------------------------------------------------------

\subsection{Definizione di base}

\begin{defbox}
\begin{definizione}[Base]
Una lista $\{v_i\}_{i \in I}$ di vettori di $V$ è una \emph{base} di $V$ se:
\begin{enumerate}
  \item i vettori della lista sono linearmente indipendenti;
  \item $V$ è generato dai vettori della lista.
\end{enumerate}
\end{definizione}
\end{defbox}

Equivalentemente (per la Proposizione 2.6.7 degli originali), $\{v_1,\ldots,v_n\}$ è una base se e solo se l'applicazione $K^n \to V$, $(x_1,\ldots,x_n) \mapsto \sum x_i v_i$, è una biiezione.

\begin{exbox}
\begin{esempio}
\begin{itemize}
  \item \textbf{Base standard di $K^n$}: $\{e_1, \ldots, e_n\}$ (ogni vettore ha una e una sola espressione come combinazione di $e_i$).
  \item \textbf{Base di $K[x]$}: la lista $\{1, x, x^2, \ldots, x^n, \ldots\}$ è una base (infinita).
  \item \textbf{Base non standard di $\R^2$}: $B = \{(1,1), (1,-1)\}$. Le coordinate di $(x_1, x_2)$ nella base $B$ si trovano risolvendo $t_1(1,1) + t_2(1,-1) = (x_1,x_2)$, che dà $t_1 = (x_1+x_2)/2$ e $t_2 = (x_1-x_2)/2$.
\end{itemize}
\end{esempio}
\end{exbox}

\begin{osservazione}[Coordinate]
Fissata una base ordinata $B = (v_1, \ldots, v_n)$ di $V$, ogni vettore $v \in V$ si scrive in modo \emph{unico} come $v = x_1 v_1 + \cdots + x_n v_n$. La $n$-pla $(x_1, \ldots, x_n) \in K^n$ è il \emph{vettore delle coordinate} di $v$ nella base $B$, denotato $X_B(v)$. Cambiare base cambia le coordinate: lo stesso vettore geometrico ha coordinate diverse rispetto a sistemi di riferimento diversi.
\end{osservazione}

\subsection{Esistenza di basi e loro cardinalità}

\begin{proposizione}
Sia $V$ finitamente generato da $v_1, \ldots, v_n$. Allora esiste una base di $V$ ottenuta eliminando alcuni dei $v_i$.
\end{proposizione}

\begin{proof}
Per induzione su $n$. Se la lista è vuota, $V = \{0\}$ e la lista vuota è una base. Nel passo induttivo: se $v_1, \ldots, v_n$ sono già indipendenti, formano una base. Altrimenti, per la proposizione precedente, qualche $v_i$ è combinazione degli altri; rimuovendolo non si perde il sottoinsieme generato, e si itera.
\end{proof}

\begin{corollario}
Ogni spazio vettoriale finitamente generato ha una base.
\end{corollario}

Il seguente teorema è il cuore di tutta la teoria.

\begin{teorema}[Steinitz]
Sia $V$ uno spazio vettoriale su $K$. Se $r_1, \ldots, r_p$ sono linearmente indipendenti e $z_1, \ldots, z_q$ generano $V$, allora $p \leq q$.
\end{teorema}

\begin{proof}[Idea della dimostrazione]
Per induzione su $p$. Il caso $p=0$ è banale. Per il passo induttivo: per ipotesi induttiva si può sostituire $p$ dei $z_j$ con $r_1, \ldots, r_p$, ottenendo ancora un insieme generante. Aggiungendo $r_{p+1}$ e usando il Lemma di scambio (uno dei $z$ rimanenti può essere rimpiazzato da $r_{p+1}$), si conclude.
\end{proof}

\begin{corollario}[Invarianza della dimensione]
In uno spazio vettoriale finitamente generato, tutte le basi hanno la stessa cardinalità.
\end{corollario}

\begin{proof}
Siano $\{r_1,\ldots,r_p\}$ e $\{z_1,\ldots,z_q\}$ due basi. La prima è indipendente e la seconda genera: $p \leq q$. Viceversa $q \leq p$. Quindi $p = q$.
\end{proof}

\begin{defbox}
\begin{definizione}[Dimensione]
La \emph{dimensione} di uno spazio vettoriale finitamente generato $V$ è la cardinalità di una qualsiasi sua base. Si denota $\dim V$ o $\dim_K V$.
\end{definizione}
\end{defbox}

\begin{exbox}
\begin{esempio}
\begin{itemize}
  \item $\dim K^n = n$.
  \item $\dim_\R \C^n = 2n$ (base reale: $\{e_1, ie_1, e_2, ie_2, \ldots, e_n, ie_n\}$).
  \item $\dim K[x]_{\leq d} = d+1$ (base: $\{1, x, \ldots, x^d\}$).
  \item $\dim \mathcal{V}(E^2) = 2$ (qualsiasi coppia di vettori non paralleli è una base).
\end{itemize}
\end{esempio}
\end{exbox}

\begin{proposizione}
Sia $\dim V = n$ e $v_1, \ldots, v_m \in V$.
\begin{enumerate}
  \item Se $v_1,\ldots,v_m$ sono linearmente indipendenti, allora $m \leq n$; se $m = n$, formano una base.
  \item Se $\langle v_1,\ldots,v_m\rangle = V$, allora $m \geq n$; se $m = n$, formano una base.
\end{enumerate}
\end{proposizione}

\begin{osservazione}[Applicazione: sistemi con più incognite che equazioni]
Se un sistema lineare omogeneo $m \times n$ ha più incognite che equazioni ($n > m$), allora ha sempre una soluzione non banale. Infatti i vettori colonna del sistema vivono in $K^m$, che ha dimensione $m < n$: per la proposizione, $n$ vettori in uno spazio di dimensione $m < n$ non possono essere indipendenti, quindi esiste una relazione lineare non banale, che corrisponde a una soluzione non nulla.
\end{osservazione}

% --------------------------------------------------------
\section{Formula di Grassmann}
% --------------------------------------------------------

La \emph{formula di Grassmann} (o \emph{formula di inclusione-esclusione per le dimensioni}) è l'analogo per gli spazi vettoriali della formula $|A \cup B| = |A| + |B| - |A \cap B|$ per gli insiemi finiti.

\begin{proposizione}[Formula di Grassmann]
Sia $V$ uno spazio vettoriale su $K$ e $U, W \subseteq V$ sottospazi finitamente generati. Allora
\[
  \dim(U + W) + \dim(U \cap W) = \dim U + \dim W.
\]
\end{proposizione}

\begin{proof}
Sia $\{z_1,\ldots,z_a\}$ una base di $U \cap W$. La si estende a una base $\{z_1,\ldots,z_a, u_1,\ldots,u_m\}$ di $U$ e a una base $\{z_1,\ldots,z_a, w_1,\ldots,w_n\}$ di $W$. Si dimostra che
\[
  \mathcal{B} = \{z_1,\ldots,z_a,\, u_1,\ldots,u_m,\, w_1,\ldots,w_n\}
\]
è una base di $U + W$. La generazione è ovvia. Per l'indipendenza: se $\sum \lambda_i z_i + \sum \mu_j u_j + \sum \theta_k w_k = 0$, il termine $\sum \mu_j u_j$ appartiene sia a $U$ che a $W$ (per manipolazione algebrica), quindi a $U \cap W$, da cui è combinazione dei $z_i$; l'indipendenza di $\{z_i, w_k\}$ forza $\theta_k = 0$; poi l'indipendenza di $\{z_i, u_j\}$ forza $\lambda_i = \mu_j = 0$. Quindi
$\dim(U+W) = a+m+n = (a+m)+(a+n)-a = \dim U + \dim W - \dim(U \cap W)$.
\end{proof}

\begin{corollario}
Se $\dim U + \dim W > \dim V$, allora $U \cap W \neq \{0\}$.
\end{corollario}

\begin{exbox}
\begin{esempio}
In $\R^3$, siano $U$ e $W$ due piani (sottospazi di dimensione $2$) passanti per l'origine. Allora $\dim U + \dim W = 4 > 3 = \dim \R^3$, quindi $\dim(U \cap W) \geq 1$: i due piani si intersecano lungo almeno una retta.
\end{esempio}
\end{exbox}

% --------------------------------------------------------
\section{Costruzioni di spazi vettoriali}
% --------------------------------------------------------

\subsection{Somma diretta}

Dati due spazi vettoriali $V$ e $W$ su $K$, il \emph{prodotto cartesiano} $V \times W$ diventa uno spazio vettoriale con le operazioni componente per componente:
\[
  (v_1, w_1) + (v_2, w_2) \coloneqq (v_1+v_2, w_1+w_2), \quad \lambda(v,w) \coloneqq (\lambda v, \lambda w).
\]
Questa struttura si chiama \emph{somma diretta} e si denota $V \oplus W$.

\begin{proposizione}
$\dim(V \oplus W) = \dim V + \dim W$.
\end{proposizione}

\begin{proof}
Se $\{v_1,\ldots,v_m\}$ è una base di $V$ e $\{w_1,\ldots,w_n\}$ di $W$, allora
$\{(v_1,0),\ldots,(v_m,0),(0,w_1),\ldots,(0,w_n)\}$
è una base di $V \oplus W$ (si verifica generazione e indipendenza).
\end{proof}

\subsection{Spazio quoziente}

Sia $W \subseteq V$ un sottospazio. Definiamo la relazione di congruenza $v_1 \equiv v_2 \pmod{W}$ se $(v_1 - v_2) \in W$. Questa è una relazione di equivalenza.

\begin{defbox}
\begin{definizione}[Spazio quoziente]
Il \emph{quoziente} $V/W$ è l'insieme delle classi di equivalenza modulo $W$, con operazioni
\[
  [v] + [u] \coloneqq [v+u], \quad \lambda[v] \coloneqq [\lambda v].
\]
Queste operazioni sono ben definite (non dipendono dal rappresentante scelto).
\end{definizione}
\end{defbox}

\begin{exbox}
\begin{esempio}
Se $V = K^n$ e $W = \langle e_{j+1}, \ldots, e_n \rangle$, allora $X \equiv Y \pmod{W}$ se e solo se le prime $j$ coordinate coincidono. Il quoziente $K^n/W$ si identifica naturalmente con $K^j$.
\end{esempio}
\end{exbox}

\begin{proposizione}
$\dim(V/W) = \dim V - \dim W$.
\end{proposizione}

\begin{proof}
Si estende una base $\{w_1,\ldots,w_a\}$ di $W$ a una base $\{w_1,\ldots,w_a,u_1,\ldots,u_b\}$ di $V$. Si verifica che $\{[u_1],\ldots,[u_b]\}$ è una base di $V/W$.
\end{proof}

% --------------------------------------------------------
\section*{Esercizi}
\addcontentsline{toc}{section}{Esercizi del Capitolo 2}
% --------------------------------------------------------

\begin{esercizio}
Siano $X = (1,2,-3)$, $Y=(3,-5,2)$, $Z=(1,1,-2) \in \R^3$. Calcolate $2X - Y + Z$. Trovate $\lambda, \mu, \nu$ non tutti nulli con $\lambda X + \mu Y + \nu Z = 0$.
\end{esercizio}

\begin{esercizio}
Determinare quali dei seguenti sottoinsiemi sono sottospazi:
\begin{enumerate}
  \item $W_1 = \{(x,y,z) \in \R^3 \mid x+y+z=0\}$.
  \item $W_2 = \{(x,y,z) \in \R^3 \mid x+y+z \leq 1\}$.
  \item $W_3 = \{(x,y,z) \in \C^3 \mid x^2+y^2+z^2=0\}$.
  \item $W_4 = \{(x,y,z) \in \mathbb{F}_2^3 \mid x^2+y^2+z^2=0\}$.
\end{enumerate}
\end{esercizio}

\begin{esercizio}
Dimostrate che se $W_1 \cup W_2$ è un sottospazio di $V$, allora $W_1 \subseteq W_2$ oppure $W_2 \subseteq W_1$.
\end{esercizio}

\begin{esercizio}
Nello spazio $\R^4$, siano $v_1 = (1,1,0,-1)$, $v_2=(1,-2,3,2)$, $v_3=(1,-1,0,0)$, $v_4=(0,1,0,1)$. Stabilite quali tra $\{v_1,v_2,v_3\}$, $\{v_2,v_3,v_4\}$, $\{v_3,v_4,v_1\}$ e $\{v_4,v_1,v_2\}$ sono linearmente dipendenti.
\end{esercizio}

\begin{esercizio}
Sia $\alpha \in K$. Dimostrate che $B_\alpha = \{1, (x+\alpha), (x+\alpha)^2, \ldots, (x+\alpha)^d\}$ è una base di $K[x]_{\leq d}$. Calcolate le coordinate di $x^d$ nella base $B_\alpha$.
\end{esercizio}

\begin{esercizio}
Sia $W \subseteq K[x]_{\leq d}$ il sottospazio $W = \{p \in K[x]_{\leq d} \mid p(0)=p(1)=p(-1)=0\}$. Calcolate $\dim W$ (trattate separatamente il caso $\mathrm{char}\, K = 2$).
\end{esercizio}

\begin{esercizio}[Formula di Grassmann applicata]
In $V = \R^5$, siano $U$ e $W$ sottospazi di dimensione rispettivamente $3$ e $4$. Quali valori può assumere $\dim(U \cap W)$? E $\dim(U + W)$?
\end{esercizio}

\begin{esercizio}
Sia $V$ finitamente generato e $W_1, \ldots, W_p \subseteq V$ sottospazi. Dimostrate per induzione (usando Grassmann) che
\[
  \mathrm{cod}(W_1 \cap \cdots \cap W_p, V) \leq \sum_{i=1}^p \mathrm{cod}(W_i, V).
\]
\end{esercizio}

\begin{esercizio}
Sia $V$ uno spazio vettoriale su $\Q$ finitamente generato di dimensione $n$, con generatori $w_1, \ldots, w_m$. Dimostrate che $m \geq n$ senza usare il Teorema di Steinitz, ma riducendosi al caso dei campi finiti $\mathbb{F}_p$ per $p$ primo grande.
\end{esercizio}

% ============================================================
\chapter{Applicazioni lineari}
\label{cap:applin}
% ============================================================

\section*{Introduzione al capitolo}

Il Capitolo~2 ha costruito la scena: gli spazi vettoriali. Questo capitolo introduce i protagonisti dell'azione: le \emph{applicazioni lineari}, ovvero le funzioni tra spazi vettoriali che rispettano la struttura lineare.

Perché limitarsi alle applicazioni lineari? La risposta è duplice. Da un lato, sono le applicazioni \emph{naturali} per la struttura degli spazi vettoriali, nel senso che preservano somme e prodotti per scalari. Dall'altro, esse sono straordinariamente ricche: come vedremo, ogni applicazione lineare tra $K^n$ e $K^m$ è completamente descritta da una matrice $m \times n$, e le operazioni tra applicazioni lineari (composizione, somma) corrispondono ad operazioni sulle matrici.

Il filo conduttore del capitolo è il seguente: partire da una definizione astratta, svilupparne le proprietà strutturali (nucleo, immagine, teorema della dimensione), poi costruire il ponte con le matrici, e infine studiare questioni più fini come la diagonalizzazione e il duale.

Il \emph{teorema fondamentale delle applicazioni lineari} (Teorema~\ref{thm:iso}) afferma che ogni applicazione lineare induce un isomorfismo tra il quoziente dello spazio di partenza per il nucleo e l'immagine. Questo risultato, di per sé elegante, è anche uno strumento potente: permette di calcolare dimensioni, classificare applicazioni e costruire nuovi spazi.

% --------------------------------------------------------
\section{Definizione e prime proprietà}
\label{sec:defapplin}
% --------------------------------------------------------

\subsection{La definizione}

\begin{defbox}
\begin{definizione}[Applicazione lineare]
\label{def:applin}
Siano $V, W$ spazi vettoriali sullo stesso campo $K$. Un'applicazione $f \colon V \to W$ è \emph{lineare} se per ogni $v, w \in V$ e $\lambda \in K$ vale:
\[
  f(v + w) = f(v) + f(w), \qquad f(\lambda v) = \lambda f(v).
\]
\end{definizione}
\end{defbox}

Le due condizioni possono essere riunite in una sola: $f$ è lineare se e solo se per ogni $\lambda, \mu \in K$ e $v, w \in V$ vale
\[
  f(\lambda v + \mu w) = \lambda f(v) + \mu f(w).
\]

\begin{exbox}
\begin{esempio}[Forme lineari su $K^n$]
Fissati $a_1, \ldots, a_n \in K$, l'applicazione
\[
  f \colon K^n \to K, \quad (x_1, \ldots, x_n) \mapsto a_1 x_1 + \cdots + a_n x_n
\]
è lineare. Si verifica immediatamente: $f(X+Y) = \sum a_i(x_i+y_i) = f(X)+f(Y)$ e $f(\lambda X) = \lambda f(X)$.

Vedremo in seguito che \emph{tutte} le applicazioni lineari $K^n \to K$ hanno questa forma.
\end{esempio}
\end{exbox}

\begin{exbox}
\begin{esempio}[Coordinate rispetto a una base]
Sia $V$ finitamente generato con base $\mathcal{B} = \{v_1, \ldots, v_n\}$. L'applicazione
\[
  X_\mathcal{B} \colon V \to K^n, \quad v \mapsto (x_1, \ldots, x_n) \text{ dove } v = \sum_{i=1}^n x_i v_i
\]
è lineare. Questo segue dal fatto che le coordinate dell'elemento $v+w$ sono la somma delle coordinate di $v$ e di $w$.
\end{esempio}
\end{exbox}

\begin{exbox}
\begin{esempio}[Rotazione del piano]
Sia $\rho \colon E^2 \to E^2$ la rotazione intorno a un punto fissato di un angolo $\theta$. Essa manda rette parallele in rette parallele, quindi induce un'applicazione lineare $V(\rho) \colon \mathcal{V}(E^2) \to \mathcal{V}(E^2)$ sullo spazio dei vettori del piano. La matrice di $V(\rho)$ rispetto a una base ortonormale $\{i, j\}$ è
\[
  \begin{pmatrix} \cos\theta & -\sin\theta \\ \sin\theta & \cos\theta \end{pmatrix}.
\]
\end{esempio}
\end{exbox}

\begin{exbox}
\begin{esempio}[Applicazione quoziente]
Se $U \subseteq V$ è un sottospazio, l'applicazione quoziente $\pi \colon V \to V/U$, $v \mapsto [v]$, è lineare.
\end{esempio}
\end{exbox}

\subsection{Proprietà fondamentali}

\begin{lemma}
\label{lem:lin0}
Sia $f \colon V \to W$ lineare. Allora:
\begin{enumerate}
  \item $f(0) = 0$.
  \item $f(-v) = -f(v)$ per ogni $v \in V$.
  \item $f(\lambda_1 v_1 + \cdots + \lambda_n v_n) = \lambda_1 f(v_1) + \cdots + \lambda_n f(v_n)$ per ogni $v_i \in V$, $\lambda_i \in K$.
\end{enumerate}
\end{lemma}

\begin{proof}
(1) Ponendo $\lambda = 0$ nella condizione $f(\lambda v) = \lambda f(v)$ si ottiene $f(0) = 0$. (2) Si pone $\lambda = -1$. (3) Per induzione, usando entrambe le condizioni della Definizione~\ref{def:applin}.
\end{proof}

Il punto~(3) del Lemma~\ref{lem:lin0} è il fatto chiave: un'applicazione lineare è completamente determinata dai suoi valori su una base.

\begin{proposizione}[Unicità e esistenza]
\label{prop:uniq-lin}
Siano $V, W$ spazi vettoriali su $K$, con $V$ finitamente generato.
\begin{enumerate}
  \item \emph{(Unicità)} Se $\{v_i\}_{i \in I}$ genera $V$ e $f, g \colon V \to W$ sono lineari con $f(v_i) = g(v_i)$ per ogni $i$, allora $f = g$.
  \item \emph{(Esistenza)} Se $\{v_1, \ldots, v_m\}$ sono linearmente indipendenti in $V$ e $w_1, \ldots, w_m \in W$ sono vettori qualsiasi, esiste un'applicazione lineare $f \colon V \to W$ tale che $f(v_j) = w_j$ per ogni $j$.
\end{enumerate}
\end{proposizione}

\begin{proof}
(1) Ogni $v \in V$ si scrive $v = \sum_j x_j v_{i_j}$; per linearità $f(v) = \sum x_j f(v_{i_j}) = \sum x_j g(v_{i_j}) = g(v)$.

(2) Si estende $\{v_1, \ldots, v_m\}$ a una base $\{v_1, \ldots, v_n\}$ di $V$ (possibile perché $V$ è finitamente generato) e si scelgono arbitrariamente $w_{m+1}, \ldots, w_n \in W$. Si definisce
\[
  f\!\left(\sum_{j=1}^n x_j v_j\right) \coloneqq \sum_{j=1}^n x_j w_j.
\]
Questa è ben definita (l'espressione nella base è unica), lineare (si verifica direttamente) e soddisfa $f(v_j) = w_j$.
\end{proof}

\begin{corollario}
\label{cor:lin-base}
Con le notazioni della Proposizione~\ref{prop:uniq-lin}, se $\mathcal{B} = \{v_1, \ldots, v_n\}$ è una base di $V$, allora per ogni scelta di $w_1, \ldots, w_n \in W$ esiste una e una sola applicazione lineare $f \colon V \to W$ con $f(v_j) = w_j$.
\end{corollario}

In altre parole: \emph{un'applicazione lineare è determinata (e può essere prescritta liberamente) sui vettori di una base}.

% --------------------------------------------------------
\section{Matrici e applicazioni lineari}
\label{sec:matrici}
% --------------------------------------------------------

\subsection{Calcolo matriciale}

\begin{defbox}
\begin{definizione}[Matrice]
Una \emph{matrice} $m \times n$ a valori in un anello $R$ è un'applicazione $f \colon \{1,\ldots,m\} \times \{1,\ldots,n\} \to R$. Si scrive $A = (a_{ij})$, dove $a_{ij} = f(i,j)$. L'insieme delle matrici $m \times n$ a valori in $R$ si denota $M_{m,n}(R)$.
\end{definizione}
\end{defbox}

Su $M_{m,n}(R)$ sono definite la \emph{somma} $(A+B)_{ij} = a_{ij}+b_{ij}$ e il \emph{prodotto per scalare} $(\lambda A)_{ij} = \lambda a_{ij}$. Se $K$ è un campo, con queste operazioni $M_{m,n}(K)$ è uno spazio vettoriale su $K$ isomorfo a $K^{mn}$.

\begin{defbox}
\begin{definizione}[Prodotto righe per colonne]
Date $A \in M_{m,n}(R)$ e $B \in M_{n,p}(R)$, il \emph{prodotto} $A \cdot B \in M_{m,p}(R)$ è definito da
\[
  (A \cdot B)_{ih} \coloneqq \sum_{j=1}^n a_{ij} b_{jh}.
\]
\end{definizione}
\end{defbox}

\begin{osservazione}
Il prodotto di matrici \emph{non} è commutativo in generale. Per esempio:
\[
  \begin{pmatrix} \lambda & 0 \\ 0 & \mu \end{pmatrix} \cdot \begin{pmatrix} 0 & 1 \\ 0 & 0 \end{pmatrix} = \begin{pmatrix} 0 & \lambda \\ 0 & 0 \end{pmatrix}, \quad
  \begin{pmatrix} 0 & 1 \\ 0 & 0 \end{pmatrix} \cdot \begin{pmatrix} \lambda & 0 \\ 0 & \mu \end{pmatrix} = \begin{pmatrix} 0 & \mu \\ 0 & 0 \end{pmatrix}.
\]
Se $\lambda \neq \mu$ i due prodotti sono diversi. Inoltre il quadrato della seconda matrice è $0$, anche se la matrice stessa non è nulla.
\end{osservazione}

Il prodotto di matrici gode tuttavia dell'associatività e della distributività (Proposizione~3.2.7 degli appunti originali), e $M_{n,n}(R)$ con somma e prodotto è un \emph{anello con unità} (non commutativo per $n > 1$).

\subsection{Dalla matrice all'applicazione lineare}

Identifichiamo $K^n$ con $M_{n,1}(K)$ (vettori colonna). Data $A \in M_{m,n}(K)$, l'applicazione
\[
  L_A \colon K^n \to K^m, \quad X \mapsto A \cdot X
\]
è lineare. Esplicitamente, la riga $i$-esima di $L_A(X)$ è $\sum_{j=1}^n a_{ij} x_j$.

\begin{proposizione}
\label{prop:ogni-lin-matrice}
Ogni applicazione lineare $f \colon K^n \to K^m$ è uguale a $L_A$ per un'unica matrice $A \in M_{m,n}(K)$.
\end{proposizione}

\begin{proof}
Poniamo $A_j \coloneqq f(e_j) \in K^m$ (colonna $j$ di $A$). Per il Corollario~\ref{cor:lin-base} questa scelta determina $f$ in modo unico, e per costruzione $L_A(e_j) = A_j = f(e_j)$.
\end{proof}

La corrispondenza $A \mapsto L_A$ è dunque una biiezione $M_{m,n}(K) \leftrightarrow \mathcal{L}(K^n, K^m)$, ed è un isomorfismo di spazi vettoriali (come verificato nella Proposizione~3.4.23 degli appunti).

\begin{exbox}
\begin{esempio}
La matrice della rotazione $\rho_\theta$ del piano (rispetto alla base standard) è
$\begin{pmatrix} \cos\theta & -\sin\theta \\ \sin\theta & \cos\theta \end{pmatrix}$.
Il prodotto di due matrici di rotazione corrisponde alla composizione: $\rho_\alpha \circ \rho_\beta = \rho_{\alpha+\beta}$. Sviluppando esplicitamente questo prodotto si ottengono le formule di addizione di seno e coseno:
\[
  \cos(\alpha+\beta) = \cos\alpha\cos\beta - \sin\alpha\sin\beta, \quad
  \sin(\alpha+\beta) = \sin\alpha\cos\beta + \cos\alpha\sin\beta.
\]
\end{esempio}
\end{exbox}

% --------------------------------------------------------
\section{Nucleo e immagine}
\label{sec:kerIm}
% --------------------------------------------------------

\begin{defbox}
\begin{definizione}[Nucleo e immagine]
Sia $f \colon V \to W$ lineare. Il \emph{nucleo} di $f$ è
\[
  \ker f \coloneqq f^{-1}(0) = \{v \in V \mid f(v) = 0\}.
\]
L'\emph{immagine} di $f$ è
\[
  \im f \coloneqq f(V) = \{f(v) \mid v \in V\}.
\]
\end{definizione}
\end{defbox}

\begin{proposizione}
$\ker f$ è un sottospazio di $V$ e $\im f$ è un sottospazio di $W$.
\end{proposizione}

\begin{proof}
Se $v_1, v_2 \in \ker f$ e $\lambda_1, \lambda_2 \in K$, allora $f(\lambda_1 v_1 + \lambda_2 v_2) = \lambda_1 f(v_1) + \lambda_2 f(v_2) = 0$, quindi $\lambda_1 v_1 + \lambda_2 v_2 \in \ker f$. Analogamente per $\im f$: se $w_i = f(v_i)$, allora $\lambda_1 w_1 + \lambda_2 w_2 = f(\lambda_1 v_1 + \lambda_2 v_2) \in \im f$.
\end{proof}

\begin{proposizione}
\label{prop:ker-iniett}
$f$ è iniettiva se e solo se $\ker f = \{0\}$.
\end{proposizione}

\begin{proof}
Se $f$ è iniettiva e $f(v) = 0 = f(0)$, allora $v = 0$. Viceversa, se $\ker f = \{0\}$ e $f(v) = f(w)$, allora $f(v-w) = 0$, quindi $v-w \in \ker f = \{0\}$, da cui $v = w$.
\end{proof}

\begin{exbox}
\begin{esempio}
Sia $A \in M_{m,n}(K)$ e $L_A \colon K^n \to K^m$. Allora:
\begin{itemize}
  \item $\ker L_A$ è il sottospazio delle soluzioni del sistema lineare omogeneo $AX = 0$.
  \item $\im L_A$ è il sottospazio di $K^m$ generato dalle colonne di $A$.
\end{itemize}
\end{esempio}
\end{exbox}

\begin{exbox}
\begin{esempio}
Sia $g \colon K[x] \to K$, $g(p) = p(c)$, la valutazione nel punto $c \in K$. Allora $\ker g = \{(x-c)q \mid q \in K[x]\}$, cioè l'ideale dei polinomi con radice $c$.
\end{esempio}
\end{exbox}

\subsection{Il teorema della dimensione}

Il risultato seguente è uno dei pilastri dell'algebra lineare. Mette in relazione le dimensioni di nucleo e immagine con quella dello spazio di partenza.

\begin{teorema}[della dimensione, o del rango e della nullità]
\label{thm:dim}
Siano $V, W$ spazi vettoriali su $K$, con $V$ finitamente generato. Allora
\[
  \dim V = \dim(\ker f) + \dim(\im f).
\]
\end{teorema}

\begin{proof}
Sia $\{v_1, \ldots, v_a\}$ una base di $\ker f$ e $\{w_1, \ldots, w_b\}$ una base di $\im f$. Scegliamo $u_1, \ldots, u_b \in V$ tali che $f(u_i) = w_i$.

\emph{Generazione.} Sia $v \in V$; poiché $f(v) \in \im f$, esistono $\alpha_i$ tali che $f(v) = \sum \alpha_i w_i = f(\sum \alpha_i u_i)$. Quindi $v - \sum \alpha_i u_i \in \ker f$, da cui $v - \sum \alpha_i u_i = \sum \beta_k v_k$, e si ottiene $v = \sum \alpha_i u_i + \sum \beta_k v_k$.

\emph{Indipendenza.} Supponiamo $\sum \lambda_k v_k + \sum \mu_i u_i = 0$. Applicando $f$: $\sum \mu_i f(u_i) = \sum \mu_i w_i = 0$. Poiché i $w_i$ sono indipendenti, $\mu_i = 0$. Resta $\sum \lambda_k v_k = 0$, da cui $\lambda_k = 0$.

Dunque $\{v_1, \ldots, v_a, u_1, \ldots, u_b\}$ è una base di $V$, e $\dim V = a + b$.
\end{proof}

\begin{corollario}
\label{cor:ker-dim}
Se $\dim V = \dim W$, allora $f$ è biunivoca se e solo se $\ker f = \{0\}$ (equivalentemente, se e solo se $f$ è suriettiva).
\end{corollario}

\begin{proof}
Dal Teorema~\ref{thm:dim}: $\ker f = \{0\}$ implica $\dim(\im f) = \dim V = \dim W$, quindi $\im f = W$. La proposizione~\ref{prop:ker-iniett} dà l'iniettività.
\end{proof}

\begin{corollario}
$\dim(\ker f) \geq \dim V - \dim W$.
\end{corollario}

\begin{exbox}
\begin{esempio}[Successione di Fibonacci e diagonalizzazione]
Sia $A = \begin{pmatrix} 1 & 1 \\ 1 & 0 \end{pmatrix}$. La successione di Fibonacci $\{x_n\}$ verifica $(x_{n+1}, x_n) = L_A^n(1,1)$. La matrice $A$ è diagonalizzabile su $\R$: gli autovalori sono $\phi = \frac{1+\sqrt{5}}{2}$ e $\hat\phi = \frac{1-\sqrt{5}}{2}$, con autovettori $(2, -1+\sqrt{5})$ e $(2, -1-\sqrt{5})$. Dalla diagonalizzazione si ottiene la \emph{formula di Binet}:
\[
  x_n = \frac{1}{\sqrt{5}}\left[\phi^{n+1} - \hat\phi^{n+1}\right].
\]
\end{esempio}
\end{exbox}

% --------------------------------------------------------
\section{Isomorfismi e operazioni tra applicazioni lineari}
\label{sec:iso}
% --------------------------------------------------------

\subsection{Composizione e isomorfismi}

\begin{proposizione}
La composizione di applicazioni lineari è lineare.
\end{proposizione}

\begin{defbox}
\begin{definizione}[Isomorfismo]
Un'applicazione lineare $f \colon V \to W$ è un \emph{isomorfismo} se è biunivoca. Si scrive $f \colon V \xrightarrow{\sim} W$ e si dice che $V$ e $W$ sono \emph{isomorfi} ($V \cong W$).
\end{definizione}
\end{defbox}

\begin{lemma}
Se $f \colon V \to W$ è un isomorfismo, allora $f^{-1} \colon W \to V$ è anch'essa lineare (quindi un isomorfismo).
\end{lemma}

\begin{proof}
Siano $w_1, w_2 \in W$ e $\lambda_1, \lambda_2 \in K$. Poiché $f$ è biunivoca esistono $v_i$ con $f(v_i) = w_i$; per linearità di $f$, $f(\lambda_1 v_1 + \lambda_2 v_2) = \lambda_1 w_1 + \lambda_2 w_2$, quindi $f^{-1}(\lambda_1 w_1 + \lambda_2 w_2) = \lambda_1 v_1 + \lambda_2 v_2 = \lambda_1 f^{-1}(w_1) + \lambda_2 f^{-1}(w_2)$.
\end{proof}

\begin{proposizione}
Due spazi vettoriali finitamente generati su $K$ sono isomorfi se e solo se hanno la stessa dimensione.
\end{proposizione}

\begin{proof}
Se $f \colon V \to W$ è un isomorfismo, esso manda basi in basi (poiché preserva indipendenza lineare e generazione), quindi $\dim V = \dim W$. Viceversa, se $\dim V = \dim W = n$, scegliendo basi $\mathcal{B}$ di $V$ e $\mathcal{C}$ di $W$, le mappe coordinate $X_\mathcal{B} \colon V \to K^n$ e $X_\mathcal{C} \colon W \to K^n$ sono isomorfismi, quindi $X_\mathcal{C}^{-1} \circ X_\mathcal{B} \colon V \to W$ è un isomorfismo.
\end{proof}

\subsection{Matrici invertibili e gruppo lineare}

\begin{defbox}
\begin{definizione}[Matrice invertibile e gruppo lineare]
Una matrice $A \in M_{n,n}(K)$ è \emph{invertibile} se esiste $B \in M_{n,n}(K)$ con $AB = 1_n = BA$. Il \emph{gruppo lineare generale} è
\[
  \GL_n(K) \coloneqq \{A \in M_{n,n}(K) \mid A \text{ invertibile}\}.
\]
\end{definizione}
\end{defbox}

\begin{proposizione}
$A \in M_{n,n}(K)$ è invertibile se e solo se $L_A \colon K^n \to K^n$ è un isomorfismo.
\end{proposizione}

In altri termini, $\GL_n(K)$ parametrizza gli automorfismi di $K^n$.

\subsection{Lo spazio delle applicazioni lineari}

\begin{defbox}
\begin{definizione}
Siano $V, W$ spazi vettoriali su $K$. Denotiamo con $\mathcal{L}(V, W)$ l'insieme delle applicazioni lineari $f \colon V \to W$. Esso è uno spazio vettoriale su $K$ con le operazioni
\[
  (f+g)(v) \coloneqq f(v) + g(v), \qquad (\lambda f)(v) \coloneqq \lambda f(v).
\]
Il \emph{duale} di $V$ è $V^\vee \coloneqq \mathcal{L}(V, K)$.
\end{definizione}
\end{defbox}

% --------------------------------------------------------
\section{Il primo teorema di isomorfismo}
\label{sec:primoisomorf}
% --------------------------------------------------------

Sia $f \colon V \to W$ lineare e sia $\pi \colon V \to V/\ker f$ l'applicazione quoziente. Poiché $\ker f$ ``collassa'' a zero sotto $f$, è naturale aspettarsi che $f$ discenda a un'applicazione su $V/\ker f$.

\begin{proposizione}
Esiste una e una sola applicazione lineare $\bar{f} \colon V/\ker f \to W$ tale che $\bar{f} \circ \pi = f$.
\end{proposizione}

\begin{proof}
Definiamo $\bar{f}([v]) \coloneqq f(v)$. Questa è ben definita: se $[v'] = [v]$, allora $v' - v \in \ker f$, quindi $f(v') = f(v)$. La linearità di $\bar{f}$ si verifica direttamente.
\end{proof}

\begin{teorema}[Primo teorema di isomorfismo]
\label{thm:iso}
L'applicazione $\bar{f} \colon V/\ker f \to \im f$ è un isomorfismo.
\end{teorema}

\begin{proof}
$\bar{f}$ è suriettiva su $\im f$ per definizione. Se $\bar{f}([v]) = 0$, allora $f(v) = 0$, quindi $v \in \ker f$, cioè $[v] = 0$: $\bar{f}$ è iniettiva.
\end{proof}

\begin{osservazione}
Il Teorema~\ref{thm:iso} implica il Teorema~\ref{thm:dim}: poiché $V/\ker f \cong \im f$, si ha $\dim(V/\ker f) = \dim(\im f)$; e $\dim(V/\ker f) = \dim V - \dim(\ker f)$ per la proposizione sulla dimensione del quoziente del Capitolo~2.
\end{osservazione}

\begin{exbox}
\begin{esempio}
Sia $f \colon \R[x]_{\leq 2} \to \R^2$, $f(p) = (p(0), p'(0))$. Il nucleo è $\{p \mid p(0) = 0, p'(0) = 0\} = \langle x^2 \rangle$ (polinomi della forma $cx^2$). L'immagine è tutto $\R^2$. Quindi $\R[x]_{\leq 2} / \langle x^2 \rangle \cong \R^2$.
\end{esempio}
\end{exbox}

% --------------------------------------------------------
\section{Dalla matrice all'applicazione: la matrice associata}
\label{sec:matrice-associata}
% --------------------------------------------------------

Quando $V$ e $W$ sono spazi astratti (non necessariamente $K^n$), possiamo ancora associare a ogni applicazione lineare una matrice, a patto di aver scelto basi.

\begin{defbox}
\begin{definizione}[Matrice di un'applicazione lineare]
Siano $V, W$ finitamente generati con basi $\mathcal{B} = \{v_1, \ldots, v_n\}$ e $\mathcal{C} = \{w_1, \ldots, w_m\}$. La \emph{matrice associata} a $f \colon V \to W$ è la matrice $M^\mathcal{B}_\mathcal{C}(f) \in M_{m,n}(K)$ la cui $j$-esima colonna è il vettore delle $\mathcal{C}$-coordinate di $f(v_j)$:
\[
  f(v_j) = \sum_{i=1}^m a_{ij} w_i.
\]
\end{definizione}
\end{defbox}

La formula fondamentale è:
\[
  X_\mathcal{C}(f(v)) = M^\mathcal{B}_\mathcal{C}(f) \cdot X_\mathcal{B}(v) \quad \forall v \in V.
\]

\begin{proposizione}
L'applicazione $f \mapsto M^\mathcal{B}_\mathcal{C}(f)$ è un isomorfismo di spazi vettoriali $\mathcal{L}(V,W) \xrightarrow{\sim} M_{m,n}(K)$.
\end{proposizione}

\begin{proposizione}[Compatibilità con la composizione]
Siano $U, V, W$ finitamente generati con basi $\mathcal{B}, \mathcal{C}, \mathcal{D}$ rispettivamente. Per $g \colon U \to V$ e $f \colon V \to W$ lineari:
\[
  M^\mathcal{B}_\mathcal{D}(f \circ g) = M^\mathcal{C}_\mathcal{D}(f) \cdot M^\mathcal{B}_\mathcal{C}(g).
\]
\end{proposizione}

\begin{proof}
Per ogni $u \in U$: $X_\mathcal{D}(f(g(u))) = M^\mathcal{C}_\mathcal{D}(f) \cdot X_\mathcal{C}(g(u)) = M^\mathcal{C}_\mathcal{D}(f) \cdot M^\mathcal{B}_\mathcal{C}(g) \cdot X_\mathcal{B}(u)$.
\end{proof}

\begin{exbox}
\begin{esempio}[Derivazione su polinomi]
Su $V = \R[x]_{\leq 2}$ con base $\mathcal{B} = \{1, x, x^2\}$, l'applicazione di derivazione $f(p) = p + p'$ ha matrice
\[
  M^\mathcal{B}_\mathcal{B}(f) = \begin{pmatrix} 1 & 1 & 0 \\ 0 & 1 & 2 \\ 0 & 0 & 1 \end{pmatrix}.
\]
(Colonne: $f(1)=1$, $f(x)=x+1$, $f(x^2)=x^2+2x$.)
\end{esempio}
\end{exbox}

% --------------------------------------------------------
\section{Operazioni elementari ed eliminazione di Gauss}
\label{sec:gauss}
% --------------------------------------------------------

\subsection{Matrici a scala}

\begin{defbox}
\begin{definizione}[Matrice a scala per colonne]
Una matrice $A \in M_{m,n}(K)$ è \emph{a scala per colonne} se, denotando con $p_A(j) = \min\{i \mid a_{ij} \neq 0\}$ la ``profondità'' della colonna $j$ (e $p_A(j) = \infty$ se la colonna è nulla), vale $p_A(1) < p_A(2) < \cdots < p_A(n)$. L'entrata non nulla di minimo indice di riga in ogni colonna non nulla si chiama \emph{pivot}.
\end{definizione}
\end{defbox}

Analogamente si definisce una \emph{matrice a scala per righe}.

\begin{proposizione}
Se $S$ è a scala per colonne, le sue colonne non nulle formano una base di $\im L_S$.
\end{proposizione}

\begin{proof}
Le colonne non nulle generano $\im L_S$ (per definizione). Sono linearmente indipendenti: in una combinazione lineare nulla $\sum \lambda_j S^{(j)} = 0$, la riga del pivot della prima colonna non nulla ha contributo solo da quella colonna, quindi $\lambda_1 = 0$; per induzione tutti i coefficienti sono nulli.
\end{proof}

\subsection{Operazioni elementari}

\begin{defbox}
\begin{definizione}[Operazioni elementari]
Le \emph{operazioni elementari} su una lista di vettori $v_1, \ldots, v_n$ di uno spazio vettoriale sono:
\begin{enumerate}
  \item Scambiare $v_i$ e $v_j$.
  \item Sostituire $v_i$ con $v_i + \lambda v_j$ ($i \neq j$, $\lambda \in K$).
  \item Moltiplicare $v_i$ per un non-zero $\lambda \in K$.
\end{enumerate}
\end{definizione}
\end{defbox}

\begin{proposizione}
Un'operazione elementare non cambia il sottospazio generato: $\langle v_1, \ldots, v_n \rangle = \langle w_1, \ldots, w_n \rangle$.
\end{proposizione}

\begin{teorema}[Riduzione a scala per colonne]
Data $A \in M_{m,n}(K)$, con operazioni elementari sulle colonne (di tipo~(1) e~(2)) si può ridurre $A$ a una matrice a scala $S$. Le colonne non nulle di $S$ formano una base di $\im L_A$, e $\rg(A) = $ numero di colonne non nulle di $S$.
\end{teorema}

\begin{exbox}
\begin{esempio}
Determiniamo una base di $\im L_A$ per
\[
  A = \begin{pmatrix} 3 & 2 & 1 & 5 \\ 1 & 0 & 2 & 11 \\ 2 & -3 & -1 & -1 \\ 6 & -1 & 2 & 15 \end{pmatrix}.
\]
Scambiando prima e terza colonna e applicando operazioni di tipo~(2) (si veda l'Esempio~3.7.13 degli appunti originali) si arriva alla forma a scala
\[
  S = \begin{pmatrix} 1 & 0 & 0 & 0 \\ 2 & 1 & 0 & 0 \\ -1 & 4 & 1 & 0 \\ 2 & 5 & 1 & 0 \end{pmatrix}.
\]
Una base di $\im L_A$ è $\{(1,2,-1,2), (0,1,4,5), (0,0,1,1)\}$, e $\rg(A) = 3$.
\end{esempio}
\end{exbox}

\subsection{Eliminazione di Gauss}

L'eliminazione di Gauss risolve i sistemi $AX = B$ operando sulle righe della matrice aumentata $[A|B]$.

\begin{defbox}
\begin{definizione}[Matrice trasposta]
La \emph{trasposta} di $A \in M_{m,n}(K)$ è la matrice $A^t \in M_{n,m}(K)$ con $(A^t)_{ij} = a_{ji}$.
\end{definizione}
\end{defbox}

\begin{proposizione}[Invarianza del rango per trasposizione]
$\rg(A) = \rg(A^t)$.
\end{proposizione}

Questo risultato, non ovvio a prima vista, afferma che il massimo numero di righe linearmente indipendenti di $A$ è uguale al massimo numero di colonne linearmente indipendenti.

\begin{teorema}[Eliminazione di Gauss]
Con operazioni elementari sulle righe si può portare $A$ a una matrice a scala per righe $S$, conservando l'insieme delle soluzioni di $AX = 0$ (e più in generale di $AX = B$). Il sistema $AX = 0$ ha $n - \rg(A)$ gradi di libertà.
\end{teorema}

\begin{exbox}
\begin{esempio}
Risolviamo in $\Q^3$ il sistema
\[
  x_1 - 4x_2 + 5x_3 = b_1, \quad 2x_2 + 3x_3 = b_2.
\]
La matrice è già a scala per righe. Dalla seconda equazione: $x_2 = -\tfrac{3}{2}x_3 - \tfrac{1}{2}b_2$. Sostituendo: $x_1 = -11x_3 + b_1 - 2b_2$. Le soluzioni formano la famiglia
\[
  \{(-11t + b_1 - 2b_2,\; -\tfrac{3}{2}t - \tfrac{1}{2}b_2,\; t) \mid t \in \Q\}.
\]
\end{esempio}
\end{exbox}

\subsection{Calcolo dell'inversa}

\begin{proposizione}
Sia $A \in \GL_n(K)$. Con operazioni elementari sulle righe di $[A \,|\, 1_n]$ si ottiene $[1_n \,|\, A^{-1}]$.
\end{proposizione}

\begin{proof}
Ogni operazione elementare sulle righe equivale alla moltiplicazione a sinistra per una matrice invertibile $L_i$. Al termine, $L_r \cdots L_1 A = 1_n$, da cui $L_r \cdots L_1 = A^{-1}$; la stessa sequenza applicata a $1_n$ dà $A^{-1}$.
\end{proof}

\begin{exbox}
\begin{esempio}
Calcoliamo l'inversa di $A = \begin{pmatrix} 2 & 1 & 3 \\ -1 & 0 & 1 \\ 3 & 2 & 8 \end{pmatrix}$ con il metodo di Gauss-Jordan. Partendo da $[A\,|\,I_3]$ e applicando le operazioni elementari descritte nell'Esempio~3.9.4 degli appunti originali si ottiene:
\[
  A^{-1} = \begin{pmatrix} -2 & -2 & 1 \\ 11 & 7 & -5 \\ -2 & -1 & 1 \end{pmatrix}.
\]
\end{esempio}
\end{exbox}

% --------------------------------------------------------
\section{Cambiamenti di base e coniugio}
\label{sec:cambiobase}
% --------------------------------------------------------

\subsection{Matrice del cambiamento di base}

Sia $V$ di dimensione $n$ con basi $\mathcal{B} = \{u_1, \ldots, u_n\}$ e $\mathcal{C} = \{w_1, \ldots, w_n\}$. La \emph{matrice del cambiamento di base da $\mathcal{B}$ a $\mathcal{C}$} è
\[
  M^\mathcal{B}_\mathcal{C}(\Id_V) \in \GL_n(K).
\]
Essa permette di convertire $\mathcal{B}$-coordinate in $\mathcal{C}$-coordinate:
\[
  X_\mathcal{C}(v) = M^\mathcal{B}_\mathcal{C}(\Id_V) \cdot X_\mathcal{B}(v).
\]

\begin{osservazione}
La matrice del cambiamento di base è invertibile e la sua inversa è la matrice del cambiamento inverso:
\[
  \bigl(M^\mathcal{B}_\mathcal{C}(\Id_V)\bigr)^{-1} = M^\mathcal{C}_\mathcal{B}(\Id_V).
\]
\end{osservazione}

\begin{exbox}
\begin{esempio}
In $\Q^3$, siano $\mathcal{B} = \{(2,1,0),(1,0,-1),(0,-1,2)\}$ e $\mathcal{D} = \{(2,-1,3),(1,0,2),(3,1,8)\}$. Per la formula $M^\mathcal{B}_\mathcal{D}(\Id) = [D_1|D_2|D_3]^{-1} \cdot [B_1|B_2|B_3]$, usando l'inversa calcolata nell'esempio precedente:
\[
  M^\mathcal{B}_\mathcal{D}(\Id) = \begin{pmatrix} -6 & -3 & 4 \\ 29 & 16 & -17 \\ -5 & -3 & 3 \end{pmatrix}.
\]
Per esempio, le $\mathcal{D}$-coordinate di $(2,1,0)$ sono $(-6,29,-5)$.
\end{esempio}
\end{exbox}

\subsection{Matrici coniugate}

Quando cambiamo base, la matrice di un endomorfismo cambia per coniugio. Precisamente, se $f \colon V \to V$ è un endomorfismo e $\mathcal{B}, \mathcal{C}$ sono due basi, allora
\[
  M^\mathcal{B}_\mathcal{B}(f) = G^{-1} \cdot M^\mathcal{C}_\mathcal{C}(f) \cdot G, \quad G = M^\mathcal{B}_\mathcal{C}(\Id_V).
\]

\begin{defbox}
\begin{definizione}[Matrici coniugate]
Due matrici $M, N \in M_{n,n}(K)$ sono \emph{coniugate} (si scrive $M \sim N$) se esiste $G \in \GL_n(K)$ tale che $M = G^{-1} N G$.
\end{definizione}
\end{defbox}

Il coniugio è una relazione di equivalenza su $M_{n,n}(K)$ (verificare riflessività, simmetria, transitività è un esercizio). La classificazione delle classi di coniugio è uno dei problemi centrali dell'algebra lineare, intimamente legato alla diagonalizzazione.

% --------------------------------------------------------
\section{Diagonalizzazione}
\label{sec:diag}
% --------------------------------------------------------

\subsection{Autovalori e autovettori}

Un endomorfismo $f \colon V \to V$ è \emph{diagonalizzabile} se esiste una base $\mathcal{B}$ di $V$ rispetto alla quale $M^\mathcal{B}_\mathcal{B}(f)$ è diagonale. Questo accade se e solo se esiste una base di $V$ formata da autovettori.

\begin{defbox}
\begin{definizione}[Autovalore, autovettore, autospazio]
Sia $f \colon V \to V$ lineare. Un vettore $0 \neq v \in V$ è un \emph{autovettore} di $f$ con \emph{autovalore} $\lambda \in K$ se $f(v) = \lambda v$. L'\emph{autospazio} associato a $\lambda$ è
\[
  V_\lambda(f) \coloneqq \ker(f - \lambda\, \Id_V).
\]
\end{definizione}
\end{defbox}

\begin{osservazione}
$\lambda$ è un autovalore se e solo se $V_\lambda(f) \neq \{0\}$, ovvero se e solo se $\ker(f - \lambda \Id) \neq \{0\}$, ovvero se e solo se $f - \lambda \Id$ non è iniettiva.
\end{osservazione}

\begin{exbox}
\begin{esempio}
Sia $A = \begin{pmatrix} 1 & -1 \\ 2 & 4 \end{pmatrix}$. Verifichiamo che $L_A(1,-1) = (2,-2) = 2(1,-1)$ e $L_A(1,-2) = (3,-6) = 3(1,-2)$. Quindi $2$ e $3$ sono autovalori con autovettori $(1,-1)$ e $(1,-2)$. La base $\{(1,-1),(1,-2)\}$ diagonalizza $A$:
\[
  M^\mathcal{B}_\mathcal{B}(L_A) = \begin{pmatrix} 2 & 0 \\ 0 & 3 \end{pmatrix}.
\]
\end{esempio}
\end{exbox}

\begin{exbox}
\begin{esempio}[Rotazioni non diagonalizzabili su $\R$]
La matrice di rotazione $R_\theta = \begin{pmatrix} \cos\theta & -\sin\theta \\ \sin\theta & \cos\theta \end{pmatrix}$ con $\theta \notin \pi\Z$ non ha autovalori reali (nessuna retta è invariante sotto una rotazione non banale). Tuttavia, vista su $\C$, gli autovalori sono $e^{i\theta}$ e $e^{-i\theta}$, con autovettori $(1, -i)$ e $(1, i)$. Questo esempio mostra come la diagonalizzabilità dipenda dal campo.
\end{esempio}
\end{exbox}

\begin{exbox}
\begin{esempio}[Matrice non diagonalizzabile su qualsiasi campo]
La matrice $N = \begin{pmatrix} 0 & 1 \\ 0 & 0 \end{pmatrix}$ ha $0$ come unico autovalore (con autovettori $(t, 0)$, $t \neq 0$). Non esiste una base di autovettori in $K^2$, quindi $N$ non è diagonalizzabile su nessun campo.
\end{esempio}
\end{exbox}

\subsection{Perché diagonalizzare?}

Se $A = G \Lambda G^{-1}$ con $\Lambda$ diagonale (con entrate $\lambda_1, \ldots, \lambda_n$), allora per ogni $r \in \Z$:
\[
  A^r = G \Lambda^r G^{-1}, \quad \Lambda^r = \mathrm{diag}(\lambda_1^r, \ldots, \lambda_n^r).
\]

Ancor più importante: l'esponenziale di matrice $e^A = \sum_{r=0}^\infty \frac{1}{r!} A^r$ si calcola agevolmente:
\[
  e^{tA} = G \cdot \mathrm{diag}(e^{t\lambda_1}, \ldots, e^{t\lambda_n}) \cdot G^{-1}.
\]
Questa formula risolve il sistema di equazioni differenziali $Y'(t) = A Y(t)$, $Y(0) = B$: la soluzione è $Y(t) = e^{tA} B$.

\begin{exbox}
\begin{esempio}
Sia $A = \begin{pmatrix} 1 & -1 \\ 2 & 4 \end{pmatrix}$ come sopra, con autovalori $2, 3$. Allora
\[
  e^{tA} = \begin{pmatrix} 2e^{2t} - e^{3t} & e^{2t} - e^{3t} \\ -2e^{2t} + 2e^{3t} & -e^{2t} + 2e^{3t} \end{pmatrix}.
\]
\end{esempio}
\end{exbox}

% --------------------------------------------------------
\section{Il duale di uno spazio vettoriale}
\label{sec:duale}
% --------------------------------------------------------

\subsection{Base duale}

Ricordiamo che il \emph{duale} di $V$ è $V^\vee = \mathcal{L}(V, K)$.

\begin{proposizione}
Se $\dim V = n$, allora $\dim V^\vee = n$.
\end{proposizione}

\begin{defbox}
\begin{definizione}[Base duale]
Sia $\mathcal{B} = \{v_1, \ldots, v_n\}$ una base di $V$. La \emph{base duale} $\mathcal{B}^\vee = \{v_1^\vee, \ldots, v_n^\vee\}$ è definita da
\[
  v_i^\vee(v_j) = \delta_{ij} \coloneqq \begin{cases} 1 & i = j \\ 0 & i \neq j \end{cases}.
\]
\end{definizione}
\end{defbox}

\begin{proposizione}
$\mathcal{B}^\vee$ è effettivamente una base di $V^\vee$.
\end{proposizione}

Attenzione: l'isomorfismo $\phi_\mathcal{B} \colon V \to V^\vee$, $\sum x_i v_i \mapsto \sum x_i v_i^\vee$, dipende dalla scelta di $\mathcal{B}$: \emph{non esiste un isomorfismo naturale (indipendente dalle scelte) tra $V$ e $V^\vee$}.

\subsection{Il biduale e la naturalità}

\begin{proposizione}
L'applicazione di valutazione
\[
  \mathrm{Val} \colon V \to V^{\vee\vee}, \quad v \mapsto [\mathrm{Val}(v) \colon f \mapsto f(v)]
\]
è un isomorfismo (naturale, ovvero indipendente da scelte), per ogni $V$ finitamente generato.
\end{proposizione}

\begin{proof}
$\mathrm{Val}$ è lineare (verifica diretta). Per la Proposizione~\ref{prop:ker-iniett} basta mostrare che $\ker\mathrm{Val} = \{0\}$: se $v \neq 0$, completando $v$ a una base si trova $f \in V^\vee$ con $f(v) = 1 \neq 0$.
\end{proof}

Questo risultato giustifica l'identificazione $V \cong V^{\vee\vee}$ come identificazione \emph{canonica}, a differenza di $V \cong V^\vee$ che richiede la scelta di una base.

\subsection{Applicazione duale}

\begin{defbox}
\begin{definizione}[Applicazione duale]
Sia $\varphi \colon V \to W$ lineare. L'\emph{applicazione duale} è
\[
  \varphi^\vee \colon W^\vee \to V^\vee, \quad f \mapsto f \circ \varphi.
\]
\end{definizione}
\end{defbox}

\begin{proposizione}
Se $\mathcal{B}, \mathcal{C}$ sono basi di $V, W$, allora la matrice di $\varphi^\vee$ rispetto alle basi duali è la trasposta della matrice di $\varphi$:
\[
  M^{\mathcal{C}^\vee}_{\mathcal{B}^\vee}(\varphi^\vee) = \bigl(M^\mathcal{B}_\mathcal{C}(\varphi)\bigr)^t.
\]
\end{proposizione}

\begin{proposizione}
$\varphi$ è iniettiva se e solo se $\varphi^\vee$ è suriettiva; $\varphi$ è suriettiva se e solo se $\varphi^\vee$ è iniettiva.
\end{proposizione}

\begin{proof}[Idea della dimostrazione]
Per la fattorizzazione $\varphi = \iota \circ \psi$ (con $\iota \colon \im\varphi \hookrightarrow W$ e $\psi \colon V \twoheadrightarrow \im\varphi$) si mostra che $\iota^\vee$ è suriettiva (usando la base duale di $\im\varphi$ estesa a quella di $W$) e si conclude.
\end{proof}

\begin{corollario}
$\rg(\varphi) = \rg(\varphi^\vee)$.
\end{corollario}

\begin{exbox}
\begin{esempio}[Annullatore]
Dato un sottospazio $W \subseteq V$, l'\emph{annullatore} di $W$ è
\[
  \Ann(W) \coloneqq \{\phi \in V^\vee \mid \phi|_W = 0\}.
\]
Questo è un sottospazio di $V^\vee$, isomorfo a $(V/W)^\vee$ tramite l'applicazione duale della proiezione $\pi \colon V \to V/W$. In particolare $\dim \Ann(W) = \dim V - \dim W$.
\end{esempio}
\end{exbox}

% --------------------------------------------------------
\section*{Esercizi}
\addcontentsline{toc}{section}{Esercizi del Capitolo 3}
% --------------------------------------------------------

\begin{esercizio}
Stabilire quali delle seguenti applicazioni sono lineari su $K[x]$:
\begin{enumerate}
  \item $\Phi(p) = p_0 \cdot p$ per un fissato $p_0 \in K[x]$.
  \item $\Psi(p) = p^2$.
  \item $\Theta(p) = p(x^2)$.
  \item $F(p) = p(0) + p(1)$.
\end{enumerate}
Per (2), la risposta dipende dalla caratteristica di $K$ (si pensi a $\mathbb{F}_2$).
\end{esercizio}

\begin{esercizio}
Sia $f \colon \R^3 \to \R^2$ definita da $f(x_1, x_2, x_3) = (2x_1 + x_2, x_1 - x_3)$. Trovare una base di $\ker f$ e calcolare $\dim(\im f)$.
\end{esercizio}

\begin{esercizio}[Guidato]
Sia $A = \begin{pmatrix} 1 & 2 & -1 \\ -1 & 1 & -1 \\ 0 & -3 & 2 \end{pmatrix}$.
\begin{enumerate}
  \item Applicate l'eliminazione di Gauss ad $A$ per trovare una matrice a scala per righe $S$.
  \item Determinare $\ker L_A$ e $\im L_A$ (basi esplicite).
  \item Verificare il Teorema~\ref{thm:dim}: $\dim\ker + \dim\im = 3$.
\end{enumerate}
\end{esercizio}

\begin{esercizio}
Calcolare $A \cdot B$ per
\[
  A = \begin{pmatrix} 2 & 3 \\ 3 & 4 \end{pmatrix}, \quad B = \begin{pmatrix} 1 & -1 \\ 1 & 1 \end{pmatrix}.
\]
\end{esercizio}

\begin{esercizio}
Sia $C = \begin{pmatrix} 1 & 2 \\ 2 & 1 \end{pmatrix}$ e $\mathcal{B} = \{(1,1),(1,-1)\}$.
\begin{enumerate}
  \item Calcolare $M^\mathcal{B}_\mathcal{B}(L_C)$.
  \item Dedurre una formula chiusa per $C^n$ (per ogni $n \in \N$).
\end{enumerate}
\end{esercizio}

\begin{esercizio}
Dimostrare che due matrici $A, B \in M_{n,n}(K)$ coniugate hanno lo stesso rango.
\end{esercizio}

\begin{esercizio}
Siano $U, W \subseteq \R^4$ i sottospazi
\[
  U = \langle (1,2,3,-1),(3,5,0,2)\rangle, \quad W = \langle (-1,0,3,2),(1,-1,1,-1),(1,-2,5,0)\rangle.
\]
Trovare equazioni cartesiane di $U$ e di $W$ (suggerimento: usare il teorema del rango e la trasposizione).
\end{esercizio}

\begin{esercizio}[Successione di Fibonacci]
Sia $A = \begin{pmatrix} 1 & 1 \\ 1 & 0 \end{pmatrix}$.
\begin{enumerate}
  \item Diagonalizzare $A$ su $\R$ trovando esplicitamente autovettori e autovalori.
  \item Dedurre la formula di Binet per $x_n$ (vedi l'Esempio nel testo).
\end{enumerate}
\end{esercizio}

\begin{esercizio}
Sia $f \colon V \to V$ un endomorfismo di uno spazio finitamente generato su $K$ con $\mathrm{char}(K) \neq 2$, e supponiamo $f \circ f = \Id_V$.
\begin{enumerate}
  \item Dimostrare che $V = V_1(f) \oplus V_{-1}(f)$ (cioè $V$ si decompone in somma diretta degli autospazi relativi a $+1$ e $-1$).
  \item Dare un esempio in cui $\mathrm{char}(K) = 2$ e la decomposizione fallisce.
\end{enumerate}
\end{esercizio}

\begin{esercizio}
Sia $V \subseteq K^n$ il sottospazio $\{X \mid x_1 + \cdots + x_n = 0\}$.
\begin{enumerate}
  \item Trovare una base di $V^\vee$.
  \item Descrivere $\Ann(V)$ come sottospazio di $(K^n)^\vee$.
\end{enumerate}
\end{esercizio}

\begin{esercizio}
Calcolare l'inversa della matrice
\[
  A = \begin{pmatrix} 1 & 1 & 1 \\ 1 & 2 & 3 \\ 1 & 4 & 9 \end{pmatrix}
\]
con il metodo di Gauss-Jordan.
\end{esercizio}


% ============================================================
% CAPITOLO 4 — SPAZI AFFINI
% ============================================================

% ============================================================
\chapter{Spazi affini}
\label{cap:affini}
% ============================================================

\section*{Introduzione al capitolo}

Finora abbiamo lavorato esclusivamente con spazi vettoriali, strutture in cui ogni elemento è un vettore e l'origine $\mathbf{0}$ svolge un ruolo privilegiato. Eppure la geometria elementare del piano e dello spazio --- quella di Euclide --- non ha un'origine privilegiata: tutti i punti sono equivalenti, e ha senso parlare di triangoli, rette e parallelismo indipendentemente dalla scelta di un sistema di riferimento.

Gli \emph{spazi affini} formalizzano questa idea. Uno spazio affine è un insieme di \emph{punti} su cui uno spazio vettoriale (il \emph{gruppo delle traslazioni}) agisce in modo coerente: possiamo «traslare» un punto lungo un vettore, ma non possiamo «sommare» due punti. La struttura è più povera di quella vettoriale, ma sufficiente per parlare di rette, piani, parallelismo, baricentri e applicazioni affini --- trasformazioni che preservano queste strutture.

Il capitolo è organizzato come segue. Nella Sezione~\ref{sec:spazi-affini} introduciamo la definizione assiomatica di spazio affine, motivata dal comportamento del piano euclideo. La Sezione~\ref{sec:comb-lin-punti} spiega come combinare \emph{punti} (non vettori) con coefficienti la cui somma valga~$1$. La Sezione~\ref{sec:sottospazi-affini} studia i sottospazi affini e il parallelismo. La Sezione~\ref{sec:appl-affini} è dedicata alle applicazioni affini e ai sistemi di riferimento. Infine, la Sezione~\ref{sec:ginnastica-affine} sviluppa la «ginnastica» concreta: equazioni parametriche e cartesiane di sottospazi, e fasci di rette e piani.

% --------------------------------------------------------
\section{Spazi affini}
\label{sec:spazi-affini}
% --------------------------------------------------------

\subsection{La struttura del piano euclideo come modello}

Il piano euclideo $E^2$ è l'insieme dei punti del piano della geometria elementare. A esso è associato lo spazio vettoriale $\VV(E^2)$ dei \emph{vettori geometrici}: classi di equipollenza di segmenti orientati, che si sommano con la regola del parallelogramma. Esiste un'applicazione naturale
\[
  \VV(E^2) \times E^2 \longrightarrow E^2, \quad (v, P) \longmapsto P + v,
\]
dove $P + v$ è il punto $Q$ tale che $\overrightarrow{PQ}$ rappresenta il vettore $v$. Questa applicazione soddisfa tre proprietà fondamentali:
\begin{enumerate}
  \item $P + \mathbf{0} = P$ per ogni $P \in E^2$ (traslare di zero non sposta nessun punto),
  \item $P + (v + w) = (P + v) + w$ per ogni $P \in E^2$, $v, w \in \VV(E^2)$ (le traslazioni si compongono),
  \item dati $P, Q \in E^2$, esiste un unico $v \in \VV(E^2)$ tale che $P + v = Q$ (ogni coppia di punti è collegata da un unico vettore).
\end{enumerate}
Queste tre proprietà catturano l'essenza geometrica del piano euclideo e servono come assiomi per la definizione astratta di spazio affine.

\begin{defbox}
\begin{definizione}[Spazio affine]
\label{def:spazio-affine}
Sia $V$ uno spazio vettoriale sul campo $\K$. Uno \emph{spazio affine con gruppo delle traslazioni $V$} è un insieme non vuoto $S$ munito di un'applicazione
\[
  V \times S \longrightarrow S, \quad (v, P) \longmapsto P + v,
\]
che soddisfa le seguenti proprietà:
\begin{enumerate}
  \item $P + \mathbf{0} = P$ per ogni $P \in S$,
  \item $P + (v + w) = (P + v) + w$ per ogni $P \in S$ e $v, w \in V$,
  \item per ogni $P, Q \in S$ esiste un unico $v \in V$ tale che $P + v = Q$.
\end{enumerate}
Gli elementi di $S$ si chiamano \emph{punti}. Si pone $\VV(S) \coloneqq V$ e si dice che $S$ è uno spazio affine \emph{su $\K$} se $\VV(S)$ è uno spazio vettoriale su $\K$.
\end{definizione}
\end{defbox}

Notate subito la differenza con gli spazi vettoriali: in uno spazio affine \emph{non} è definita la somma di due punti né il prodotto di un punto per uno scalare. Solo la «traslazione» di un punto lungo un vettore è ammessa.

\begin{exbox}
\begin{esempio}[Il modello fondamentale: soluzioni di sistemi lineari]
\label{es:spazio-affine-sistemi}
Siano $A \in M_{m,n}(\K)$ e $B \in M_{m,1}(\K)$. Consideriamo
\[
  S \coloneqq \{ Y \in \K^n \mid A \cdot Y = B \}, \qquad V \coloneqq \{ X \in \K^n \mid A \cdot X = \mathbf{0} \}.
\]
Se $S \neq \emptyset$, allora $S$ è uno spazio affine con gruppo delle traslazioni $V$. Infatti: se $Y \in S$ e $X \in V$, allora $A(Y + X) = AY + AX = B + \mathbf{0} = B$, quindi $Y + X \in S$; le proprietà (1) e (2) sono ovvie; per~(3), se $Y, Z \in S$ il vettore $Z - Y$ soddisfa $A(Z - Y) = 0$, dunque $Z - Y \in V$ e $Y + (Z - Y) = Z$.

Questo esempio chiarisce perché lo studio delle soluzioni di sistemi lineari è intrinsecamente «affine»: l'insieme delle soluzioni non è mai un sottospazio vettoriale (a meno che $B = \mathbf{0}$), ma è sempre uno spazio affine traslato del nucleo.
\end{esempio}
\end{exbox}

\begin{exbox}
\begin{esempio}[Ogni spazio vettoriale è uno spazio affine]
\label{es:V-come-affine}
Sia $V$ uno spazio vettoriale su $\K$. L'applicazione $(u, v) \mapsto v + u$ dà a $V$ una struttura di spazio affine con gruppo delle traslazioni $V$ stesso. Lo \emph{spazio affine standard}
\[
  \mathbb{A}^n(\K) \coloneqq \K^n
\]
è $\K^n$ visto come spazio affine con questa struttura.
\end{esempio}
\end{exbox}

\subsection{Traslazioni e vettore tra due punti}

Sia $S$ uno spazio affine con gruppo delle traslazioni $V$. Dato $v \in V$, la \emph{traslazione di vettore $v$} è l'applicazione
\[
  \tau_v \colon S \to S, \quad P \mapsto P + v.
\]
Le proprietà della Definizione~\ref{def:spazio-affine} si riscrivono come:
\begin{itemize}
  \item[(I)] $\tau_{\mathbf{0}} = \Id_S$,
  \item[(II)] $\tau_w \circ \tau_v = \tau_{w+v}$ per ogni $v, w \in V$,
  \item[(III)] per ogni $P, Q \in S$ esiste un unico $v \in V$ tale che $\tau_v(P) = Q$.
\end{itemize}

\begin{proposizione}
\label{prop:tau-biunivoca}
Sia $S$ uno spazio affine con gruppo delle traslazioni $V$ e $v \in V$.
\begin{enumerate}[label=(\alph*)]
  \item L'applicazione $\tau_v$ è biunivoca.
  \item Se esiste $P \in S$ tale che $\tau_v(P) = P$, allora $v = \mathbf{0}$. In altre parole, $\tau_v$ è priva di punti fissi se $v \neq \mathbf{0}$.
\end{enumerate}
\end{proposizione}

\begin{proof}
(a) Per (II) e (I) si ha $\tau_{-v} \circ \tau_v = \tau_v \circ \tau_{-v} = \tau_{\mathbf{0}} = \Id_S$, dunque $\tau_v$ è biunivoca con inversa $\tau_{-v}$.

(b) Se $\tau_v(P) = P = \tau_{\mathbf{0}}(P)$, per l'unicità in (III) si conclude $v = \mathbf{0}$.
\end{proof}

\begin{osservazione}
La scelta di un punto $P \in S$ determina una biiezione $V \to S$, $v \mapsto P + v$ (biunivoca per (III)). Pertanto $S$ è, come insieme, in corrispondenza biunivoca con $V$; la differenza è che $S$ non ha un'origine privilegiata. Tutti i punti di $S$ sono equivalenti.
\end{osservazione}

\begin{defbox}
\begin{definizione}[Vettore $\overrightarrow{PQ}$]
\label{def:vett-pq}
Sia $S$ uno spazio affine con gruppo delle traslazioni $V$. Dati $P, Q \in S$, il \emph{vettore} $\overrightarrow{PQ} \in V$ è l'unico vettore tale che $P + \overrightarrow{PQ} = Q$ (esiste e unico per la proprietà (3) della Definizione~\ref{def:spazio-affine}).
\end{definizione}
\end{defbox}

\begin{proposizione}[Regola di Chasles]
\label{prop:chasles}
Sia $S$ uno spazio affine e $P, Q, R \in S$. Valgono:
\begin{align}
  \overrightarrow{PP} &= \mathbf{0}, \label{eq:PP}\\
  \overrightarrow{PQ} + \overrightarrow{QR} &= \overrightarrow{PR}, \label{eq:chasles}\\
  \overrightarrow{PQ} &= -\overrightarrow{QP}. \label{eq:inverso}
\end{align}
\end{proposizione}

\begin{proof}
\eqref{eq:PP}: poiché $P + \mathbf{0} = P$, per unicità $\overrightarrow{PP} = \mathbf{0}$.

\eqref{eq:chasles}: si ha
\[
  P + (\overrightarrow{PQ} + \overrightarrow{QR}) = (P + \overrightarrow{PQ}) + \overrightarrow{QR} = Q + \overrightarrow{QR} = R,
\]
dunque $\overrightarrow{PQ} + \overrightarrow{QR} = \overrightarrow{PR}$ per unicità.

\eqref{eq:inverso}: da \eqref{eq:PP} e \eqref{eq:chasles} si ricava $\mathbf{0} = \overrightarrow{PP} = \overrightarrow{PQ} + \overrightarrow{QP}$, quindi $\overrightarrow{PQ} = -\overrightarrow{QP}$.
\end{proof}

\begin{defbox}
\begin{definizione}[Dimensione di uno spazio affine]
La \emph{dimensione} di uno spazio affine $S$ è
\[
  \dim S \coloneqq \dim \VV(S).
\]
\end{definizione}
\end{defbox}

Una \emph{retta affine} è uno spazio affine di dimensione~$1$; un \emph{piano affine} ha dimensione~$2$; un \emph{solido affine} ha dimensione~$3$. Per $m \in \{1,2,3\}$, la dimensione di $E^m$ è $m$.

\begin{osservazione}
La definizione di spazio affine cattura esattamente le proprietà di $E^2$ ed $E^3$ che permettono di parlare di rette e piani. Una retta in $E^m$ è data dai punti $P_0 + sv_1$ con $P_0 \in E^m$, $v_1 \neq \mathbf{0}$ e $s \in \R$ arbitrario; un piano da $P_0 + sv_1 + tv_2$ con $v_1, v_2$ lin.\ ind. In uno spazio affine generale rimane vero l'assioma di Euclide: per due punti distinti passa una e una sola retta. E vale l'assioma delle parallele: per un punto esterno a una retta data passa un'unica retta parallela.
\end{osservazione}

% --------------------------------------------------------
\section{Combinazioni lineari di punti}
\label{sec:comb-lin-punti}
% --------------------------------------------------------

In uno spazio affine $S$ non ha senso scrivere $\lambda P + \mu Q$ per coefficienti $\lambda, \mu \in \K$ arbitrari. Se per esempio $S = \{Y \in \K^n \mid AY = B\}$ con $B \neq \mathbf{0}$, e $\lambda + \mu \neq 1$, allora $\lambda X + \mu Z \notin S$ in generale. Tuttavia, quando i coefficienti hanno somma~$1$, la costruzione diventa intrinseca e indipendente dal punto di riferimento.

\begin{lemma}
\label{lem:comb-affine-ben-def}
Sia $S$ uno spazio affine su $\K$ e $P_0, \ldots, P_d, Q, R \in S$. Siano $\lambda_0, \ldots, \lambda_d \in \K$ con $\sum_{i=0}^d \lambda_i = 1$. Allora
\[
  Q + \sum_{i=0}^d \lambda_i \,\overrightarrow{QP_i} = R + \sum_{i=0}^d \lambda_i \,\overrightarrow{RP_i}.
\]
\end{lemma}

\begin{proof}
Per la regola di Chasles~\eqref{eq:chasles}, $\overrightarrow{QP_i} = \overrightarrow{QR} + \overrightarrow{RP_i}$. Sostituendo:
\[
  Q + \sum_{i=0}^d \lambda_i \,\overrightarrow{QP_i} = Q + \Bigl(\sum_{i=0}^d \lambda_i\Bigr)\overrightarrow{QR} + \sum_{i=0}^d \lambda_i \,\overrightarrow{RP_i} = (Q + \overrightarrow{QR}) + \sum_{i=0}^d \lambda_i \,\overrightarrow{RP_i} = R + \sum_{i=0}^d \lambda_i \,\overrightarrow{RP_i},
\]
dove si è usato che $\sum \lambda_i = 1$.
\end{proof}

\begin{defbox}
\begin{definizione}[Combinazione affine di punti]
\label{def:comb-affine}
Sia $S$ uno spazio affine su $\K$ e $P_0, \ldots, P_d \in S$. Siano $\lambda_0, \ldots, \lambda_d \in \K$ con $\sum_{i=0}^d \lambda_i = 1$. La \emph{combinazione affine} di $P_0, \ldots, P_d$ con pesi $\lambda_0, \ldots, \lambda_d$ è
\[
  \sum_{i=0}^d \lambda_i P_i \;\coloneqq\; Q + \sum_{i=0}^d \lambda_i \,\overrightarrow{QP_i},
\]
dove $Q \in S$ è un punto qualsiasi. Il risultato è indipendente da $Q$ per il Lemma~\ref{lem:comb-affine-ben-def}.
\end{definizione}
\end{defbox}

\begin{exbox}
\begin{esempio}[Il caso degli spazi affini di sistemi lineari]
Sia $S \subset \K^n$ lo spazio affine dell'Esempio~\ref{es:spazio-affine-sistemi}. Se $X, Y \in S$ e $\lambda + \mu = 1$, la combinazione affine $\lambda X + \mu Y$ coincide con la combinazione lineare di vettori $\lambda X + \mu Y \in \K^n$, e appartiene a $S$.
\end{esempio}
\end{exbox}

\begin{exbox}
\begin{esempio}[Punti su una retta]
Nello spazio affine $E^m$, siano $P, Q \in E^m$ distinti. Al variare di $\lambda \in \R$, le combinazioni affini $\lambda P + (1-\lambda)Q$ descrivono esattamente la retta per $P$ e $Q$: per $\lambda = 0$ si ottiene $Q$, per $\lambda = 1$ si ottiene $P$, per $\lambda = 1/2$ il punto medio. Se $P = Q$ tutte le combinazioni affine danno $P$.
\end{esempio}
\end{exbox}

% --------------------------------------------------------
\section{Sottospazi affini}
\label{sec:sottospazi-affini}
% --------------------------------------------------------

\begin{defbox}
\begin{definizione}[Sottospazio affine]
\label{def:sottospazio-affine}
Sia $S$ uno spazio affine. Un sottoinsieme $T \subseteq S$ è un \emph{sottospazio affine} se esistono $P_0 \in S$ e un sottospazio vettoriale $W \subseteq \VV(S)$ tali che
\[
  T = P_0 + W \coloneqq \{ P_0 + w \mid w \in W \}.
\]
Il sottospazio vettoriale $W = \VV(T)$ si chiama \emph{giacitura} di $T$.
\end{definizione}
\end{defbox}

\begin{exbox}
\begin{esempio}
In uno spazio vettoriale $V$ visto come spazio affine (Esempio~\ref{es:V-come-affine}), i sottospazi affini sono le traslate $v_0 + W$ dei sottospazi vettoriali $W \subseteq V$. Un sottospazio vettoriale è un caso speciale di sottospazio affine (con $v_0 \in W$). Per esempio ogni punto $v_0 \in V$ è un sottospazio affine di dimensione~$0$, ma è un sottospazio vettoriale solo se $v_0 = \mathbf{0}$.
\end{esempio}
\end{exbox}

La proposizione seguente mostra che la giacitura è ben definita e non dipende dal punto di base $P_0$.

\begin{proposizione}
\label{prop:giacitura}
Siano $S$ uno spazio affine e $T = P_0 + W \subseteq S$ un sottospazio affine. Per ogni $P \in T$:
\[
  \{ \overrightarrow{PQ} \mid Q \in T \} = W.
\]
In particolare, l'applicazione $W \times T \to T$, $(v, P) \mapsto P + v$, dà a $T$ una struttura di spazio affine con gruppo delle traslazioni $W$.
\end{proposizione}

\begin{proof}
Sia $P = P_0 + v$ con $v \in W$. Se $Q = P_0 + u \in T$ (con $u \in W$), allora $\overrightarrow{PQ} = u - v \in W$. Viceversa, se $w \in W$, ponendo $u \coloneqq v + w \in W$ e $Q \coloneqq P_0 + u$, si ottiene $Q \in T$ e $\overrightarrow{PQ} = w$.
\end{proof}

\subsection{Parallelismo}

\begin{defbox}
\begin{definizione}[Sottospazi paralleli]
Sia $S$ uno spazio affine. Due sottospazi affini $T_1, T_2 \subseteq S$ sono \emph{paralleli} se $\VV(T_1) \subseteq \VV(T_2)$ oppure $\VV(T_1) \supseteq \VV(T_2)$.
\end{definizione}
\end{defbox}

\begin{osservazione}
Se $T_1, T_2$ sono paralleli e $T_1 \cap T_2 \neq \emptyset$, allora $T_1 \subseteq T_2$ o $T_1 \supseteq T_2$. Infatti, se $P \in T_1 \cap T_2$ e $\VV(T_1) \subseteq \VV(T_2)$, allora $T_1 = P + \VV(T_1) \subseteq P + \VV(T_2) = T_2$.

Attenzione: in $\mathbb{A}^3(\K)$ due rette senza punti in comune non sono necessariamente parallele (rette \emph{sghembe}). Per esempio, le rette $T_1 = \{(x,0,0)\}$ e $T_2 = \{(x,0,1+x)\}$ non si incontrano, ma $\VV(T_1) = \langle(1,0,0)\rangle$ e $\VV(T_2) = \langle(1,0,1)\rangle$ sono distinti e nessuno contiene l'altro. Invece, in un \emph{piano} affine, due rette senza punti in comune sono sempre parallele (si dimostra usando che $\VV(S)$ è bidimensionale).
\end{osservazione}

\subsection{Intersezione e sottospazio generato}

\begin{proposizione}
\label{prop:intersezione-affine}
Sia $S$ uno spazio affine e $\{T_i\}_{i \in I}$ una collezione di sottospazi affini di $S$. Se $\bigcap_{i \in I} T_i \neq \emptyset$, allora è un sottospazio affine di $S$ con giacitura $\bigcap_{i \in I} \VV(T_i)$.
\end{proposizione}

\begin{proof}
Sia $P_0 \in \bigcap_{i \in I} T_i$. Allora $T_i = P_0 + \VV(T_i)$ per ogni $i$. Un punto $Q \in S$ appartiene a $\bigcap_i T_i$ se e solo se $\overrightarrow{P_0Q} \in \VV(T_i)$ per ogni $i$, ovvero $\overrightarrow{P_0Q} \in \bigcap_i \VV(T_i)$. Poiché quest'ultima è un'intersezione di sottospazi vettoriali, è un sottospazio vettoriale, e $\bigcap_i T_i = P_0 + \bigcap_i \VV(T_i)$.
\end{proof}

\begin{defbox}
\begin{definizione}[Sottospazio affine generato]
\label{def:span-affine}
Sia $S$ uno spazio affine e $Z \subseteq S$ un sottoinsieme non vuoto. Il \emph{sottospazio affine generato da $Z$} è
\[
  \langle Z \rangle \coloneqq \bigcap \{ T \subseteq S \mid T \text{ sottospazio affine},\; Z \subseteq T \}.
\]
Per $Z = \{P_0, \ldots, P_d\}$ finito, si scrive $\langle P_0, \ldots, P_d \rangle$.
\end{definizione}
\end{defbox}

È ben definito (l'intersezione è non vuota perché $S$ stesso contiene $Z$) ed è il più piccolo sottospazio affine contenente $Z$.

\begin{proposizione}
\label{prop:generato-finito}
Siano $S$ uno spazio affine e $P_0, \ldots, P_d \in S$. Allora
\[
  \langle P_0, \ldots, P_d \rangle = P_0 + \langle \overrightarrow{P_0P_1}, \ldots, \overrightarrow{P_0P_d} \rangle
\]
e inoltre
\[
  \langle P_0, \ldots, P_d \rangle = \Bigl\{ \sum_{i=0}^d \lambda_i P_i \;\Big|\; \lambda_i \in \K,\; \sum_{i=0}^d \lambda_i = 1 \Bigr\}.
\]
In particolare, $\dim \langle P_0, \ldots, P_d \rangle \leq d$.
\end{proposizione}

\begin{proof}
Il membro $P_0 + \langle \overrightarrow{P_0P_1}, \ldots, \overrightarrow{P_0P_d}\rangle$ è un sottospazio affine contenente $P_0, \ldots, P_d$. Ogni sottospazio affine $T$ contenente $P_0, \ldots, P_d$ soddisfa $\overrightarrow{P_0P_i} \in \VV(T)$ (per la Proposizione~\ref{prop:giacitura}), dunque $T \supseteq P_0 + \langle \overrightarrow{P_0P_1},\ldots,\overrightarrow{P_0P_d}\rangle$. Questo dà la prima uguaglianza. La seconda si ricava osservando che $P_0 + \sum_{i=1}^d \mu_i \overrightarrow{P_0P_i} = (1 - \sum_{i=1}^d \mu_i)P_0 + \sum_{i=1}^d \mu_i P_i$.
\end{proof}

\begin{defbox}
\begin{definizione}[Indipendenza affine]
\label{def:ind-affine}
Sia $S$ uno spazio affine. Una lista di punti $P_0, \ldots, P_d \in S$ è \emph{affinemente indipendente} se $\dim \langle P_0, \ldots, P_d \rangle = d$, e \emph{affinemente dipendente} se $\dim \langle P_0, \ldots, P_d \rangle < d$.
\end{definizione}
\end{defbox}

\begin{lemma}
\label{lem:ind-affine-car}
La lista $P_0, \ldots, P_d$ è affinemente indipendente se e solo se l'applicazione
\[
  \Bigl\{ (\lambda_0,\ldots,\lambda_d) \in \K^{d+1} \;\Big|\; \sum_{i=0}^d \lambda_i = 1 \Bigr\} \longrightarrow \langle P_0,\ldots,P_d \rangle, \quad (\lambda_i) \mapsto \sum_{i=0}^d \lambda_i P_i
\]
è iniettiva.
\end{lemma}

\begin{exbox}
\begin{esempio}
In uno spazio vettoriale $V$, i punti $v_0, \ldots, v_d \in V$ sono affinemente indipendenti se e solo se i vettori $v_1 - v_0, \ldots, v_d - v_0$ sono linearmente indipendenti in $V$. Per esempio, in $\R^2$ tre punti non collineari sono affinemente indipendenti.

Si noti che l'indipendenza lineare (vettoriale) implica quella affine, ma non vale il viceversa: $\mathbf{0}$ e $v \neq \mathbf{0}$ sono affinemente indipendenti (determinano una retta), ma vettorialmente dipendenti.
\end{esempio}
\end{exbox}

% --------------------------------------------------------
\section{Applicazioni affini}
\label{sec:appl-affini}
% --------------------------------------------------------

\subsection{Definizione e prime proprietà}

\begin{defbox}
\begin{definizione}[Applicazione affine]
\label{def:appl-affine}
Un'applicazione $F \colon S \to T$ tra spazi affini sullo stesso campo $\K$ è \emph{affine} se esiste un'applicazione lineare $f \colon \VV(S) \to \VV(T)$ tale che
\[
  \overrightarrow{F(P)F(Q)} = f\bigl(\overrightarrow{PQ}\bigr) \qquad \forall\, P, Q \in S.
\]
L'applicazione lineare $f$, univocamente determinata da $F$, si chiama \emph{parte lineare} di $F$ e si denota $\VV(F)$.
\end{definizione}
\end{defbox}

In altre parole, $F$ manda segmenti orientati equipollenti in segmenti orientati equipollenti, e l'applicazione indotta sui vettori è lineare.

\begin{exbox}
\begin{esempio}[Traslazioni]
La traslazione $\tau_v \colon S \to S$ è affine con $\VV(\tau_v) = \Id_{\VV(S)}$. Viceversa, ogni applicazione affine $F \colon S \to S$ con $\VV(F) = \Id$ è una traslazione di vettore $\overrightarrow{P_0F(P_0)}$ per qualsiasi $P_0 \in S$.
\end{esempio}
\end{exbox}

\begin{exbox}
\begin{esempio}[Applicazioni affini tra spazi standard]
\label{es:appl-affine-matriciale}
Siano $A \in M_{m,n}(\K)$ e $B \in M_{m,1}(\K)$. L'applicazione
\[
  F \colon \mathbb{A}^n(\K) \to \mathbb{A}^m(\K), \quad X \mapsto A \cdot X + B,
\]
è affine con $\VV(F) = L_A$ (la moltiplicazione per $A$). Viceversa, ogni applicazione affine $\mathbb{A}^n(\K) \to \mathbb{A}^m(\K)$ è di questa forma: esistono un unico $A \in M_{m,n}(\K)$ e un unico $B \in M_{m,1}(\K)$ tali che $F(X) = A \cdot X + B$.
\end{esempio}
\end{exbox}

\begin{exbox}
\begin{esempio}[Proiezione parallela]
Siano $T \subset S$ un sottospazio affine e $D \subset \VV(S)$ un sottospazio vettoriale tale che $\VV(S) = D \oplus \VV(T)$. Per ogni $P \in S$, il sottospazio $L_P \coloneqq P + D$ interseca $T$ in esattamente un punto. Si definisce così la \emph{proiezione su $T$ parallela a $D$}:
\[
  \pi \colon S \to T, \quad P \mapsto L_P \cap T.
\]
Questa è un'applicazione affine.
\end{esempio}
\end{exbox}

\begin{osservazione}
\label{oss:F-determinata}
Un'applicazione affine $F \colon S \to T$ è completamente determinata dalla sua parte lineare $\VV(F)$ e dall'immagine di un punto $P_0 \in S$. Infatti, per ogni $P \in S$:
\[
  F(P) = F(P_0) + \VV(F)\bigl(\overrightarrow{P_0P}\bigr).
\]
\end{osservazione}

\begin{proposizione}
\label{prop:costruzione-affine}
Siano $S, T$ spazi affini sullo stesso campo $\K$, $P_0 \in S$, $Q_0 \in T$ e $f \colon \VV(S) \to \VV(T)$ lineare. L'applicazione
\[
  F \colon S \to T, \quad F(P) \coloneqq Q_0 + f\bigl(\overrightarrow{P_0P}\bigr),
\]
è l'unica applicazione affine tale che $F(P_0) = Q_0$ e $\VV(F) = f$.
\end{proposizione}

\begin{proof}
Per $P_1, P_2 \in S$, con $v_i = \overrightarrow{P_0P_i}$:
\[
  \overrightarrow{F(P_1)F(P_2)} = (Q_0 + f(v_2)) - (Q_0 + f(v_1)) = f(v_2 - v_1) = f\bigl(\overrightarrow{P_1P_2}\bigr).
\]
Dunque $F$ è affine con $\VV(F) = f$. L'unicità segue dall'Osservazione~\ref{oss:F-determinata}.
\end{proof}

\begin{proposizione}
\label{prop:immagini-affini}
Siano $S, T$ spazi affini sullo stesso campo $\K$ e $F \colon S \to T$ affine.
\begin{enumerate}[label=(\alph*)]
  \item Se $D \subseteq S$ è un sottospazio affine, allora $F(D)$ è un sottospazio affine di $T$ con giacitura $\VV(F)(\VV(D))$.
  \item Se $E \subseteq T$ è un sottospazio affine, allora $F^{-1}(E)$ è vuoto oppure è un sottospazio affine di $S$ con giacitura $\VV(F)^{-1}(\VV(E))$.
  \item Se $\sum_{i=0}^d \lambda_i = 1$, allora $F\bigl(\sum_{i=0}^d \lambda_i P_i\bigr) = \sum_{i=0}^d \lambda_i F(P_i)$.
\end{enumerate}
\end{proposizione}

\begin{proof}
(a): Sia $P_0 \in D$. Per la definizione di applicazione affine, $F(D) = F(P_0) + \VV(F)(\VV(D))$.

(b): Se $P_0 \in F^{-1}(E)$, allora $F^{-1}(E) = P_0 + \VV(F)^{-1}(\VV(E))$.

(c): Sia $Q \in S$ arbitrario. Usando la definizione~\ref{def:comb-affine}:
\[
  F\Bigl(\sum_{i=0}^d \lambda_i P_i\Bigr) = F\Bigl(Q + \sum_{i=0}^d \lambda_i \overrightarrow{QP_i}\Bigr) = F(Q) + \sum_{i=0}^d \lambda_i \VV(F)(\overrightarrow{QP_i}) = \sum_{i=0}^d \lambda_i F(P_i). \qedhere
\]
\end{proof}

\subsection{Composizione e isomorfismi affini}

\begin{proposizione}
Siano $S, T, U$ spazi affini sullo stesso campo e $F \colon S \to T$, $G \colon T \to U$ affini. Allora $G \circ F$ è affine con $\VV(G \circ F) = \VV(G) \circ \VV(F)$.
\end{proposizione}

\begin{proof}
Per ogni $P, Q \in S$: $\overrightarrow{G(F(P))\,G(F(Q))} = \VV(G)(\overrightarrow{F(P)F(Q)}) = \VV(G)(\VV(F)(\overrightarrow{PQ}))$.
\end{proof}

\begin{defbox}
\begin{definizione}[Isomorfismo e automorfismo affine]
Un'applicazione affine $F \colon S \to T$ è un \emph{isomorfismo affine} se ammette un'inversa affine. Un isomorfismo $F \colon S \to S$ si chiama \emph{automorfismo} o \emph{affinità} di $S$. Il gruppo delle affinità di $S$ è
\[
  \Aut(S) \coloneqq \{ F \colon S \to S \mid F \text{ è un'affinità} \}.
\]
\end{definizione}
\end{defbox}

\begin{proposizione}
\label{prop:iso-affine}
Un'applicazione affine $F \colon S \to T$ è un isomorfismo se e solo se $\VV(F)$ è un isomorfismo di spazi vettoriali.
\end{proposizione}

\begin{proof}
Se $F$ è un isomorfismo con inversa affine $G$, allora $\VV(G) \circ \VV(F) = \VV(G \circ F) = \VV(\Id_S) = \Id_{\VV(S)}$ e analogamente $\VV(F) \circ \VV(G) = \Id_{\VV(T)}$, dunque $\VV(F)$ è invertibile. Viceversa, se $\VV(F)$ è un isomorfismo, scelti $P_0 \in S$ e $Q_0 = F(P_0)$, l'applicazione $G(Q) \coloneqq P_0 + \VV(F)^{-1}(\overrightarrow{Q_0Q})$ è affine con $G \circ F = \Id_S$ e $F \circ G = \Id_T$.
\end{proof}

\begin{exbox}
\begin{esempio}
Per lo spazio affine standard $\mathbb{A}^n(\K)$:
\[
  \Aut(\mathbb{A}^n(\K)) = \{ F \colon X \mapsto A \cdot X + B \mid A \in \GL_n(\K),\; B \in \K^n \}.
\]
L'applicazione $\Aut(\mathbb{A}^n(\K)) \to \GL_n(\K)$, $F \mapsto A$, è un omomorfismo suriettivo di gruppi. Il suo nucleo è il gruppo delle traslazioni $\{\tau_B \mid B \in \K^n\}$.
\end{esempio}
\end{exbox}

\subsection{Funzioni affini ed equazioni cartesiane}

\begin{defbox}
\begin{definizione}
Una \emph{funzione affine} su uno spazio affine $S$ è un'applicazione affine $F \colon S \to \mathbb{A}^1(\K)$.
\end{definizione}
\end{defbox}

Esplicitamente, $F$ è affine se esiste $f \colon \VV(S) \to \K$ lineare tale che $F(P) = F(P_0) + f(\overrightarrow{P_0P})$ per ogni $P \in S$.

\begin{proposizione}
\label{prop:eq-cartesiane}
Sia $S$ uno spazio affine su $\K$ di dimensione finita $n$ e $T \subseteq S$ un sottospazio affine. Posto $m \coloneqq n - \dim T$, esistono $m$ funzioni affini $F_1, \ldots, F_m \colon S \to \K$ tali che
\[
  T = \{ P \in S \mid F_1(P) = \cdots = F_m(P) = 0 \}.
\]
\end{proposizione}

\begin{proof}
Si ha $T = P_0 + W$ con $m = \dim \VV(S) - \dim W$. Per l'Esercizio relativo all'annullatore (Esercizio~\ref{es:ann-eq-cart} del Capitolo~3), esistono $g_1, \ldots, g_m \in \Ann(W)$ tali che $W = \{v \mid g_1(v) = \cdots = g_m(v) = 0\}$. Si pone $F_i(P) \coloneqq g_i(\overrightarrow{P_0P})$: ogni $F_i$ è affine con $\VV(F_i) = g_i$, e $F_i(P) = 0 \Leftrightarrow \overrightarrow{P_0P} \in W \Leftrightarrow P \in T$.
\end{proof}

\begin{exbox}
\begin{esempio}[Equazioni cartesiane in $\mathbb{A}^4(\Q)$]
Sia $T = P_0 + U \subset \mathbb{A}^4(\Q)$ con $P_0 = (-2,-1,1,2)$ e
\[
  U = \langle (1,-2,3,-4),\, (2,0,3,1) \rangle.
\]
Una funzione lineare $g(x_1,x_2,x_3,x_4) = \lambda_1 x_1 + \cdots + \lambda_4 x_4$ appartiene a $\Ann(U)$ se e solo se
\[
  \lambda_1 - 2\lambda_2 + 3\lambda_3 - 4\lambda_4 = 0 \quad\text{e}\quad 2\lambda_1 + 3\lambda_3 + \lambda_4 = 0.
\]
Risolvendo: $\Ann(U) = \langle 2e_1^\vee + 9e_2^\vee - 4e_4^\vee,\; 6e_1^\vee - 3e_2^\vee - 4e_3^\vee \rangle$. Le equazioni cartesiane di $T$ sono:
\begin{align*}
  F_1(X) &= 2(x_1+2) + 9(x_2+1) - 4(x_4-2) = 2x_1 + 9x_2 - 4x_4 + 21 = 0,\\
  F_2(X) &= 6(x_1+2) - 3(x_2+1) - 4(x_3-1) = 6x_1 - 3x_2 - 4x_3 + 13 = 0.
\end{align*}
\end{esempio}
\end{exbox}

\subsection{Coordinate affini}

\begin{defbox}
\begin{definizione}[Sistema di riferimento affine]
\label{def:rif-affine}
Un \emph{sistema di riferimento affine} $\mathcal{R} = \RA(O;\mathcal{B})$ su uno spazio affine $S$ di dimensione $n$ consiste di un punto $O \in S$ (\emph{origine}) e di una base $\mathcal{B} = \{v_1, \ldots, v_n\}$ di $\VV(S)$.
\end{definizione}
\end{defbox}

Dato $\mathcal{R} = \RA(O;\mathcal{B})$, l'applicazione
\[
  X_\mathcal{R} \colon S \to \mathbb{A}^n(\K), \quad P \mapsto X_\mathcal{B}(\overrightarrow{OP})
\]
è un isomorfismo affine (per le Proposizioni~\ref{prop:costruzione-affine} e~\ref{prop:iso-affine}). Le \emph{coordinate affini} di $P$ nel riferimento $\mathcal{R}$ sono le entrate di $X_\mathcal{R}(P)$. Scriviamo $P(a_1, \ldots, a_n)$ per indicare che $P = O + \sum_{i=1}^n a_i v_i$.

\begin{exbox}
\begin{esempio}[Calcolo di coordinate in un sottospazio affine]
Sia $S \subset \mathbb{A}^3(\R)$ lo spazio affine di equazione $x_1 + x_2 + x_3 = 3$, con gruppo delle traslazioni $V = \ker(x_1+x_2+x_3)$. Sia $O = (1,1,1)$ e $\mathcal{B} = \{(1,-1,0),(1,0,-1)\}$ base di $V$. Per il punto $P = (5,-2,0)$:
\[
  \overrightarrow{OP} = (4,-3,-1) = \lambda(1,-1,0) + \mu(1,0,-1) \implies \lambda = 3,\; \mu = 1.
\]
Dunque $P$ ha coordinate $(3,1)$ nel riferimento $\RA(O;\mathcal{B})$.
\end{esempio}
\end{exbox}

\begin{osservazione}[Cambio di riferimento affine]
Siano $\mathcal{R} = \RA(O;\mathcal{B})$ e $\mathcal{R}' = \RA(O';\mathcal{B}')$ due riferimenti affini su $S$. La composizione $X_{\mathcal{R}'} \circ X_\mathcal{R}^{-1}$ è un isomorfismo affine di $\mathbb{A}^n(\K)$, dunque esistono $A \in \GL_n(\K)$ e $B \in \K^n$ tali che
\[
  X_{\mathcal{R}'}(P) = A \cdot X_\mathcal{R}(P) + B \qquad \forall\, P \in S.
\]
La matrice $A$ è la matrice del cambio di base da $\mathcal{B}$ a $\mathcal{B}'$; il vettore $B$ tiene conto dello spostamento dell'origine.
\end{osservazione}

% --------------------------------------------------------
\section{Ginnastica affine}
\label{sec:ginnastica-affine}
% --------------------------------------------------------

In tutta questa sezione $S$ è uno spazio affine di dimensione finita $n$, con riferimento affine fisso $\mathcal{R} = \RA(O;\mathcal{B})$, $\mathcal{B} = \{v_1, \ldots, v_n\}$. Identifichiamo $S \cong \mathbb{A}^n(\K)$ tramite $X_\mathcal{R}$.

\subsection{Equazioni parametriche di sottospazi affini}

Sia $T = P_0 + W \subseteq S$ un sottospazio affine. Siano $P_0(a_1,\ldots,a_n)$ e $\{w_1, \ldots, w_r\}$ una base di $W$, con $w_i$ di coordinate $(c_{i1},\ldots,c_{in})$ in $\mathcal{B}$. Le \emph{equazioni parametriche} di $T$ sono:
\[
  x_j = a_j + \sum_{i=1}^r t_i\, c_{ij}, \quad j = 1,\ldots,n, \quad t_1,\ldots,t_r \in \K.
\]

\begin{exbox}
\begin{esempio}[Retta parametrica in $\mathbb{A}^n$]
\label{es:retta-param}
Se $T$ è una retta ($r = 1$, base $\{w\} = \{(l_1,\ldots,l_n)\}$), le equazioni parametriche sono
\[
  x_j = a_j + l_j t, \quad j = 1,\ldots,n, \quad t \in \K,
\]
con $(l_1,\ldots,l_n) \neq \mathbf{0}$. Gli scalari $l_1,\ldots,l_n$ si chiamano \emph{parametri direttori} della retta.
\end{esempio}
\end{exbox}

\begin{exbox}
\begin{esempio}[Retta per due punti]
Siano $P_0(a_1,\ldots,a_n)$ e $P_1(b_1,\ldots,b_n)$ distinti. La retta $\langle P_0, P_1 \rangle$ ha equazioni parametriche
\[
  x_j = a_j + (b_j - a_j)\,t, \quad j = 1,\ldots,n, \quad t \in \K,
\]
poiché $\overrightarrow{P_0P_1}$ ha coordinate $(b_1 - a_1, \ldots, b_n - a_n)$.
\end{esempio}
\end{exbox}

\begin{exbox}
\begin{esempio}[Piano per tre punti]
Siano $P_0(a_j)$, $P_1(b_j)$, $P_2(c_j)$ affinemente indipendenti. Il piano $\langle P_0,P_1,P_2 \rangle$ ha equazioni parametriche
\[
  x_j = a_j + (b_j - a_j)\,s + (c_j - a_j)\,t, \quad j = 1,\ldots,n, \quad s, t \in \K.
\]
\end{esempio}
\end{exbox}

\begin{exbox}
\begin{esempio}[Non unicità delle equazioni parametriche]
Le rette in $\mathbb{A}^2(\K)$:
\[
  L\colon \begin{cases} x_1 = 3 + 2t \\ x_2 = -1 + 3t \end{cases} \qquad M\colon \begin{cases} x_1 = 5 + 4t \\ x_2 = 2 + 6t \end{cases}
\]
hanno giacitura $\langle (2,3) \rangle = \langle (4,6) \rangle$, dunque sono parallele. Il punto $(5,2)$ appartiene a $L$ (con $t=1$) e a $M$ (con $t=0$), quindi $L = M$ pur avendo equazioni parametriche diverse.
\end{esempio}
\end{exbox}

\subsection{Equazioni cartesiane e passaggio da un tipo all'altro}

\begin{proposizione}[Unicità dell'equazione cartesiana per iperpiani]
Sia $T \subseteq S$ un iperpiano ($\dim T = n-1$). Se $a_1 x_1 + \cdots + a_n x_n + b = 0$ e $c_1 x_1 + \cdots + c_n x_n + d = 0$ sono entrambe equazioni cartesiane di $T$, allora esiste $\lambda \in \K^*$ tale che $a_i = \lambda c_i$ per ogni $i$ e $b = \lambda d$.
\end{proposizione}

\begin{proof}
Le applicazioni lineari $\phi = \sum a_i e_i^\vee$ e $\psi = \sum c_i e_i^\vee$ appartengono entrambe ad $\Ann(\VV(T))$, che ha dimensione~$1$ (poiché $\dim \VV(T) = n - 1$). Dunque $\phi = \lambda \psi$ per qualche $\lambda \neq 0$, da cui $a_i = \lambda c_i$. Sostituendo un punto $P \in T$ nell'equazione si ottiene $b = \lambda d$.
\end{proof}

\bigskip
\noindent\textbf{Da parametriche a cartesiane: caso di un piano in $\mathbb{A}^3(\K)$.}
Dato il piano di equazioni parametriche
\[
  x_j = a_j + l_j s + m_j t, \quad j = 1,2,3,
\]
un'equazione cartesiana $c_1 x_1 + c_2 x_2 + c_3 x_3 = b$ richiede che $(c_1,c_2,c_3)$ sia soluzione del sistema
\[
  l_1 y_1 + l_2 y_2 + l_3 y_3 = 0, \qquad m_1 y_1 + m_2 y_2 + m_3 y_3 = 0.
\]
Se il piano è proprio (cioè $l = (l_1,l_2,l_3)$ e $m = (m_1,m_2,m_3)$ sono lin.\ ind.), tale sistema ha spazio delle soluzioni di dimensione~$1$ generato da
\[
  (l_2 m_3 - l_3 m_2,\; -(l_1 m_3 - l_3 m_1),\; l_1 m_2 - l_2 m_1).
\]

\begin{exbox}
\begin{esempio}[Da parametriche a cartesiane in $\mathbb{A}^3$]
Sia il piano di equazioni parametriche $x_1 = 1 + s + 2t$, $x_2 = 2 + s$, $x_3 = 3 + s - t$. Vettori di direzione: $l = (1,1,1)$ e $m = (2,0,-1)$. Il vettore normale è $(1\cdot(-1) - 1\cdot 0,\, -(1\cdot(-1)-1\cdot 2),\, 1\cdot 0 - 1\cdot 2) = (-1, 3, -2)$. Poiché $P_0 = (1,2,3)$ appartiene al piano, l'equazione cartesiana è
\[
  -(x_1-1) + 3(x_2-2) - 2(x_3-3) = 0 \quad\Longleftrightarrow\quad -x_1 + 3x_2 - 2x_3 + 1 = 0.
\]
\end{esempio}
\end{exbox}

\begin{exbox}
\begin{esempio}[Da cartesiane a parametriche in $\mathbb{A}^3(\R)$]
\label{es:cart-param-a3}
Sia $L \subset \mathbb{A}^3(\R)$ la retta di equazioni cartesiane
\[
  \begin{cases} x_1 + 3x_2 - x_3 = 1 \\ 2x_1 + 7x_2 + x_3 = 5. \end{cases}
\]
Con eliminazione di Gauss: si ottiene $x_1 = -8 + 10x_3$ e $x_2 = 3 - 3x_3$. Ponendo $x_3 = t$:
\[
  \begin{cases} x_1 = -8 + 10t \\ x_2 = 3 - 3t \\ x_3 = t. \end{cases}
\]
\end{esempio}
\end{exbox}

\subsection{Fasci di rette e piani}

\begin{defbox}
\begin{definizione}[Fascio proprio di rette]
Sia $S$ un piano affine e $P_0 \in S$. Il \emph{fascio proprio di rette di centro $P_0$} è l'insieme delle rette di $S$ contenenti $P_0$.
\end{definizione}
\end{defbox}

\begin{proposizione}
\label{prop:fascio-proprio}
Sia $S$ un piano affine e $P_0 \in S$ di equazioni cartesiane
\[
  \{a_1 x_1 + a_2 x_2 + b = 0\} \cap \{c_1 x_1 + c_2 x_2 + d = 0\}.
\]
Una retta $L \subseteq S$ contiene $P_0$ se e solo se esistono $\lambda, \mu \in \K$ non entrambi nulli tali che $L$ abbia equazione
\[
  \lambda(a_1 x_1 + a_2 x_2 + b) + \mu(c_1 x_1 + c_2 x_2 + d) = 0.
\]
\end{proposizione}

\begin{proof}
Se $L$ ha equazione di tale forma, $P_0 \in L$ perché le sue coordinate soddisfano il sistema. Viceversa, sia $L$ una retta contenente $P_0$ e sia $Q \in L \setminus \{P_0\}$. Esiste $(\lambda_0, \mu_0) \neq (0,0)$ tale che $\lambda_0(a_1 x_1(Q) + a_2 x_2(Q) + b) + \mu_0(c_1 x_1(Q) + c_2 x_2(Q) + d) = 0$. La retta di equazione $\lambda_0(\cdots) + \mu_0(\cdots) = 0$ contiene $P_0$ e $Q$, dunque è uguale a $L$.
\end{proof}

\begin{defbox}
\begin{definizione}[Fascio improprio di rette]
Sia $S$ un piano affine e $W \subseteq \VV(S)$ un sottospazio di dimensione~$1$. Il \emph{fascio improprio di rette di direzione $W$} è l'insieme delle rette di $S$ con giacitura $W$.
\end{definizione}
\end{defbox}

\begin{proposizione}
Una retta $L \subseteq S$ appartiene al fascio improprio di direzione $W$ se e solo se, detto $a_1 x_1 + a_2 x_2 = 0$ un'equazione cartesiana di $W \subseteq \VV(S)$, esiste $b \in \K$ tale che $L$ abbia equazione $a_1 x_1 + a_2 x_2 + b = 0$.
\end{proposizione}

\begin{proof}
Se $L$ ha equazione del tipo indicato, la sua giacitura è $W$ per la proprietà~\eqref{eq:chasles} della giacitura. Viceversa, se $VpLq = W$ e $P(x_1, x_2) \in L$, la retta di equazione $a_1(x - x_1) + a_2(x_2 - x_2) = 0$ ha giacitura $W$ e contiene $P$, dunque coincide con $L$.
\end{proof}

Definizioni e proposizioni analoghe si hanno per i \emph{fasci di piani} in un solido affine: il fascio proprio di centro $L$ (piani contenenti la retta $L$) e il fascio improprio di direzione $W$ (piani paralleli tra loro, con giacitura il piano vettoriale $W$).

\begin{exbox}
\begin{esempio}[Fascio di piani in $\mathbb{A}^3(\R)$]
Sia $L \subset \mathbb{A}^3(\R)$ la retta di equazioni cartesiane
\[
  \begin{cases} 3x_1 + 2x_2 + x_3 - 6 = 0 \\ x_1 + 2x_2 + 3x_3 - 1 = 0. \end{cases}
\]
Sia $J$ la retta di equazioni parametriche $x_1 = 1+2t$, $x_2 = -2+t$, $x_3 = 3t$, con direzione $(2,1,3)$. Il vettore $(2,1,3)$ non soddisfa il sistema omogeneo associato a $L$, dunque $J$ non è parallela a $L$: esiste un unico piano del fascio di asse $L$ parallelo a $J$.

Per la Proposizione~\ref{prop:fascio-proprio} adattata ai piani, cercando $\lambda(3x_1 + 2x_2 + x_3 - 6) + \mu(x_1 + 2x_2 + 3x_3 - 1) = 0$ con $\VV(J) \subset \VV(T)$:
\[
  (3\lambda + \mu) \cdot 2 + (2\lambda + 2\mu) \cdot 1 + (\lambda + 3\mu) \cdot 3 = 11\lambda + 13\mu = 0.
\]
Prendendo $\lambda = 13$, $\mu = -11$, l'equazione cartesiana del piano cercato è
\[
  28x_1 + 4x_2 - 20x_3 - 67 = 0.
\]
\end{esempio}
\end{exbox}

% --------------------------------------------------------
\section*{Esercizi}
\addcontentsline{toc}{section}{Esercizi del Capitolo 4}
\label{sec:es-cap4}
% --------------------------------------------------------

\begin{esercizio}
Sia $S$ un piano affine con riferimento $\RA(O;\mathcal{B})$. Siano $P_0(1,2)$ e $P_1(-1,1)$ punti di $S$. Scrivere equazioni parametriche e cartesiane della retta $\langle P_0, P_1 \rangle$.
\end{esercizio}

\begin{esercizio}
Sia $S$ un piano affine reale con riferimento $\RA(O;\mathcal{B})$, e $L, J \subseteq S$ le rette di equazioni parametriche
\[
  L\colon \begin{cases} x_1 = 1 + 3t \\ x_2 = -2 + t \end{cases} \qquad J\colon \begin{cases} x_1 = s \\ x_2 = 1 - s. \end{cases}
\]
Determinare le coordinate del punto d'intersezione tra $L$ e $J$.
\end{esercizio}

\begin{esercizio}[Guidato: cambio di riferimento]
Sia $\mathcal{B} \coloneqq \{\mathbf{i}, \mathbf{j}\}$ una base di $\VV(E^2)$ e $\mathbf{k} \coloneqq \mathbf{i} + 2\mathbf{j}$, $\mathbf{h} \coloneqq \mathbf{i} + \mathbf{j}$. Sia $Q \in E^2$ di coordinate $(1,-1)$ nel riferimento $\RA(O;\mathcal{B})$.
\begin{enumerate}
  \item Verificare che $\mathcal{C} \coloneqq \{\mathbf{k}, \mathbf{h}\}$ è una base di $\VV(E^2)$.
  \item Determinare le coordinate di $O$ nel riferimento $\RA(Q;\mathcal{C})$. [\emph{Suggerimento}: $\overrightarrow{QO} = -\overrightarrow{OQ}$; esprimere $\overrightarrow{OQ} = \mathbf{i} - \mathbf{j}$ in termini di $\mathbf{k}$ e $\mathbf{h}$.]
\end{enumerate}
\end{esercizio}

\begin{esercizio}
Sia $S$ un solido affine reale con riferimento $\RA(O;\mathcal{B})$. Siano $P_0(1,1,-1)$, $P_1(3,0,2)$, $P_2(4,2,3)$.
\begin{enumerate}
  \item Verificare che $P_0, P_1, P_2$ non sono collineari.
  \item Determinare equazioni parametriche e cartesiane del piano $\Lambda = \langle P_0, P_1, P_2 \rangle$.
\end{enumerate}
\end{esercizio}

\begin{esercizio}
Sia $S$ un solido affine reale e siano i piani $T_1 \colon 3x_1 + 2x_2 + x_3 = 1$ e $T_2 \colon x_1 - x_2 - 2x_3 = 2$. Verificare che $T_1 \cap T_2$ è una retta $L$ e determinarne equazioni parametriche.
\end{esercizio}

\begin{esercizio}
Sia $S$ un solido affine reale e siano le rette
\[
  L\colon \begin{cases} x_1 = 1+2t \\ x_2 = -t \\ x_3 = 2+5t \end{cases} \qquad J\colon \begin{cases} x_1 = s \\ x_2 = 1+2s \\ x_3 = 3. \end{cases}
\]
Determinare un'equazione cartesiana del piano $T$ contenente $L$ e parallelo a $J$.
\end{esercizio}

\begin{esercizio}
Siano $T_1 \colon x_1 + 2x_2 - x_3 + 1 = 0$ e $T_2 \colon 2x_1 + x_3 - 3 = 0$ due piani in un solido affine reale. Determinare equazioni parametriche della retta contenente $P(1,1,1)$ e parallela a entrambi $T_1$ e $T_2$.
\end{esercizio}

\begin{esercizio}[Guidato: chi ha sbagliato?]
Sia $S$ un piano affine reale con riferimento $\RA(O;\mathcal{B})$. Siano $P_0(1,1)$, $P_1(2,3)$, $P_2(3,5)$. Tre studenti misurano le coordinate di questi punti in un nuovo riferimento $\RA(Q;\mathcal{C})$:
\begin{itemize}
  \item[(A)] $P_0(0,1)$, $P_1(-1,0)$, $P_2(-2,-1)$.
  \item[(M)] $P_0(0,2)$, $P_1(1,2)$, $P_2(0,3)$.
  \item[(L)] $P_0(0,0)$, $P_1(1,1)$, $P_2(3,3)$.
\end{itemize}
Due tra loro hanno sicuramente sbagliato: determinare chi. [\emph{Suggerimento}: osservare che $P_2 = 2P_1 - P_0$ in $S$ (relazione affine con $\lambda_0 = -1$, $\lambda_1 = 2$, $\lambda_2 = 1$); le applicazioni affini preservano le combinazioni affini.]
\end{esercizio}

\begin{esercizio}[Baricentro]
Sia $S$ uno spazio affine su un campo di caratteristica zero. Siano $P_0,\ldots,P_d \in S$ affinemente indipendenti. Il \emph{baricentro} è
\[
  B(P_0,\ldots,P_d) \coloneqq \frac{1}{d+1}\sum_{i=0}^d P_i.
\]
Sia $r$ la retta $\langle P_d,\, B(P_0,\ldots,P_d) \rangle$. Dimostrare che $r \cap \langle P_0,\ldots,P_{d-1} \rangle = \{B(P_0,\ldots,P_{d-1})\}$.
\end{esercizio}

\begin{esercizio}
Stabilire se esiste un'applicazione affine $F \colon \mathbb{A}^2(\R) \to \mathbb{A}^2(\R)$ tale che:
\begin{enumerate}
  \item $F(1,1) = (1,2)$, $F(3,2) = (-1,-2)$, $F(2,3/2) = (0,1)$.
  \item $F(0,0) = (1,1)$, $F(1,1) = (2,1)$, $F(1,-1) = (1,2)$.
\end{enumerate}
Nel caso esista, determinare $F$ esplicitamente. [\emph{Suggerimento}: per (1) verificare se $(2,3/2)$ è combinazione affine di $(1,1)$ e $(3,2)$ in modo incompatibile con i valori assegnati; per (2) usare la Proposizione~\ref{prop:costruzione-affine}.]
\end{esercizio}

\begin{esercizio}
\label{es:ann-eq-cart}
Sia $S = P + U \subset \mathbb{A}^4(\R)$ con $P = (1,0,3,-1)$ e $U = \langle (1,1,1,1), (3,2,-1,5), (4,3,0,6) \rangle$. Determinare equazioni cartesiane di $S$. [\emph{Suggerimento}: verificare prima che $\dim U < 3$, poi calcolare $\Ann(U)$ applicando Gauss alla matrice trasposta di quella dei generatori.]
\end{esercizio}

\begin{esercizio}
Sia $F \colon S \to S$ un'affinità di uno spazio affine. Il \emph{luogo dei punti fissi} è $\mathrm{Fix}(F) \coloneqq \{P \in S \mid F(P) = P\}$.
\begin{enumerate}
  \item Dimostrare che se $\mathrm{Fix}(F) \neq \emptyset$, allora è un sottospazio affine di $S$ con giacitura $\ker(\VV(F) - \Id_{\VV(S)})$.
  \item Dare esempi con $\dim S = n$ e $\dim \mathrm{Fix}(F) = m$ per ogni $0 \leq m \leq n$. [\emph{Suggerimento}: per $m = 0$ usare una traslazione non nulla; per $m = n$ prendere $F = \Id_S$; per $0 < m < n$ usare applicazioni di $\mathbb{A}^n(\R)$ con matrice $A$ tale che $\ker(A - I_n)$ abbia dimensione $m$.]
\end{enumerate}
\end{esercizio}


% ============================================================
\chapter{Determinanti}
\label{cap:determinanti}
% ============================================================

\section*{Introduzione al capitolo}

Nei capitoli precedenti abbiamo visto che una matrice quadrata $A \in M_{n,n}(K)$ è invertibile se e solo se le sue colonne (o righe) sono linearmente indipendenti. Ma come si verifica concretamente questa condizione? Il \emph{determinante} è la risposta: è un'unica funzione scalare
\[
  \det \colon M_{n,n}(K) \longrightarrow K
\]
tale che $\det A \neq 0$ se e soltanto se $A$ è invertibile.

Il determinante ha però un significato molto più ricco di un semplice criterio di invertibilità. Su $\R$, esso misura il \emph{volume orientato} del parallelepipedo generato dalle colonne della matrice: il segno dice se l'orientazione è preservata o invertita, il valore assoluto è il volume. In Algebra lineare astratta, il determinante è l'unica forma multilineare alternante normalizzata — una caratterizzazione che porta a comprenderne le proprietà in modo profondo e conciso.

Il capitolo è organizzato come segue. Prima diamo la definizione ricorsiva (sviluppo di Laplace lungo l'ultima riga) e ne deriviamo le proprietà fondamentali per le colonne. Poi studiamo le applicazioni multilineari alternanti, che forniscono il quadro teorico corretto. Infine sviluppiamo le applicazioni: formula di Binet per il prodotto, sviluppo di Laplace per righe e colonne, formula di Cramer per l'inversa, polinomio caratteristico e diagonalizzazione.

% --------------------------------------------------------
\section{Definizione ricorsiva}
\label{sec:det-def}
% --------------------------------------------------------

\subsection{La definizione per induzione}

Sia $R$ un anello commutativo con unità e sia $A \in M_{n,n}(R)$. Per $1 \leq i,j \leq n$ denotiamo con $A^i_j$ la matrice $(n-1)\times(n-1)$ ottenuta eliminando la riga $i$-esima e la colonna $j$-esima di $A$.

\begin{defbox}
\begin{definizione}[Determinante]
\label{def:determinante}
La funzione $\det_n \colon M_{n,n}(R) \to R$ è definita ricorsivamente:
\begin{enumerate}
  \item $\det_1\bigl((a)\bigr) \coloneqq a$.
  \item Per $n \geq 2$, posto $\det_{n-1}$ già definito:
  \[
    \det_n(A) \coloneqq \sum_{j=1}^{n} (-1)^{n+j}\, a_{nj}\, \det_{n-1}(A^n_j).
  \]
\end{enumerate}
Questa formula si chiama \emph{sviluppo di Laplace lungo l'ultima riga}. Quando non vi è ambiguità scriviamo $\det$ in luogo di $\det_n$, e anche $|A|$ in luogo di $\det A$.
\end{definizione}
\end{defbox}

\begin{osservazione}
La costruzione è ben fondata per induzione: $\det_1$ è data dal punto (1), da essa si ricava $\det_2$ tramite il punto (2), poi $\det_3$, ecc.
\end{osservazione}

\subsection{Calcolo esplicito nei casi piccoli}

\begin{exbox}
\begin{esempio}[Caso $n=2$]
\[
  \det\begin{pmatrix} a_{11} & a_{12} \\ a_{21} & a_{22} \end{pmatrix}
  = (-1)^{2+1}a_{21}\det(a_{12}) + (-1)^{2+2}a_{22}\det(a_{11})
  = -a_{21}a_{12} + a_{22}a_{11}
  = a_{11}a_{22} - a_{12}a_{21}.
\]
È la regola di Sarrus per $2\times 2$: prodotto della diagonale principale meno prodotto della diagonale secondaria.
\end{esempio}
\end{exbox}

\begin{exbox}
\begin{esempio}[Caso $n=3$]
\label{es:det3}
Per $A \in M_{3,3}(R)$, espandendo lungo la terza riga:
\begin{align*}
  \det A &= a_{31}(a_{12}a_{23} - a_{13}a_{22}) - a_{32}(a_{11}a_{23} - a_{13}a_{21}) + a_{33}(a_{11}a_{22} - a_{12}a_{21}) \\
         &= a_{11}a_{22}a_{33} + a_{12}a_{23}a_{31} + a_{13}a_{21}a_{32}
            - a_{13}a_{22}a_{31} - a_{12}a_{21}a_{33} - a_{11}a_{23}a_{32}.
\end{align*}
Questa è la \emph{regola di Sarrus}: i tre prodotti con segno $+$ corrispondono alle diagonali che scendono da sinistra a destra, i tre con segno $-$ alle diagonali che scendono da destra a sinistra.
\end{esempio}
\end{exbox}

\begin{exbox}
\begin{esempio}
\[
  \left|\begin{matrix}1&2\\3&4\end{matrix}\right| = 1\cdot4 - 2\cdot3 = -2.
\]
\end{esempio}
\end{exbox}

\subsection{Prime proprietà per induzione}

La definizione ricorsiva si presta naturalmente a dimostrazioni per induzione. Ne diamo due come prototipo.

\begin{proposizione}
\label{prop:col-nulla}
Se $A \in M_{n,n}(R)$ ha una colonna nulla, allora $\det A = 0$.
\end{proposizione}

\begin{proof}
Per induzione su $n$. Se $n=1$, $A=(0)$ e $\det A = 0$. Per $n \geq 2$, supponiamo vera l'asserzione per $(n-1)\times(n-1)$.

Sia $A^n_j$ la colonna nulla di $A$. Nella formula di definizione
\[
  \det A = \sum_{j=1}^n (-1)^{n+j} a_{nj}\,\det(A^n_j),
\]
il termine $j=j_0$ si annulla poiché $a_{n,j_0}=0$. Per ogni $j \neq j_0$, la matrice $A^n_j$ ha ancora una colonna nulla (la colonna $j_0-1$ se $j < j_0$, o la colonna $j_0$ se $j > j_0$), e per ipotesi induttiva $\det(A^n_j)=0$. Quindi tutti i termini sono nulli.
\end{proof}

\begin{proposizione}
\label{prop:scala-colonne}
Se $B \in M_{n,n}(R)$ è a scala per colonne (cioè $b_{ij}=0$ per $i < j$), allora
\[
  \det B = b_{11} \cdot b_{22} \cdots b_{nn}.
\]
\end{proposizione}

\begin{proof}
Per induzione su $n$. Se $n=1$ è ovvio. Per $n \geq 2$, nella formula $\det B = \sum_{j=1}^n(-1)^{n+j}b_{nj}\det(B^n_j)$: per $j < n$, la matrice $B^n_j$ ha l'ultima colonna nulla (poiché $B$ è a scala e la colonna $n$ di $B$ ha tutte le entrate nulle tranne l'ultima), quindi per la Proposizione~\ref{prop:col-nulla} il suo determinante è $0$. Resta solo il termine $j=n$: $B^n_n$ è ancora a scala per colonne con entrate diagonali $b_{11},\ldots,b_{n-1,n-1}$, e per ipotesi induttiva $\det(B^n_n) = b_{11}\cdots b_{n-1,n-1}$. Pertanto $\det B = (-1)^{2n}b_{nn}\cdot b_{11}\cdots b_{n-1,n-1} = b_{11}\cdots b_{nn}$.
\end{proof}

% --------------------------------------------------------
\section{Proprietà algebriche}
\label{sec:det-proprieta}
% --------------------------------------------------------

D'ora in avanti $K$ è un campo. Vediamo il determinante come funzione delle $n$ colonne:
\[
  \det_n \colon K^n \times \cdots \times K^n \longrightarrow K,
  \quad (A_1,\ldots,A_n) \mapsto \det[A_1 \mid \cdots \mid A_n].
\]

\begin{proposizione}[Multilinearità nelle colonne]
\label{prop:multilineare}
La funzione $\det_n$ è lineare in ciascuna colonna: fissate tutte le colonne tranne la $j_0$-esima, l'applicazione $D \mapsto \det[A_1,\ldots,A_{j_0-1},D,A_{j_0+1},\ldots,A_n]$ è lineare.
\end{proposizione}

\begin{proof}
Per induzione su $n$. Il caso $n=1$ è banale. Per $n \geq 2$ e $j \neq j_0$, sia $X_j \in K^{n-1}$ la colonna ottenuta eliminando l'ultima entrata di $A_j$, e sia $\tilde D \in K^{n-1}$ analoga per $D$. Dalla formula di definizione si vede che $\det_n$ è somma di $n$ applicazioni, ciascuna lineare in $D$ (quella con $j = j_0$ lo è esplicitamente, le altre per ipotesi induttiva). Una somma di applicazioni lineari è lineare.
\end{proof}

\begin{proposizione}[Colonne uguali $\Rightarrow$ determinante nullo]
\label{prop:col-uguali}
Se $A \in M_{n,n}(K)$ ha due colonne uguali, $A_{j_0} = A_{j_1}$ con $j_0 \neq j_1$, allora $\det A = 0$.
\end{proposizione}

\begin{proof}
Per induzione su $n$. Il passo base $n=1$ è banale. Mostriamo prima il caso di colonne \emph{adiacenti}: $j_1 = j_0+1$.

Per la multilinearità, $\det[\ldots, A_{j_0}+A_{j_0+1}, A_{j_0}+A_{j_0+1},\ldots]$ si espande come somma di quattro determinanti in cui le posizioni $j_0, j_0+1$ sono rispettivamente $(A_{j_0},A_{j_0})$, $(A_{j_0},A_{j_0+1})$, $(A_{j_0+1},A_{j_0})$, $(A_{j_0+1},A_{j_0+1})$. Per ipotesi induttiva (le sottomatrici $(n-1)\times(n-1)$ hanno due colonne uguali), i termini con colonne ripetute sono nulli. Quindi
\[
  0 = \det[\ldots,A_{j_0},A_{j_0+1},\ldots] + \det[\ldots,A_{j_0+1},A_{j_0},\ldots],
\]
il che dimostra che scambiare due colonne adiacenti cambia segno al determinante:
\begin{equation}
  \label{eq:scambio-adiacenti}
  \det[\ldots,A_{j_0+1},A_{j_0},\ldots] = -\det[\ldots,A_{j_0},A_{j_0+1},\ldots].
\end{equation}
Il caso generale $j_1 - j_0 = k > 1$ si riduce al caso $k=1$ per induzione su $k$, spostando la colonna $j_1$ verso sinistra passo per passo.
\end{proof}

\begin{corollario}[Scambio di colonne]
\label{cor:scambio-colonne}
Scambiando due colonne di $A$ di indici distinti, il determinante cambia segno.
\end{corollario}

\begin{corollario}[Operazione elementare di tipo 2]
\label{cor:op-elem}
$\det[\ldots, A_{j_0} + \lambda A_{j_1},\ldots] = \det A$ per ogni $\lambda \in K$ e $j_0 \neq j_1$.
\end{corollario}

\begin{proof}
Per la multilinearità, il determinante a sinistra è $\det A + \lambda \det[\ldots,A_{j_1},\ldots]$ dove la seconda matrice ha la colonna $j_1$ ripetuta in posizione $j_0$. Per la Proposizione~\ref{prop:col-uguali}, quel determinante è zero.
\end{proof}

\begin{osservazione}[Algoritmo per il calcolo]
\label{oss:algoritmo-det}
I Corollari~\ref{cor:scambio-colonne} e \ref{cor:op-elem} mostrano che:
\begin{itemize}
  \item un'operazione elementare di tipo (1) sulle colonne (scambio) moltiplica il determinante per $-1$;
  \item un'operazione elementare di tipo (2) sulle colonne (somma di multiplo di un'altra colonna) non cambia il determinante.
\end{itemize}
Quindi, portando $A$ a forma a scala per colonne $B$ con $s$ scambi:
\[
  \det A = (-1)^s \cdot b_{11} \cdots b_{nn}.
\]
Questo fornisce un algoritmo efficiente per il calcolo.
\end{osservazione}

\begin{proposizione}[Invertibilità e determinante]
\label{prop:det-invertibilita}
Una matrice quadrata $A \in M_{n,n}(K)$ è invertibile se e solo se $\det A \neq 0$.
\end{proposizione}

\begin{proof}
Sia $B$ la forma a scala per colonne di $A$ ottenuta con operazioni elementari di tipo (1) e (2). Le operazioni elementari non cambiano il rango, quindi $\rg(A) = \rg(B)$.

Per l'Osservazione~\ref{oss:algoritmo-det}, $\det A = 0 \Leftrightarrow \det B = 0 \Leftrightarrow b_{nn}=0$ (poiché $B$ è a scala, $b_{nn}=0$ implica che l'ultima colonna è nulla). Ma $b_{nn}=0 \Leftrightarrow \rg(B) < n \Leftrightarrow A$ è singolare.
\end{proof}

% --------------------------------------------------------
\section{Determinanti e rango}
\label{sec:det-rango}
% --------------------------------------------------------

Per la Proposizione~\ref{prop:det-invertibilita}, una matrice quadrata $n\times n$ ha rango $n$ se e solo se il suo determinante è non nullo. Il rango di una matrice rettangolare è caratterizzato dai determinanti delle sue sottomatrici quadrate.

\begin{defbox}
\begin{definizione}[Sottomatrice]
Sia $A \in M_{m,n}(K)$ e siano $1 \leq i_1 < \cdots < i_p \leq m$, $1\leq j_1 < \cdots < j_q \leq n$. La \emph{sottomatrice di $A$} corrispondente alle righe $i_1,\ldots,i_p$ e alle colonne $j_1,\ldots,j_q$ è la matrice $B \in M_{p,q}(K)$ con $b_{rs} = a_{i_r, j_s}$. Se $p=q$, il suo determinante si chiama \emph{minore} di $A$.

Una sottomatrice $C$ di $A$ \emph{contiene} $B$ se corrisponde a insiemi di righe e colonne che contengono rispettivamente $\{i_1,\ldots,i_p\}$ e $\{j_1,\ldots,j_q\}$.
\end{definizione}
\end{defbox}

\begin{teorema}[degli Orlati]
\label{thm:orlati}
Una matrice $A \in M_{m,n}(K)$ ha rango $r$ se e solo se esiste una sua sottomatrice $B$ di dimensione $r\times r$ con $\det B \neq 0$, e ogni sottomatrice di $A$ di dimensione $(r+1)\times(r+1)$ contenente $B$ ha determinante nullo.
\end{teorema}

\begin{proof}
Mostriamo prima un lemma utile.

\emph{Lemma.} Se $B$ è una sottomatrice $r\times r$ di $A$ con $\det B \neq 0$, allora le $r$ colonne (e le $r$ righe) di $A$ corrispondenti sono linearmente indipendenti, e in particolare $\rg A \geq r$.

\emph{Dimostrazione del Lemma.} Sia $\sum_{k=1}^r \lambda_k A_{j_k} = 0$ una combinazione nulla. Selezionando le righe $i_1,\ldots,i_r$ otteniamo $\sum_{k=1}^r \lambda_k B_k = 0$ (combinazione nulla delle colonne di $B$). Poiché $\det B \neq 0$, $B$ è invertibile, quindi $\lambda_1 = \cdots = \lambda_r = 0$.

\emph{Dimostrazione del Teorema.} Supponiamo $\rg A = r$. Esistono $r$ colonne linearmente indipendenti $A_{j_1},\ldots,A_{j_r}$, tali che ogni altra colonna è loro combinazione lineare. Sia $M$ la matrice $m\times r$ di queste colonne; $\rg M = r$, quindi esistono anche $r$ righe linearmente indipendenti di $M$, con indici $i_1,\ldots,i_r$. La sottomatrice $B = (a_{i_s,j_k})$ ha $\det B \neq 0$ (per il Lemma applicato a $M$). D'altra parte, se esistesse una sottomatrice $(r+1)\times(r+1)$ con determinante non nullo, per il Lemma il rango di $A$ sarebbe $\geq r+1$, contraddizione.

Il viceversa: data $B$ con le proprietà enunciate, per il Lemma $\rg A \geq r$. Basta quindi mostrare che ogni colonna di $A$ di indice diverso da $j_1,\ldots,j_r$ è combinazione lineare di $A_{j_1},\ldots,A_{j_r}$. Sia $A_j$ tale colonna. La matrice $M'$ con colonne $A_{j_1},\ldots,A_{j_r},A_j$ ha rango $< r+1$ se e solo se ogni sua sottomatrice $(r+1)\times(r+1)$ ha determinante nullo, il che vale per ipotesi. Quindi $\rg M' \leq r$, e poiché le prime $r$ colonne di $M'$ sono indipendenti, $A_j$ è loro combinazione lineare.
\end{proof}

\begin{corollario}
\label{cor:rango-minori}
$\rg A = r$ se e solo se esiste un minore $r\times r$ non nullo di $A$ e tutti i minori $(r+1)\times(r+1)$ di $A$ sono nulli.
\end{corollario}

\begin{exbox}
\begin{esempio}
Calcoliamo il rango di
\[
  A \coloneqq \begin{pmatrix} 2 & 1 & 3 & 5 \\ 3 & 0 & 3 & 6 \\ 0 & -1 & -1 & -1 \end{pmatrix} \in M_{3,4}(\Q).
\]
La sottomatrice delle righe $2,3$ e colonne $1,2$ è
$B = \begin{pmatrix}3&0\\0&-1\end{pmatrix}$
con $\det B = -3 \neq 0$. Le due sottomatrici $3\times 3$ contenenti $B$ sono
$\begin{pmatrix}2&1&3\\3&0&3\\0&-1&-1\end{pmatrix}$ e $\begin{pmatrix}2&1&5\\3&0&6\\0&-1&-1\end{pmatrix}$;
entrambe hanno determinante nullo (verifica diretta con l'Esempio~\ref{es:det3}). Per il Teorema degli Orlati, $\rg A = 2$.
\end{esempio}
\end{exbox}

% --------------------------------------------------------
\section{Applicazioni multilineari alternanti}
\label{sec:multilineare-alt}
% --------------------------------------------------------

Il determinante non è solo un numero: è \emph{il} prototipo di una classe di funzioni con proprietà speciali. Questa sezione sviluppa la teoria astratta che ne spiega l'unicità.

\subsection{Applicazioni multilineari}

\begin{defbox}
\begin{definizione}[Applicazione multilineare]
Siano $V_1,\ldots,V_n$ e $W$ spazi vettoriali su $K$. Un'applicazione $\Phi \colon V_1 \times \cdots \times V_n \to W$ è \emph{multilineare} se è lineare in ciascun argomento:
\[
  \Phi(v_1,\ldots,\lambda u + \mu w,\ldots,v_n) = \lambda\Phi(v_1,\ldots,u,\ldots,v_n) + \mu\Phi(v_1,\ldots,w,\ldots,v_n)
\]
per ogni posizione $j$-esima e per tutti $\lambda,\mu \in K$, $u,w \in V_j$.

Se $n=2$ si dice \emph{bilineare}. Se $W = K$ si dice \emph{forma multilineare}.
\end{definizione}
\end{defbox}

\begin{exbox}
\begin{esempio}
\leavevmode
\begin{enumerate}
  \item Il determinante $\det_n \colon (K^n)^n \to K$ è multilineare (Proposizione~\ref{prop:multilineare}).
  \item Il prodotto scalare standard $\R^n \times \R^n \to \R$, $(X,Y) \mapsto X^T Y$, è bilineare.
  \item Il prodotto vettoriale $\mathcal{V}(\mathbf{E}^3) \times \mathcal{V}(\mathbf{E}^3) \to \mathcal{V}(\mathbf{E}^3)$, $(v,w) \mapsto v \wedge w$, è bilineare.
  \item L'applicazione di valutazione $V^* \times V \to K$, $(f,v) \mapsto f(v)$, è bilineare.
\end{enumerate}
\end{esempio}
\end{exbox}

Le applicazioni multilineari formano uno spazio vettoriale: se $\Phi_1,\Phi_2$ sono multilineari, anche $\Phi_1 + \Phi_2$ e $\lambda\Phi_1$ lo sono.

\subsection{Applicazioni alternanti}

\begin{defbox}
\begin{definizione}[Applicazione alternante]
Sia $\Phi \colon V^n \to W$. Diciamo che $\Phi$ è \emph{alternante} (nelle entrate $j$ e $h$) se $\Phi(v_1,\ldots,v_n) = 0$ ogni qualvolta $v_j = v_h$. Diciamo che $\Phi$ è \emph{alternante} se lo è per ogni $1 \leq j < h \leq n$.
\end{definizione}
\end{defbox}

\begin{exbox}
\begin{esempio}
Il prodotto vettoriale $v \wedge w$ è bilineare e alternante: $v \wedge v = 0$ per ogni $v$.
\end{esempio}
\end{exbox}

\begin{proposizione}
\label{prop:alt-proprieta}
Sia $\Phi \colon V^n \to W$ multilineare e alternante.
\begin{enumerate}
  \item Se $v_1,\ldots,v_n$ sono linearmente dipendenti, allora $\Phi(v_1,\ldots,v_n) = 0$.
  \item Scambiando due argomenti, il valore di $\Phi$ cambia segno:
  \[
    \Phi(v_1,\ldots,v_h,\ldots,v_j,\ldots,v_n) = -\Phi(v_1,\ldots,v_j,\ldots,v_h,\ldots,v_n).
  \]
\end{enumerate}
\end{proposizione}

\begin{proof}
\emph{(1).} Per la Proposizione~2.5.10 (esistenza di un vettore esprimibile come combinazione lineare degli altri), esiste $j_0$ tale che $v_{j_0} = \sum_{i \neq j_0}\mu_i v_i$. Per la multilinearità:
\[
  \Phi(v_1,\ldots,v_n) = \sum_{i \neq j_0}\mu_i\,\Phi(v_1,\ldots,v_{j_0-1},v_i,v_{j_0+1},\ldots,v_n).
\]
Ogni termine ha due argomenti uguali (quello in posizione $j_0$ e quello in posizione $i$), quindi è zero per l'alternanza.

\emph{(2).} Per $n=2$: poiché $\Phi$ è alternante e bilineare,
\[
  0 = \Phi(u+w,\,u+w) = \Phi(u,u)+\Phi(w,w)+\Phi(u,w)+\Phi(w,u) = \Phi(u,w)+\Phi(w,u).
\]
Il caso $n>2$ si riduce a $n=2$ fissando gli altri argomenti.
\end{proof}

\subsection{Unicità del determinante}

Il risultato chiave di questa sezione afferma che ogni forma multilineare alternante su uno spazio $n$-dimensionale è un multiplo del determinante.

\begin{defbox}
\begin{definizione}
Se $V$ è uno spazio vettoriale su $K$, denotiamo con $\mathrm{Alt}(V^n,K)$ l'insieme delle forme multilineari alternanti $\Phi \colon V^n \to K$. Questo è un sottospazio di $\mathrm{MultLin}(V^n,K)$.
\end{definizione}
\end{defbox}

\begin{proposizione}[Unicità]
\label{prop:unicita-det}
Sia $V$ uno spazio vettoriale su $K$ di dimensione $n$, con base $\mathcal{B} = \{v_1,\ldots,v_n\}$. L'applicazione di valutazione
\[
  \mathrm{ev}_{\mathcal{B}} \colon \mathrm{Alt}(V^n,K) \longrightarrow K, \quad \Phi \mapsto \Phi(v_1,\ldots,v_n)
\]
è un isomorfismo di spazi vettoriali. In particolare, $\dim\,\mathrm{Alt}(V^n,K) = 1$.
\end{proposizione}

\begin{proof}
\emph{Iniettività.} Sia $\Phi \in \ker(\mathrm{ev}_{\mathcal{B}})$, cioè $\Phi(v_1,\ldots,v_n)=0$. Dobbiamo mostrare $\Phi \equiv 0$. Per $u_1,\ldots,u_n \in V$, scriviamo $u_i = \sum_{j=1}^n a_{ij}v_j$. Per la multilinearità:
\[
  \Phi(u_1,\ldots,u_n) = \sum_{j_1,\ldots,j_n} a_{1,j_1}\cdots a_{n,j_n}\,\Phi(v_{j_1},\ldots,v_{j_n}).
\]
Se tra gli indici $j_1,\ldots,j_n$ c'è una ripetizione, il termine è zero per l'alternanza. Se non ci sono ripetizioni, $(j_1,\ldots,j_n)$ è una permutazione di $(1,\ldots,n)$ e, scambiando argomenti opportunamente (con cambio di segno per la Proposizione~\ref{prop:alt-proprieta}), si ottiene $\pm\Phi(v_1,\ldots,v_n) = 0$. Dunque $\Phi(u_1,\ldots,u_n)=0$.

\emph{Suriettività.} Basta esibire una $\Phi \in \mathrm{Alt}(V^n,K)$ con $\Phi(v_1,\ldots,v_n) \neq 0$. Tramite l'isomorfismo $\chi_{\mathcal{B}} \colon V \xrightarrow{\sim} K^n$, si trasporta $\det_n$ su $V$: la funzione $\Phi(u_1,\ldots,u_n) \coloneqq \det_n(\chi_{\mathcal{B}}(u_1),\ldots,\chi_{\mathcal{B}}(u_n))$ è multilineare alternante e $\Phi(v_1,\ldots,v_n) = \det(1_n) = 1 \neq 0$.
\end{proof}

\begin{corollario}
Ogni forma multilineare alternante $\Phi \colon (K^n)^n \to K$ è della forma $\Phi = c\cdot\det_n$ per un unico $c \in K$.
\end{corollario}

% --------------------------------------------------------
\section{Formula di Binet e sue conseguenze}
\label{sec:binet}
% --------------------------------------------------------

\subsection{La formula di Binet}

\begin{teorema}[Binet]
\label{thm:binet}
Siano $A, B \in M_{n,n}(K)$. Allora
\[
  \det(A \cdot B) = \det A \cdot \det B.
\]
\end{teorema}

\begin{proof}
Definiamo $\Phi \colon M_{n,n}(K) \to K$ da $\Phi(M) \coloneqq \det(A \cdot M)$, vista come funzione delle colonne di $M$. Per la multilinearità di $\det$ e il fatto che $A \cdot [X_1|\cdots|X_n] = [AX_1|\cdots|AX_n]$, la funzione $\Phi$ è multilineare nelle colonne di $M$. È alternante: se due colonne di $M$ sono uguali, due colonne di $A\cdot M$ sono uguali.

Anche $\Psi(M) \coloneqq \det A \cdot \det M$ è multilineare e alternante nelle colonne di $M$.

Poiché $\Phi(1_n) = \det(A\cdot 1_n) = \det A = \det A \cdot \det 1_n = \Psi(1_n)$, per la Proposizione~\ref{prop:unicita-det} (applicata con $V=K^n$ e base standard) si ottiene $\Phi = \Psi$.
\end{proof}

\begin{corollario}
\label{cor:det-inversa}
Se $A \in M_{n,n}(K)$ è invertibile, allora $\det(A^{-1}) = (\det A)^{-1}$.
\end{corollario}

\begin{proof}
$1 = \det(1_n) = \det(A \cdot A^{-1}) = \det A \cdot \det(A^{-1})$.
\end{proof}

\begin{osservazione}
La formula di Binet mostra che l'applicazione $\det \colon \GL_n(K) \to K^*$ è un \emph{omomorfismo di gruppi}.
\end{osservazione}

\begin{corollario}[Invarianza per coniugazione]
\label{cor:det-coniugazione}
Sia $V$ uno spazio vettoriale finitamente generato e $f \colon V \to V$ un endomorfismo. Se $\mathcal{B}$ e $\mathcal{C}$ sono due basi di $V$, allora
\[
  \det(M^{\mathcal{B}}_{\mathcal{B}}(f)) = \det(M^{\mathcal{C}}_{\mathcal{C}}(f)).
\]
\end{corollario}

\begin{proof}
$M^{\mathcal{B}}_{\mathcal{B}}(f) = (M^{\mathcal{B}}_{\mathcal{C}}(\Id_V))^{-1}\cdot M^{\mathcal{C}}_{\mathcal{C}}(f) \cdot M^{\mathcal{B}}_{\mathcal{C}}(\Id_V)$. Applicando Binet e il Corollario~\ref{cor:det-inversa}:
\[
  \det(M^{\mathcal{B}}_{\mathcal{B}}(f)) = (\det M^{\mathcal{B}}_{\mathcal{C}}(\Id_V))^{-1}\cdot\det M^{\mathcal{C}}_{\mathcal{C}}(f)\cdot\det M^{\mathcal{B}}_{\mathcal{C}}(\Id_V) = \det M^{\mathcal{C}}_{\mathcal{C}}(f). \qedhere
\]
\end{proof}

\begin{defbox}
\begin{definizione}[Determinante di un endomorfismo]
Sia $f \colon V \to V$ un endomorfismo di uno spazio finitamente generato. Il \emph{determinante di $f$} è
\[
  \det f \coloneqq \det(M^{\mathcal{B}}_{\mathcal{B}}(f))
\]
per un'arbitraria base $\mathcal{B}$ di $V$. È ben definito per il Corollario~\ref{cor:det-coniugazione}.
\end{definizione}
\end{defbox}

\subsection{Sviluppo di Laplace}

\begin{proposizione}[Sviluppo lungo la riga $i$-esima]
\label{prop:laplace-righe}
Per ogni $A \in M_{n,n}(K)$ e ogni $1 \leq i \leq n$:
\[
  \det A = \sum_{j=1}^n (-1)^{i+j}\,a_{ij}\,\det(A^i_j).
\]
\end{proposizione}

\begin{proof}
Sia $\Phi^i(A)$ il membro di destra. Si verifica che $\Phi^i$ è multilineare e alternante nelle colonne (argomenti analoghi a quelli delle Proposizioni~\ref{prop:multilineare} e~\ref{prop:col-uguali}), e che $\Phi^i(1_n) = 1$. Per la Proposizione~\ref{prop:unicita-det}, $\Phi^i = \det$.
\end{proof}

\begin{proposizione}[Determinante della trasposta]
\label{prop:det-trasposta}
Per ogni $A \in M_{n,n}(K)$, $\det(A^T) = \det A$.
\end{proposizione}

\begin{proof}
L'applicazione $\Phi(A) \coloneqq \det(A^T)$, vista come funzione delle colonne di $A$ (che sono le righe di $A^T$), è multilineare e alternante per la Proposizione~\ref{prop:laplace-righe} (il determinante è multilineare e alternante anche nelle righe). Poiché $\Phi(1_n) = \det(1_n) = 1$, si ha $\Phi = \det$.
\end{proof}

\begin{corollario}[Sviluppo lungo la colonna $j$-esima]
Per ogni $A \in M_{n,n}(K)$ e ogni $1 \leq j \leq n$:
\[
  \det A = \sum_{i=1}^n (-1)^{i+j}\,a_{ij}\,\det(A^i_j).
\]
\end{corollario}

\begin{proof}
Espandere $\det(A^T)$ lungo la riga $j$-esima e usare $\det(A^T) = \det A$.
\end{proof}

\begin{osservazione}
La Proposizione~\ref{prop:det-trasposta} mostra che il determinante è multilineare e alternante anche nelle righe. Pertanto l'Osservazione~\ref{oss:algoritmo-det} vale anche riducendo $A$ a scala per righe (con operazioni elementari sulle righe).
\end{osservazione}

\subsection{Formula di Cramer}

\begin{defbox}
\begin{definizione}[Cofattore e matrice aggiunta]
Sia $A \in M_{n,n}(K)$. Il \emph{cofattore} (o \emph{complemento algebrico}) di $A$ di indici $(i,j)$ è
\[
  A_{ij} \coloneqq (-1)^{i+j}\det(A^i_j).
\]
La \emph{matrice dei cofattori} (o \emph{matrice aggiunta classica}) è
\[
  A^c \coloneqq (A_{ji})_{i,j} \in M_{n,n}(K),
\]
cioè la trasposta della matrice degli $A_{ij}$.
\end{definizione}
\end{defbox}

\begin{exbox}
\begin{esempio}
Per $A = \begin{pmatrix}1&1&1\\1&-1&2\\1&1&4\end{pmatrix}$, si calcola $\det A = -6$ e
\[
  A^c = \begin{pmatrix}-6 & -3 & 3\\ -2 & 3 & -1 \\ 2 & 0 & -2\end{pmatrix}.
\]
\end{esempio}
\end{exbox}

\begin{teorema}[Formula di Cramer]
\label{thm:cramer}
Sia $A \in M_{n,n}(K)$. Allora
\[
  A \cdot A^c = A^c \cdot A = (\det A)\,1_n.
\]
In particolare, se $\det A \neq 0$, allora $A^{-1} = (\det A)^{-1} A^c$.
\end{teorema}

\begin{proof}
Calcoliamo l'entrata $(i,j)$ di $A \cdot A^c$:
\[
  (A \cdot A^c)_{ij} = \sum_{s=1}^n a_{is} A_{js} = \sum_{s=1}^n (-1)^{j+s} a_{is}\det(A^j_s).
\]
Se $i = j$: questa è esattamente lo sviluppo di $\det A$ lungo la riga $i$-esima, quindi $(A\cdot A^c)_{ii} = \det A$.

Se $i \neq j$: la somma è lo sviluppo lungo la riga $j$-esima della matrice $B$ ottenuta da $A$ sostituendo alla riga $j$-esima la riga $i$-esima. Poiché $B$ ha due righe uguali, $\det B = 0$.

Dunque $A \cdot A^c = (\det A)1_n$. La formula $A^c \cdot A = (\det A)1_n$ si dimostra analogamente.
\end{proof}

\begin{exbox}
\begin{esempio}
Con la matrice dell'esempio precedente ($\det A = -6$):
\[
  A^{-1} = \frac{1}{-6}\begin{pmatrix}-6&-3&3\\-2&3&-1\\2&0&-2\end{pmatrix}
  = \begin{pmatrix}1&\tfrac{1}{2}&-\tfrac{1}{2}\\\tfrac{1}{3}&-\tfrac{1}{2}&\tfrac{1}{6}\\-\tfrac{1}{3}&0&\tfrac{1}{3}\end{pmatrix}.
\]
\end{esempio}
\end{exbox}

% --------------------------------------------------------
\section{Il determinante come volume orientato}
\label{sec:det-area}
% --------------------------------------------------------

Questa sezione dà un'interpretazione geometrica del determinante su $\R$.

\begin{defbox}
\begin{definizione}[Parallelogramma]
Siano $P_0 \in \mathbf{E}^2$ e $v, w \in \mathcal{V}(\mathbf{E}^2)$. Il \emph{parallelogramma} generato da $v$ e $w$ con punto base $P_0$ è
\[
  \Pi(P_0,v,w) \coloneqq \{P_0 + sv + tw \mid (s,t) \in [0,1]^2\}.
\]
\end{definizione}
\end{defbox}

\begin{proposizione}
\label{prop:det-area}
Sia $\mathcal{B}=\{i,j\}$ una base di $\mathcal{V}(\mathbf{E}^2)$ tale che $\Pi(P_0,i,j)$ abbia area $1$. Se $v = a_{11}i+a_{21}j$ e $w=a_{12}i+a_{22}j$, allora
\[
  A(\Pi(P_0,v,w)) = \left|\det\begin{pmatrix}a_{11}&a_{12}\\a_{21}&a_{22}\end{pmatrix}\right|.
\]
\end{proposizione}

\begin{proof}[Idea della dimostrazione]
Si definisce l'applicazione $\Phi \colon \mathcal{V}(\mathbf{E}^2)^2 \to \R$ come l'area con segno $\epsilon(v,w)\cdot A(\Pi(P_0,v,w))$, dove $\epsilon(v,w) = +1$ se $\{v,w\}$ è una base con la stessa orientazione di $\mathcal{B}$, $-1$ altrimenti. Si dimostra che $\Phi$ è bilineare, alternante e $\Phi(i,j)=1$. Per la Proposizione~\ref{prop:unicita-det}, $\Phi$ coincide con $\Psi(v,w) = \det\begin{pmatrix}a_{11}&a_{12}\\a_{21}&a_{22}\end{pmatrix}$. Prendendo il valore assoluto si ottiene la formula.
\end{proof}

\begin{osservazione}
Il risultato si estende in ogni dimensione: il volume del parallelepipedo in $\R^n$ con lati $A_1,\ldots,A_n$ è $|\det[A_1|\cdots|A_n]|$.

Se $F \colon \mathbf{E}^n \to \mathbf{E}^n$ è un'applicazione affine con parte lineare $f = \mathcal{V}(F)$, allora per ogni regione misurabile $T$:
\[
  \mathrm{Vol}(F(T)) = |\det f|\cdot\mathrm{Vol}(T).
\]
Questo è il teorema di cambiamento di variabile nel calcolo integrale.
\end{osservazione}

% --------------------------------------------------------
\section{Determinante e permutazioni}
\label{sec:det-permutazioni}
% --------------------------------------------------------

\subsection{La formula di Leibniz}

Denotiamo $\mathbf{n} \coloneqq \{1,\ldots,n\}$ e con $S_n$ il \emph{gruppo simmetrico} delle permutazioni di $\mathbf{n}$.

\begin{defbox}
\begin{definizione}
Per $\sigma \in S_n$, la \emph{matrice di permutazione} $M_\sigma \in M_{n,n}(K)$ ha $(i,j)$-entrata pari a $1$ se $i=\sigma(j)$ e $0$ altrimenti. Equivalentemente, la colonna $j$ di $M_\sigma$ è $e_{\sigma(j)}$.

Il \emph{segno} di $\sigma$ è $\varepsilon(\sigma) \coloneqq \det(M_\sigma) \in \{+1,-1\}$.
\end{definizione}
\end{defbox}

\begin{teorema}[Formula di Leibniz]
\label{thm:leibniz}
Sia $A = (a_{ij}) \in M_{n,n}(K)$. Allora
\[
  \det A = \sum_{\sigma \in S_n} \varepsilon(\sigma)\,a_{1,\sigma(1)}\,a_{2,\sigma(2)}\cdots a_{n,\sigma(n)}.
\]
\end{teorema}

\begin{proof}
Dalla multilinearità nelle colonne: $\det A = \det\Bigl[\sum_i a_{i1}e_i,\ldots,\sum_i a_{in}e_i\Bigr] = \sum_\phi \det(M_\phi)\,a_{\phi(1),1}\cdots a_{\phi(n),n}$, dove la somma è su tutte le $\phi \colon \mathbf{n}\to\mathbf{n}$. Per la Proposizione~\ref{prop:col-uguali}, $\det(M_\phi)=0$ se $\phi$ non è biunivoca. Per $\sigma \in S_n$, si verifica che $M_{\sigma\circ\tau} = M_\sigma \cdot M_\tau$ (prodotto di permutazioni corrisponde a prodotto di matrici), quindi $\varepsilon$ è un omomorfismo $S_n \to K^*$. Riordinando gli indici $a_{\sigma(j),j} = a_{j,\sigma^{-1}(j)}$ e sostituendo $\sigma \to \sigma^{-1}$ (che lascia invariata la somma su $S_n$) si ottiene la formula.
\end{proof}

\subsection{Segno di una permutazione}

\begin{defbox}
\begin{definizione}[Trasposizione]
Una \emph{trasposizione} $\tau_{ab} \in S_n$ ($a \neq b$) è la permutazione che scambia $a$ e $b$ e fissa tutti gli altri elementi.
\end{definizione}
\end{defbox}

\begin{proposizione}
\label{prop:segno-trasposizione}
Se $\tau$ è una trasposizione, allora $\varepsilon(\tau) = -1$.
\end{proposizione}

\begin{proof}
La matrice $M_\tau$ si ottiene da $1_n$ scambiando le colonne $a$ e $b$. Per il Corollario~\ref{cor:scambio-colonne}, $\det M_\tau = -\det 1_n = -1$.
\end{proof}

\begin{proposizione}
\label{prop:sn-trasposizioni}
Ogni $\sigma \in S_n$ è prodotto di trasposizioni. Se $\sigma = \tau_1\circ\cdots\circ\tau_k = \rho_1\circ\cdots\circ\rho_l$ sono due scritture come prodotto di trasposizioni, allora $(-1)^k = (-1)^l$, ovvero la parità del numero di fattori è ben definita.
\end{proposizione}

\begin{proof}[Idea]
L'esistenza della decomposizione si dimostra per induzione su $n$: se $\sigma(n)=m\neq n$, la composizione $\tau_{mn}\circ\sigma$ fissa $n$ e si riduce a $S_{n-1}$. La unicità della parità segue dall'omomorfismo $\varepsilon \colon S_n \to \{+1,-1\}$: poiché $\varepsilon(\tau)=-1$ per ogni trasposizione, $\varepsilon(\sigma) = (-1)^k$.
\end{proof}

\begin{defbox}
\begin{definizione}
$\sigma \in S_n$ è \emph{pari} se $\varepsilon(\sigma)=+1$, \emph{dispari} se $\varepsilon(\sigma)=-1$.
\end{definizione}
\end{defbox}

\begin{exbox}
\begin{esempio}
Sia $\sigma \in S_4$ con $\sigma(1)=3,\sigma(2)=4,\sigma(3)=1,\sigma(4)=2$. Partiamo dalla successione $(3,4,1,2)$ e applichiamo scambi per raggiungere $(1,2,3,4)$:
\[
  (3,4,1,2) \xrightarrow{\text{scambio pos.\,1,3}} (1,4,3,2) \xrightarrow{\text{scambio pos.\,2,4}} (1,2,3,4).
\]
Due scambi: $\sigma$ è pari, $\varepsilon(\sigma) = +1$.
\end{esempio}
\end{exbox}

% --------------------------------------------------------
\section{Determinanti su anelli e polinomio caratteristico}
\label{sec:det-anello-polcar}
% --------------------------------------------------------

\subsection{Determinanti con entrate in un anello}

Tutte le formule dimostrate (Binet, Laplace, Cramer, Leibniz) valgono per matrici $A \in M_{n,n}(R)$ purché $R$ sia un \emph{sottoanello di un campo} $K$. In particolare valgono per $R = \Z$ (sottoanello di $\Q$) e per $R = K[x]$ (sottoanello di $K(x)$).

\begin{proposizione}
\label{prop:invertibile-anello}
Sia $R$ un sottoanello di un campo e $A \in M_{n,n}(R)$. Esiste $B \in M_{n,n}(R)$ con $A\cdot B = B\cdot A = 1_n$ se e solo se $\det A$ è invertibile in $R$.
\end{proposizione}

\begin{proof}
Se $B$ è l'inversa, Binet dà $1 = \det(1_n) = \det A \cdot \det B$, quindi $\det A$ è invertibile in $R$. Viceversa, se $\det A$ è invertibile, allora $(\det A)^{-1} A^c \in M_{n,n}(R)$ e Cramer dà $A\cdot (\det A)^{-1}A^c = 1_n$.
\end{proof}

\begin{osservazione}
Per $R=\Z$: $A \in M_{n,n}(\Z)$ è invertibile in $M_{n,n}(\Z)$ se e solo se $\det A = \pm 1$.
\end{osservazione}

\subsection{Polinomio caratteristico}

Sia $A \in M_{n,n}(K)$. La matrice $\lambda 1_n - A$ ha entrate nel campo $K[\lambda]$ dei polinomi in $\lambda$. Per la Proposizione~\ref{prop:invertibile-anello}, il suo determinante è un elemento ben definito di $K[\lambda]$.

\begin{defbox}
\begin{definizione}[Polinomio caratteristico]
Il \emph{polinomio caratteristico} di $A \in M_{n,n}(K)$ è
\[
  p_A(\lambda) \coloneqq \det(\lambda 1_n - A) \in K[\lambda].
\]
\end{definizione}
\end{defbox}

\begin{proposizione}
\label{prop:polcar-monico}
$p_A$ ha grado $n$ ed è monico (il coefficiente di $\lambda^n$ è $1$). Più precisamente:
\[
  p_A(\lambda) = \lambda^n - (\mathrm{Tr}\, A)\,\lambda^{n-1} + \cdots + (-1)^n \det A,
\]
dove $\mathrm{Tr}\, A \coloneqq \sum_{i=1}^n a_{ii}$ è la \emph{traccia} di $A$.
\end{proposizione}

\begin{proof}
Dalla formula di Leibniz, $p_A(\lambda) = \sum_{\sigma \in S_n}\varepsilon(\sigma)\prod_{i=1}^n(\lambda\delta_{i,\sigma(i)} - a_{i,\sigma(i)})$. L'unico addendo di grado $n$ in $\lambda$ è quello con $\sigma = \mathrm{Id}$: $\prod_{i=1}^n(\lambda - a_{ii})$, che contribuisce $\lambda^n - (\sum_i a_{ii})\lambda^{n-1} + \ldots$. Il termine costante ($\lambda=0$) è $\det(-A) = (-1)^n\det A$.
\end{proof}

\begin{defbox}
\begin{definizione}[Traccia]
La \emph{traccia} di $A \in M_{n,n}(K)$ è $\mathrm{Tr}\,A \coloneqq \sum_{i=1}^n a_{ii}$.

Se $f \colon V \to V$ è un endomorfismo, la \emph{traccia di $f$} è $\mathrm{Tr}\,f \coloneqq \mathrm{Tr}(M^{\mathcal{B}}_{\mathcal{B}}(f))$, indipendente dalla base $\mathcal{B}$ (poiché $p_f$ non dipende dalla base per la Proposizione~\ref{prop:polcar-coniugazione} che segue).
\end{definizione}
\end{defbox}

\begin{proposizione}
\label{prop:polcar-coniugazione}
Se $A$ e $B$ sono coniugate ($B = G^{-1}AG$ per qualche $G \in \GL_n(K)$), allora $p_A = p_B$.
\end{proposizione}

\begin{proof}
$p_B(\lambda) = \det(\lambda 1_n - G^{-1}AG) = \det(G^{-1}(\lambda 1_n - A)G) = \det(G^{-1})\det(\lambda 1_n - A)\det G = p_A(\lambda)$.
\end{proof}

\begin{osservazione}
Il collegamento con gli autovalori: $\lambda_0$ è autovalore di $f$ se e solo se $p_f(\lambda_0) = 0$. Infatti, $\lambda_0$ è autovalore $\Leftrightarrow$ $\ker(\lambda_0 \Id - f) \neq \{0\}$ $\Leftrightarrow$ $\lambda_0 \Id - f$ non è invertibile $\Leftrightarrow$ $\det(\lambda_0 1_n - A) = 0$.
\end{osservazione}

% --------------------------------------------------------
\section{Diagonalizzazione}
\label{sec:diagonalizzazione}
% --------------------------------------------------------

Il polinomio caratteristico è lo strumento principale per studiare la diagonalizzabilità di un endomorfismo. Richiamiamo le definizioni dal Capitolo~\ref{cap:applin}.

\begin{defbox}
\begin{definizione}
Sia $f \colon V \to V$ un endomorfismo e $\lambda_0 \in K$ un autovalore. La \emph{molteplicità algebrica} di $\lambda_0$ è $\mathrm{mult}_{\lambda_0}(p_f)$ (molteplicità di $\lambda_0$ come radice di $p_f$). La \emph{molteplicità geometrica} di $\lambda_0$ è $\dim V_{\lambda_0}(f)$.
\end{definizione}
\end{defbox}

\begin{proposizione}
\label{prop:mult-geom-alg}
Sia $\lambda_0$ autovalore di $f \colon V \to V$ con $\dim V = n$. Allora
\[
  1 \leq \dim V_{\lambda_0}(f) \leq \mathrm{mult}_{\lambda_0}(p_f).
\]
\end{proposizione}

\begin{proof}
Il limite inferiore è la definizione di autovalore. Per il limite superiore: sia $r = \dim V_{\lambda_0}(f)$ ed estendiamo una base $\{v_1,\ldots,v_r\}$ di $V_{\lambda_0}(f)$ a una base $\mathcal{B}$ di $V$. La matrice $M^{\mathcal{B}}_{\mathcal{B}}(\lambda_0\Id-f)$ ha le prime $r$ colonne del tipo $(\lambda-\lambda_0)e_i$ e le rimanenti colonne di grado $\leq 1$ in $\lambda$. Sviluppando il determinante si ottiene $p_f(\lambda) = (\lambda-\lambda_0)^r \cdot q(\lambda)$ con $q \in K[\lambda]$, quindi $\mathrm{mult}_{\lambda_0}(p_f) \geq r$.
\end{proof}

\begin{lemma}
\label{lemma:autovettori-indipendenti}
Se $v_1,\ldots,v_d \in V$ sono autovettori con autovalori distinti, allora sono linearmente indipendenti.
\end{lemma}

\begin{proof}
Per induzione su $d$. Se $d=1$, un autovettore è non nullo. Supponiamo $d>1$ e $\alpha_1 v_1 + \cdots + \alpha_d v_d = 0$ con coefficienti non tutti nulli; per ipotesi induttiva, ciascun $\alpha_i \neq 0$ (altrimenti avremmo una dipendenza tra $d-1$ autovettori con autovalori distinti). Applicando $f$: $\alpha_1\lambda_1 v_1 + \cdots + \alpha_d\lambda_d v_d = 0$. Moltiplicando la prima relazione per $\lambda_d$ e sottraendo:
\[
  \sum_{i=1}^{d-1}\alpha_i(\lambda_i - \lambda_d)v_i = 0.
\]
Poiché $\lambda_i \neq \lambda_d$ e $\alpha_i \neq 0$, per ipotesi induttiva i $v_1,\ldots,v_{d-1}$ sono indipendenti: contraddizione.
\end{proof}

\begin{teorema}[Criterio di diagonalizzabilità]
\label{thm:diagonalizzabilita}
Sia $f \colon V \to V$ un endomorfismo con $\dim V = n$. Le seguenti condizioni sono equivalenti:
\begin{enumerate}
  \item $f$ è diagonalizzabile.
  \item $p_f$ ha esattamente $n$ radici in $K$ (contate con molteplicità) e per ogni autovalore $\lambda$:
  \[
    \dim V_\lambda(f) = \mathrm{mult}_\lambda(p_f).
  \]
  \item $\displaystyle\sum_{\lambda \in K}\dim V_\lambda(f) = n$.
\end{enumerate}
\end{teorema}

\begin{proof}
\emph{(2)$\Leftrightarrow$(3).} Seguono dalla Proposizione~\ref{prop:mult-geom-alg} e dalla relazione $\sum_\lambda\mathrm{mult}_\lambda(p_f) \leq \deg p_f = n$.

\emph{(1)$\Rightarrow$(2).} Se $\mathcal{B}=\{v_1,\ldots,v_n\}$ è una base di autovettori con autovalori $\lambda_1,\ldots,\lambda_n$, allora $p_f(\lambda) = \prod_{i=1}^n(\lambda-\lambda_i)$. Per ogni autovalore distinto $\lambda_j$, la molteplicità algebrica è $|\{i \mid \lambda_i = \lambda_j\}|$, ma tutti i vettori $v_i$ con $\lambda_i = \lambda_j$ stanno in $V_{\lambda_j}(f)$, quindi $\dim V_{\lambda_j}(f) \geq \mathrm{mult}_{\lambda_j}(p_f)$, e per la Proposizione~\ref{prop:mult-geom-alg} si ha l'uguaglianza.

\emph{(2)$\Rightarrow$(1).} Siano $\lambda_1,\ldots,\lambda_d$ gli autovalori distinti. Per ogni $i$, sia $\mathcal{B}_i$ una base di $V_{\lambda_i}(f)$, di cardinalità $\mathrm{mult}_{\lambda_i}(p_f)$. Per il Lemma~\ref{lemma:autovettori-indipendenti}, l'unione $\mathcal{B}_1 \cup \cdots \cup \mathcal{B}_d$ è linearmente indipendente, e ha cardinalità $\sum_i \mathrm{mult}_{\lambda_i}(p_f) = n$. È quindi una base di $V$ formata da autovettori.
\end{proof}

\begin{corollario}
Se $p_f$ ha $n$ radici distinte in $K$, allora $f$ è diagonalizzabile.
\end{corollario}

\begin{exbox}
\begin{esempio}[Matrice diagonalizzabile]
Sia $A = \begin{pmatrix}1 & 5\\ -\tfrac{1}{4} & -2\end{pmatrix} \in M_{2,2}(\Q)$. Il polinomio caratteristico è
\[
  p_A(\lambda) = \lambda^2 + \lambda - \frac{3}{4} = \left(\lambda + \frac{3}{2}\right)\left(\lambda - \frac{1}{2}\right).
\]
Le radici $\lambda_1 = -\tfrac{3}{2}$ e $\lambda_2 = \tfrac{1}{2}$ sono distinte, quindi $A$ è diagonalizzabile. Gli autovettori si trovano risolvendo $(A - \lambda_i 1_2)X = 0$: per $\lambda_1$ si ottiene $(2,-1)^T$, per $\lambda_2$ si ottiene $(10,-1)^T$. La base $\mathcal{B} = \{(2,-1),(10,-1)\}$ diagonalizza $A$:
\[
  M^{\mathcal{B}}_{\mathcal{B}}(L_A) = \begin{pmatrix}-\tfrac{3}{2}&0\\0&\tfrac{1}{2}\end{pmatrix}.
\]
\end{esempio}
\end{exbox}

\begin{exbox}
\begin{esempio}[Decidere la diagonalizzabilità con molteplicità]
Sia $A = \begin{pmatrix}0&-4&-1\\3&13&3\\-12&-48&-11\end{pmatrix} \in M_{3,3}(\Q)$. Si calcola $p_A(\lambda) = \lambda(\lambda-1)^2$. Le radici sono $0$ (molteplicità $1$) e $1$ (molteplicità $2$). $A$ è diagonalizzabile se e solo se $\dim V_1(L_A) = 2$. Il sistema $(A-1_3)X=0$ ha matrice dei coefficienti
$\begin{pmatrix}-1&-4&-1\\3&12&3\\-12&-48&-12\end{pmatrix}$
di rango $1$, quindi $\dim V_1 = 2$. Sì, $A$ è diagonalizzabile.
\end{esempio}
\end{exbox}

% --------------------------------------------------------
\section*{Esercizi}
\addcontentsline{toc}{section}{Esercizi del Capitolo 5}
% --------------------------------------------------------

\begin{esercizio}
Calcolate il determinante delle seguenti matrici:
\[
  A \coloneqq \begin{pmatrix}2&3&1\\3&5&0\\2&4&2\end{pmatrix}, \qquad
  B \coloneqq \begin{pmatrix}1&2&3\\2&3&4\\3&4&5\end{pmatrix}.
\]
[\emph{Suggerimento}: per $B$, notare che le righe non sono indipendenti.]
\end{esercizio}

\begin{esercizio}
Calcolate la matrice dei cofattori di $A$ e $B$ dell'esercizio precedente. Per quale delle due si può calcolare l'inversa? Calcolatela.
\end{esercizio}

\begin{esercizio}[Determinante di Vandermonde]
Sia $K$ un campo e $x_1,\ldots,x_n \in K$. Dimostrate che
\[
  \det\begin{pmatrix}1&1&\cdots&1\\x_1&x_2&\cdots&x_n\\x_1^2&x_2^2&\cdots&x_n^2\\\vdots&\vdots&\ddots&\vdots\\x_1^{n-1}&x_2^{n-1}&\cdots&x_n^{n-1}\end{pmatrix}
  = \prod_{1\leq i < j \leq n}(x_j - x_i).
\]
[\emph{Suggerimento}: per induzione su $n$; sottrarre a ciascuna colonna $j\geq 2$ la colonna precedente moltiplicata per $x_1$, poi espandere lungo la prima riga.]
\end{esercizio}

\begin{esercizio}[Guidato: potenza di una matrice tramite diagonalizzazione]
Sia $A = \begin{pmatrix}22&-12\\35&-19\end{pmatrix} \in M_{2,2}(\Z)$.
\begin{enumerate}
  \item Calcolate $p_A(\lambda)$ e le sue radici in $\Q$.
  \item Trovate una base di autovettori e la matrice di passaggio $P$.
  \item Usate la diagonalizzazione per scrivere $A^m$ in forma chiusa per $m \in \Z$.
\end{enumerate}
\end{esercizio}

\begin{esercizio}
Sia $A \in M_{n,n}(K)$ e sia $B$ la matrice ottenuta riscrivendo le colonne di $A$ nell'ordine inverso. Quale relazione esiste tra $\det A$ e $\det B$? [\emph{Suggerimento}: quanti scambi di colonne adiacenti servono per invertire l'ordine di $n$ elementi?]
\end{esercizio}

\begin{esercizio}
Sia $A \in M_{n,n}(\Z)$ e sia $p$ un primo. Denotiamo $\bar A \in M_{n,n}(\Z/p\Z)$ la riduzione modulo $p$. Dimostrate che se $\rg A = r$ allora $p^{n-r} \mid \det A$.
\end{esercizio}

\begin{esercizio}[Ricorrenza tridiagonale]
Per $n \geq 1$ sia $A^{(n)} \in M_{n,n}(\Z)$ la matrice tridiagonale con $2$ sulla diagonale principale e $-1$ sulle due diagonali adiacenti. Dimostrate che $\det A^{(n)} = n+1$ per ogni $n \geq 1$. [\emph{Suggerimento}: sviluppare lungo l'ultima riga e ricavare la ricorrenza $d_n = 2d_{n-1} - d_{n-2}$ con $d_0=1$, $d_1=2$.]
\end{esercizio}

\begin{esercizio}
Sia $A \in M_{5,1}(K)$ e $B \in M_{1,5}(K)$. Calcolare $\det(A\cdot B)$. [\emph{Suggerimento}: pensare al rango di $A\cdot B$.]
\end{esercizio}

\begin{esercizio}[Autovettori e diagonalizzazione]
Sia $A = \begin{pmatrix}3&-2&1\\0&1&1\\4&-1&-1\end{pmatrix} \in M_{3,3}(\R)$.
\begin{enumerate}
  \item Calcolate $p_A(\lambda)$.
  \item Determinate autovalori, autospazi e molteplicità geometriche e algebriche.
  \item $A$ è diagonalizzabile? Se sì, trovate una base diagonalizzante.
\end{enumerate}
\end{esercizio}

\begin{esercizio}
Sia $K$ un campo con $\mathrm{char}\,K \neq 2$ e $n$ dispari. Dimostrate che se $A \in M_{n,n}(K)$ è antisimmetrica ($A^T = -A$) allora $\det A = 0$. [\emph{Suggerimento}: usare $\det(A^T) = \det A$ e la relazione tra $\det A$ e $\det(-A)$.]
\end{esercizio}

\begin{esercizio}[Formula determinantale per sottospazi]
\label{es:ann-det}
Siano $A^1,\ldots,A^{n-1} \in K^n$ vettori linearmente indipendenti (scritti come vettori riga) e $A \in M_{n-1,n}(K)$ la matrice con queste righe. Per $1 \leq j \leq n$ sia $M_j$ la matrice $(n-1)\times(n-1)$ ottenuta eliminando la colonna $j$-esima da $A$. Dimostrate che il vettore
\[
  f \colon K^n \to K, \quad X \mapsto \sum_{j=1}^n (-1)^j (\det M_j)\,x_j
\]
genera $\Ann\langle A^1,\ldots,A^{n-1}\rangle$. [\emph{Suggerimento}: mostrare che $f$ è non nulla, che $f(A^i)=0$ per ogni $i$, e usare la formula $\dim\Ann(V) = n - \dim V$ con $\dim V = n-1$.]
\end{esercizio}

\begin{esercizio}[PageRank elementare]
Siano $4$ pagine web con link come in Figura: la pagina $1$ ha link verso $2$; la pagina $2$ verso $3$ e $4$; la pagina $3$ verso $1$; la pagina $4$ verso $1$, $2$ e $3$.
\begin{enumerate}
  \item Scrivere la matrice stocastica $A$ del PageRank.
  \item Verificare che $1$ è autovalore di $A$ (usando la Proposizione~\ref{prop:det-invertibilita}).
  \item Determinare l'autospazio $V_1(L_A)$ e il ranking delle pagine.
\end{enumerate}
\end{esercizio}

\begin{esercizio}[Guidato: $\det f$ tramite base adattata]
Sia $V = \{X \in \R^4 \mid x_1+x_2+x_3+x_4=0\}$ e
\[
  M \coloneqq \begin{pmatrix}2&1&0&-1\\1&0&-1&2\\0&-1&2&1\\-1&2&1&0\end{pmatrix}.
\]
\begin{enumerate}
  \item Verificate che $L_M(e^\vee_1+e^\vee_2+e^\vee_3+e^\vee_4) = 2(e^\vee_1+e^\vee_2+e^\vee_3+e^\vee_4)$ e deducete che $L_M(V)\subset V$.
  \item Sia $f = L_M\vert_V \colon V \to V$. Calcolate $\det f$ usando la definizione (trovare una base di $V$ e la matrice di $f$ rispetto ad essa).
  \item Calcolate invece $\det M$ tramite l'algoritmo di riduzione a scala, poi deducete $\det f$ usando la decomposizione $\R^4 = \langle(1,1,1,1)\rangle \oplus V$ e la moltiplicatività del determinante.
\end{enumerate}
\end{esercizio}

\begin{esercizio}
Sia $V$ uno spazio vettoriale reale di dimensione $n$ finita. Sia $g \in \GL(V)$ con $\det g < 0$. Dimostrate che il gruppo $\GL^+(V) = \{h \in \GL(V) \mid \det h > 0\}$ è un sottogruppo normale di $\GL(V)$ di indice $2$, e che ogni elemento di $\GL(V) \setminus \GL^+(V)$ è coniugato a $g$.
\end{esercizio}

% ============================================================
\chapter{Spazi vettoriali euclidei ed hermitiani}
\label{cap:euclidei}
% ============================================================

\section*{Introduzione al capitolo}

Fino a questo punto gli spazi vettoriali sono stati strutture puramente algebriche: somme, combinazioni lineari, basi. Mancava qualcosa di essenziale per fare geometria: la nozione di \emph{lunghezza} e di \emph{angolo}. Questo capitolo introduce esattamente questi concetti, formalizzandoli nella nozione di \emph{prodotto scalare euclideo}.

L'idea di partenza è semplice. Nel piano euclideo $\mathbf{E}^2$ sappiamo misurare lunghezze di segmenti e angoli tra rette. Il prodotto scalare $v \cdot w = \|v\|\,\|w\|\cos\theta$ codifica entrambe le informazioni in un'unica operazione bilineare. Le sue proprietà — bilinearità, simmetria, positività — sono la chiave per generalizzarla a spazi astratti.

Una volta fissato un prodotto scalare euclideo, lo spazio si arricchisce enormemente: si possono definire angoli, proiezioni, basi ortonormali. L'algoritmo di Gram-Schmidt mostra come costruire sistematicamente basi ortonormali a partire da basi qualsiasi. La decomposizione ortogonale $V = U \oplus U^\perp$ è lo strumento geometrico fondamentale del capitolo: ogni vettore si scompone univocamente in una parte parallela e una perpendicolare a un dato sottospazio.

La seconda parte del capitolo estende la teoria al caso complesso (spazi hermitiani) e studia le isometrie: le applicazioni lineari che preservano il prodotto scalare. Il gruppo ortogonale $O(V)$ e il sottogruppo speciale $SO(V)$ sono i protagonisti, e il teorema spettrale per le isometrie — ogni elemento di $O(V)$ decompone $V$ in sottospazi invarianti di dimensione al più $2$, su ciascuno dei quali agisce come rotazione — è uno dei risultati più belli dell'algebra lineare.

% --------------------------------------------------------
\section{Prodotto scalare euclideo}
\label{sec:prodotto-scalare}
% --------------------------------------------------------

\subsection{Definizione e primi esempi}

\begin{defbox}
\begin{definizione}[Prodotto scalare euclideo]
\label{def:prodotto-scalare}
Sia $V$ uno spazio vettoriale reale. Un \emph{prodotto scalare euclideo} su $V$ è una forma bilineare
\[
  \langle\cdot,\cdot\rangle \colon V \times V \longrightarrow \R
\]
che soddisfa:
\begin{enumerate}
  \item \textbf{Simmetria:} $\langle v,w\rangle = \langle w,v\rangle$ per ogni $v,w \in V$.
  \item \textbf{Positività:} $\langle v,v\rangle \geq 0$ per ogni $v \in V$, con uguaglianza se e solo se $v = 0$.
\end{enumerate}
Uno \emph{spazio vettoriale euclideo} è una coppia $(V,\langle\cdot,\cdot\rangle)$ dove $V$ è uno spazio vettoriale reale e $\langle\cdot,\cdot\rangle$ è un prodotto scalare euclideo.
\end{definizione}
\end{defbox}

\begin{osservazione}
Un prodotto scalare euclideo è in particolare una forma bilineare simmetrica. La condizione aggiuntiva di positività lo rende uno strumento geometrico: dice che ogni vettore non nullo ha ``lunghezza'' positiva.
\end{osservazione}

\begin{exbox}
\begin{esempio}[Prodotto euclideo standard su $\R^n$]
\label{es:prodotto-standard}
L'applicazione
\[
  \langle X, Y\rangle \coloneqq X^T \cdot Y = \sum_{i=1}^n x_i y_i
\]
è un prodotto scalare euclideo su $\R^n$, detto \emph{prodotto euclideo standard}. Infatti è bilineare, simmetrica, e $\langle X,X\rangle = \sum x_i^2 \geq 0$ con uguaglianza solo se $X=0$.
\end{esempio}
\end{exbox}

\begin{exbox}
\begin{esempio}[Prodotto scalare su funzioni continue]
\label{es:prodotto-funzioni}
Sia $C^0([-\pi,\pi])$ lo spazio vettoriale delle funzioni continue $f \colon [-\pi,\pi] \to \R$. L'applicazione
\[
  \langle f,g\rangle \coloneqq \frac{1}{\pi}\int_{-\pi}^{\pi} f(t)g(t)\,dt
\]
è un prodotto scalare euclideo. La bilinearità e la simmetria sono ovvie. La positività segue dal fatto che $\langle f,f\rangle = \frac{1}{\pi}\int_{-\pi}^\pi f(t)^2\,dt \geq 0$, con uguaglianza solo se $f^2 \equiv 0$, cioè $f \equiv 0$ (per continuità).
\end{esempio}
\end{exbox}

\begin{defbox}
\begin{definizione}[Norma]
Sia $(V,\langle\cdot,\cdot\rangle)$ uno spazio vettoriale euclideo. La \emph{norma} di $v \in V$ è
\[
  \|v\| \coloneqq \langle v,v\rangle^{1/2} \geq 0.
\]
\end{definizione}
\end{defbox}

\begin{exbox}
\begin{esempio}
In $C^0([-\pi,\pi])$ con il prodotto dell'Esempio~\ref{es:prodotto-funzioni}, per $m \in \Z^*$:
\[
  \|\cos mt\| = \|\sin mt\| = 1, \qquad \|1\| = \sqrt{2}.
\]
\end{esempio}
\end{exbox}

La norma determina completamente il prodotto scalare, grazie alla seguente identità.

\begin{proposizione}[Identità di polarizzazione]
\label{prop:polarizzazione}
In uno spazio vettoriale euclideo $(V,\langle\cdot,\cdot\rangle)$:
\[
  \langle v,w\rangle = \frac{1}{2}\bigl(\|v+w\|^2 - \|v\|^2 - \|w\|^2\bigr).
\]
\end{proposizione}

\begin{proof}
Espandere $\|v+w\|^2 = \langle v+w,v+w\rangle = \|v\|^2 + 2\langle v,w\rangle + \|w\|^2$.
\end{proof}

% --------------------------------------------------------
\section{Cauchy-Schwarz e disuguaglianza triangolare}
\label{sec:cauchy-schwarz}
% --------------------------------------------------------

\begin{teorema}[Cauchy-Schwarz]
\label{thm:cauchy-schwarz}
Sia $(V,\langle\cdot,\cdot\rangle)$ uno spazio vettoriale euclideo. Per ogni $v,w \in V$:
\[
  \langle v,w\rangle^2 \leq \|v\|^2 \cdot \|w\|^2,
\]
con uguaglianza se e solo se $v,w$ sono linearmente dipendenti.
\end{teorema}

\begin{proof}
Se $v=0$ la tesi è banale. Sia $v \neq 0$ e consideriamo il polinomio $p(x) \coloneqq \|xv+w\|^2 = \|v\|^2 x^2 + 2\langle v,w\rangle x + \|w\|^2$. Poiché $p(x) \geq 0$ per ogni $x \in \R$ (è una norma al quadrato), il discriminante è $\leq 0$:
\[
  4\langle v,w\rangle^2 - 4\|v\|^2\|w\|^2 \leq 0.
\]
Il discriminante è $0$ se e solo se $p(x)=0$ ha una radice reale $x_0$, cioè $x_0 v + w = 0$: poiché $v \neq 0$, questo equivale alla dipendenza lineare di $v$ e $w$.
\end{proof}

\begin{corollario}[Disuguaglianza triangolare]
\label{cor:triangolare}
$\|v+w\| \leq \|v\| + \|w\|$ per ogni $v,w \in V$, con uguaglianza se e solo se esiste $\lambda \geq 0$ tale che $v = \lambda w$ o $w = \lambda v$.
\end{corollario}

\begin{proof}
$\|v+w\|^2 = \|v\|^2 + 2\langle v,w\rangle + \|w\|^2 \leq \|v\|^2 + 2\|v\|\|w\| + \|w\|^2 = (\|v\|+\|w\|)^2$.
\end{proof}

\begin{osservazione}
Per $\R^n$ con il prodotto standard, Cauchy-Schwarz diventa la classica disuguaglianza:
\[
  \left(\sum_{i=1}^n x_i y_i\right)^2 \leq \left(\sum_{i=1}^n x_i^2\right)\left(\sum_{j=1}^n y_j^2\right).
\]
La disuguaglianza triangolare in $\mathbf{E}^m$ afferma $d(P,R) \leq d(P,Q) + d(Q,R)$: da qui il nome.
\end{osservazione}

La disuguaglianza di Cauchy-Schwarz permette di definire l'angolo in ogni spazio euclideo.

\begin{defbox}
\begin{definizione}[Angolo e ortogonalità]
\label{def:angolo}
Sia $(V,\langle\cdot,\cdot\rangle)$ uno spazio vettoriale euclideo.
\begin{enumerate}
  \item Per $v,w \neq 0$, l'\emph{angolo} tra $v$ e $w$ è l'unico $\theta \in [0,\pi]$ tale che
  $\displaystyle\cos\theta = \frac{\langle v,w\rangle}{\|v\|\,\|w\|}$.
  \item Vettori $v,w \in V$ sono \emph{ortogonali} (o \emph{perpendicolari}) se $\langle v,w\rangle = 0$. In simboli: $v \perp w$.
  \item Per $S \subseteq V$, il \emph{complemento ortogonale} di $S$ è $S^\perp \coloneqq \{w \in V \mid \langle v,w\rangle = 0\;\forall\,v \in S\}$.
\end{enumerate}
\end{definizione}
\end{defbox}

\begin{osservazione}
La definizione di angolo ha senso perché Cauchy-Schwarz garantisce $-1 \leq \frac{\langle v,w\rangle}{\|v\|\|w\|} \leq 1$.
\end{osservazione}

\begin{exbox}
\begin{esempio}
Sia $(C^0([-\pi,\pi]),\langle\cdot,\cdot\rangle)$ l'Esempio~\ref{es:prodotto-funzioni}. Per $m,n \in \N$, un calcolo con le formule di Eulero dà:
\[
  \int_{-\pi}^{\pi}\cos mt\,\cos nt\,dt = \pi\delta_{mn},\quad
  \int_{-\pi}^{\pi}\sin mt\,\sin nt\,dt = \pi\delta_{mn},\quad
  \int_{-\pi}^{\pi}\cos mt\,\sin nt\,dt = 0.
\]
Quindi le funzioni $\cos mt$ e $\sin mt$ per $m \in \N$ sono a due a due ortogonali.
\end{esempio}
\end{exbox}

% --------------------------------------------------------
\section{Basi ortonormali}
\label{sec:basi-ON}
% --------------------------------------------------------

\begin{defbox}
\begin{definizione}[Lista ortonormale e base ON]
\label{def:base-ON}
Sia $(V,\langle\cdot,\cdot\rangle)$ uno spazio vettoriale euclideo. Una lista $\{v_1,\ldots,v_n\}$ è \emph{ortonormale} (ON) se $\langle v_i,v_j\rangle = \delta_{ij}$ per ogni $i,j$. Una \emph{base ortonormale} (base ON) è una base che è anche ON.
\end{definizione}
\end{defbox}

\begin{exbox}
\begin{esempio}
La base standard $\{e_1,\ldots,e_n\}$ di $\R^n$ è una base ON per il prodotto euclideo standard.
\end{esempio}
\end{exbox}

\begin{proposizione}
\label{prop:ON-indipendenti}
Una lista ortonormale è linearmente indipendente.
\end{proposizione}

\begin{proof}
Sia $\sum_{i=1}^n \lambda_i v_i = 0$. Per ogni $k$, facendo il prodotto scalare con $v_k$: $\lambda_k = \langle \sum_i \lambda_i v_i, v_k\rangle = 0$.
\end{proof}

\begin{proposizione}[Calcolo del prodotto scalare in una base ON]
\label{prop:calcolo-ON}
Se $\{v_1,\ldots,v_n\}$ è una lista ON, allora
\[
  \left\langle\sum_{i=1}^n x_i v_i,\; \sum_{i=1}^n y_i v_i\right\rangle = \sum_{i=1}^n x_i y_i.
\]
\end{proposizione}

\begin{proof}
Per bilinearità: $\langle \sum_i x_i v_i, \sum_j y_j v_j\rangle = \sum_{i,j} x_i y_j \langle v_i,v_j\rangle = \sum_{i,j} x_i y_j \delta_{ij} = \sum_i x_i y_i$.
\end{proof}

\begin{proposizione}[Esistenza di basi ON]
\label{prop:esistenza-base-ON}
Ogni spazio vettoriale euclideo finitamente generato ammette una base ortonormale.
\end{proposizione}

\begin{proof}
Per induzione su $\dim V$. Se $\dim V = 0$, banale. Se $\dim V = 1$, sia $\{v_1\}$ una base qualsiasi e poniamo $w_1 \coloneqq v_1/\|v_1\|$. Per il passo induttivo: sia $w_1 \coloneqq v_1/\|v_1\|$ per un qualsiasi $v_1 \neq 0$. L'applicazione $\varphi \colon V \to \R$, $v \mapsto \langle v,w_1\rangle$, è lineare e suriettiva (poiché $\varphi(w_1)=1$), quindi $\ker\varphi = w_1^\perp$ ha dimensione $\dim V - 1$. La restrizione di $\langle\cdot,\cdot\rangle$ a $w_1^\perp$ è un prodotto scalare euclideo, e per ipotesi induttiva ammette una base ON $\{w_2,\ldots,w_n\}$. La lista $\{w_1,\ldots,w_n\}$ è ON (ogni $w_i$ per $i\geq 2$ è ortogonale a $w_1$ per costruzione) e ha $n$ elementi: è una base ON di $V$.
\end{proof}

% --------------------------------------------------------
\section{Decomposizione ortogonale}
\label{sec:decomp-ortogonale}
% --------------------------------------------------------

\begin{lemma}
$S^\perp$ è un sottospazio di $V$ per ogni $S \subseteq V$. In particolare, se $U \subseteq V$ è un sottospazio finitamente generato da $u_1,\ldots,u_m$, allora $U^\perp = \bigcap_{i=1}^m u_i^\perp$.
\end{lemma}

\begin{proof}
$S^\perp = \bigcap_{v \in S}\ker(\langle\cdot,v\rangle)$ è intersezione di nuclei di applicazioni lineari, quindi è un sottospazio. La seconda parte segue dall'uguaglianza $u^\perp_1 \cap \cdots \cap u^\perp_m = \{v \mid \langle v,u_i\rangle = 0\;\forall i\}$, che vale per bilinearità.
\end{proof}

\begin{proposizione}[Decomposizione ortogonale]
\label{prop:decomp-ortogonale}
Sia $(V,\langle\cdot,\cdot\rangle)$ uno spazio vettoriale euclideo e $U \subseteq V$ un sottospazio finitamente generato. Allora
\[
  V = U \oplus U^\perp.
\]
\end{proposizione}

\begin{proof}
Mostriamo $U \cap U^\perp = \{0\}$ e $V = U + U^\perp$.

\emph{Intersezione banale.} Se $v \in U \cap U^\perp$, allora $\langle v,v\rangle = 0$, quindi $v=0$.

\emph{Generazione.} Sia $\{u_1,\ldots,u_m\}$ una base ON di $U$ (esiste per la Proposizione~\ref{prop:esistenza-base-ON}). Definiamo la \emph{proiezione ortogonale} su $U$:
\[
  \pi_U(v) \coloneqq \sum_{i=1}^m \langle v,u_i\rangle\, u_i.
\]
Per ogni $j$: $\langle v - \pi_U(v), u_j\rangle = \langle v,u_j\rangle - \langle v,u_j\rangle = 0$, quindi $v - \pi_U(v) \in U^\perp$. La decomposizione $v = \pi_U(v) + (v - \pi_U(v))$ dà $v \in U + U^\perp$.
\end{proof}

\begin{osservazione}
La proiezione $\pi_U$ è lineare e non dipende dalla scelta della base ON di $U$ (è l'unico elemento di $U$ tale che $v - \pi_U(v) \in U^\perp$).
\end{osservazione}

\begin{defbox}
\begin{definizione}[Proiezione ortogonale]
Sia $U \subseteq V$ un sottospazio finitamente generato. La \emph{proiezione ortogonale} su $U$ è l'applicazione lineare $\pi_U \colon V \to U$ definita da $\pi_U(v) \coloneqq \sum_{i=1}^m \langle v,u_i\rangle u_i$, dove $\{u_1,\ldots,u_m\}$ è una qualsiasi base ON di $U$.
\end{definizione}
\end{defbox}

\begin{corollario}
\label{cor:ortogonale-doppio}
$(U^\perp)^\perp = U$ e $\dim U^\perp = \dim V - \dim U$.
\end{corollario}

\begin{proof}
$U \subseteq (U^\perp)^\perp$ è ovvio. Per dimensioni: $\dim U + \dim U^\perp = \dim V$ dalla Proposizione~\ref{prop:decomp-ortogonale}, e analogamente $\dim U^\perp + \dim(U^\perp)^\perp = \dim V$, quindi $\dim(U^\perp)^\perp = \dim U$. Poiché $U \subseteq (U^\perp)^\perp$ e hanno la stessa dimensione, coincidono.
\end{proof}

\begin{defbox}
\begin{definizione}[Somma diretta ortogonale]
$(V,\langle\cdot,\cdot\rangle)$ è \emph{somma diretta ortogonale} di $W_1,\ldots,W_r$ se $V = W_1 \oplus \cdots \oplus W_r$ e $W_i \perp W_j$ per ogni $i \neq j$. Notazione: $V = W_1 \oplus^\perp \cdots \oplus^\perp W_r$.
\end{definizione}
\end{defbox}

\begin{proposizione}[Minima distanza]
\label{prop:minima-distanza}
$\|v - \pi_U(v)\| \leq \|v - u\|$ per ogni $u \in U$, con uguaglianza se e solo se $u = \pi_U(v)$.
\end{proposizione}

\begin{proof}
$\|v-u\|^2 = \|v - \pi_U(v) - (u-\pi_U(v))\|^2 = \|v-\pi_U(v)\|^2 + \|u-\pi_U(v)\|^2 \geq \|v-\pi_U(v)\|^2$, con uguaglianza iff $u = \pi_U(v)$.
\end{proof}

\begin{exbox}
\begin{esempio}[Serie di Fourier e identità di Parseval]
\label{es:fourier}
In $C^0([-\pi,\pi])$, sia $U_n$ il sottospazio generato da $\frac{1}{\sqrt{2}}, \cos t, \sin t, \ldots, \cos nt, \sin nt$. Per $\varphi(t) = t$, la proiezione su $U_n$ è
\[
  \varphi_n(t) = 2\sum_{m=1}^n \frac{(-1)^{m+1}}{m}\sin mt.
\]
Dalla decomposizione ortogonale $\|\varphi\|^2 \geq \|\varphi_n\|^2$, cioè $\frac{2\pi^2}{3} \geq 4\sum_{m=1}^n \frac{1}{m^2}$. Il fatto (non banale) che la differenza tende a $0$ dà la celebre formula di Eulero:
\[
  \sum_{m=1}^\infty \frac{1}{m^2} = \frac{\pi^2}{6}.
\]
\end{esempio}
\end{exbox}

% --------------------------------------------------------
\section{Algoritmo di Gram-Schmidt}
\label{sec:gram-schmidt}
% --------------------------------------------------------

L'algoritmo di Gram-Schmidt costruisce una base ortogonale (o ON) a partire da vettori linearmente indipendenti, mantenendo invariati i sottospazi generati.

\begin{defbox}
\begin{definizione}[Algoritmo di Gram-Schmidt]
Siano $v_1,\ldots,v_m$ vettori linearmente indipendenti in $(V,\langle\cdot,\cdot\rangle)$. Poniamo $w_1 \coloneqq v_1$ e, per $s = 2,\ldots,m$:
\[
  w_s \coloneqq v_s - \pi_{U_{s-1}}(v_s) = v_s - \sum_{i=1}^{s-1}\frac{\langle v_s, w_i\rangle}{\|w_i\|^2}\,w_i,
\]
dove $U_{s-1} = \Span(v_1,\ldots,v_{s-1})$.
\end{definizione}
\end{defbox}

\begin{proposizione}
\label{prop:gram-schmidt}
Per ogni $s \in \{1,\ldots,m\}$, la lista $\{w_1,\ldots,w_s\}$ è una base ortogonale di $U_s = \Span(v_1,\ldots,v_s)$.
\end{proposizione}

\begin{proof}
Per induzione su $s$. Se $s=1$: $w_1 = v_1 \neq 0$, e $\{w_1\}$ è una base di $U_1$. Passo induttivo: per ipotesi $\{w_1,\ldots,w_s\}$ è una base ortogonale di $U_s$. Allora $w_{s+1} = v_{s+1} - \pi_{U_s}(v_{s+1})$ è ortogonale a tutti i $w_i$ per $i \leq s$ per costruzione. È non nullo perché $v_{s+1} \notin U_s$ (indipendenza lineare), e genera $U_{s+1}$ insieme a $\{w_1,\ldots,w_s\}$.
\end{proof}

Per ottenere una base ON, si normalizza: $w'_i \coloneqq w_i/\|w_i\|$.

\begin{exbox}
\begin{esempio}[Gram-Schmidt in $\R^3$]
Sia $(V,\langle\cdot,\cdot\rangle)$ uno spazio euclideo con base ON $\{u_1,u_2,u_3\}$. Siano
\[
  v_1 = u_1 - 2u_2 + 2u_3,\quad v_2 = 3u_1 - 12u_2 - 9u_3.
\]
\textbf{Passo 1.} $w_1 = v_1$, con $\|w_1\|^2 = 1+4+4=9$.

\textbf{Passo 2.} $\langle v_2,w_1\rangle = 3(-1)(1)+(-12)(-2)+(-9)(2) = \ldots$. Usando la Proposizione~\ref{prop:calcolo-ON}: $\langle v_2,w_1\rangle = 3\cdot1 + (-12)(-2) + (-9)(2) = 3+24-18 = 9$. Quindi:
\[
  w_2 = v_2 - \frac{9}{9}w_1 = (3u_1-12u_2-9u_3)-(u_1-2u_2+2u_3) = 2u_1-10u_2-11u_3.
\]
Normalizzando: $w'_1 = \frac{1}{3}(u_1-2u_2+2u_3)$, $w'_2 = \frac{1}{15}(2u_1-10u_2-11u_3)$.
\end{esempio}
\end{exbox}

% --------------------------------------------------------
\section{Matrice di Gram e forme bilineari}
\label{sec:gram}
% --------------------------------------------------------

\begin{defbox}
\begin{definizione}[Matrice di Gram]
\label{def:gram}
Sia $(V,\langle\cdot,\cdot\rangle)$ uno spazio vettoriale euclideo e $\mathcal{B} = \{v_1,\ldots,v_n\}$ una lista di vettori. La \emph{matrice di Gram} associata a $\mathcal{B}$ è la matrice $M_{\mathcal{B}}(\langle\cdot,\cdot\rangle) \in M_{n,n}(\R)$ con entrata $(i,j)$ uguale a $\langle v_i,v_j\rangle$.
\end{definizione}
\end{defbox}

\begin{osservazione}
La matrice di Gram è simmetrica (per la simmetria del prodotto scalare) e, se $\mathcal{B}$ è una base ON, coincide con $1_n$.
\end{osservazione}

\subsection{Corrispondenza tra matrici simmetriche e forme bilineari}

Sia $V$ uno spazio vettoriale finitamente generato su $K$ con base $\mathcal{B} = \{v_1,\ldots,v_n\}$. Per ogni $A \in M_{n,n}(K)$ la forma
\[
  \Phi^{\mathcal{B}}_A(v,w) \coloneqq X_{\mathcal{B}}(v)^T \cdot A \cdot X_{\mathcal{B}}(w)
\]
è bilineare, e si ha $\Phi^{\mathcal{B}}_A(v_i,v_j) = a_{ij}$.

\begin{proposizione}
\label{prop:bilineare-matrice}
L'applicazione $A \mapsto \Phi^{\mathcal{B}}_A$ è una biiezione $M_{n,n}(K) \xrightarrow{\sim} \mathrm{Bil}(V)$, e manda le matrici simmetriche nelle forme bilineari simmetriche.
\end{proposizione}

\begin{proof}
L'\emph{iniettività} segue da $a_{ij} = \Phi^{\mathcal{B}}_A(v_i,v_j)$. La \emph{suriettività}: data $F \in \mathrm{Bil}(V)$, poniamo $a_{ij} \coloneqq F(v_i,v_j)$ e si verifica $\Phi^{\mathcal{B}}_A = F$ per bilinearità. La parte sulle simmetriche: $\Phi^{\mathcal{B}}_A$ è simmetrica iff $a_{ij} = a_{ji}$ per ogni $i,j$ iff $A = A^T$.
\end{proof}

La matrice di una forma bilineare dipende dalla base. La legge di trasformazione è:

\begin{proposizione}
\label{prop:cambio-base-bilineare}
Se $\mathcal{B}$ e $\mathcal{C}$ sono basi di $V$ e $F \in \mathrm{Bil}(V)$, allora
\[
  M_{\mathcal{C}}(F) = (M^{\mathcal{C}}_{\mathcal{B}}(\Id))^T \cdot M_{\mathcal{B}}(F) \cdot M^{\mathcal{C}}_{\mathcal{B}}(\Id).
\]
In particolare, $\det M_{\mathcal{C}}(F) = (\det M^{\mathcal{C}}_{\mathcal{B}}(\Id))^2 \cdot \det M_{\mathcal{B}}(F)$.
\end{proposizione}

\begin{proof}
Per ogni $v,w \in V$: $F(v,w) = X_{\mathcal{B}}(v)^T \cdot M_{\mathcal{B}}(F) \cdot X_{\mathcal{B}}(w) = (M^{\mathcal{C}}_{\mathcal{B}}(\Id)\cdot X_{\mathcal{C}}(v))^T \cdot M_{\mathcal{B}}(F) \cdot M^{\mathcal{C}}_{\mathcal{B}}(\Id)\cdot X_{\mathcal{C}}(w)$, da cui la formula. La formula per i determinanti segue da Binet e $\det(M^T) = \det M$.
\end{proof}

\begin{corollario}
\label{cor:gram-det-positivo}
Sia $(V,\langle\cdot,\cdot\rangle)$ uno spazio vettoriale euclideo finitamente generato. Per ogni base $\mathcal{C}$ di $V$, $\det M_{\mathcal{C}}(\langle\cdot,\cdot\rangle) > 0$.
\end{corollario}

\begin{proof}
Sia $\mathcal{B}$ una base ON (esiste per la Proposizione~\ref{prop:esistenza-base-ON}). Allora $M_{\mathcal{B}}(\langle\cdot,\cdot\rangle) = 1_n$, quindi $\det M_{\mathcal{B}}(\langle\cdot,\cdot\rangle) = 1 > 0$. Per la Proposizione~\ref{prop:cambio-base-bilineare}, $\det M_{\mathcal{C}}(\langle\cdot,\cdot\rangle) = (\det M^{\mathcal{C}}_{\mathcal{B}}(\Id))^2 > 0$.
\end{proof}

\begin{osservazione}
Da questo corollario si ricava un'ulteriore dimostrazione della disuguaglianza di Cauchy-Schwarz: per $v,w$ linearmente indipendenti, la matrice di Gram $\begin{pmatrix}\langle v,v\rangle & \langle v,w\rangle\\\langle w,v\rangle & \langle w,w\rangle\end{pmatrix}$ ha determinante $>0$, cioè $\|v\|^2\|w\|^2 - \langle v,w\rangle^2 > 0$.
\end{osservazione}

\subsection{Criterio di Sylvester}

\begin{defbox}
\begin{definizione}[Minori principali]
Per $A \in M_{n,n}(K)$ e $p \in \{1,\ldots,n\}$, indichiamo con $A^{(p)}$ la sottomatrice $p \times p$ formata dalle prime $p$ righe e colonne di $A$.
\end{definizione}
\end{defbox}

\begin{proposizione}[Criterio di Sylvester]
\label{prop:sylvester}
Sia $V$ uno spazio vettoriale reale finitamente generato con base $\mathcal{B} = \{v_1,\ldots,v_n\}$ e $A \in M^+_{n,n}(\R)$ simmetrica. La forma bilineare simmetrica $\Phi^{\mathcal{B}}_A$ è un prodotto scalare euclideo se e solo se $\det A^{(p)} > 0$ per ogni $p \in \{1,\ldots,n\}$.
\end{proposizione}

\begin{proof}
\emph{Necessità.} Se $\Phi^{\mathcal{B}}_A$ è un prodotto scalare euclideo, allora la sua restrizione a $U_p = \Span(v_1,\ldots,v_p)$ è un prodotto scalare euclideo con matrice di Gram $A^{(p)}$, e per il Corollario~\ref{cor:gram-det-positivo} $\det A^{(p)} > 0$.

\emph{Sufficienza.} Per induzione su $n$. Se $n=1$: $\det A^{(1)} = a_{11} > 0$ e $\Phi_A(v,v) = a_{11}x_1^2 > 0$ per $v \neq 0$. Passo induttivo: la restrizione di $F = \Phi^{\mathcal{B}}_A$ a $U_{n-1}$ è un prodotto scalare euclideo per ipotesi induttiva, con decomposizione ortogonale $V = U_{n-1} \oplus^\perp U_{n-1}^\perp$. Poiché $\dim U_{n-1}^\perp = 1$, sia $w$ una base di $U_{n-1}^\perp$. Dalla struttura a blocchi della matrice di Gram in una base adattata: $\det M_{\mathcal{D}}(F) = \det A^{(n-1)} \cdot F(w,w)$. Questo determante ha lo stesso segno di $\det A^{(n)} > 0$, e poiché $\det A^{(n-1)} > 0$, si ha $F(w,w) > 0$. Per ogni $v = u + \lambda w$ non nullo: $F(v,v) = F(u,u) + \lambda^2 F(w,w) > 0$.
\end{proof}

\begin{exbox}
\begin{esempio}[Esercizio svolto]
\label{es:sylvester}
Sia $A = \begin{pmatrix}3&1&0\\1&1&2\\0&2&7\end{pmatrix}$.

\textbf{(a) Verificare che $\langle X,Y\rangle \coloneqq X^T A Y$ è un prodotto scalare su $\R^3$.}

I minori principali: $\det A^{(1)} = 3 > 0$, $\det A^{(2)} = 3\cdot1-1\cdot1 = 2 > 0$, $\det A^{(3)} = 3(7-4)-1(7-0)+0 = 9-7 = 2 > 0$. Per il criterio di Sylvester, sì.

\textbf{(b) Trovare la proiezione ortogonale di $v=(1,1,1)$ su $U = \{x_1+x_2+x_3=0\}$.}

Una base di $U$: $\{(1,-1,0),(1,0,-1)\}$. Applichiamo Gram-Schmidt con il prodotto $\langle\cdot,\cdot\rangle_A$:
\[
  \langle(1,-1,0),(1,-1,0)\rangle_A = (1,-1,0)A(1,-1,0)^T = 2,
\]
\[
  \langle(1,0,-1),(1,-1,0)\rangle_A = 4, \quad w_2 = (1,0,-1) - \tfrac{4}{2}(1,-1,0) = (-1,2,-1).
\]
Quindi $\langle v,w_1\rangle_A = \langle(1,1,1),(1,-1,0)\rangle_A = 2$, $\langle v,w_2\rangle_A = \langle(1,1,1),(-1,2,-1)\rangle_A = -5$. La proiezione è:
\[
  \pi_U(v) = \frac{2}{2}(1,-1,0) + \frac{-5}{\|w_2\|^2_A}(-1,2,-1).
\]
\end{esempio}
\end{exbox}

% --------------------------------------------------------
\section{Isometrie di spazi vettoriali euclidei}
\label{sec:isometrie}
% --------------------------------------------------------

\begin{defbox}
\begin{definizione}[Isometria]
\label{def:isometria}
Siano $(V,\langle\cdot,\cdot\rangle_V)$ e $(W,\langle\cdot,\cdot\rangle_W)$ spazi vettoriali euclidei. Un'applicazione lineare $f \colon V \to W$ è un'\emph{isometria} se
\[
  \langle v_1,v_2\rangle_V = \langle f(v_1),f(v_2)\rangle_W \quad\forall\,v_1,v_2 \in V.
\]
Un'\emph{isometria suriettiva} si dice \emph{isomorfismo di spazi vettoriali euclidei}.
\end{definizione}
\end{defbox}

\begin{proposizione}
\label{prop:isometria-norma}
$f$ è un'isometria se e solo se $\|v\|_V = \|f(v)\|_W$ per ogni $v \in V$.
\end{proposizione}

\begin{proof}
La necessità è ovvia. La sufficienza segue dall'identità di polarizzazione (Proposizione~\ref{prop:polarizzazione}): $\langle v_1,v_2\rangle_V = \frac{1}{2}(\|v_1+v_2\|^2 - \|v_1\|^2 - \|v_2\|^2)$, e il membro di destra coincide con $\langle f(v_1),f(v_2)\rangle_W$ applicando l'ipotesi a $v_1$, $v_2$, $v_1+v_2$.
\end{proof}

\begin{corollario}
Ogni isometria è iniettiva. Ogni isometria di uno spazio finitamente generato in sé è suriettiva (quindi un isomorfismo).
\end{corollario}

\begin{exbox}
\begin{esempio}
Se $\{v_1,\ldots,v_n\}$ è una base ON di $(V,\langle\cdot,\cdot\rangle)$, l'applicazione $f \colon \R^n \to V$, $X \mapsto \sum_i x_i v_i$, è un'isometria (con il prodotto standard su $\R^n$): $\|f(X)\|^2 = \|\sum x_i v_i\|^2 = \sum x_i^2 = \|X\|^2$.
\end{esempio}
\end{exbox}

\begin{corollario}
Spazi vettoriali euclidei finitamente generati della stessa dimensione sono isomorfi (come spazi vettoriali euclidei).
\end{corollario}

\begin{defbox}
\begin{definizione}[Gruppo ortogonale]
\label{def:gruppo-ortogonale}
Sia $(V,\langle\cdot,\cdot\rangle)$ uno spazio vettoriale euclideo. Il \emph{gruppo ortogonale} è
\[
  O(V,\langle\cdot,\cdot\rangle) \coloneqq \{f \colon V \to V \mid f \text{ è un'isometria}\}.
\]
Il \emph{gruppo ortogonale speciale} è $SO(V) \coloneqq \{f \in O(V) \mid \det f = 1\}$.
\end{definizione}
\end{defbox}

\begin{proposizione}
\label{prop:ortogonale-matrice}
$A \in \GL_n(\R)$ appartiene a $O_n(\R)$ se e solo se $A^T A = 1_n$. In tal caso $\det A = \pm 1$.
\end{proposizione}

\begin{proof}
$\langle AX,AY\rangle = (AX)^T(AY) = X^T A^T A Y$. Questo è uguale a $X^T Y$ per ogni $X,Y$ iff $A^T A = 1_n$. Per Binet: $1 = \det(A^T A) = (\det A)^2$, quindi $\det A = \pm 1$.
\end{proof}

% --------------------------------------------------------
\section{Spazi vettoriali hermitiani}
\label{sec:hermitiani}
% --------------------------------------------------------

Passando a spazi vettoriali complessi, la nozione di prodotto scalare si modifica: la bilinearità viene sostituita dalla sesquilinearità.

\begin{defbox}
\begin{definizione}[Prodotto scalare hermitiano]
\label{def:hermitiano}
Sia $V$ uno spazio vettoriale complesso. Un \emph{prodotto scalare hermitiano} su $V$ è un'applicazione $\langle\cdot,\cdot\rangle \colon V \times V \to \C$ tale che:
\begin{enumerate}
  \item $v \mapsto \langle v,w\rangle$ è lineare per ogni $w$ fisso (\emph{lineare nella prima variabile}).
  \item $\langle w,v\rangle = \overline{\langle v,w\rangle}$ per ogni $v,w \in V$ (\emph{simmetria hermitiana}).
  \item $\langle v,v\rangle > 0$ per ogni $v \neq 0$ (positività; $\langle v,v\rangle \in \R$ per la proprietà (2)).
\end{enumerate}
Uno \emph{spazio vettoriale hermitiano} è una coppia $(V,\langle\cdot,\cdot\rangle)$.
\end{definizione}
\end{defbox}

\begin{osservazione}
Dalla simmetria hermitiana: $\langle v,\lambda w\rangle = \overline{\langle \lambda w,v\rangle} = \bar\lambda\,\overline{\langle w,v\rangle} = \bar\lambda\langle v,w\rangle$. Quindi $\langle\cdot,\cdot\rangle$ è \emph{coniugato-lineare} nella seconda variabile. Un'applicazione lineare nella prima e coniugato-lineare nella seconda si chiama \emph{forma sesquilineare}.
\end{osservazione}

\begin{exbox}
\begin{esempio}
Il \emph{prodotto hermitiano standard} su $\C^n$: $\langle X,Y\rangle \coloneqq X^T \bar Y = \sum_{j=1}^n x_j \bar y_j$. La norma indotta è $\|X\| = \sqrt{\sum|x_j|^2}$.
\end{esempio}
\end{exbox}

\begin{teorema}[Cauchy-Schwarz hermitiana]
\label{thm:cs-hermitiana}
$|\langle v,w\rangle|^2 \leq \|v\|^2\|w\|^2$, con uguaglianza iff $v,w$ dipendenti.
\end{teorema}

\begin{proof}
Se $\langle v,w\rangle = 0$, banale. Altrimenti sia $\theta = \mathrm{Arg}(\langle v,w\rangle)$ e $v' = e^{-i\theta}v$: allora $\langle v',w\rangle = e^{-i\theta}\langle v,w\rangle \in \R_{\geq 0}$ e $\|v'\|=\|v\|$. Si procede come nel caso reale con $p(x) = \|xv'+w\|^2 \geq 0$.
\end{proof}

\begin{osservazione}
La teoria delle basi ON, dell'ortogonalità, della decomposizione ortogonale, di Gram-Schmidt si estende senza modifiche agli spazi hermitiani.
\end{osservazione}

\subsection{Matrici hermitiane e gruppo unitario}

\begin{defbox}
\begin{definizione}
Una matrice $A \in M_{n,n}(\C)$ è \emph{hermitiana} se $\bar A^T = A$. Il \emph{gruppo unitario} è $U_n(\C) \coloneqq \{A \in \GL_n(\C) \mid \bar A^T A = 1_n\}$, e $SU_n(\C) \coloneqq \{A \in U_n(\C) \mid \det A = 1\}$.
\end{definizione}
\end{defbox}

\begin{proposizione}
\label{prop:unitario-autovalori}
Se $f \in U(V)$ (gruppo unitario di uno spazio hermitiano), allora gli autovalori di $f$ hanno modulo $1$ e $f$ è diagonalizzabile con una base ON.
\end{proposizione}

\begin{proof}
Se $f(v) = \lambda v$ con $\|v\| = 1$: $1 = \|v\|^2 = \|f(v)\|^2 = |\lambda|^2$. Per induzione su $\dim V$: poiché $K = \C$ esiste un autovalore $\lambda_n$ con autovettore $v_n$. Si verifica che $f(v_n^\perp) \subseteq v_n^\perp$ (se $\langle w,v_n\rangle = 0$ allora $\langle f(w),f(v_n)\rangle = \langle w,v_n\rangle = 0$, e poiché $f(v_n) = \lambda_n v_n$ con $|\lambda_n|=1$, questo dà $\langle f(w),v_n\rangle = 0$). Per ipotesi induttiva $f|_{v_n^\perp}$ è diagonalizzabile con una base ON.
\end{proof}

% --------------------------------------------------------
\section{Il gruppo ortogonale}
\label{sec:gruppo-ortogonale}
% --------------------------------------------------------

\subsection{Decomposizione canonica di un'isometria}

Il teorema principale di questa sezione descrive la struttura di ogni elemento di $O(V)$.

\begin{teorema}[Decomposizione canonica di isometrie reali]
\label{thm:decomp-isometrie}
Sia $(V,\langle\cdot,\cdot\rangle)$ uno spazio vettoriale euclideo finitamente generato e $f \in O(V)$. Esiste una decomposizione ortogonale
\[
  V = V_1(f) \oplus^\perp V_{-1}(f) \oplus^\perp A_1 \oplus^\perp \cdots \oplus^\perp A_c,
\]
dove ogni $A_j$ ha dimensione $2$, è $f$-invariante e la restrizione di $f$ ad $A_j$ è una rotazione (di angolo $\theta_j \notin \pi\Z$).
\end{teorema}

\begin{proof}[Idea della dimostrazione]
Poiché $O_n(\R) \subset U_n(\C)$, per la Proposizione~\ref{prop:unitario-autovalori} l'estensione $F$ di $f$ a $\C^n$ è diagonalizzabile con autovalori di modulo $1$. Gli autovalori reali sono $\pm 1$; gli autovalori complessi non reali vengono in coppie coniugate $\lambda, \bar\lambda$. Per una coppia $\{T_j, \bar T_j\}$ di autovettori di $F$ relativi a $\lambda_j = e^{i\theta_j}$, i vettori reali $U_j = \frac{1}{\sqrt{2}}(T_j+\bar T_j)$ e $W_j = \frac{1}{\sqrt{2}i}(T_j-\bar T_j)$ generano un sottospazio $A_j$ di dimensione $2$ sul quale $f$ agisce con matrice di rotazione $R_{\theta_j}$. La base $X_1,\ldots,Y_1,\ldots,U_1,W_1,\ldots$ così costruita è ON e realizza la decomposizione.
\end{proof}

\begin{defbox}
\begin{definizione}[Rotazione]
Sia $(V,\langle\cdot,\cdot\rangle)$ euclideo di dimensione $2$. Un'isometria $f \in O(V)$ è una \emph{rotazione} se, scelta una base ON $\mathcal{B}$, la matrice $M^{\mathcal{B}}_{\mathcal{B}}(f) = R_\theta \coloneqq \begin{pmatrix}\cos\theta & -\sin\theta\\\sin\theta & \cos\theta\end{pmatrix}$ per qualche $\theta \in \R$.
\end{definizione}
\end{defbox}

\begin{corollario}
\label{cor:so-decomp}
Se $f \in SO(V)$, allora $V = V_1(f) \oplus^\perp B_1 \oplus^\perp \cdots \oplus^\perp B_d$ con ogni $B_j$ bidimensionale e $f|_{B_j}$ una rotazione. Se $f \in O(V) \setminus SO(V)$, si aggiunge un addendo $W$ monodimensionale con $f|_W = -\Id_W$.
\end{corollario}

\begin{exbox}
\begin{esempio}
In dimensione $3$: ogni $f \in SO(\R^3)$ ha un asse fisso (uno spazio $V_1(f)$ di dimensione $\geq 1$) e ruota il complemento ortogonale — è una rotazione attorno ad un asse. Ogni $f \in O(\R^3) \setminus SO(\R^3)$ è la composizione di una rotazione e una riflessione rispetto al piano ortogonale all'asse.
\end{esempio}
\end{exbox}

\subsection{Connessione per archi in $O(V)$}

\begin{defbox}
\begin{definizione}
Un cammino continuo in $O(V)$ è un'applicazione continua $\gamma \colon [a,b] \to O(V)$ (dove la continuità si intende nel senso delle entrate della matrice rispetto a una base). Due isometrie $g,h \in O(V)$ sono \emph{deformabili} l'una nell'altra se esiste un tale cammino con $\gamma(a)=g$, $\gamma(b)=h$.
\end{definizione}
\end{defbox}

\begin{proposizione}
\label{prop:connessione-O}
$g,h \in O(V)$ sono deformabili l'uno nell'altro se e solo se $\det g = \det h$, ovvero appartengono entrambi a $SO(V)$ o entrambi a $O(V)\setminus SO(V)$.
\end{proposizione}

\begin{proof}
Se $\det g \neq \det h$, non sono deformabili perché $t \mapsto \det\gamma(t)$ è continua e a valori in $\{+1,-1\}$, quindi costante (Teorema di Bolzano).

Se $g,h \in SO(V)$: basta deformare ciascuno in $\Id_V$. Per il Corollario~\ref{cor:so-decomp}, ogni $g \in SO(V)$ è prodotto di rotazioni piane; ogni rotazione $R_\theta$ si connette a $\Id$ tramite $t \mapsto R_{t\theta}$.

Se $g,h \in O(V)\setminus SO(V)$: $g^{-1}\cdot h \in SO(V)$, e si usa il caso precedente come nella dimostrazione della Proposizione~5.12.8.
\end{proof}

% --------------------------------------------------------
\section*{Esercizi}
\addcontentsline{toc}{section}{Esercizi del Capitolo 6}
% --------------------------------------------------------

\begin{esercizio}
Siano $A = \begin{pmatrix}2&-1\\-1&2\end{pmatrix}$ e $B = \begin{pmatrix}2&3\\3&5\end{pmatrix}$. Verificate che $\langle X,Y\rangle_A = X^T A Y$ e $\langle X,Y\rangle_B = X^T B Y$ definiscono prodotti scalari euclidei su $\R^2$ (usare il criterio di Sylvester).
\end{esercizio}

\begin{esercizio}
Sia $U = \{X \in \R^3 \mid x_1+2x_2+3x_3=0\}$ e $v=(1,1,1)$. Determinate la proiezione ortogonale di $v$ su $U$ rispetto al prodotto euclideo standard, e verificate che $v - \pi_U(v) \perp U$.
\end{esercizio}

\begin{esercizio}[Riflessione ortogonale]
Sia $(V,\langle\cdot,\cdot\rangle)$ uno spazio vettoriale euclideo e $0 \neq v \in V$. Definiamo
\[
  R_v \colon V \to V, \quad w \mapsto w - 2\frac{\langle w,v\rangle}{\|v\|^2}\,v.
\]
\begin{enumerate}
  \item Dimostrate che $R_v \in O(V)$ e $\det R_v = -1$.
  \item Descrivete $R_v$ geometricamente in $V = \mathcal{V}(\mathbf{E}^2)$ e in $V = \mathcal{V}(\mathbf{E}^3)$. [\emph{Suggerimento}: che cosa fa $R_v$ ai vettori paralleli e perpendicolari a $v$?]
  \item Dimostrate che ogni elemento di $O_n(\R)$ è prodotto di al più $n$ riflessioni del tipo $R_v$.
\end{enumerate}
\end{esercizio}

\begin{esercizio}[Guidato: Gram-Schmidt e proiezione]
\label{es:gram-schmidt-guidato}
Su $\R^3$ con il prodotto scalare standard, siano $v_1=(1,1,0)$, $v_2=(1,0,1)$, $v_3=(0,1,1)$.
\begin{enumerate}
  \item Verificate che $\{v_1,v_2,v_3\}$ è una base di $\R^3$.
  \item Applicate Gram-Schmidt per ottenere una base ON $\{w_1,w_2,w_3\}$.
  \item Calcolate la proiezione ortogonale di $u=(1,2,3)$ sul sottospazio $\Span(v_1,v_2)$.
\end{enumerate}
\end{esercizio}

\begin{esercizio}
Sia $(V,\langle\cdot,\cdot\rangle)$ uno spazio vettoriale euclideo finitamente generato e $U,W \subseteq V$ sottospazi. Dimostrate che $(U+W)^\perp = U^\perp \cap W^\perp$ e $(U \cap W)^\perp = U^\perp + W^\perp$.
\end{esercizio}

\begin{esercizio}
Sia $\{u_1,\ldots,u_m\}$ una lista ON di uno spazio vettoriale euclideo $(V,\langle\cdot,\cdot\rangle)$ e $v \in V$. Dimostrate la \emph{disuguaglianza di Bessel}:
\[
  \sum_{i=1}^m \langle v,u_i\rangle^2 \leq \|v\|^2,
\]
con uguaglianza se e solo se $v \in \Span(u_1,\ldots,u_m)$.
\end{esercizio}

\begin{esercizio}
Sia $A \in M_{n,n}(\R)$ con $A^T A = 1_n$ (cioè $A \in O_n(\R)$). Dimostrate che le colonne di $A$ formano una base ON di $\R^n$, e analogamente per le righe.
\end{esercizio}

\begin{esercizio}[Matrice di Gram e volume]
Siano $v_1,\ldots,v_n \in \R^n$. Dimostrate che il quadrato del volume del parallelepipedo generato da $v_1,\ldots,v_n$ è uguale al determinante della matrice di Gram $G = (\langle v_i,v_j\rangle)_{i,j}$. [\emph{Suggerimento}: se $M$ è la matrice con colonne $v_1,\ldots,v_n$, allora $G = M^T M$.]
\end{esercizio}

\begin{esercizio}
Sia $(V,\langle\cdot,\cdot\rangle)$ uno spazio vettoriale euclideo di dimensione $n$ e $f \in O(V)$ con $\det f = -1$.
\begin{enumerate}
  \item Dimostrate che $f$ ha $-1$ come autovalore (cioè $V_{-1}(f) \neq \{0\}$).
  \item Dedurre che se $n$ è dispari, ogni $f \in O(V)$ ha almeno un autovalore reale.
\end{enumerate}
\end{esercizio}

\begin{esercizio}[Proiezione e sistema lineare]
Sia $U = \Span(A_1,\ldots,A_m) \subseteq \R^n$ con $A_i \in \R^n$ vettori colonna linearmente indipendenti, e $M \in M_{n,m}(\R)$ la matrice con colonne $A_1,\ldots,A_m$. Dimostrate che la proiezione ortogonale di $b \in \R^n$ su $U$ è $\pi_U(b) = M(M^T M)^{-1}M^T b$ (il \emph{proiettore ortogonale}), e che $\pi_U(b)$ è l'unica soluzione del sistema delle equazioni normali $M^T M x = M^T b$.
\end{esercizio}

\begin{esercizio}[Guidato: isometrie di $\R^3$]
Sia $f \in SO_3(\R)$ con
\[
  M = \begin{pmatrix}0&-1&0\\1&0&0\\0&0&1\end{pmatrix}.
\]
\begin{enumerate}
  \item Verificate che $M \in SO_3(\R)$.
  \item Trovate l'asse di rotazione (cioè $V_1(f)$).
  \item Trovate l'angolo di rotazione $\theta$ dal Teorema~\ref{thm:decomp-isometrie}.
  \item Deformate $M$ con continuità in $1_3$ all'interno di $SO_3(\R)$.
\end{enumerate}
\end{esercizio}

\begin{esercizio}
Sia $(V,\langle\cdot,\cdot\rangle)$ uno spazio vettoriale hermitiano finitamente generato e $f \in U(V)$. Sia $W \subseteq V$ un sottospazio $f$-invariante. Dimostrate che anche $W^\perp$ è $f$-invariante. [\emph{Suggerimento}: se $w \in W^\perp$ e $u \in W$, calcolate $\langle f(w), u\rangle$.]
\end{esercizio}

% ============================================================
\chapter{Spazi affini euclidei}
\label{cap:affini-euclidei}
% ============================================================

\section*{Introduzione al capitolo}

Nei capitoli precedenti abbiamo separato due strutture: la geometria affine (spazi affini, sottospazi, applicazioni affini) e la geometria metrica (prodotto scalare, norme, angoli). Questo capitolo le fonde in un'unica nozione: lo \emph{spazio affine euclideo}.

Intuitivamente, uno spazio affine euclideo è uno spazio in cui si possono misurare distanze tra punti e angoli tra rette — esattamente come nel piano euclideo $\mathbf{E}^2$ o nello spazio $\mathbf{E}^3$ delle scuole medie, ma in dimensione arbitraria. La struttura è semplice: si prende uno spazio affine reale $S$ e si dota il suo spazio delle traslazioni $\mathcal{V}(S)$ di un prodotto scalare euclideo.

Il capitolo si sviluppa attorno a tre temi. \textbf{Distanza e proiezione:} la distanza tra due punti è la norma del vettore che li connette; la proiezione ortogonale di un punto su un sottospazio affine generalizza il piede dell'altezza. \textbf{Isometrie:} le applicazioni che preservano le distanze sono affini per il Teorema~\ref{thm:isometrie-affini}, e il gruppo $\mathrm{Isom}(S)$ si descrive esplicitamente. \textbf{Classificazione:} in dimensione $1$, $2$, $3$ le isometrie hanno nomi familiari --- traslazioni, rotazioni, riflessioni, glissoriflessioni, movimenti elicoidali --- e il Teorema di Chasles descrive completamente i movimenti rigidi di $\mathbf{E}^3$.

% --------------------------------------------------------
\section{Definizione e prime proprietà}
\label{sec:affine-euclideo-def}
% --------------------------------------------------------

\begin{defbox}
\begin{definizione}[Spazio affine euclideo]
\label{def:spazio-affine-euclideo}
Uno \emph{spazio affine euclideo} è una coppia $(S, \langle\cdot,\cdot\rangle)$ dove $S$ è uno spazio affine reale e $\langle\cdot,\cdot\rangle$ è un prodotto scalare euclideo su $\mathcal{V}(S)$.

La \emph{distanza} tra $P, Q \in S$ è $d(P,Q) \coloneqq \|\overrightarrow{PQ}\|$.
\end{definizione}
\end{defbox}

\begin{exbox}
\begin{esempio}
Il piano euclideo $\mathbf{E}^2$ e lo spazio euclideo $\mathbf{E}^3$, con il prodotto scalare definito dopo aver fissato un'unità di misura, sono spazi affini euclidei.
\end{esempio}
\end{exbox}

\begin{exbox}
\begin{esempio}[Spazio affine euclideo standard]
\label{es:En}
$\mathbb{A}^n(\R)$, cioè $\R^n$ visto come spazio affine, con il prodotto scalare standard $\langle X,Y\rangle = \sum_i x_i y_i$, è uno spazio affine euclideo. Lo denotiamo $\mathbf{E}^n(\R)$, e la distanza è
\[
  d(X,Y) = \left(\sum_{i=1}^n (x_i - y_i)^2\right)^{1/2}.
\]
\end{esempio}
\end{exbox}

\begin{proposizione}[Proprietà della distanza]
\label{prop:distanza}
Sia $(S,\langle\cdot,\cdot\rangle)$ uno spazio affine euclideo. Per ogni $P,Q,R \in S$:
\begin{enumerate}
  \item $d(P,Q) \geq 0$, con $d(P,Q) = 0 \Leftrightarrow P = Q$.
  \item $d(P,Q) = d(Q,P)$.
  \item $d(P,R) \leq d(P,Q) + d(Q,R)$, con uguaglianza sse $Q$ appartiene all'inviluppo convesso di $P$ e $R$.
\end{enumerate}
\end{proposizione}

\begin{proof}
(1) e (2) seguono dalla definizione di norma. Per (3): $\|\overrightarrow{PR}\| = \|\overrightarrow{PQ}+\overrightarrow{QR}\| \leq \|\overrightarrow{PQ}\|+\|\overrightarrow{QR}\|$ per la disuguaglianza triangolare vettoriale. L'uguaglianza si ha iff $\overrightarrow{PQ} = \lambda\overrightarrow{QR}$ per qualche $\lambda \geq 0$, che equivale a $Q = \frac{1}{1+\lambda}P + \frac{\lambda}{1+\lambda}R$.
\end{proof}

\begin{defbox}
\begin{definizione}[Riferimento ortonormale]
\label{def:riferimento-ON}
Un \emph{riferimento affine ortonormale} $\mathrm{RO}(O,\mathcal{B})$ è un riferimento affine in cui la base $\mathcal{B}$ di $\mathcal{V}(S)$ è ortonormale.
\end{definizione}
\end{defbox}

\begin{proposizione}
In un riferimento ortonormale, se $P(x_1,\ldots,x_n)$ e $Q(y_1,\ldots,y_n)$, allora $d(P,Q) = \bigl(\sum_i(x_i-y_i)^2\bigr)^{1/2}$.
\end{proposizione}

\begin{proof}
$d(P,Q) = \|\overrightarrow{PQ}\| = \|\sum_i(y_i-x_i)v_i\| = \bigl(\sum_i(y_i-x_i)^2\bigr)^{1/2}$ per la Proposizione~\ref{prop:calcolo-ON}.
\end{proof}

% --------------------------------------------------------
\section{Ginnastica affine euclidea}
\label{sec:ginnastica-affine}
% --------------------------------------------------------

\subsection{Ortogonalità tra sottospazi affini e interpretazione delle equazioni cartesiane}

Sia $S$ uno spazio affine euclideo di dimensione finita e $T \subset S$ con equazioni cartesiane (in un riferimento ON)
$a_{i,1}x_1 + \cdots + a_{i,n}x_n + b_i = 0$, $i=1,\ldots,m$. Posto $u_i = \sum_j a_{ij}v_j$ e fissato $P_0 \in T$:
\[
  T = \{P \in S \mid \langle\overrightarrow{P_0 P},u_i\rangle = 0\;\forall\,i\}, \qquad \mathcal{V}(T) = \{u_1,\ldots,u_m\}^\perp.
\]

\begin{defbox}
\begin{definizione}
Sottospazi affini $A, B \subset S$ sono \emph{ortogonali} se $\mathcal{V}(A) \perp \mathcal{V}(B)$.
\end{definizione}
\end{defbox}

\subsection{Proiezione ortogonale su un sottospazio affine}

\begin{proposizione}
\label{prop:proiezione-affine}
Siano $S$ uno spazio affine euclideo di dimensione finita e $T \subset S$ un sottospazio affine. Per ogni $P \in S$ esiste un unico $Q_0 \in T$ che minimizza $d(P,Q)$ al variare di $Q \in T$.
\end{proposizione}

\begin{proof}
Sia $P_0 \in T$ e $Q_0 \coloneqq P_0 + \pi_{\mathcal{V}(T)}(\overrightarrow{P_0 P})$. Per definizione di proiezione ortogonale vettoriale, $\overrightarrow{Q_0 P} \perp \mathcal{V}(T)$. Per $Q \in T$: $\|\overrightarrow{PQ}\|^2 = \|\overrightarrow{PQ_0}\|^2 + \|\overrightarrow{Q_0 Q}\|^2 \geq \|\overrightarrow{PQ_0}\|^2$.
\end{proof}

\begin{defbox}
\begin{definizione}
La \emph{proiezione ortogonale} di $P$ su $T$ è il punto $Q_0$ della proposizione. La \emph{distanza} $d(P,T) \coloneqq d(P,Q_0)$.
\end{definizione}
\end{defbox}

\begin{exbox}
\begin{esempio}[Distanza punto-iperpiano]
\label{es:distanza-iperpiano}
Se $T$ ha equazione $f(x_1,\ldots,x_n)=a_1x_1+\cdots+a_nx_n+b=0$ e $u=\sum_i a_i v_i$, allora per $P(c_1,\ldots,c_n)$:
\[
  d(P,T) = \frac{|f(c_1,\ldots,c_n)|}{\|u\|} = \frac{|a_1c_1+\cdots+a_nc_n+b|}{\sqrt{\sum_i a_i^2}}.
\]
La proiezione ortogonale è $Q_0 = \bigl(c_i - \tfrac{f(\mathbf{c})\,a_i}{\|u\|^2}\bigr)_i$.

\emph{Derivazione.} Si cerca $x \in \R$ con $\overrightarrow{P_0 P} = \overrightarrow{P_0 Q_0} + xu$ e $\overrightarrow{P_0 Q_0} \perp u$. Allora $x = \langle\overrightarrow{P_0 P},u\rangle/\|u\|^2 = f(\mathbf{c})/\|u\|^2$, e $d(P,T)=|x|\|u\|$.
\end{esempio}
\end{exbox}

\begin{exbox}
\begin{esempio}
Distanza di $P(1,2,3)$ dal piano $2x_1-x_2+2x_3-1=0$: $d = |2-2+6-1|/\sqrt{4+1+4} = 5/3$.
\end{esempio}
\end{exbox}

% --------------------------------------------------------
\section{Applicazioni che preservano le distanze}
\label{sec:preserva-distanze}
% --------------------------------------------------------

\begin{defbox}
\begin{definizione}
$\Phi \colon S \to T$ \emph{preserva le distanze} se $d(\Phi(P),\Phi(Q)) = d(P,Q)$ per ogni $P,Q$. È un'\emph{isometria} se preserva le distanze ed è biunivoca.
\end{definizione}
\end{defbox}

\begin{proposizione}
\label{prop:iso-lineare}
Sia $\varphi \colon V \to W$ tra spazi euclidei con $\dim V < \infty$, con $\varphi(0)=0$ e $\|\varphi(v)-\varphi(w)\|=\|v-w\|$ per ogni $v,w$. Allora $\varphi$ è lineare e isometrica.
\end{proposizione}

\begin{proof}
$\|\varphi(v)\|=\|v\|$ (da $\varphi(0)=0$). Espandendo le norme: $\langle\varphi(v),\varphi(w)\rangle = \langle v,w\rangle$. Se $\{v_i\}$ è una base ON di $V$, allora $\{\varphi(v_i)\}$ è ON in $W$. Dall'identità del prodotto scalare: $\varphi(v) = \sum_i\langle v,v_i\rangle\varphi(v_i)$ per ogni $v$, quindi $\varphi$ è lineare.
\end{proof}

\begin{teorema}
\label{thm:isometrie-affini}
Sia $S$ uno spazio affine euclideo di dimensione finita. Se $\Phi \colon S \to T$ preserva le distanze, allora $\Phi$ è affine e $\mathcal{V}(\Phi)$ è un'isometria di spazi vettoriali euclidei.
\end{teorema}

\begin{proof}
Fissiamo $P_0 \in S$; poniamo $\varphi(\overrightarrow{P_0 P}) \coloneqq \overrightarrow{\Phi(P_0)\Phi(P)}$. Le ipotesi della Proposizione~\ref{prop:iso-lineare} sono soddisfatte: $\varphi(0)=0$ e $\|\varphi(u)-\varphi(w)\| = d(\Phi(P),\Phi(Q)) = d(P,Q) = \|u-w\|$ dove $u=\overrightarrow{P_0 P}$, $w=\overrightarrow{P_0 Q}$. Quindi $\varphi$ è lineare, e $\Phi(P) = \Phi(P_0)+\varphi(\overrightarrow{P_0 P})$ è affine.
\end{proof}

\begin{corollario}
$\Phi$ preserva le distanze sse è affine con $\mathcal{V}(\Phi) \in O(\mathcal{V}(S),\mathcal{V}(T))$.
\end{corollario}

\begin{exbox}
\begin{esempio}[Riflessione ortogonale]
\label{es:riflessione}
Sia $T \subset S$ un sottospazio affine. La \emph{riflessione} in $T$ è $R_T(P) \coloneqq Q_0 + \overrightarrow{PQ_0}$ dove $Q_0 = \pi_T(P)$. In un riferimento ON con $T: x_{m+1}=\cdots=x_n=0$:
\[
  R_T(x_1,\ldots,x_n) = (x_1,\ldots,x_m,-x_{m+1},\ldots,-x_n).
\]
È un'isometria con $\det\mathcal{V}(R_T)=(-1)^{n-m}$ e $R_T\circ R_T = \Id_S$.
\end{esempio}
\end{exbox}

% --------------------------------------------------------
\section{Isometrie}
\label{sec:isometrie-affini}
% --------------------------------------------------------

\subsection{Il gruppo delle isometrie}

\begin{defbox}
\begin{definizione}
$\mathrm{Isom}(S)$ è il gruppo delle isometrie di $S$. Il \emph{sottogruppo dei movimenti rigidi} è
\[
  \mathrm{Isom}^+(S) \coloneqq \{\Phi \in \mathrm{Isom}(S) \mid \det\mathcal{V}(\Phi) = 1\}.
\]
\end{definizione}
\end{defbox}

\begin{proposizione}
Spazi affini euclidei di uguale dimensione finita sono isometrici.
\end{proposizione}

\begin{proof}
Tramite riferimenti ON si costruiscono isometrie $S \to \mathbf{E}^n(\R) \leftarrow T$, la cui composizione dà $S \xrightarrow{\sim} T$.
\end{proof}

\begin{osservazione}
$\mathrm{Isom}(\mathbf{E}^n(\R)) = \{X \mapsto AX+B \mid A \in O_n(\R),\, B \in \R^n\}$ e $\mathrm{Isom}^+$ corrisponde a $A \in SO_n(\R)$.
\end{osservazione}

\subsection{Deformabilità}

\begin{proposizione}
\label{prop:deformabilita-isom}
$\Phi,\Psi \in \mathrm{Isom}(S)$ sono deformabili sse $\det\mathcal{V}(\Phi) = \det\mathcal{V}(\Psi)$.
\end{proposizione}

\begin{proof}
Il determinante è costante lungo ogni cammino continuo. Ogni $\Phi \in \mathrm{Isom}^+$ con $\Phi(X)=AX+B$ si deforma in $\Id$ tramite $\mu(t)(X) = \gamma_A(t)\cdot X + tB$, dove $\gamma_A$ è un cammino in $SO_n(\R)$ da $1_n$ ad $A$ (Proposizione~\ref{prop:connessione-O}).
\end{proof}

\subsection{Isometrie in dimensione 1}

Le isometrie sono $x \mapsto \pm x + b$: traslazioni ($+1$, movimenti rigidi) e simmetrie di punto ($-1$, non movimenti rigidi).

\subsection{Isometrie in dimensione 2}

\begin{defbox}
\begin{definizione}
Sia $S$ uno spazio affine euclideo di dimensione $2$.
\begin{enumerate}
  \item Una \emph{rotazione} di centro $P$ e angolo $\theta$ è $X \mapsto R_\theta X$ in coordinate ON con origine $P$.
  \item Una \emph{glissoriflessione} di asse $L$ e vettore $v \in \mathcal{V}(L)$ è $\tau_v \circ R_L$ (riflessione in $L$ poi traslazione). Se $v=0$: riflessione.
\end{enumerate}
\end{definizione}
\end{defbox}

\begin{proposizione}
\label{prop:isom-dim2}
$\mathrm{Isom}^+(S)$ = traslazioni $\cup$ rotazioni. $\mathrm{Isom}(S)\setminus\mathrm{Isom}^+(S)$ = glissoriflessioni.
\end{proposizione}

\begin{proof}
\emph{Movimenti rigidi.} $A \in SO_2(\R)$: se $A=1_2$ è una traslazione; se $A=R_\theta$ ($\theta \notin 2\pi\Z$), l'equazione $AX+B=X$ ha soluzione unica (centro di rotazione) poiché $\det(A-1_2)\neq 0$.

\emph{Non movimenti rigidi.} $A \in O_2(\R)\setminus SO_2(\R)$ ha autovettori ortogonali $v$ ($\lambda=+1$) e $w$ ($\lambda=-1$). La retta $L$ con $\mathcal{V}(L)=\Span(w)$ soddisfa $\Phi(L)=L$; si decompone $\Phi=\tau_v\circ R_L$.
\end{proof}

\subsection{Movimenti rigidi in dimensione 3: Teorema di Chasles}

\begin{defbox}
\begin{definizione}
Sia $S$ uno spazio affine euclideo di dimensione $3$.
\begin{enumerate}
  \item Una \emph{rotazione del piano} $T$ è $\Psi(x,y,z) = (x\cos\theta-y\sin\theta,\, x\sin\theta+y\cos\theta,\, z)$ in coordinate ON con $T: z=0$.
  \item Un \emph{movimento elicoidale} è $\Phi = \tau_v \circ \Psi$ dove $\Psi$ è rotazione di $T$ e $v \in \mathcal{V}(T)^\perp$. Traslazioni e rotazioni di piani sono casi degeneri.
\end{enumerate}
\end{definizione}
\end{defbox}

\begin{osservazione}
Per $v\neq 0$ e $\theta \notin 2\pi\Z$: $\tau_v \circ \Psi = \Psi \circ \tau_v$ (la traslazione è nella direzione dell'asse, ortogonale al piano).
\end{osservazione}

\begin{teorema}[Chasles]
\label{thm:chasles}
I movimenti rigidi di uno spazio affine euclideo di dimensione $3$ sono esattamente i movimenti elicoidali.
\end{teorema}

\begin{proof}
I movimenti elicoidali sono in $\mathrm{Isom}^+$ (parte lineare in $SO_3$). Sia $\Phi \in \mathrm{Isom}^+$. Se $\mathcal{V}(\Phi)=\Id$: traslazione. Altrimenti per il Corollario~\ref{cor:so-decomp}: $\mathcal{V}(S)=U\oplus^\perp W$ con $\mathcal{V}(\Phi)|_U=\Id$ e $\mathcal{V}(\Phi)|_W$ rotazione non banale. Sia $T$ il piano con $\mathcal{V}(T)=W$. Per ogni $P\in T$, la retta $L_P=P+U$ si mappa su una retta $\perp T$ che incontra $T$ in $\Psi(P)$: $\Psi \colon T\to T$ è una rotazione con unico punto fisso $P_0$. La retta $L_{P_0}$ è $\Phi$-invariante con restrizione $\tau_v$ ($v\in U$). Allora $\Phi = \tau_v \circ \Psi_T$.
\end{proof}

% --------------------------------------------------------
\section*{Esercizi}
\addcontentsline{toc}{section}{Esercizi del Capitolo 7}
% --------------------------------------------------------

\begin{esercizio}
Sia $T \colon \mathrm{Isom}(\mathbf{E}^n(\R)) \to (\R^n,+)$ definita da $T(\Phi)=B$ dove $\Phi(X)=AX+B$. Dimostrate che $T$ non è un omomorfismo di gruppi. [\emph{Suggerimento}: calcolate $T(\Phi_1\circ\Phi_2)$ vs $T(\Phi_1)+T(\Phi_2)$.]
\end{esercizio}

\begin{esercizio}
Siano $A = \begin{pmatrix}119/169&-120/169\\-120/169&-119/169\end{pmatrix}$ e $B = \begin{pmatrix}3\\2\end{pmatrix}$.
\begin{enumerate}
  \item Verificate $A \in O_2(\R)\setminus SO_2(\R)$, quindi $\Phi(X)=AX+B$ è una glissoriflessione.
  \item Determinate l'asse della glissoriflessione. [\emph{Suggerimento}: trovate l'autovettore di $A$ con autovalore $+1$.]
\end{enumerate}
\end{esercizio}

\begin{esercizio}
Sia $A$ un piano affine euclideo. Siano $\Phi_1,\Phi_2$ glissoriflessioni con assi $S_1,S_2$ ortogonali e vettori di traslazione $v_1,v_2$.
\begin{enumerate}
  \item Dimostrate che $\Psi \coloneqq \Phi_1\circ\Phi_2$ è una rotazione di angolo $\pi$.
  \item Se $\|v_1\|=\|v_2\|$, dimostrate che il centro di $\Psi$ appartiene a una delle bisettrici di $S_1$ e $S_2$.
\end{enumerate}
\end{esercizio}

\begin{esercizio}
Sia $G \coloneqq \{\Phi \in \mathrm{Isom}(\mathbf{E}^2(\R)) \mid \Phi(\Z^2)=\Z^2\}$. Descrivete i movimenti rigidi in $G$ e gli assi delle glissoriflessioni in $G$.
\end{esercizio}

\begin{esercizio}[Guidato: distanza e proiezione in $\mathbf{E}^3$]
Sia $L \colon (1+t,\,2t,\,1-t)$ e $P(2,1,3)$.
\begin{enumerate}
  \item Trovate la proiezione ortogonale $Q_0$ di $P$ su $L$.
  \item Calcolate $d(P,L)$ direttamente e con la formula $d(P,L)^2 = \|\overrightarrow{P_0 P}\|^2 - \langle\overrightarrow{P_0 P},\hat{v}\rangle^2$.
\end{enumerate}
\end{esercizio}

\begin{esercizio}
Dimostrate che la composizione di due riflessioni in rette $L_1,L_2$ di $\mathbf{E}^2$ è: una traslazione se $L_1\parallel L_2$; una rotazione di angolo $2\alpha$ se $L_1\cap L_2\neq\emptyset$ con angolo $\alpha$ tra esse.
\end{esercizio}

\begin{esercizio}[Guidato: Chasles in pratica]
Sia $\Phi(X)=AX+B$ con $A=\begin{pmatrix}0&-1&0\\1&0&0\\0&0&1\end{pmatrix}$, $B=\begin{pmatrix}1\\0\\2\end{pmatrix}$.
\begin{enumerate}
  \item Verificate $\Phi \in \mathrm{Isom}^+(\mathbf{E}^3(\R))$.
  \item Trovate $T$, $v \in \mathcal{V}(T)^\perp$ e $P_0$ del Teorema di Chasles.
  \item Descrivete $\Phi$: asse, angolo e passo elicoidale.
\end{enumerate}
\end{esercizio}

\begin{esercizio}
Sia $T \subset S$ di dimensione $m$ in uno spazio euclideo di dimensione $n$. Dimostrate che $\det\mathcal{V}(R_T)=(-1)^{n-m}$. Quando è $R_T$ un movimento rigido?
\end{esercizio}


% ============================================================
\chapter{Forme quadratiche e forme bilineari simmetriche}
\label{cap:forme-quadratiche}
% ============================================================

\section*{Introduzione al capitolo}

Le \emph{forme quadratiche} sono le funzioni polinomiali omogenee di grado~$2$ su uno spazio vettoriale: oggetti semplicissimi da definire, eppure al centro di una teoria ricchissima. Si presentano ovunque: la norma quadratica $\|v\|^2$ di un vettore, la forma associata a una conica o a una quadrica, la matrice Hessiana di una funzione, l'energia cinetica in meccanica, l'intervallo spaziotemporale nella relatività speciale.

Il problema centrale del capitolo è la \emph{classificazione a meno di congruenza}: due forme quadratiche $f, g$ su $V$ sono congruenti se esiste un automorfismo $\varphi \in \GL(V)$ tale che $f(v) = g(\varphi^{-1}v)$ per ogni $v$. In termini di matrici, questo corrisponde alla relazione $A \sim G^t B G$ tra matrici simmetriche. La complessità della classificazione dipende fortemente dal campo:
\begin{itemize}
  \item Su $\C$ basta un invariante: il \emph{rango} (Proposizione~\ref{prop:class-complesse}).
  \item Su $\R$ servono due invarianti: la \emph{segnatura positiva} $s^+$ e quella \emph{negativa} $s^-$ --- il celebre \textbf{teorema di Sylvester} (Proposizione~\ref{prop:sylvester}).
  \item Su $\Q$ o su campi finiti la classificazione richiede strumenti più sottili (teoria dei numeri).
\end{itemize}

Il filo conduttore è il legame tra forme quadratiche e forme bilineari simmetriche (§\ref{sec:polarizzazione}): ogni forma quadratica, quando $\car\K \neq 2$, ha un'unica \emph{polarizzazione} bilineare simmetrica, e i due linguaggi sono interscambiabili.

% --------------------------------------------------------
\section{Funzioni polinomiali omogenee}
\label{sec:funz-pol-omog}
% --------------------------------------------------------

In tutto il capitolo gli spazi vettoriali sono finitamente generati. Richiamiamo la nozione di funzione polinomiale omogenea introdotta nella Sottosezione~1.7.

\begin{defbox}
\begin{definizione}[Funzione polinomiale omogenea]
\label{def:pol-omog}
Sia $V$ uno spazio vettoriale su $\K$. Una funzione $P \colon V \to \K$ è \emph{polinomiale omogenea di grado $d$} se, scelta una base $\{v_1,\ldots,v_n\}$ di $V$, la funzione $f \colon \K^n \to \K$ definita da $f(x_1,\ldots,x_n) \coloneqq P(x_1v_1+\cdots+x_nv_n)$ è un polinomio omogeneo di grado $d$. Si denota $P^d(V)$ l'insieme delle tali funzioni (un sottospazio vettoriale di $\K^V$).
\end{definizione}
\end{defbox}

La definizione è indipendente dalla base: se $\{w_1,\ldots,w_n\}$ è un'altra base con $w_i = \sum_j a_{ij}v_j$, allora $P(y_1w_1+\cdots+y_nw_n)$ è la composizione del polinomio omogeneo di grado $d$ con una trasformazione lineare, dunque è ancora polinomiale omogenea di grado $d$.

\begin{exbox}
\begin{esempio}
Le funzioni in $P^1(V)$ sono le forme lineari $V^\vee$. Le funzioni in $P^0(V)$ sono le costanti. La norma quadratica $v \mapsto \|v\|^2$ su uno spazio euclideo è in $P^2(V)$: se $\{v_i\}$ è una base, $\|{\textstyle\sum} x_iv_i\|^2 = \sum_{i,j}\langle v_i,v_j\rangle x_ix_j$.
\end{esempio}
\end{exbox}

\begin{osservazione}
\label{oss:omog-proprieta}
Se $P \in P^d(V)$, allora $P(tv) = t^d P(v)$ per ogni $t \in \K$, $v \in V$ (omotetia di grado $d$). In particolare $P(\mathbf{0}) = 0$ se $d \geq 1$.
\end{osservazione}

\subsection{Congruenza}

Il gruppo $\GL(V)$ agisce su $P^d(V)$ via $(g \cdot P)(v) \coloneqq P(g^{-1}v)$. Si verifica che questa è un'azione di gruppo: $\Id_V \cdot P = P$ e $(gh) \cdot P = g \cdot (h \cdot P)$.

\begin{defbox}
\begin{definizione}[Congruenza]
Due funzioni $P, Q \in P^d(V)$ sono \emph{congruenti} ($P \sim Q$) se esiste $g \in \GL(V)$ tale che $Q = g \cdot P$.
\end{definizione}
\end{defbox}

La congruenza è una relazione di equivalenza. In termini di basi: $P \sim Q$ se e solo se esistono basi $\{v_i\}$ e $\{w_i\}$ di $V$ tali che $P({\textstyle\sum}x_iv_i) = Q({\textstyle\sum}x_iw_i)$ per ogni $(x_i) \in \K^n$. Classificare $P^d(V)$ a meno di congruenza equivale quindi a classificare i polinomi omogenei di grado $d$ in $n$ variabili a meno di cambiamenti di coordinate lineari invertibili.

\begin{exbox}
\begin{esempio}
In $P^1(V) = V^\vee$ ci sono esattamente due classi di congruenza: $\{0\}$ e $V^\vee \setminus \{0\}$. Questo perché ogni due forme lineari non nulle sono legate da un automorfismo di $V$ (si portano entrambe nella stessa forma coordinata).
\end{esempio}
\end{exbox}

% --------------------------------------------------------
\section{Forme quadratiche}
\label{sec:forme-quadratiche}
% --------------------------------------------------------

\begin{defbox}
\begin{definizione}[Forma quadratica]
\label{def:forma-quadratica}
Sia $V$ uno spazio vettoriale su $\K$. Una \emph{forma quadratica} su $V$ è una funzione $Q \in P^2(V)$. Lo spazio delle forme quadratiche si denota $\mathcal{Q}(V) \coloneqq P^2(V)$.
\end{definizione}
\end{defbox}

\begin{osservazione}
A differenza di quanto succede per polinomi su campi finiti (dove polinomi diversi possono indurre la stessa funzione), una forma quadratica è completamente determinata dalla funzione $Q \colon V \to \K$. Questo segue dal fatto che se $Q(\sum x_iv_i) = \sum_{i \le j} c_{ij}x_ix_j$ e $Q \equiv 0$, allora $c_{ii} = Q(v_i) = 0$ e $c_{ij} = Q(v_i+v_j) - Q(v_i) - Q(v_j) = 0$ per $i < j$, dunque $Q = 0$ come polinomio.
\end{osservazione}

\subsection{Nucleo (radicale) e non-degenerazione}

\begin{proposizione}
\label{prop:nucleo-qf-sott}
Sia $Q \in \mathcal{Q}(V)$. L'insieme
\[
  \ker(Q) \coloneqq \{ u \in V \mid Q(u + v) = Q(v)\;\forall\, v \in V \}
\]
è un sottospazio vettoriale di $V$.
\end{proposizione}

\begin{proof}
Se $Q(u_1+v) = Q(v)$ e $Q(u_2+v) = Q(v)$ per ogni $v$, allora $Q((u_1+u_2)+v) = Q(u_1+(u_2+v)) = Q(u_2+v) = Q(v)$: la somma è chiusa. Per lo scalare: se $t = 0$ è ovvio; se $t \neq 0$ e $u \in \ker(Q)$, usando $Q(tv) = t^2Q(v)$ (Osservazione~\ref{oss:omog-proprieta}):
\[
  Q(tu+v) = Q\!\bigl(t(u+t^{-1}v)\bigr) = t^2 Q(u+t^{-1}v) = t^2 Q(t^{-1}v) = Q(v). \qedhere
\]
\end{proof}

\begin{defbox}
\begin{definizione}[Nucleo, nullità, non-degenerazione]
\label{def:nucleo-nondeg}
Sia $Q \in \mathcal{Q}(V)$.
\begin{itemize}
  \item Il \emph{nucleo} (o \emph{radicale}) di $Q$ è il sottospazio $\ker(Q)$ della Proposizione~\ref{prop:nucleo-qf-sott}.
  \item La \emph{nullità} di $Q$ è $\dim\ker(Q)$; il \emph{rango} di $Q$ è $\rg(Q) \coloneqq \dim V - \dim\ker(Q)$.
  \item $Q$ è \emph{non degenere} se $\ker(Q) = \{\mathbf{0}\}$, \emph{degenere} altrimenti.
\end{itemize}
\end{definizione}
\end{defbox}

\begin{osservazione}
\label{oss:ker-ne-zero}
Se $u \in \ker(Q)$ allora $Q(u) = Q(u+\mathbf{0}) = Q(\mathbf{0}) = 0$. Il viceversa è falso: $Q(u)=0$ non implica $u \in \ker(Q)$. Per esempio, $Q(x,y,z) = x^2-y^2$ soddisfa $Q(1,1,0)=0$ ma $(1,1,0) \notin \ker(Q)$ se $\car\K\neq 2$.
\end{osservazione}

\begin{exbox}
\begin{esempio}[Nucleo di $Q(x,y,z) = x^2-y^2$]
\label{es:nucleo-x2y2}
Un vettore $(a,b,c)$ è in $\ker(Q)$ se $(x+a)^2-(y+b)^2 = x^2-y^2$ per ogni $x,y$, cioè se $2ax+a^2-2by-b^2 = 0$ per ogni $x,y$.
\begin{itemize}
  \item Se $\car\K \neq 2$: necessariamente $a = b = 0$, dunque $\ker(Q) = \{(0,0,c)\}$ e la nullità è $1$.
  \item Se $\car\K = 2$: l'uguaglianza diventa $(a+b)^2 = 0$, cioè $a = b$, dunque $\ker(Q) = \{(a,a,c)\}$ e la nullità è $2$.
\end{itemize}
\end{esempio}
\end{exbox}

\begin{exbox}
\begin{esempio}
Per $r \in \{1,\ldots,n\}$, sia $Q_r(x_1,\ldots,x_n) \coloneqq x_1^2+\cdots+x_r^2 \in \mathcal{Q}(\K^n)$. Se $\car\K \neq 2$, la nullità di $Q_r$ è $n-r$, quindi $Q_r$ è congruente a $Q_s$ solo se $r = s$. Se invece $\car\K = 2$, tutte le $Q_r$ sono congruenti tra loro.
\end{esempio}
\end{exbox}

\begin{proposizione}
\label{prop:ker-congruenza}
Se $Q, Q' \in \mathcal{Q}(V)$ sono congruenti, le loro nullità (e i loro ranghi) coincidono.
\end{proposizione}

\begin{proof}
Se $Q' = g \cdot Q$, si verifica che $u \in \ker(Q')$ se e solo se $g^{-1}(u) \in \ker(Q)$, cioè $\ker(Q') = g(\ker(Q))$. Poiché $g$ è un isomorfismo, $\dim\ker(Q') = \dim\ker(Q)$.
\end{proof}

% --------------------------------------------------------
\section{Polarizzazione e forme bilineari simmetriche}
\label{sec:polarizzazione}
% --------------------------------------------------------

\textbf{Ipotesi.} Da questa sezione in poi supponiamo $\car\K \neq 2$.

Richiamiamo dalla Sezione~6.7: una \emph{forma bilineare} su $V$ è un'applicazione $F \colon V \times V \to \K$ bilineare; è \emph{simmetrica} se $F(v,w) = F(w,v)$. Si denota $\mathrm{Bil}^+(V)$ il sottospazio delle forme bilineari simmetriche.

\subsection{L'isomorfismo tra $\mathrm{Bil}^+(V)$ e $\mathcal{Q}(V)$}

Ogni $F \in \mathrm{Bil}^+(V)$ determina una forma quadratica: $q_F(v) \coloneqq F(v,v)$. Se $\{v_1,\ldots,v_n\}$ è una base:
\[
  q_F\!\Bigl(\sum_i x_iv_i\Bigr) = \sum_{i} F(v_i,v_i)\,x_i^2 + 2\!\sum_{i<j} F(v_i,v_j)\,x_ix_j.
\]

\begin{proposizione}
\label{prop:iso-bil-qf}
L'applicazione $\mathrm{Bil}^+(V) \to \mathcal{Q}(V)$, $F \mapsto q_F$, è un isomorfismo di $\K$-spazi vettoriali.
\end{proposizione}

\begin{proof}
\textit{Linearità}: banale. \textit{Iniettività}: se $q_F = 0$, per ogni $v,w$: $0 = q_F(v+w) = q_F(v)+2F(v,w)+q_F(w) = 2F(v,w)$, e $\car\K\neq 2$ dà $F = 0$. \textit{Suriettività}: data $q \in \mathcal{Q}(V)$ con coefficienti $c_{ij}$ ($i \le j$) in una base $\mathcal{B}$, si pone $a_{ij} = c_{ij}$ per $i=j$, $a_{ij} = c_{ij}/2$ per $i < j$, $a_{ij} = c_{ji}/2$ per $i > j$, e $F(v,w) \coloneqq X_\mathcal{B}(v)^t (a_{ij}) X_\mathcal{B}(w)$. Allora $F \in \mathrm{Bil}^+(V)$ e $q_F = q$.
\end{proof}

\begin{defbox}
\begin{definizione}[Polarizzazione]
\label{def:polarizzazione}
La \emph{polarizzazione} di $q \in \mathcal{Q}(V)$ è l'unica $F \in \mathrm{Bil}^+(V)$ tale che $q_F = q$.
\end{definizione}
\end{defbox}

I valori di $F$ si recuperano tramite l'\emph{identità di polarizzazione}:

\begin{proposizione}[Identità di polarizzazione]
\label{prop:id-pol}
Se $F$ è la polarizzazione di $q$, allora per ogni $v,w \in V$:
\[
  F(v,w) = \tfrac{1}{2}\bigl(q(v+w) - q(v) - q(w)\bigr).
\]
\end{proposizione}

\begin{proof}
$q(v+w) = F(v+w,v+w) = q(v)+2F(v,w)+q(w)$. Risolvendo per $F(v,w)$ e usando $\car\K \neq 2$.
\end{proof}

\begin{exbox}
\begin{esempio}
Il prodotto scalare euclideo $\langle\cdot,\cdot\rangle$ è la polarizzazione di $\|v\|^2$.
\end{esempio}
\end{exbox}

\begin{exbox}
\begin{esempio}
\label{es:pol-matrice}
Se $B \in M_{n,n}(\K)$ (non simmetrica) e $q(X) = X^t B X$, la polarizzazione è $F(X,Y) = \frac{1}{2}X^t(B^t+B)Y$. Infatti $F(X,X) = \frac{1}{2}X^t(B^t+B)X = X^tBX = q(X)$, perché $X^tB^tX = (X^tBX)^t = X^tBX$ (scalare).
\end{esempio}
\end{exbox}

\subsection{Matrice di Gram e congruenza di matrici}

Data $F \in \mathrm{Bil}^+(V)$ e una base $\mathcal{B} = \{v_1,\ldots,v_n\}$, la \emph{matrice di Gram} è $M_\mathcal{B}(F) \coloneqq (F(v_i,v_j)) \in M_n^+(\K)$ (matrice simmetrica). Per una forma quadratica $q$ con polarizzazione $F$, si pone $M_\mathcal{B}(q) \coloneqq M_\mathcal{B}(F)$.

L'azione di $\GL(V)$ su $\mathrm{Bil}^+(V)$ è $(g \cdot F)(v,w) \coloneqq F(g^{-1}v,\,g^{-1}w)$. Se $G = M^\mathcal{B}_\mathcal{B}(g)$, allora $M_\mathcal{B}(g\cdot F) = (G^{-1})^t\,M_\mathcal{B}(F)\,G^{-1}$.

\begin{defbox}
\begin{definizione}[Congruenza di matrici]
\label{def:matrici-congruenti}
Due matrici $A, B \in M_n(\K)$ sono \emph{congruenti} se esiste $G \in \GL_n(\K)$ tale che $A = G^t B G$.
\end{definizione}
\end{defbox}

\begin{osservazione}
Classificare le forme quadratiche su $V$ (di dimensione $n$) a meno di congruenza equivale a classificare le matrici simmetriche $n \times n$ a meno di congruenza.
\end{osservazione}

\subsection{Discriminante}

Il determinante di $M_\mathcal{B}(F)$ dipende dalla base: se $\mathcal{C}$ è un'altra base con matrice di cambio $P = M^\mathcal{B}_\mathcal{C}(\Id)$, allora $\det M_\mathcal{C}(F) = \det M_\mathcal{B}(F) \cdot (\det P)^{-2}$. Questo mostra che la classe di $\det M_\mathcal{B}(F)$ modulo i quadrati di $\K^*$ è un invariante.

Sia $\mathrm{Sqf}(\K)$ l'insieme delle classi di equivalenza per $x \sim y \Leftrightarrow \exists\,\mu \in \K^*: x = \mu^2 y$ (aggiungendo la classe $[0]$).

\begin{exbox}
\begin{esempio}
$\mathrm{Sqf}(\C) = \{[0],[1]\}$;\quad $\mathrm{Sqf}(\R) = \{[0],[1],[-1]\}$.
\end{esempio}
\end{exbox}

\begin{defbox}
\begin{definizione}[Discriminante]
Il \emph{discriminante} di $F \in \mathrm{Bil}^+(V)$ è $\mathrm{Disc}(F) \coloneqq [\det M_\mathcal{B}(F)] \in \mathrm{Sqf}(\K)$ (indipendente dalla base).
\end{definizione}
\end{defbox}

Il discriminante è un invariante di congruenza.

\begin{exbox}
\begin{esempio}
Le forme quadratiche su $\Q^2$:
\[
  f(x_1,x_2) = 2x_1^2-2x_1x_2+3x_2^2, \quad h(x_1,x_2) = 5x_1^2+6x_1x_2+3x_2^2,
\]
hanno $M_S(f) = \bigl(\begin{smallmatrix}2&-1\\-1&3\end{smallmatrix}\bigr)$, $M_S(h) = \bigl(\begin{smallmatrix}5&3\\3&3\end{smallmatrix}\bigr)$, dunque $\mathrm{Disc}(f) = [5]$ e $\mathrm{Disc}(h) = [6]$. Poiché $5/6$ non è un quadrato in $\Q^*$, $f \not\sim h$ su $\Q$. Le stesse $f, h$ viste su $\R$ sono invece congruenti (stesso rango e stessa segnatura, come vedremo).
\end{esempio}
\end{exbox}

\subsection{Ortogonalit\`a}

\begin{defbox}
\begin{definizione}[Ortogonalit\`a]
\label{def:ortogonalita-F}
Sia $F \in \mathrm{Bil}^+(V)$. I vettori $v,w\in V$ sono \emph{$F$-ortogonali} ($v \perp w$) se $F(v,w)=0$. L'\emph{$F$-ortogonale} di $S \subseteq V$ è $S^\perp \coloneqq \{w \in V \mid F(v,w)=0\;\forall v\in S\}$.
\end{definizione}
\end{defbox}

Per la simmetria di $F$: $v\perp w \Leftrightarrow w \perp v$. Inoltre, l'applicazione $L_F \colon V \to V^\vee$, $v \mapsto F(v,\cdot)$, è lineare e $V^\perp = \ker(L_F) = \ker(q_F)$ (il nucleo della forma quadratica).

\begin{lemma}
\label{lem:ortogonale-sottospazio}
$S^\perp$ è un sottospazio di $V$. Se $U = \langle u_1,\ldots,u_m\rangle$, allora $U^\perp = \bigcap_{i=1}^m u_i^\perp$.
\end{lemma}

\begin{proposizione}
\label{prop:dim-ortogonale}
Sia $F \in \mathrm{Bil}^+(V)$ non degenere e $U \subseteq V$ un sottospazio. Allora $\dim U^\perp = \dim V - \dim U$.
\end{proposizione}

\begin{proof}
Sia $\{u_1,\ldots,u_m\}$ base di $U$. Allora $U^\perp = \ker\Phi$ dove $\Phi(v) = (F(u_1,v),\ldots,F(u_m,v))$. Poiché $F$ è non degenere, $L_F$ è un isomorfismo $V \xrightarrow{\sim} V^\vee$, quindi $F(u_1,\cdot),\ldots,F(u_m,\cdot)$ sono linearmente indipendenti in $V^\vee$, $\Phi$ è suriettiva e $\dim\ker\Phi = \dim V - m$.
\end{proof}

\begin{exbox}
\begin{esempio}
In geometria euclidea, $S^\perp$ è il complemento ortogonale standard (introdotto nel Cap.~6): il Lemma~\ref{lem:ortogonale-sottospazio} e la Proposizione~\ref{prop:dim-ortogonale} generalizzano le proprietà già note.
\end{esempio}
\end{exbox}

\subsection{Diagonalizzazione}

\begin{defbox}
\begin{definizione}[Forma diagonale]
$q \in \mathcal{Q}(V)$ è \emph{diagonale} nella base $\mathcal{B} = \{v_1,\ldots,v_n\}$ se $q({\textstyle\sum} x_iv_i) = \sum_i c_i x_i^2$ per certi $c_i \in \K$. In questo caso $\mathcal{B}$ \emph{diagonalizza} $q$.
\end{definizione}
\end{defbox}

\begin{teorema}[Diagonalizzazione]
\label{thm:diagonalizzazione}
Sia $F \in \mathrm{Bil}^+(V)$. Esiste una base di $V$ che diagonalizza $F$.
\end{teorema}

\begin{proof}
Per induzione su $n = \dim V$. Se $n \le 1$ non c'è nulla. Se $F = 0$ ogni base va bene. Se $F \neq 0$, esiste $v_0$ con $q_F(v_0) = F(v_0,v_0) \neq 0$. Poniamo $U \coloneqq v_0^\perp$: si ha $\dim U = n-1$ (perché $F(v_0,v_0) \neq 0$ implica $v_0 \notin v_0^\perp$) e $V = \K v_0 \oplus U$. Per ipotesi induttiva, $F|_{U\times U}$ ha una base diagonalizzante $\{w_1,\ldots,w_{n-1}\}$ di $U$. Poiché $F(w_i,v_0) = 0$ per ogni $i$, la base $\{w_1,\ldots,w_{n-1},v_0\}$ diagonalizza $F$ su $V$.
\end{proof}

\begin{corollario}[Lagrange]
\label{cor:lagrange}
Ogni $q \in \mathcal{Q}(V)$ è diagonalizzabile, cioè esiste una base di $V$ in cui $q = \sum_i c_i x_i^2$.
\end{corollario}

\begin{corollario}
Ogni matrice simmetrica $A \in M_n^+(\K)$ è congruente a una matrice diagonale.
\end{corollario}

\subsection{Algoritmo di Lagrange}

Il Teorema~\ref{thm:diagonalizzazione} è costruttivo. Supposto $q \neq 0$, si sceglie $u_n \in V$ con $q(u_n) = c_n \neq 0$ e si completa a una base. Un calcolo esplicito mostra che, detti $a_{ij}$ le entrate di $M_\mathcal{B}(q)$:
\[
  q\!\Bigl(\sum_{i=1}^n z_i u_i\Bigr) = r(z_1,\ldots,z_{n-1}) + c_n\!\left(z_n + c_n^{-1}\sum_{i=1}^{n-1}a_{in}z_i\right)^2,
\]
dove $r$ è una forma quadratica in $n-1$ variabili. Il cambio di variabili $y_i = z_i$ ($i < n$), $y_n = z_n + c_n^{-1}\sum_{i<n}a_{in}z_i$ (matrice triangolare inferiore con $1$ sulla diagonale, quindi invertibile) riduce $q$ a $r(y_1,\ldots,y_{n-1}) + c_ny_n^2$. Si itera su $r$.

Se $K = \R$: dopo la diagonalizzazione, rescalando ogni $v_i$ con $|c_i|^{-1/2}$ (se $c_i \neq 0$), si riduce ogni coefficiente non nullo a $\pm 1$.

\begin{exbox}
\begin{esempio}[Lagrange su $\R^3$]
\label{es:lagrange}
Sia $f(x_1,x_2,x_3) = x_1x_2 + x_2x_3 + x_3x_1$. Poiché $f(0,1,1) = 1 \neq 0$, scegliamo la base $\{u_1,u_2,u_3\} = \{(1,0,0),(0,1,0),(0,1,1)\}$ e coordinate $y_i$ via $(x_1,x_2,x_3) = (y_1,y_2+y_3,y_3)$:
\[
  f = y_1y_2 + 2y_1y_3 + y_2y_3 + y_3^2.
\]
Completiamo il quadrato in $y_3$ (coefficiente $c_3 = 1$ di $y_3^2$ dopo la manipolazione):
\[
  f = -y_1^2 - \tfrac{1}{4}y_2^2 + \bigl(y_3+y_1+\tfrac{1}{2}y_2\bigr)^2.
\]
Poniamo $z_1 = y_1$, $z_2 = y_2$, $z_3 = y_1+\tfrac{1}{2}y_2+y_3$. La base corrispondente è $\{w_1,w_2,w_3\} = \{(1,0,-1),(0,1,-\tfrac{1}{2}),(0,0,1)\}$. Infine rescaliamo: $t_1 = z_1$, $t_2 = z_2/2$, $t_3 = z_3$, con base $\{r_1,r_2,r_3\} = \{w_1,2w_2,w_3\}$. Risultato:
\[
  f(t_1r_1+t_2r_2+t_3r_3) = -t_1^2 - t_2^2 + t_3^2.
\]
\end{esempio}
\end{exbox}

% --------------------------------------------------------
\section{Forme quadratiche reali e complesse}
\label{sec:forme-reali-complesse}
% --------------------------------------------------------

\subsection{Forme complesse: classificazione per rango}

Su $\C$ ogni non-zero è un quadrato (radice quadrata complessa sempre esistente), quindi i coefficienti diagonali non nulli si possono ridurre a $1$.

\begin{proposizione}
\label{prop:class-complesse}
Sia $V$ uno spazio vettoriale su $\C$. Due forme quadratiche $f, g \in \mathcal{Q}(V)$ sono congruenti se e solo se $\rg(f) = \rg(g)$.
\end{proposizione}

\begin{proof}
($\Rightarrow$) Per la Proposizione~\ref{prop:ker-congruenza}. ($\Leftarrow$) Per il Corollario~\ref{cor:lagrange}, $f = \sum_{i=1}^r \lambda_i x_i^2$ e $g = \sum_{i=1}^r \mu_i x_i^2$ in basi opportune, con $\lambda_i, \mu_i \neq 0$. Sostituendo $v_i \leftarrow \lambda_i^{-1/2}v_i$ e $w_i \leftarrow \mu_i^{-1/2}w_i$ (radici quadrate complesse), si ottiene $f = g = x_1^2+\cdots+x_r^2$, dunque $f \sim g$.
\end{proof}

\subsection{Forme reali: segnatura di Sylvester}

Su $\R$ i quadrati sono solo i reali non negativi, quindi $f(x) = x^2$ e $f(x) = -x^2$ non sono congruenti: occorre un ulteriore invariante. Anche rango e discriminante insieme non bastano: $f = x_1^2+x_2^2$ e $f = -x_1^2-x_2^2$ hanno rango $2$ e discriminante $[1]$ ma non sono congruenti.

\begin{defbox}
\begin{definizione}[Definitezza e segnatura]
\label{def:segnatura}
Siano $V$ uno spazio vettoriale reale e $f \in \mathcal{Q}(V)$.
\begin{itemize}
  \item $f$ è \emph{definita positiva} ($f > 0$) se $f(v) > 0$ per ogni $v \neq \mathbf{0}$.
  \item $f$ è \emph{definita negativa} ($f < 0$) se $f(v) < 0$ per ogni $v \neq \mathbf{0}$.
  \item $f$ è \emph{semidefinita positiva} (risp.\ negativa) se $f(v) \geq 0$ (risp.\ $\leq 0$) per ogni $v$.
  \item La \emph{segnatura positiva} $s^+(f)$ è $\max\{\dim U \mid U \subseteq V,\,f|_U > 0\}$.
  \item La \emph{segnatura negativa} $s^-(f)$ è $\max\{\dim U \mid U \subseteq V,\,f|_U < 0\}$.
\end{itemize}
\end{definizione}
\end{defbox}

\begin{osservazione}
Una forma quadratica reale è definita positiva se e solo se la sua polarizzazione è un prodotto scalare euclideo (Capitolo~6).
\end{osservazione}

\begin{lemma}
\label{lem:segnatura-invariante}
Segnatura positiva e negativa sono invarianti di congruenza: se $g \in \GL(V)$, allora $s^\pm(g\cdot f) = s^\pm(f)$.
\end{lemma}

\begin{proof}
L'applicazione $U \mapsto g(U)$ è una biiezione tra $\{U \mid f|_U > 0\}$ e $\{W \mid (g\cdot f)|_W > 0\}$: se $v \in U\setminus\{\mathbf{0}\}$ allora $(g\cdot f)(gv) = f(v) > 0$.
\end{proof}

\begin{proposizione}[Sylvester]
\label{prop:sylvester}
Siano $V$ uno spazio vettoriale reale e $f \in \mathcal{Q}(V)$. Se in una base $\mathcal{B}$ si ha
\[
  f = c_1x_1^2+\cdots+c_ax_a^2 - d_{a+1}x_{a+1}^2-\cdots-d_{a+b}x_{a+b}^2, \quad c_i,d_j > 0,
\]
allora $s^+(f) = a$ e $s^-(f) = b$.
\end{proposizione}

\begin{proof}
Siano $V_+ = \langle v_1,\ldots,v_a\rangle$, $V_- = \langle v_{a+1},\ldots,v_{a+b}\rangle$, $V_0 = \langle v_{a+b+1},\ldots,v_n\rangle$. Chiaramente $f|_{V_+} > 0$ e $f|_{V_-} < 0$, quindi $s^+(f) \geq a$ e $s^-(f) \geq b$.

Mostriamo $s^+(f) \leq a$ per assurdo. Se esistesse $U$ con $\dim U > a$ e $f|_U > 0$, per la formula di Grassmann:
\[
  \dim(U \cap (V_-+V_0)) \geq \dim U + \dim(V_-+V_0) - n = \dim U - a > 0.
\]
Quindi esiste $\mathbf{0} \neq v \in U \cap (V_-+V_0)$: da $v \in U$ segue $f(v) > 0$, ma da $v \in V_-+V_0$ (coordinate dei primi $a$ termini nulle) segue $f(v) \leq 0$: contraddizione. Analogo argomento per $s^-(f)$.
\end{proof}

\begin{teorema}[Classificazione delle forme quadratiche reali]
\label{thm:class-reali}
Siano $V$ reale e $f, h \in \mathcal{Q}(V)$. Allora $f \sim h$ se e solo se $s^+(f) = s^+(h)$ e $s^-(f) = s^-(h)$.
\end{teorema}

\begin{proof}
($\Rightarrow$) Lemma~\ref{lem:segnatura-invariante}. ($\Leftarrow$) Per il Corollario~\ref{cor:lagrange} si diagonalizzano $f$ e $h$; per la Proposizione~\ref{prop:sylvester} entrambe hanno $a = s^+(f)$ coefficienti positivi e $b = s^-(f)$ negativi. Rescalando con le radici quadrate inverse dei coefficienti, entrambe si riducono a $x_1^2+\cdots+x_a^2-x_{a+1}^2-\cdots-x_{a+b}^2$, dunque $f \sim h$.
\end{proof}

\begin{exbox}
\begin{esempio}[La metrica di Lorentz]
La forma quadratica $q(x_1,x_2,x_3,t) = c^2t^2 - x_1^2 - x_2^2 - x_3^2$ su $\R^4$ ha $s^+ = 1$, $s^- = 3$ e rango $4$. Il suo gruppo ortogonale $O(\R^4,q)$ è il \emph{gruppo di Lorentz}: le sue trasformazioni collegano i sistemi di riferimento inerziali preservando l'intervallo spaziotemporale $c^2(\Delta t)^2 - (\Delta x_1)^2 - (\Delta x_2)^2 - (\Delta x_3)^2$.
\end{esempio}
\end{exbox}

% --------------------------------------------------------
\section{Il gruppo ortogonale di una forma quadratica}
\label{sec:gruppo-ortogonale-qf}
% --------------------------------------------------------

\begin{defbox}
\begin{definizione}[Gruppo ortogonale di una forma quadratica]
\label{def:gruppo-ort-qf}
Sia $q \in \mathcal{Q}(V)$. Il \emph{gruppo ortogonale} di $q$ è
\[
  O(V,q) \coloneqq \{ g \in \GL(V) \mid g \cdot q = q \} = \{ g \in \GL(V) \mid q(g^{-1}v) = q(v)\;\forall v \}.
\]
\end{definizione}
\end{defbox}

\begin{lemma}
$O(V,q)$ è un sottogruppo di $\GL(V)$ (lo stabilizzatore di $q$ per l'azione di $\GL(V)$).
\end{lemma}

\begin{proof}
$\Id_V \cdot q = q$ \checkmark. Se $g \in O(V,q)$ e $w = g^{-1}v$: $q(w) = q(gw)$ per ogni $w$, cioè $g^{-1} \in O(V,q)$. La chiusura per composizione è analoga.
\end{proof}

In termini di matrici: se $A = M_\mathcal{B}(q)$ e $G = M^\mathcal{B}_\mathcal{B}(g)$, allora $g \in O(V,q)$ se e solo se $G^t A G = A$.

\begin{exbox}
\begin{esempio}
$O(V,0) = \GL(V)$. Se $(V,\langle\cdot,\cdot\rangle)$ è euclideo e $q(v) = \|v\|^2$, allora $O(V,q) = O(V,\langle\cdot,\cdot\rangle)$ è il gruppo ortogonale euclideo del Capitolo~6.
\end{esempio}
\end{exbox}

\begin{defbox}
\begin{definizione}[Prodotto scalare su $\R$]
Un \emph{prodotto scalare} su uno spazio vettoriale reale $V$ è una forma bilineare simmetrica \emph{non degenere} $\langle\cdot,\cdot\rangle$ su $V$. Il suo gruppo ortogonale è $O(V,q)$ con $q(v) = \langle v,v\rangle$. Un prodotto scalare \emph{definito positivo} è un prodotto scalare euclideo.
\end{definizione}
\end{defbox}

Se il prodotto scalare è indefinito (né positivo né negativo), il gruppo ortogonale cambia radicalmente carattere: per esempio, il gruppo di Lorentz ha quattro componenti connesse (a differenza di $O(n)$ che ne ha due).

\begin{exbox}
\begin{esempio}[Gruppo di Lorentz in dimensione $1+1$]
Per il prodotto scalare $\langle X,X\rangle = 2x_1x_2$ su $\R^2$ (segnatura $(1,1)$), il gruppo ortogonale è formato da quattro insiemi di matrici:
$\bigl\{\bigl(\begin{smallmatrix}t&0\\0&t^{-1}\end{smallmatrix}\bigr)\bigr\}_{t>0}$, $\bigl\{\bigl(\begin{smallmatrix}t&0\\0&t^{-1}\end{smallmatrix}\bigr)\bigr\}_{t<0}$, $\bigl\{\bigl(\begin{smallmatrix}0&t\\t^{-1}&0\end{smallmatrix}\bigr)\bigr\}_{t>0}$, $\bigl\{\bigl(\begin{smallmatrix}0&t\\t^{-1}&0\end{smallmatrix}\bigr)\bigr\}_{t<0}$, ciascuno costituente una componente connessa.
\end{esempio}
\end{exbox}

% --------------------------------------------------------
\section*{Esercizi}
\addcontentsline{toc}{section}{Esercizi del Capitolo 8}
\label{sec:es-cap8}
% --------------------------------------------------------

\begin{esercizio}
Sia $q \in \mathcal{Q}(\R^2)$ data da $q(X) = X^t\!AX$ con $A = \bigl(\begin{smallmatrix}2&-1\\-1&2\end{smallmatrix}\bigr)$.
\begin{enumerate}[label=(\arabic*)]
  \item Con la base $\mathcal{B} = \{(1,1),(1,-1)\}$ di $\R^2$, calcolate $M_\mathcal{B}(q)$.
  \item Verificate che $q$ è definita positiva. [\emph{Suggerimento}: diagonalizzate e controllate i segni.]
\end{enumerate}
\end{esercizio}

\begin{esercizio}
Sia $V \subset C^0([0,1])$ il sottospazio generato da $f \equiv 1$, $g = \cos\pi t$, $h = \sin\pi t$, e $F(\varphi,\psi) = \int_0^1\varphi(t)\psi(t)\,dt$.
\begin{enumerate}[label=(\arabic*)]
  \item Verificate che $f, g, h$ sono linearmente indipendenti.
  \item Calcolate $M_\mathcal{B}(F)$ con $\mathcal{B} = \{f,g,h\}$.
\end{enumerate}
\end{esercizio}

\begin{esercizio}
Siano $f,g \in \mathcal{Q}(\R^3)$ date da $f = -7x_1^2+2x_1x_2-6x_1x_3+5x_2^2-x_3^2$ e $g = x_1^2-4x_1x_2-2x_1x_3+8x_2^2+6x_2x_3+2x_3^2$. Basi: $\mathcal{B} = \{(1,1,3),(2,-1,0),(0,0,1)\}$, $\mathcal{C} = \{(0,1,2),(3,0,1),(1,-1,0)\}$.
\begin{enumerate}[label=(\arabic*)]
  \item Calcolate $M_\mathcal{B}(f)$ e $M_\mathcal{C}(g)$.
  \item Determinate se $f \sim g$ (cioè se $(\R^3,f) \cong (\R^3,g)$).
\end{enumerate}
\end{esercizio}

\begin{esercizio}
Sia $f \in \mathcal{Q}(\R^{2n})$ definita da $f(x_1,\ldots,x_n,y_1,\ldots,y_n) = x_1y_1+\cdots+x_ny_n$. Determinate $s^+(f)$ e $s^-(f)$. [\emph{Suggerimento}: considerate i sottospazi $U^\pm = \langle e_i \pm e_{n+i}\rangle_{i=1}^n$.]
\end{esercizio}

\begin{esercizio}
Sia $A = \bigl(\begin{smallmatrix}2&1&0\\1&0&2\\0&2&3\end{smallmatrix}\bigr)$. Trovate una base di $\R^3$ che diagonalizza la forma quadratica $q(X) = X^tAX$.
\end{esercizio}

\begin{esercizio}
(a) Siano $A = \bigl(\begin{smallmatrix}3&2\\2&5\end{smallmatrix}\bigr)$, $B = \bigl(\begin{smallmatrix}2&1\\1&4\end{smallmatrix}\bigr) \in M_2(\Q)$ e $f(X)=X^tAX$, $g(X)=X^tBX$. Dimostrate che $f \not\sim g$ su $\Q$.

(b) Stesse matrici viste su $\R$: dimostrate che $\varphi(X) = X^tAX \sim \psi(X) = X^tBX$.
\end{esercizio}

\begin{esercizio}
Sia $f \in \mathcal{Q}(V)$ con $V$ reale. Un sottospazio $U \subseteq V$ è \emph{isotropo} se $f|_U \equiv 0$. Dimostrate che $\max\{\dim U \mid U \text{ isotropo}\} = \dim\ker(f) + \min\{s^+(f),s^-(f)\}$.
\end{esercizio}

\begin{esercizio}
Per $a\in\C$ indichiamo $\mathrm{Re}(a)$ e $\mathrm{Im}(a)$. Siano $f(w,z) = \mathrm{Re}(w\bar{z})$ e $g(w,z) = \mathrm{Im}(w\bar{z})$ su $\C\times\C$ (con $\C$ visto come $\R$-spazio vettoriale). Verificate che $f$ e $g$ sono bilineari e non degeneri; stabilite quale è simmetrica.
\end{esercizio}

\begin{esercizio}
Sia $\Phi(A,B) = \Tr(AB)$ su $M_{2,2}(\R)\times M_{2,2}(\R)$ (dove $\Tr$ è la traccia). Verificate che $\Phi$ è bilineare simmetrica e trovate una base di $M_{2,2}(\R)$ che diagonalizza $\Phi$.
\end{esercizio}

\begin{esercizio}
Sia $q(A) = \Tr(A^2)$ su $M_{n,n}(\R)$. Determinate rango e segnatura di $q$. [\emph{Suggerimento}: esaminate la restrizione di $q$ al sottospazio $\mathrm{Sym}_n(\R)$ delle matrici simmetriche e al sottospazio $\mathrm{Alt}_n(\R)$ delle antisimmetriche.]
\end{esercizio}

\begin{esercizio}[Guidato: diagonalizzazione in base ortonormale]
Sia $f(x_1,x_2) = x_1^2+2x_1x_2+3x_2^2$.
\begin{enumerate}[label=(\arabic*)]
  \item Calcolate la matrice di Gram $A$ di $f$ nella base standard di $\R^2$.
  \item Trovate una base \emph{ortonormale per il prodotto euclideo standard} che diagonalizza $f$. [\emph{Suggerimento}: diagonalizzate $A$ per similitudine ortogonale (teorema spettrale, Cap.~9); gli autovettori normalizzati formano la base cercata.]
\end{enumerate}
\end{esercizio}

\begin{esercizio}
\label{es:base-diag-catena}
Siano $U_1 \subset U_2 \subset \cdots \subset U_n = V$ sottospazi con $\dim U_i = i$, e $f \in \mathcal{Q}(V)$ tale che $f|_{U_i}$ sia non degenere per ogni $i$. Dimostrate che esiste una base $\{v_1,\ldots,v_n\}$ che diagonalizza $f$ con $\langle v_1,\ldots,v_i\rangle = U_i$ per ogni $i$.
\end{esercizio}

\begin{esercizio}[Criterio di Sylvester--Jacobi]
Sia $A \in M_n^+(\R)$ e $f(X) = X^tAX$. Per $p \in \{1,\ldots,n\}$ sia $A^{(p)}$ la sottomatrice $p\times p$ superiore sinistra di $A$. Supponete che $\det A^{(p)} \neq 0$ per ogni $p$. Dimostrate che $s^-(f)$ è il numero di cambi di segno nella sequenza $1,\det A^{(1)},\det A^{(2)},\ldots,\det A^{(n)}$. In particolare $f > 0$ se e solo se tutti i $\det A^{(p)}$ sono positivi. [\emph{Suggerimento}: usate l'Esercizio~\ref{es:base-diag-catena} con la catena canonica $U_i = \langle e_1,\ldots,e_i\rangle$.]
\end{esercizio}

\begin{esercizio}
Sia $\langle X,Y\rangle = X^t\bigl(\begin{smallmatrix}0&1\\1&0\end{smallmatrix}\bigr)Y$ il prodotto scalare su $\R^2$ (segnatura $(1,1)$, $\langle X,X\rangle = 2x_1x_2$). Dimostrate che
\[
  O(\R^2,\langle\cdot,\cdot\rangle) = \left\{\begin{pmatrix}t&0\\0&t^{-1}\end{pmatrix}\bigg|\,t\in\R^*\right\} \cup \left\{\begin{pmatrix}0&t\\t^{-1}&0\end{pmatrix}\bigg|\,t\in\R^*\right\}.
\]
Dimostrate inoltre che ciascuno dei due sottoinsiemi si suddivide in due componenti connesse ($t>0$ e $t<0$), per un totale di quattro componenti di $O(\R^2,\langle\cdot,\cdot\rangle)$.
\end{esercizio}


% ============================================================
\chapter{Teoremi spettrali}
\label{cap:spettrale}
% ============================================================

\section*{Introduzione al capitolo}

I \emph{teoremi spettrali} sono tra i risultati più profondi e utili dell'algebra lineare. In poche parole dicono che le «migliori» basi per diagonalizzare una forma quadratica o un operatore simmetrico non sono basi qualunque, ma basi \emph{ortonormali}. Questo impone vincoli molto più forti sulla struttura dell'operatore.

Due sono i risultati principali:

\begin{enumerate}
  \item \textbf{Teorema spettrale reale.} Se $(V,\langle\cdot,\cdot\rangle)$ è uno spazio vettoriale euclideo e $A \colon V \to V$ è un endomorfismo \emph{simmetrico} (cioè $\langle Av,w\rangle = \langle v,Aw\rangle$), allora $V$ ammette una base ortonormale di autovettori di $A$. In termini di matrici: ogni matrice simmetrica reale è diagonalizzabile tramite una matrice ortogonale.

  \item \textbf{Teorema spettrale complesso.} Se $(V,\langle\cdot,\cdot\rangle)$ è uno spazio vettoriale hermitiano e $T \colon V \to V$ è \emph{autoaggiunto} ($\langle Tv,w\rangle = \langle v,Tw\rangle$), allora $T$ ha autovalori tutti reali e $V$ ammette una base ortonormale di autovettori di $T$. In termini di matrici: ogni matrice hermitiana complessa è diagonalizzabile tramite una matrice unitaria, con autovalori reali.
\end{enumerate}

L'aggettivo «spettrale» si riferisce allo spettro di un operatore (l'insieme dei suoi autovalori): il teorema afferma, tra le altre cose, che lo spettro di una matrice simmetrica reale o hermitiana complessa è contenuto in $\R$.

Il capitolo è organizzato come segue. La Sezione~\ref{sec:spettrale-reale} sviluppa il teorema spettrale reale in due formulazioni equivalenti --- quella per endomorfismi simmetrici e quella per forme quadratiche --- e ne dà la dimostrazione geometrica. La Sezione~\ref{sec:forme-hermitiane} introduce le forme hermitiane e la loro teoria, del tutto analoga a quella delle forme bilineari simmetriche reali. La Sezione~\ref{sec:spettrale-complesso} enuncia e dimostra il teorema spettrale complesso.

% --------------------------------------------------------
\section{Il teorema spettrale reale}
\label{sec:spettrale-reale}
% --------------------------------------------------------

In tutta questa sezione $(V,\langle\cdot,\cdot\rangle)$ è uno spazio vettoriale euclideo (reale) di dimensione finita $n$.

\subsection{Endomorfismi simmetrici}

\begin{defbox}
\begin{definizione}[Endomorfismo simmetrico]
\label{def:endo-simm}
Un endomorfismo $A \colon V \to V$ è \emph{simmetrico} se
\[
  \langle A(u), w \rangle = \langle u, A(w) \rangle \qquad \forall\, u, w \in V.
\]
Denotiamo $\mathrm{End}^+(V) \coloneqq \{ A \in \mathrm{End}(V) \mid A \text{ è simmetrico} \}$.
\end{definizione}
\end{defbox}

\begin{osservazione}
\label{oss:simm-matrice}
Sia $\mathcal{B} = \{v_1,\ldots,v_n\}$ una base ON di $V$. Allora $A$ è simmetrico se e solo se $M^\mathcal{B}_\mathcal{B}(A)$ è una matrice simmetrica. Questo segue dal fatto che l'entrata $(i,j)$ di $M^\mathcal{B}_\mathcal{B}(A)$ è $\langle A(v_j),v_i\rangle$, e per un endomorfismo simmetrico $\langle A(v_j),v_i\rangle = \langle v_j,A(v_i)\rangle = \langle A(v_i),v_j\rangle$, che è l'entrata $(j,i)$.
\end{osservazione}

L'insieme $\mathrm{End}^+(V)$ è un sottospazio vettoriale di $\mathrm{End}(V)$, e l'applicazione $A \mapsto M^\mathcal{B}_\mathcal{B}(A)$ (con $\mathcal{B}$ base ON) è un isomorfismo $\mathrm{End}^+(V) \xrightarrow{\sim} M^+_{n,n}(\R)$.

\subsection{Due formulazioni del teorema spettrale}

\begin{proposizione}
\label{prop:autospazi-ortogonali}
Sia $A \colon V \to V$ un endomorfismo simmetrico. Se $\lambda \neq \mu$ sono autovalori di $A$, allora $V_\lambda(A) \perp V_\mu(A)$.
\end{proposizione}

\begin{proof}
Siano $u \in V_\lambda(A)$ e $w \in V_\mu(A)$. Allora
\[
  \lambda\langle u,w\rangle = \langle A(u),w\rangle = \langle u,A(w)\rangle = \mu\langle u,w\rangle.
\]
Poiché $\lambda \neq \mu$, si ha $\langle u,w\rangle = 0$.
\end{proof}

\begin{corollario}
\label{cor:diag-ON}
Se $A \colon V \to V$ è simmetrico e diagonalizzabile su $\R$, allora esiste una base ON di $V$ che diagonalizza $A$.
\end{corollario}

\begin{proof}
Sia $V = V_{\lambda_1}(A) \oplus \cdots \oplus V_{\lambda_m}(A)$. Per la Proposizione~\ref{prop:autospazi-ortogonali} gli autospazi sono a due a due ortogonali. Scegliamo una base ON $\mathcal{B}_i$ di ciascun $V_{\lambda_i}(A)$: l'unione $\mathcal{B} = \mathcal{B}_1 \cup \cdots \cup \mathcal{B}_m$ è una base ON di $V$ di autovettori.
\end{proof}

\begin{teorema}[Teorema spettrale reale --- prima formulazione]
\label{thm:spettrale-reale-I}
Sia $(V,\langle\cdot,\cdot\rangle)$ uno spazio vettoriale euclideo e $A \colon V \to V$ un endomorfismo simmetrico. Allora esiste una base ON di $V$ che diagonalizza $A$.
\end{teorema}

La dimostrazione è rinviata alla Sottosezione~\ref{ssec:dim-spettrale}. Vediamo prima la formulazione equivalente in termini di forme quadratiche.

Ogni endomorfismo $A \colon V \to V$ determina una forma bilineare
\[
  \Phi_A \colon V \times V \to \R, \quad \Phi_A(u,w) \coloneqq \langle A(u),w\rangle.
\]
$A$ è simmetrico se e solo se $\Phi_A \in \mathrm{Bil}^+(V)$.

\begin{proposizione}
\label{prop:iso-end-bil}
L'applicazione $\Phi \colon \mathrm{End}^+(V) \to \mathrm{Bil}^+(V)$, $A \mapsto \Phi_A$, è un isomorfismo di $\R$-spazi vettoriali.
\end{proposizione}

\begin{proof}
Sia $\mathcal{B}$ una base ON. L'isomorfismo $M^\mathcal{B}_\mathcal{B} \colon \mathrm{End}^+(V) \xrightarrow{\sim} M^+_{n,n}(\R)$ e l'isomorfismo $M_\mathcal{B} \colon \mathrm{Bil}^+(V) \xrightarrow{\sim} M^+_{n,n}(\R)$ soddisfano $M^\mathcal{B}_\mathcal{B}(A) = M_\mathcal{B}(\Phi_A)$ (per l'Osservazione~\ref{oss:simm-matrice}), dunque $\Phi = M_\mathcal{B}^{-1} \circ M^\mathcal{B}_\mathcal{B}$ è un isomorfismo.
\end{proof}

Componendo con l'isomorfismo $\mathrm{Bil}^+(V) \xrightarrow{\sim} \mathcal{Q}(V)$ della Proposizione~\ref{prop:iso-bil-qf} otteniamo:

\begin{defbox}
\begin{definizione}
Data $A \in \mathrm{End}^+(V)$, la \emph{forma quadratica associata} è
\[
  q_A \colon V \to \R, \quad q_A(v) \coloneqq \langle A(v),v\rangle.
\]
\end{definizione}
\end{defbox}

L'applicazione $A \mapsto q_A$ è un isomorfismo $\mathrm{End}^+(V) \xrightarrow{\sim} \mathcal{Q}(V)$. Una base $\mathcal{B}$ diagonalizza $A$ se e solo se diagonalizza $q_A$.

\begin{teorema}[Teorema spettrale reale --- seconda formulazione]
\label{thm:spettrale-reale-II}
Sia $(V,\langle\cdot,\cdot\rangle)$ uno spazio vettoriale euclideo e $f \colon V \to \R$ una forma quadratica. Allora esiste una base ON $\mathcal{B} = \{v_1,\ldots,v_n\}$ di $V$ che diagonalizza $f$:
\[
  f(x_1v_1+\cdots+x_nv_n) = \sum_{i=1}^n \lambda_i x_i^2.
\]
\end{teorema}

\begin{osservazione}
Il gruppo ortogonale $O(V)$ agisce per coniugio su $\mathrm{End}^+(V)$: se $A \in \mathrm{End}^+(V)$ e $g \in O(V)$, allora $g \circ A \circ g^{-1}$ è simmetrico. Questa azione corrisponde (tramite gli isomorfismi) all'azione di $O(V)$ per congruenza su $\mathcal{Q}(V)$. Il Teorema spettrale reale implica che $A, B \in M^+_{n,n}(\R)$ sono nella stessa orbita per l'azione di $O_n(\R)$ (per coniugio) se e solo se hanno lo stesso polinomio caratteristico.
\end{osservazione}

\subsection{Dimostrazione del teorema spettrale reale}
\label{ssec:dim-spettrale}

Il passaggio cruciale è trovare un autovettore. Daremo una dimostrazione geometrica che non usa il Teorema fondamentale dell'Algebra.

\begin{proposizione}
\label{prop:autovett-crit}
Sia $A \colon V \to V$ simmetrico. Un vettore non nullo $v \in V$ è un autovettore di $A$ se e solo se
\[
  v^\perp \subseteq \{ w \in V \mid \langle A(v), w\rangle = 0 \}.
\]
\end{proposizione}

\begin{proof}
($\Rightarrow$) Se $A(v) = \lambda v$ e $w \perp v$, allora $\langle A(v),w\rangle = \lambda\langle v,w\rangle = 0$.

($\Leftarrow$) L'ipotesi dice che $A(v) \perp v^\perp$. Poiché l'ortogonale di $v^\perp$ è $\mathrm{Span}(v)$, si ha $A(v) \in \mathrm{Span}(v)$.
\end{proof}

\begin{corollario}
\label{cor:A-preserva-ort}
Se $v$ è un autovettore di $A$, allora $A(v^\perp) \subseteq v^\perp$.
\end{corollario}

\begin{proof}
Se $w \perp v$, allora per simmetria di $A$: $\langle A(w),v\rangle = \langle w,A(v)\rangle = \lambda\langle w,v\rangle = 0$.
\end{proof}

\begin{proposizione}
\label{prop:esistenza-autovett}
Sia $A \colon V \to V$ simmetrico. Esiste un vettore non nullo $v \in V$ che è autovettore di $A$.
\end{proposizione}

\begin{proof}
Sia $f = q_A$ e consideriamo
\[
  \rho \colon V \setminus \{\mathbf{0}\} \to \R, \quad \rho(v) = \frac{f(v)}{\|v\|^2}.
\]
\textbf{Passo 1: $\rho$ ha un massimo.} Tramite un'isometria, possiamo assumere $V = \R^n$ con prodotto standard. Sia $F_n = \{X \in \R^n \mid |x_i| \leq 1\;\forall i,\;\exists i_0 : |x_{i_0}|=1\}$ la frontiera del cubo unitario. La funzione $\rho$ è continua, e $F_n$ è compatta (unione di iperfacce chiuse e limitate), quindi $\rho|_{F_n}$ ammette massimo (per il teorema di Bolzano--Weierstrass in $\R^n$). Per ogni $X \neq \mathbf{0}$, il vettore $m^{-1}X$ (con $m = \max_i|x_i|$) appartiene a $F_n$ e $\rho(X) = \rho(m^{-1}X)$, quindi il massimo di $\rho|_{F_n}$ è il massimo di $\rho$.

\textbf{Passo 2: il massimizzatore è un autovettore.} Sia $v$ un massimizzatore di $\rho$. Per ogni $w \perp v$ le funzioni $g(s) = f(v+sw)$ e $h(s) = \|v+sw\|^2$ sono differenziabili in $s$. Poiché $\xi(s) = g(s)/h(s)$ ha un massimo in $s=0$:
\[
  0 = \xi'(0) = \frac{g'(0)h(0) - g(0)h'(0)}{h(0)^2}.
\]
Si calcola $g'(0) = 2\langle A(v),w\rangle$ e $h'(0) = 2\langle v,w\rangle = 0$ (poiché $w \perp v$). Dunque $\langle A(v),w\rangle = 0$ per ogni $w \perp v$, che per la Proposizione~\ref{prop:autovett-crit} significa che $v$ è autovettore.
\end{proof}

\begin{proof}[Dimostrazione del Teorema~\ref{thm:spettrale-reale-I}]
Per induzione su $n = \dim V$. Il caso $n = 1$ è banale. Per $n \geq 2$: per la Proposizione~\ref{prop:esistenza-autovett} esiste un autovettore $v$ di $A$; normalizziamolo: $\|v\| = 1$. Per il Corollario~\ref{cor:A-preserva-ort}, $A$ manda $v^\perp$ in sé, quindi $B \coloneqq A|_{v^\perp}$ è un endomorfismo simmetrico dello spazio euclideo $(v^\perp, \langle\cdot,\cdot\rangle|_{v^\perp})$, che ha dimensione $n-1$. Per ipotesi induttiva, esiste una base ON $\mathcal{C}$ di $v^\perp$ che diagonalizza $B$. Allora $\mathcal{C} \cup \{v\}$ è una base ON di $V$ che diagonalizza $A$.
\end{proof}

\subsection{Implementazione: algoritmo}

Il teorema spettrale fornisce il seguente algoritmo per trovare una base ON che diagonalizza $A$ (o $f$):

\begin{enumerate}
  \item Scegliamo una base ON $\mathcal{C}$ di $V$ e calcoliamo $A = M_\mathcal{C}(f)$ (o $M^\mathcal{C}_\mathcal{C}(A)$).
  \item Calcoliamo le radici $\lambda_1,\ldots,\lambda_m$ del polinomio caratteristico $p_A$ (per il teorema spettrale, sono tutte reali).
  \item Per ogni $\lambda_i$, troviamo una base ON dell'autospazio $V_{\lambda_i}(A)$.
  \item L'unione delle basi degli autospazi è una base ON di $V$ che diagonalizza $A$.
\end{enumerate}

\begin{exbox}
\begin{esempio}
Sia $A = \bigl(\begin{smallmatrix}1&2&3\\2&3&1\\3&1&2\end{smallmatrix}\bigr) \in M_3(\R)$ (simmetrica). Il polinomio caratteristico è $p_A(\lambda) = \lambda^3-6\lambda^2-3\lambda+18$, con radici $\lambda_1 = 6$, $\lambda_2 = \sqrt{3}$, $\lambda_3 = -\sqrt{3}$. Gli autospazi sono
\[
  V_6(A) = \langle (1,1,1)\rangle, \quad V_{\sqrt{3}}(A) = \langle (7-3\sqrt{3},\,-5-\sqrt{3},\,-2+4\sqrt{3})\rangle,
\]
\[
  V_{-\sqrt{3}}(A) = \langle (7+3\sqrt{3},\,-5+\sqrt{3},\,-2-4\sqrt{3})\rangle.
\]
Per la Proposizione~\ref{prop:autospazi-ortogonali} questi vettori sono già ortogonali; normalizzandoli si ottiene la base ON cercata.
\end{esempio}
\end{exbox}

% --------------------------------------------------------
\section{Forme hermitiane}
\label{sec:forme-hermitiane}
% --------------------------------------------------------

Le forme hermitiane sono l'analogo complesso delle forme bilineari simmetriche reali. La differenza essenziale è che la \emph{linearità nella seconda variabile} viene sostituita dalla \emph{coniugato-linearità}.

\subsection{Forme sesquilineari ed hermitiane}

\begin{defbox}
\begin{definizione}[Forma sesquilineare]
\label{def:sesquilineare}
Una \emph{forma sesquilineare} su uno $\C$-spazio vettoriale $V$ è un'applicazione $H \colon V \times V \to \C$ che è:
\begin{itemize}
  \item lineare nella prima variabile: $H(\lambda v+v',w) = \lambda H(v,w)+H(v',w)$,
  \item \emph{coniugato-lineare} nella seconda: $H(v,\lambda w+w') = \bar{\lambda}H(v,w)+H(v,w')$.
\end{itemize}
\end{definizione}
\end{defbox}

\begin{defbox}
\begin{definizione}[Forma hermitiana]
\label{def:forma-herm}
Una forma sesquilineare $H$ su $V$ è \emph{hermitiana} se $H(w,v) = \overline{H(v,w)}$ per ogni $v,w \in V$. Si denota $\mathrm{Herm}(V)$ l'insieme delle forme hermitiane su $V$.
\end{definizione}
\end{defbox}

\begin{osservazione}
Per una forma hermitiana $H$: $H(v,v) = \overline{H(v,v)} \in \R$ per ogni $v$ (i valori diagonali sono reali). I prodotti scalari hermitiani (Capitolo~6) sono esattamente le forme hermitiane definite positive.
\end{osservazione}

$\mathrm{Herm}(V)$ è uno $\R$-spazio vettoriale (non $\C$: se $H$ è hermitiana e $\lambda \notin \R$, l'applicazione $\lambda H$ non è hermitiana).

\subsection{Matrici hermitiane}

\begin{defbox}
\begin{definizione}[Matrice aggiunta e hermitiana]
L'\emph{aggiunta} di $A \in M_{n,n}(\C)$ è $A^* \coloneqq \overline{A}^{\,t}$ (trasposta coniugata). La matrice $A$ è \emph{hermitiana} (o \emph{autoaggiunta}) se $A = A^*$.
\end{definizione}
\end{defbox}

\begin{osservazione}
La forma sesquilineare $\Psi_A(X,Y) \coloneqq X^t \cdot A \cdot \overline{Y}$ è hermitiana se e solo se $A = A^*$. La diagonale di una matrice hermitiana ha entrate reali (poiché $a_{jj} = \overline{a_{jj}}$).
\end{osservazione}

Si denota $M^*_{n,n}(\C)$ il sottospazio \emph{reale} delle matrici hermitiane.

Data una base $\mathcal{B} = \{v_1,\ldots,v_n\}$ di $V$, la \emph{matrice di Gram} di $H \in \mathrm{Herm}(V)$ è $M_\mathcal{B}(H) \coloneqq (H(v_j,v_k)) \in M^*_{n,n}(\C)$.

\begin{proposizione}
L'applicazione $M_\mathcal{B} \colon \mathrm{Herm}(V) \to M^*_{n,n}(\C)$, $H \mapsto M_\mathcal{B}(H)$, è un isomorfismo di $\R$-spazi vettoriali.
\end{proposizione}

\subsection{Equivalenza e classificazione}

Il gruppo $\GL(V)$ agisce su $\mathrm{Herm}(V)$ via $(g \cdot H)(v,w) \coloneqq H(g^{-1}v,g^{-1}w)$.

\begin{defbox}
\begin{definizione}[Equivalenza di forme hermitiane]
$H_1, H_2 \in \mathrm{Herm}(V)$ sono \emph{equivalenti} se esiste $g \in \GL(V)$ tale che $H_1 = g \cdot H_2$.
\end{definizione}
\end{defbox}

In termini di matrici: $A_1, A_2 \in M^*_{n,n}(\C)$ sono equivalenti se esiste $G \in \GL_n(\C)$ con $A_1 = G^t \cdot A_2 \cdot \overline{G} = G^t A_2 (G^*)^t$.

\begin{defbox}
\begin{definizione}[Definitezza e segnatura hermitiana]
Sia $H \in \mathrm{Herm}(V)$.
\begin{itemize}
  \item $H$ è \emph{definita positiva} (risp.\ negativa) se $H(v,v) > 0$ (risp.\ $< 0$) per ogni $v \neq \mathbf{0}$.
  \item La \emph{segnatura positiva} $s^+(H)$ è $\max\{\dim U \mid H|_U > 0\}$.
  \item La \emph{segnatura negativa} $s^-(H)$ è $\max\{\dim U \mid H|_U < 0\}$.
\end{itemize}
\end{definizione}
\end{defbox}

\begin{proposizione}
\label{prop:class-hermitiane}
$H_1, H_2 \in \mathrm{Herm}(V)$ sono equivalenti se e solo se $s^+(H_1) = s^+(H_2)$ e $s^-(H_1) = s^-(H_2)$.
\end{proposizione}

\begin{proof}[Schema della dimostrazione]
Analoga alla classificazione delle forme quadratiche reali (Teorema~\ref{thm:class-reali}).

\textit{Passo 1 (diagonalizzazione):} Se $H \neq 0$, esiste $v \in V$ con $H(v,v) \neq 0$ (si dimostra che, se esistono $v_1,v_2$ con $H(v_1,v_2) \neq 0$, allora dopo un opportuno riscalamento esiste $t_0 \in \R$ con $H(v_1+t_0v_2,v_1+t_0v_2) \neq 0$). Si pone $v^\perp_H = \{w \mid H(v,w)=0\}$ e si itera.

\textit{Passo 2 (segnatura):} In una base diagonalizzante, le entrate diagonali sono reali. La segnatura conta le entrate positive e negative, esattamente come nel Teorema di Sylvester.
\end{proof}

% --------------------------------------------------------
\section{Il teorema spettrale complesso}
\label{sec:spettrale-complesso}
% --------------------------------------------------------

\subsection{L'aggiunto di un operatore}

\begin{defbox}
\begin{definizione}[Operatore aggiunto]
\label{def:aggiunto}
Sia $(V,\langle\cdot,\cdot\rangle)$ uno spazio hermitiano (di dimensione finita). L'\emph{aggiunto} di $T \colon V \to V$ è l'unico operatore $T^* \colon V \to V$ tale che
\[
  \langle T(v),w\rangle = \langle v,T^*(w)\rangle \qquad \forall\, v,w \in V.
\]
$T$ è \emph{autoaggiunto} se $T = T^*$.
\end{definizione}
\end{defbox}

\begin{lemma}
L'aggiunto esiste ed è unico. Se $\mathcal{B}$ è una base ON di $V$, allora $M^\mathcal{B}_\mathcal{B}(T^*) = \overline{M^\mathcal{B}_\mathcal{B}(T)}^{\,t} = M^\mathcal{B}_\mathcal{B}(T)^*$.
\end{lemma}

\begin{proof}
Per ogni $v,w \in V$:
\[
  \langle T(v),w\rangle = (M^\mathcal{B}_\mathcal{B}(T)\,X_\mathcal{B}(v))^t \cdot \overline{X_\mathcal{B}(w)} = X_\mathcal{B}(v)^t \cdot M^\mathcal{B}_\mathcal{B}(T)^t \cdot \overline{X_\mathcal{B}(w)} = X_\mathcal{B}(v)^t \cdot \overline{M^\mathcal{B}_\mathcal{B}(T)^* \cdot X_\mathcal{B}(w)}.
\]
Quindi $T^*$ è l'operatore con matrice $M^\mathcal{B}_\mathcal{B}(T)^*$ in $\mathcal{B}$, ed è unico.
\end{proof}

\begin{osservazione}
$T$ è autoaggiunto se e solo se $M^\mathcal{B}_\mathcal{B}(T)$ è hermitiana (per ogni base ON $\mathcal{B}$). In particolare, su $\C^n$ con il prodotto hermitiano standard, $L_A$ è autoaggiunto se e solo se $A = A^*$. L'insieme degli operatori autoaggiunti è un sottospazio \emph{reale} di $\mathrm{End}(V)$.
\end{osservazione}

\subsection{Il teorema spettrale}

\begin{teorema}[Teorema spettrale complesso]
\label{thm:spettrale-complesso}
Sia $(V,\langle\cdot,\cdot\rangle)$ uno spazio vettoriale hermitiano finitamente generato e $T \colon V \to V$ un operatore autoaggiunto. Allora:
\begin{enumerate}[label=(\alph*)]
  \item tutti gli autovalori di $T$ sono reali;
  \item esiste una base ON di $V$ che diagonalizza $T$.
\end{enumerate}
\end{teorema}

\begin{proof}
\textbf{(a) Autovalori reali.} Sia $\lambda$ autovalore con autovettore $v \neq \mathbf{0}$. Allora
\[
  \lambda\|v\|^2 = \langle T(v),v\rangle = \langle v,T(v)\rangle = \overline{\lambda}\|v\|^2,
\]
e poiché $\|v\|^2 > 0$ si ha $\lambda = \overline{\lambda}$, cioè $\lambda \in \R$.

\textbf{(b) Base ON.} Per induzione su $n = \dim_\C V$. Il caso $n = 1$ è banale. Per $n \geq 2$: poiché $V$ è uno spazio vettoriale \emph{complesso}, il polinomio caratteristico di $T$ ha almeno una radice complessa, che per (a) è reale. Sia $\lambda$ tale autovalore e $v$ un autovettore normalizzato ($\|v\|=1$). Mostriamo che $T(v^\perp) \subseteq v^\perp$: se $w \perp v$,
\[
  \langle T(w),v\rangle = \langle w,T(v)\rangle = \langle w,\lambda v\rangle = \overline{\lambda}\langle w,v\rangle = 0.
\]
Quindi $S \coloneqq T|_{v^\perp}$ è un operatore autoaggiunto su $v^\perp$ (di dimensione $n-1$). Per ipotesi induttiva esiste una base ON $\mathcal{C}$ di $v^\perp$ che diagonalizza $S$. Allora $\mathcal{C} \cup \{v\}$ è la base ON cercata.
\end{proof}

\begin{corollario}
\label{cor:spettrale-forme-herm}
Sia $(V,\langle\cdot,\cdot\rangle)$ uno spazio vettoriale hermitiano (di dimensione finita) e $H$ una forma hermitiana su $V$. Esiste una base ON di $V$ che diagonalizza $H$.
\end{corollario}

\begin{proof}
Per l'Osservazione~\ref{oss:simm-matrice} generalizzata al caso hermitiano, esiste un unico operatore autoaggiunto $A \colon V \to V$ tale che $H(v,w) = \langle A(v),w\rangle$. Per il Teorema~\ref{thm:spettrale-complesso} esiste una base ON che diagonalizza $A$, cioè diagonalizza $H$.
\end{proof}

\begin{osservazione}
Il Teorema spettrale complesso implica quello reale: uno spazio euclideo reale $(V,\langle\cdot,\cdot\rangle)$ si può complessificare in $V_\C = V \otimes_\R \C$ con prodotto hermitiano standard; un endomorfismo simmetrico $A$ su $V$ diventa autoaggiunto su $V_\C$, e il teorema complesso dà la diagonalizzazione, che poi si restringe a $V$.

In alternativa, la dimostrazione per il caso reale (Sezione~\ref{sec:spettrale-reale}) è autonoma e non usa il Teorema fondamentale dell'Algebra.
\end{osservazione}

% --------------------------------------------------------
\section*{Esercizi}
\addcontentsline{toc}{section}{Esercizi del Capitolo 9}
\label{sec:es-cap9}
% --------------------------------------------------------

\begin{esercizio}
Sia $f(x_1,x_2) = x_1^2+2x_1x_2+3x_2^2$. Trovate una base di $\R^2$, ON per il prodotto euclideo standard, che diagonalizza $f$.
\end{esercizio}

\begin{esercizio}
Sia $A = \bigl(\begin{smallmatrix}2&3\\3&5\end{smallmatrix}\bigr)$.
\begin{enumerate}[label=(\arabic*)]
  \item Verificate che $\langle X,Y\rangle \coloneqq X^t A Y$ è un prodotto euclideo su $\R^2$.
  \item Sia $f(x_1,x_2) = x_1^2+2x_1x_2-3x_2^2$. Trovate una base di $\R^2$ che diagonalizza $f$ e che è ON per il prodotto euclideo del punto~(1).
\end{enumerate}
\end{esercizio}

\begin{esercizio}
Sia $V$ uno spazio vettoriale complesso di dimensione finita e $H \in \mathrm{Herm}(V)$. Dimostrate che
\[
  \max\{ \dim U \mid U \subseteq V,\, H|_U = 0 \} = \dim\ker(H) + \min\{s^+(H),s^-(H)\}.
\]
(Analogo hermitiano dell'Esercizio del Capitolo~8 sui sottospazi isotropi.)
\end{esercizio}

\begin{esercizio}
Sia $A \in M^*_{n,n}(\C)$ una matrice hermitiana. Dimostrate che $A$ è definita positiva se e solo se $\det A^{(p)} > 0$ per ogni $p = 1,\ldots,n$, dove $A^{(p)}$ è la sottomatrice $p\times p$ superiore sinistra di $A$. Dimostrate inoltre che $A$ è definita negativa se e solo se i determinanti $\det A^{(1)}, \det A^{(2)},\ldots$ si alternano in segno con $\det A^{(1)} < 0$.
\end{esercizio}

\begin{esercizio}
Sia $A = \bigl(\begin{smallmatrix}0&i&1\\-i&3&2-i\\1&2+i&1\end{smallmatrix}\bigr) \in M^*_{3,3}(\C)$ e $H_A(X,Y) = X^t A\,\overline{Y}$. Determinate $s^+(H_A)$ e $s^-(H_A)$.
\end{esercizio}

\begin{esercizio}
Dimostrate che il polinomio caratteristico di un operatore autoaggiunto su uno spazio vettoriale hermitiano ha tutti i coefficienti reali.
\end{esercizio}

\begin{esercizio}[Dimostrazione alternativa del teorema spettrale complesso]
Sia $(V,\langle\cdot,\cdot\rangle)$ hermitiano di dimensione finita e $A \colon V \to V$ autoaggiunto. Dimostrate il Teorema~\ref{thm:spettrale-complesso} senza usare il Teorema fondamentale dell'Algebra, seguendo i passi:
\begin{enumerate}[label=(\arabic*)]
  \item Dimostrate che $f(v) \coloneqq \langle A(v),v\rangle / \|v\|^2$ ha un massimo su $V \setminus \{\mathbf{0}\}$.
  \item Sia $v_0$ un massimizzatore: dimostrate che $v_0^\perp \subseteq \{w \mid \langle A(w),v_0\rangle = 0\}$.
  \item Conclude che $v_0$ è un autovettore di $A$.
\end{enumerate}
\end{esercizio}

\begin{esercizio}
Osservate che il Teorema spettrale complesso implica quello reale: descrivete come dedurre il Teorema~\ref{thm:spettrale-reale-I} dal Teorema~\ref{thm:spettrale-complesso}.
\end{esercizio}

\begin{esercizio}[Classificazione per l'azione di $U(V)$]
Sia $(V,\langle\cdot,\cdot\rangle)$ uno spazio hermitiano di dimensione finita. Diciamo che $H_1, H_2 \in \mathrm{Herm}(V)$ sono \emph{unitariamente equivalenti} se esiste $g \in U(V)$ tale che $H_1(v,w) = H_2(g^{-1}v,g^{-1}w)$.
\begin{enumerate}[label=(\alph*)]
  \item Per $H \in \mathrm{Herm}(V)$, definite il \emph{polinomio caratteristico di $H$} come quello dell'unico operatore autoaggiunto $A$ con $H = H_A$.
  \item Dimostrate che $H_1$ e $H_2$ sono unitariamente equivalenti se e solo se i loro polinomi caratteristici coincidono.
  \item Siano $A, B \in M^*_{n,n}(\C)$. Dimostrate che esiste $U \in U_n(\C)$ con $A = U^* B U$ se e solo se $p_A = p_B$.
\end{enumerate}
\end{esercizio}

\begin{esercizio}[Forme hermitiane e bilineari simmetriche reali]
Sia $H$ una forma hermitiana su uno spazio vettoriale complesso $V$. Definite
\[
  g_H \colon V \times V \to \R, \quad g_H(v,w) \coloneqq \tfrac{1}{2}(H(v,w)+\overline{H(v,w)}) = \mathrm{Re}\,H(v,w).
\]
\begin{enumerate}[label=(\alph*)]
  \item Dimostrate che $g_H$ è una forma bilineare simmetrica sullo $\R$-spazio vettoriale $V$.
  \item Dimostrate che $g_H(iv,iw) = g_H(v,w)$ per ogni $v,w \in V$; in particolare $g_H(v,iv) = 0$.
  \item Viceversa, sia $g \colon V \times V \to \R$ bilineare simmetrica su $V$ visto come $\R$-spazio, con $g(iv,iw) = g(v,w)$ per ogni $v,w$. Dimostrate che esiste un'unica $H \in \mathrm{Herm}(V)$ tale che $g = g_H$.
\end{enumerate}
\end{esercizio}

% ============================================================
\chapter{Coniche, quadriche e iperquadriche}
\label{cap:coniche}
% ============================================================

\section*{Introduzione al capitolo}

Abbiamo finora costruito un imponente edificio di strutture astratte: spazi vettoriali, applicazioni lineari, prodotti scalari, forme quadratiche. In questo capitolo raccogliamo i frutti di tutto questo lavoro applicandolo a uno dei temi più antichi e visivi della geometria: le \emph{curve e superfici di secondo grado}.

I Greci avevano studiato le sezioni coniche — ellisse, iperbole, parabola — geometricamente, come intersezioni di un cono con un piano. Apollonio di Perga, nel III secolo a.C., aveva già classificato tutte le possibili sezioni. Quello che offriamo qui è la versione moderna e coordinata di quella classificazione: ogni conica è (a meno di isometria) descrivibile da un'equazione canonica, e il passaggio alla forma canonica si riduce alla diagonalizzazione di una forma quadratica seguita da una traslazione.

Il salto concettuale è il seguente: un'equazione polinomiale di grado 2 in due variabili definisce una ``curva'', ma di che tipo? Il Teorema di classificazione (Teorema~\ref{thm:forme-canoniche-2d}) risponde in modo definitivo: a meno di isometrie, le uniche possibilità sono ellisse, ellisse complessa (insieme vuoto su $\R$), iperbole, parabola, e alcune configurazioni degeneri (rette, punto). Questo non è affatto ovvio a priori — l'equazione generale $ax^2 + bxy + cy^2 + dx + ey + f = 0$ ha sei parametri liberi, eppure solo un numero finito di tipi qualitativi.

Lo stesso ragionamento si estende in dimensione $n$: le \emph{iperquadriche} di uno spazio affine euclideo ammettono una classificazione completa (Teorema~\ref{thm:classeq}), e in dimensione 3 otteniamo le famose \emph{quadriche}: ellissoide, iperboloidi, paraboloidi.

% ============================================================
\section{Ellissi, iperboli, parabole}
\label{sec:coniche-def}
% ============================================================

Nel seguito, $S$ denota un \emph{piano affine euclideo}, cioè uno spazio affine euclideo di dimensione 2. Indichiamo con $d(P,Q)$ la distanza euclidea tra due punti, e con $d(P, r)$ la distanza di un punto da una retta.

\subsection{Ellissi}

\begin{defbox}
\begin{definizione}[Ellisse]
\label{def:ellisse}
Un sottoinsieme $\mathcal{C} \subseteq S$ è un'\emph{ellisse} se esiste un riferimento ortonormale di coordinate $(x,y)$ in $S$ tale che
\[
  \mathcal{C} = \left\{ P(x,y) \in S \;\middle|\; \frac{x^2}{a^2} + \frac{y^2}{b^2} = 1 \right\},
\]
con $a \geq b > 0$. I parametri $a$ e $b$ si chiamano rispettivamente \emph{semiasse maggiore} e \emph{semiasse minore}.
\end{definizione}
\end{defbox}

\begin{esempio}
Un cerchio di raggio $r$ è un'ellisse con $a = b = r$. L'equazione canonica diventa $x^2 + y^2 = r^2$.
\end{esempio}

\begin{osservazione}
\label{oss:simmetrie-ellisse}
L'ellisse di equazione $\frac{x^2}{a^2} + \frac{y^2}{b^2} = 1$ gode delle seguenti simmetrie:
\begin{enumerate}[label=(\arabic*)]
  \item Le rette $\{y=0\}$ e $\{x=0\}$ sono assi di simmetria. Se $a \neq b$ non ne esistono altri; se $a = b$ (cerchio) ogni retta per l'origine è un asse di simmetria.
  \item L'origine è un centro di simmetria: la rotazione di angolo $\pi$ attorno all'origine manda $\mathcal{C}$ in sé. Si può dimostrare che non esistono altri centri di simmetria (se una figura ha due centri distinti, contiene una traslazione non banale nel suo gruppo di simmetria, e quindi è illimitata).
\end{enumerate}
\end{osservazione}

\begin{defbox}
\begin{definizione}[Fuochi ed eccentricità dell'ellisse]
\label{def:fuochi-ellisse}
Sia $\mathcal{C}$ l'ellisse di equazione $\frac{x^2}{a^2}+\frac{y^2}{b^2}=1$ con $a \geq b > 0$. Si pone
\[
  c \coloneqq \sqrt{a^2 - b^2}.
\]
I \emph{fuochi} di $\mathcal{C}$ sono i punti $F_1 \coloneqq (c,0)$ e $F_2 \coloneqq (-c,0)$. L'\emph{eccentricità} è
\[
  e \coloneqq \frac{c}{a} = \frac{\sqrt{a^2-b^2}}{a} \in [0,1).
\]
\end{definizione}
\end{defbox}

\begin{osservazione}
L'eccentricità misura quanto l'ellisse si discosta da un cerchio: $e=0$ corrisponde al cerchio, mentre $e \to 1$ corrisponde a un'ellisse molto allungata (con $b \to 0$). Per un cerchio i fuochi coincidono con il centro.
\end{osservazione}

La proprietà geometrica fondamentale dei fuochi è la seguente.

\begin{proposizione}[Caratterizzazione focale dell'ellisse]
\label{prop:foc-ellisse}
Sia $\mathcal{C}$ l'ellisse di equazione $\frac{x^2}{a^2}+\frac{y^2}{b^2}=1$ con fuochi $F_1, F_2$. Allora
\[
  \mathcal{C} = \bigl\{ P \in S \;\big|\; d(P,F_1) + d(P,F_2) = 2a \bigr\}.
\]
Viceversa, dati due punti $F_1, F_2 \in S$ e un valore $d > d(F_1,F_2)$, l'insieme dei punti equidistanti (nel senso della somma) è un'ellisse.
\end{proposizione}

\begin{proof}[Idea della dimostrazione]
La prima asserzione si verifica con un calcolo diretto: si pone $r_1 = d(P,F_1)$, $r_2 = d(P,F_2)$ con $P=(x,y)$, e si mostra che $r_1 + r_2 = 2a$ è equivalente all'equazione cartesiana. Per il viceversa: dati $F_1, F_2$ e $d > d(F_1,F_2)$, si pone $2c = d(F_1,F_2)$ e $2a = d$; poiché $a > c$, si può definire $b \coloneqq \sqrt{a^2-c^2}$. Esiste allora un'isometria $\Phi \in \Isom(S)$ che manda i fuochi $(\pm c, 0)$ in $F_1, F_2$, e l'immagine tramite $\Phi$ dell'ellisse canonica è l'insieme cercato.
\end{proof}

\begin{defbox}
\begin{definizione}[Direttrici dell'ellisse]
\label{def:direttrici-ellisse}
Sia $\mathcal{C}$ l'ellisse con $a \geq b > 0$, $c = \sqrt{a^2-b^2}$ e $c > 0$ (cioè $\mathcal{C}$ non è un cerchio). Le \emph{direttrici} di $\mathcal{C}$ sono le rette
\[
  D_1 \colon x = \frac{a^2}{c}, \qquad D_2 \colon x = -\frac{a^2}{c}.
\]
\end{definizione}
\end{defbox}

\begin{proposizione}[Proprietà fuoco-direttrice]
\label{prop:foc-dir-ellisse}
Sia $\mathcal{C}$ l'ellisse con eccentricità $e$ e direttrici $D_1, D_2$. Per $i=1,2$,
\[
  \mathcal{C} = \bigl\{ P \in S \;\big|\; d(P, F_i) = e \cdot d(P, D_i) \bigr\}.
\]
\end{proposizione}

\begin{osservazione}
Poiché $0 < e < 1$, il punto $P$ è \emph{più vicino} al fuoco che alla direttrice corrispondente. Le direttrici si trovano al di fuori dell'ellisse (poiché $a^2/c = a/e > a$).
\end{osservazione}

\subsection{Iperboli}

\begin{defbox}
\begin{definizione}[Iperbole]
\label{def:iperbole}
Un sottoinsieme $\mathcal{C} \subseteq S$ è un'\emph{iperbole} se esiste un riferimento ortonormale $(x,y)$ tale che
\[
  \mathcal{C} = \left\{ P(x,y) \in S \;\middle|\; \frac{x^2}{a^2} - \frac{y^2}{b^2} = 1 \right\},
\]
con $a, b > 0$.
\end{definizione}
\end{defbox}

\begin{osservazione}
A differenza dell'ellisse, un'iperbole è una curva \emph{non connessa}: è l'unione di due \emph{rami},
\[
  \mathcal{C} = \left\{ \left(\sqrt{a^2 + \frac{a^2 y^2}{b^2}},\, y\right) \;\middle|\; y \in \R \right\}
  \cup
  \left\{ \left(-\sqrt{a^2 + \frac{a^2 y^2}{b^2}},\, y\right) \;\middle|\; y \in \R \right\}.
\]
Assi di simmetria e centro di simmetria sono gli stessi dell'ellisse.

Le iperbole ammettono inoltre due \emph{asintoti}: le rette $y = \pm \frac{b}{a} x$. Per $|x| \to \infty$, i punti di $\mathcal{C}$ si avvicinano sempre più a queste rette.
\end{osservazione}

\begin{defbox}
\begin{definizione}[Fuochi ed eccentricità dell'iperbole]
\label{def:fuochi-iper}
Sia $\mathcal{C}$ l'iperbole di equazione $\frac{x^2}{a^2}-\frac{y^2}{b^2}=1$. Si pone
\[
  c \coloneqq \sqrt{a^2 + b^2}.
\]
I \emph{fuochi} sono $F_1 \coloneqq (c,0)$ e $F_2 \coloneqq (-c,0)$, e l'\emph{eccentricità} è
\[
  e \coloneqq \frac{c}{a} = \frac{\sqrt{a^2+b^2}}{a} > 1.
\]
\end{definizione}
\end{defbox}

Si noti la differenza rispetto all'ellisse: per l'iperbole $c^2 = a^2 + b^2$ (con segno $+$) e $e > 1$.

\begin{proposizione}[Caratterizzazione focale dell'iperbole]
\label{prop:foc-iper}
Sia $\mathcal{C}$ l'iperbole con fuochi $F_1, F_2$. Allora
\[
  \mathcal{C} = \bigl\{ P \in S \;\big|\; |d(P,F_1) - d(P,F_2)| = 2a \bigr\}.
\]
Viceversa, dati punti distinti $F_1, F_2$ e $0 < d < d(F_1,F_2)$, l'insieme dei punti con differenza di distanze uguale a $d$ è un'iperbole.
\end{proposizione}

\begin{defbox}
\begin{definizione}[Direttrici dell'iperbole]
\label{def:direttrici-iper}
Le \emph{direttrici} dell'iperbole di equazione $\frac{x^2}{a^2}-\frac{y^2}{b^2}=1$ sono
\[
  D_1 \colon x = \frac{a^2}{c}, \qquad D_2 \colon x = -\frac{a^2}{c},
\]
con $c = \sqrt{a^2+b^2}$.
\end{definizione}
\end{defbox}

\begin{proposizione}[Proprietà fuoco-direttrice per l'iperbole]
\label{prop:foc-dir-iper}
I due rami dell'iperbole sono gli insiemi dei punti $P$ tali che $d(P, F_i) = e \cdot d(P, D_i)$ per $i=1$ e $i=2$ rispettivamente.
\end{proposizione}

\subsection{Parabole}

La parabola emerge come caso limite comune sia delle ellissi che delle iperboli, quando l'eccentricità vale esattamente $1$.

\begin{defbox}
\begin{definizione}[Parabola]
\label{def:parabola}
Un sottoinsieme $\mathcal{C} \subseteq S$ è una \emph{parabola} se esiste un riferimento ortonormale $(x,y)$ e una costante $c > 0$ tali che
\[
  \mathcal{C} = \left\{ P(x,y) \in S \;\middle|\; y = \frac{x^2}{4c} \right\}.
\]
\end{definizione}
\end{defbox}

\begin{defbox}
\begin{definizione}[Fuoco e direttrice della parabola]
\label{def:fuoco-par}
La parabola di equazione $y = \frac{x^2}{4c}$ ha:
\begin{itemize}
  \item \emph{fuoco} $F \coloneqq (0, c)$,
  \item \emph{direttrice} $D \colon y = -c$.
\end{itemize}
\end{definizione}
\end{defbox}

\begin{proposizione}[Caratterizzazione della parabola]
\label{prop:foc-par}
La parabola di equazione $y = \frac{x^2}{4c}$ è l'insieme dei punti equidistanti dal fuoco $F$ e dalla direttrice $D$:
\[
  \mathcal{C} = \bigl\{ P \in S \;\big|\; d(P,F) = d(P,D) \bigr\}.
\]
Viceversa, l'insieme dei punti equidistanti da un punto $F$ e da una retta $r$ non passante per $F$ è una parabola.
\end{proposizione}

\begin{proof}
Poniamo $P = (x,y)$. Si ha $d(P,F)^2 = x^2 + (y-c)^2$ e $d(P,D) = y+c$ (per $P$ al di sopra di $D$). L'uguaglianza $d(P,F) = d(P,D)$ diventa
\[
  x^2 + (y-c)^2 = (y+c)^2 \iff x^2 = 4cy,
\]
che è l'equazione della parabola (riscritta come $y = x^2/(4c)$). Il viceversa si dimostra usando un'isometria che porta $F$ in $(0,c)$ e $r$ in $\{y = -c\}$.
\end{proof}

% ============================================================
\section{Classificazione delle coniche a meno di isometrie}
\label{sec:class-coniche}
% ============================================================

\subsection{Funzioni polinomiali su uno spazio affine}

L'obiettivo di questa sezione è classificare tutte le figure piane definite da un'equazione polinomiale di grado 2. Il punto di partenza è osservare che tale equazione non dipende dalla scelta del sistema di riferimento, purché si considerino solo cambiamenti di coordinate affini.

\begin{defbox}
\begin{definizione}[Funzione polinomiale su uno spazio affine]
\label{def:pol-affine}
Sia $S$ uno spazio affine su $\K$ (di dimensione finita). Una funzione $f \colon S \to \K$ è \emph{polinomiale di grado $d$} se, fissato un riferimento affine $R_O(O, \mathcal{B})$ con coordinate $X_{R_O} \colon S \xrightarrow{\sim} \mathbb{A}^n(\K)$, la funzione composta
\[
  \mathbb{A}^n(\K) \xrightarrow{f \circ X_{R_O}^{-1}} \K
\]
è un polinomio di grado $d$. Il grado non dipende dalla scelta del riferimento.
\end{definizione}
\end{defbox}

\begin{osservazione}
Su uno spazio \emph{vettoriale} ha senso parlare di polinomi omogenei, poiché esiste un'origine privilegiata. Su uno spazio \emph{affine} no: la condizione di omogeneità $f(\lambda X) = \lambda^d f(X)$ non è invariante per cambio di origine.
\end{osservazione}

\begin{defbox}
\begin{definizione}[Ipersuperficie algebrica]
\label{def:iper-alg}
Sia $S$ uno spazio affine su $\K$ e $f \colon S \to \K$ polinomiale. L'\emph{ipersuperficie algebrica} di equazione $f$ è
\[
  V(f) \coloneqq \{ P \in S \mid f(P) = 0 \}.
\]
Se $\dim S = 2$, $V(f)$ si dice \emph{curva algebrica}; se $\dim S = 3$, \emph{superficie algebrica}.
\end{definizione}
\end{defbox}

\subsection{Equivalenza per isometrie}

\begin{defbox}
\begin{definizione}[Equivalenza per isometrie]
\label{def:equiv-iso}
Due funzioni polinomiali $f, g \colon S \to \R$ su un piano affine euclideo $S$ sono \emph{equivalenti per isometrie} (in simboli $f \sim g$) se esistono $\Phi \in \Isom(S)$ e $\lambda \in \R \setminus \{0\}$ tali che
\[
  g(P) = \lambda \cdot f(\Phi(P)) \qquad \forall\, P \in S.
\]
\end{definizione}
\end{defbox}

\begin{osservazione}
Il fattore $\lambda \neq 0$ è necessario perché $V(\lambda f) = V(f)$ per ogni $\lambda \neq 0$, e vogliamo che figure identiche corrispondano a funzioni equivalenti.

Si noti che $f \sim g$ implica $V(g) = \Phi^{-1}(V(f))$ (le figure sono isometricamente equivalenti), ma il viceversa è falso in generale.
\end{osservazione}

\begin{osservazione}
Le condizioni ``$f \sim g$'' e ``$V(f)$ è isometrico a $V(g)$'' non sono equivalenti. Per esempio, le funzioni $f = x^2+y^2+1$ e $g = 3x^2+y^2+1$ hanno entrambe insieme degli zeri vuoto, ma non sono equivalenti per isometrie (hanno coefficienti diversi in qualunque base ON).
\end{osservazione}

\subsection{Il teorema di classificazione}

\begin{teorema}[Forme canoniche delle coniche]
\label{thm:forme-canoniche-2d}
Sia $f \colon S \to \R$ una funzione polinomiale non nulla di grado 2 su un piano affine euclideo. Allora esiste un sistema di coordinate ortonormale nel quale $f$ è uguale (a meno di scalare non nullo) a una delle funzioni della Tabella~\ref{tab:forme-can-2d}. Inoltre, due tali funzioni sono equivalenti per isometrie se e solo se appartengono alla stessa riga della tabella con gli stessi parametri.
\end{teorema}

\begin{table}[h!]
\centering
\caption{Forme canoniche di funzioni polinomiali di grado 2 in due variabili.}
\label{tab:forme-can-2d}
\renewcommand{\arraystretch}{1.6}
\begin{tabular}{llll}
\toprule
\textbf{Equazione canonica} & \textbf{Nome} & \textbf{Tipo} & \textbf{$V(f) \subseteq \R^2$} \\
\midrule
$\dfrac{x^2}{a^2}+\dfrac{y^2}{b^2}-1,\ a\geq b>0$ & Ellisse & Non degenere & Curva chiusa \\[4pt]
$\dfrac{x^2}{a^2}+\dfrac{y^2}{b^2}+1,\ a\geq b>0$ & Ellisse complessa & Non degenere & $\emptyset$ \\[4pt]
$\dfrac{x^2}{a^2}-\dfrac{y^2}{b^2}-1,\ a,b>0$ & Iperbole & Non degenere & Due rami \\[4pt]
$\dfrac{x^2}{4c}-y,\ c>0$ & Parabola & Non degenere & Curva aperta \\[4pt]
\midrule
$\dfrac{x^2}{a^2}+\dfrac{y^2}{b^2},\ a\geq b>0$ & Coppia di rette complesse & Degenere & $\{(0,0)\}$ \\[4pt]
$\dfrac{x^2}{a^2}-\dfrac{y^2}{b^2},\ a,b>0$ & Coppia di rette incidenti & Degenere & Due rette \\[4pt]
$\dfrac{x^2}{a^2}-1,\ a>0$ & Coppia di rette parallele & Degenere & Due rette $\|$ \\[4pt]
$\dfrac{x^2}{a^2}+1,\ a>0$ & Coppia di rette complesse $\|$ & Degenere & $\emptyset$ \\[4pt]
$x^2$ & Retta doppia & Degenere & Una retta \\
\bottomrule
\end{tabular}
\end{table}

\begin{proof}[Dimostrazione (schema)]
La dimostrazione procede in due passi.

\textbf{Passo 1 (unicità delle forme canoniche).} Si mostra che funzioni con nomi diversi nella tabella non sono equivalenti (guardando alla geometria di $V(f)$), e che due funzioni dello stesso tipo, come $E_{a,b}$ e $E_{a',b'}$, sono equivalenti se e solo se $a=a'$ e $b=b'$.

Per le coniche non degeneri con zero vuoto o singolo (ellisse complessa, coppia di rette complesse), l'argomento è meno immediato. Per esempio, per l'ellisse complessa: se $E_{a,b}^+ \sim E_{a',b'}^+$, esiste $\Phi \in \Isom(S)$ e $\lambda \neq 0$ con $E_{a,b}^+(P) = \lambda E_{a',b'}^+(\Phi P)$. Poiché il minimo di $E_{a,b}^+$ su $\R^2$ è $1$ (raggiunto in $(0,0)$) e lo stesso vale per $E_{a',b'}^+$, si conclude che $\Phi(0,0)=(0,0)$, quindi $\Phi$ è lineare ortogonale. L'uguaglianza delle matrici diagonalizzate implica $(a,b)=(a',b')$.

\textbf{Passo 2 (esistenza).} Scriviamo
\[
  f(x_1, x_2) = X^t M X + \mu x_1 + \nu x_2 + \theta,
\]
con $M \in \mathrm{Sym}_2(\R)$ non nulla e $\mu, \nu, \theta \in \R$. Per il Teorema Spettrale (Teorema~\ref{thm:spettrale-reale-I}) esiste $A \in O_2(\R)$ che diagonalizza $M$. Dopo questo cambio di base, $f$ diventa $\alpha y_1^2 + \beta y_2^2 + \mu' y_1 + \nu' y_2 + \theta'$.

\textit{Caso (a): $\alpha \neq 0$ e $\beta \neq 0$.} Completando i quadrati:
\[
  \alpha\!\left(y_1+\frac{\mu'}{2\alpha}\right)^2 + \beta\!\left(y_2+\frac{\nu'}{2\beta}\right)^2 + \gamma.
\]
La traslazione $z_i = y_i + (\cdots)$ porta $f$ nella forma $\alpha z_1^2 + \beta z_2^2 + \gamma$. Se $\gamma \neq 0$, dividendo per $\gamma$ si ottiene una delle forme $E_{a,b}$, $E^+_{a,b}$ o $I_{a,b}$ (a seconda dei segni di $\alpha/\gamma$ e $\beta/\gamma$). Se $\gamma = 0$, si ottiene $W_{a,b}$ o $V_{a,b}$.

\textit{Caso (b): uno tra $\alpha, \beta$ è nullo.} Un calcolo analogo porta a una delle forme $P_c$, $L_a$, $H_a$ o $D$.
\end{proof}

\subsection{Omogeneizzazione e criteri pratici}

Un metodo efficiente per determinare il tipo di una conica — senza procedere all'esplicita riduzione — usa l'\emph{omogeneizzazione}.

\begin{defbox}
\begin{definizione}[Omogeneizzazione]
\label{def:omog}
Sia $f \in \R[x_1, x_2]$ un polinomio di grado 2. La sua \emph{omogeneizzazione} è la forma quadratica in tre variabili
\[
  F(x_0, x_1, x_2) \coloneqq x_0^2 \cdot f\!\left(\frac{x_1}{x_0}, \frac{x_2}{x_0}\right).
\]
\end{definizione}
\end{defbox}

\begin{esempio}
Se $f(x_1,x_2) = x_1^2 - 3x_1 x_2 + 5x_2^2 - 3x_1 + 4$, allora
\[
  F(x_0,x_1,x_2) = x_1^2 - 3x_1 x_2 + 5x_2^2 - 3x_0 x_1 + 4x_0^2.
\]
\end{esempio}

Il legame tra equivalenza di $f$ e equivalenza di $F$ è espresso dal seguente lemma.

\begin{lemma}
\label{lem:omog-equiv}
Se $f \sim g$, allora le corrispondenti omogeneizzazioni $F$ e $G$ sono equivalenti (a meno di scalare non nullo) sotto l'azione di $\GL_3(\R)$ che fissa il sottospazio $\{x_0 = 0\}$.
\end{lemma}

\begin{proof}
Segue dall'identità $G(y_0, y_1, y_2) = F(y_0,\, a_{11}y_1 + a_{12}y_2 + b_1 y_0,\, a_{21}y_1 + a_{22}y_2 + b_2 y_0)$ per un cambio di coordinate affine ortonormale $X = AY + B$.
\end{proof}

\begin{proposizione}[Algoritmo di classificazione tramite segnatura]
\label{prop:class-segn}
Sia $f \in \R[x_1,x_2]$ di grado 2 e sia $F$ la sua omogeneizzazione, con matrice simmetrica associata $A \in M_{3 \times 3}(\R)$.
\begin{enumerate}[label=(\Roman*)]
  \item Se $F$ è non degenere (cioè $\det A \neq 0$):
  \begin{enumerate}[label=(\alph*)]
    \item se $F|_{\{x_0=0\}}$ è definita ma $F$ non lo è ($F$ ha segnatura $(2,1)$ o $(1,2)$): $f$ è un'\emph{ellisse};
    \item se $F|_{\{x_0=0\}}$ è non degenere e non definita (segnatura $(1,1)$): $f$ è un'\emph{iperbole};
    \item se $F$ è definita (segnatura $(3,0)$ o $(0,3)$): $V(f) = \emptyset$ (ellisse complessa);
    \item se $F|_{\{x_0=0\}}$ è degenere: $f$ è una \emph{parabola}.
  \end{enumerate}
  \item Se $F$ è degenere: $f$ è una conica degenere (coppia di rette, retta doppia, o punto).
\end{enumerate}
\end{proposizione}

\subsection{Il centro di una conica}

\begin{proposizione}[Calcolo del centro]
\label{prop:centro-conica}
Sia $\mathcal{C} \subseteq S$ un'ellisse o un'iperbole con equazione $f(x_1,x_2) = 0$ in un riferimento ON, e sia $F$ l'omogeneizzazione di $f$ con matrice associata $A$. L'unico vettore $(a_0, a_1, a_2)^t$ tale che $A(a_0, a_1, a_2)^t = (1, 0, 0)^t$ soddisfa $a_0 \neq 0$, e il centro di $\mathcal{C}$ ha coordinate $\bigl(\frac{a_1}{a_0}, \frac{a_2}{a_0}\bigr)$.
\end{proposizione}

\begin{proof}
Il risultato è vero per le forme canoniche (verifica diretta). Poiché ogni conica a centro è equivalente per isometrie alla sua forma canonica, e la Proposizione~\ref{prop:class-segn} (unitamente al Lemma~\ref{lem:omog-equiv}) è invariante per cambi di coordinate ON, la formula vale in qualunque riferimento.
\end{proof}

\subsection{Esempi di riduzione alla forma canonica}

\begin{exbox}
\begin{esempio}[Un'iperbole in coordinate oblique]
\label{es:iperbole-calc}
Sia $\mathcal{C} \subseteq E^2(\R)$ la conica di equazione
\[
  \mathcal{C} \colon 2x_1^2 + 72x_1 x_2 + 23x_2^2 + 25x_1 = 0.
\]
Determiniamo il tipo di $\mathcal{C}$ e la sua forma canonica.

\textbf{Passo 1: diagonalizzazione della parte quadratica.} La matrice della forma quadratica $2x_1^2 + 72x_1x_2 + 23x_2^2$ è
\[
  M = \begin{pmatrix} 2 & 36 \\ 36 & 23 \end{pmatrix}.
\]
Gli autovalori sono $50$ e $-25$, con autovettori $(3,4)$ e $(4,-3)$. Normalizzando:
\[
  \mathcal{B} = \left\{ \left(\tfrac{3}{5}, \tfrac{4}{5}\right),\, \left(\tfrac{4}{5}, -\tfrac{3}{5}\right) \right\}.
\]
Nel riferimento $R_O(O, \mathcal{B})$ l'equazione di $\mathcal{C}$ diventa $50y_1^2 - 25y_2^2 + 15y_1 + 20y_2 = 0$, ovvero (dividendo per 5)
\[
  10y_1^2 - 5y_2^2 + 3y_1 + 4y_2 = 0.
\]

\textbf{Passo 2: completamento dei quadrati.} Scriviamo
\[
  10\!\left[\left(y_1 + \tfrac{3}{20}\right)^2 - \tfrac{9}{400}\right] - 5\!\left[\left(y_2 - \tfrac{2}{5}\right)^2 - \tfrac{4}{25}\right] = 0.
\]
Ponendo $z_1 = y_1 + \frac{3}{20}$, $z_2 = y_2 - \frac{2}{5}$:
\[
  10z_1^2 - 5z_2^2 + \frac{23}{40} = 0 \implies \frac{200}{23}z_2^2 - \frac{400}{23}z_1^2 = 1
\]
(dopo aver scambiato $w_1 = z_2$, $w_2 = z_1$). Quindi $\mathcal{C}$ è un'\textbf{iperbole} di semiassi
\[
  a^2 = \frac{23}{200}, \quad b^2 = \frac{23}{400}.
\]
\end{esempio}
\end{exbox}

\begin{exbox}
\begin{esempio}[Una parabola: fuoco e direttrice]
\label{es:parabola-calc}
Sia $\mathcal{C}$ la conica di equazione
\[
  25x_1^2 - 120x_1x_2 + 144x_2^2 - 494x_1 - 572x_2 + 169 = 0.
\]

\textbf{Passo 1: forma quadratica.} La matrice è
\[
  M = \begin{pmatrix} 25 & -60 \\ -60 & 144 \end{pmatrix}, \quad \det M = 25\cdot144 - 60^2 = 0.
\]
Poiché $\det M = 0$, la parte quadratica è degenere: un autovalore è $0$. Gli autovalori sono $169$ e $0$, con autovettori normalizzati $\left(\frac{5}{13}, -\frac{12}{13}\right)$ e $\left(\frac{12}{13}, \frac{5}{13}\right)$. Il cambio di base
\[
  \begin{pmatrix}x_1\\x_2\end{pmatrix} = \begin{pmatrix}\frac{5}{13} & \frac{12}{13} \\ -\frac{12}{13} & \frac{5}{13}\end{pmatrix}\begin{pmatrix}y_1\\y_2\end{pmatrix}
\]
porta l'equazione nella forma $169y_1^2 + 338y_1 - 676y_2 + 169 = 0$, ovvero $y_1^2 + 2y_1 - 4y_2 + 1 = 0$.

\textbf{Passo 2: completamento del quadrato.} Ponendo $z_1 = y_1 + 1$, $z_2 = y_2$:
\[
  z_1^2 = 4z_2, \quad \text{ovvero} \quad z_2 = \frac{z_1^2}{4}.
\]
Questa è una parabola con $c = 1$. Fuoco e direttrice nelle coordinate $(z_1, z_2)$: $F = (0,1)$, $D\colon z_2 = -1$. Risalendo ai cambi di variabili, nelle coordinate $(y_1, y_2)$: $F = (-1, 1)$, $D\colon y_2 = -1$.

Usando la matrice di cambio di base inversa, si trovano le coordinate originali:
\[
  F = \left(\frac{7}{13}, \frac{17}{13}\right), \qquad D \colon 12x_1 + 5x_2 + 13 = 0.
\]
\end{esempio}
\end{exbox}

\begin{exbox}
\begin{esempio}[Un'iperbole: verifica e calcolo del centro]
\label{es:iper-centro}
Sia $\mathcal{C} \subseteq E^2(\R)$ la conica di equazione
\[
  f(x_1,x_2) \coloneqq x_1^2 + 3x_1 x_2 + 2x_2^2 + x_1 - 1 = 0.
\]

\textbf{(a) Tipo di conica.} L'omogeneizzazione è
\[
  F(x_0,x_1,x_2) = x_1^2 + 3x_1x_2 + 2x_2^2 + x_0x_1 - x_0^2,
\]
con matrice simmetrica
\[
  A = \begin{pmatrix} -1 & \frac{1}{2} & 0 \\ \frac{1}{2} & 1 & \frac{3}{2} \\ 0 & \frac{3}{2} & 2 \end{pmatrix}, \quad \det A = -\frac{1}{4} \neq 0.
\]
Poiché $\det A \neq 0$, la conica è non degenere. La parte quadratica (ristretta a $x_0=0$) è $q(x_1,x_2) = x_1^2 + 3x_1x_2 + 2x_2^2$ con matrice $M = \bigl(\begin{smallmatrix}1 & 3/2 \\ 3/2 & 2\end{smallmatrix}\bigr)$; il determinante di $M$ è $2 - \frac{9}{4} = -\frac{1}{4} < 0$, quindi $q$ ha segnatura $(1,1)$. Per la Proposizione~\ref{prop:class-segn}(I)(b), $\mathcal{C}$ è un'\textbf{iperbole}.

\textbf{(b) Centro.} Per la Proposizione~\ref{prop:centro-conica}, il centro $(x_1^0, x_2^0)$ si trova risolvendo $A\mathbf{v} = (1,0,0)^t$. Il calcolo fornisce $(a_0, a_1, a_2) = (4, 16, -12)$ (a meno di normalizzazione), da cui il centro è $Q = (4, -3)$.
\end{esempio}
\end{exbox}

% ============================================================
\section{Quadriche e iperquadriche}
\label{sec:quadriche}
% ============================================================

\subsection{La classificazione generale}

La classificazione delle coniche si estende naturalmente agli spazi affini euclidei di dimensione $n$.

\begin{defbox}
\begin{definizione}[Equivalenza per isometrie in dimensione $n$]
\label{def:equiv-iso-n}
Sia $S$ uno spazio affine euclideo di dimensione $n$. Due funzioni polinomiali $f, g \colon S \to \R$ di grado 2 sono \emph{equivalenti per isometrie} se esistono $\Phi \in \Isom(S)$ e $\lambda \in \R \setminus \{0\}$ tali che $f(P) = \lambda g(\Phi(P))$ per ogni $P \in S$.
\end{definizione}
\end{defbox}

\begin{lemma}[Invarianti per equivalenza]
\label{lem:inv-eq}
Sia $F$ l'omogeneizzazione di $f \colon S \to \R$ polinomiale di grado 2 (in $n+1$ variabili).
\begin{enumerate}[label=(\alph*)]
  \item Il rango di $F$ è invariante per equivalenza; $f$ è \emph{non degenere} se $F$ è non degenere.
  \item La coppia di segnature $(s^+(F), s^-(F))$ è invariante per equivalenza a meno di riordinamento.
  \item La coppia di segnature di $F(0, x_1, \ldots, x_n)$ a meno di riordinamento è invariante.
\end{enumerate}
\end{lemma}

\begin{proof}
(a) Il rango non dipende dalle coordinate e non cambia moltiplicando per $\lambda \neq 0$. (b) La segnatura è invariante per trasformazioni ortogonali; scambia se $\lambda < 0$. (c) La matrice del cambio di coordinate affine ortonormale (della forma $\bigl(\begin{smallmatrix}1&\mathbf{0}\\B&A\end{smallmatrix}\bigr)$ con $A\in O_n(\R)$) fissa il sottospazio $\{x_0=0\}$.
\end{proof}

\begin{teorema}[Classificazione delle iperquadriche non degeneri]
\label{thm:classeq}
Sia $S$ uno spazio affine euclideo di dimensione $n$ e $f \colon S \to \R$ polinomiale di grado 2 non degenere. Allora esiste un riferimento ON con coordinate $(x_1, \ldots, x_n)$ nel quale $f$ è uguale (a meno di scalare) a una delle seguenti forme canoniche:
\begin{align}
  E_i(x_1,\ldots,x_n) &\coloneqq \frac{x_1^2}{a_1^2} + \cdots + \frac{x_i^2}{a_i^2} - \frac{x_{i+1}^2}{a_{i+1}^2} - \cdots - \frac{x_n^2}{a_n^2} + 1, \quad i \in \{0,\ldots,n\}, \label{eq:Ei} \\
  P_j(x_1,\ldots,x_n) &\coloneqq \frac{x_1^2}{a_1^2} + \cdots + \frac{x_j^2}{a_j^2} - \frac{x_{j+1}^2}{a_{j+1}^2} - \cdots - \frac{x_{n-1}^2}{a_{n-1}^2} + x_n, \quad 0 \leq j \leq \left\lfloor\frac{n-1}{2}\right\rfloor, \label{eq:Pj}
\end{align}
con $a_k > 0$. Due funzioni non degeneri di grado 2 sono equivalenti per isometrie se e solo se hanno la stessa forma canonica (con gli stessi parametri, a meno di riordinamento degli $a_k$).
\end{teorema}

\begin{proof}[Schema della dimostrazione]
\textbf{Passo 1.} Decomponiamo $f = q + l + c$ dove $q$ è una forma quadratica, $l$ è lineare e $c \in \R$. L'omogeneizzazione soddisfa $q(x_1,\ldots,x_n) = F(0, x_1, \ldots, x_n)$.

\textbf{Passo 2.} Si dimostra che $q$ ha rango almeno $n-1$ (cioè o è non degenere, oppure ha rango $n-1$). Infatti, se $q$ avesse rango $\leq n-2$, il ker di $q$ avrebbe dimensione $\geq 2$, e la restrizione di $L_F$ a ker $q$ mapperebbe nello spazio unidimensionale $\langle x_0^\vee \rangle$, il che contraddirebbe la non degenerazione di $F$.

\textbf{Caso (a): $q$ non degenere.} Per il Teorema Spettrale, si diagonalizza $q$ con una rotazione; poi si completa i quadrati e si ottiene $\sum \lambda_i z_i^2 + d$ con $d \neq 0$ (per la non degenerazione di $F$). Dividendo per $d$ e riordinando si arriva a $E_i$.

\textbf{Caso (b): $\rg q = n-1$.} Un'analisi simile porta alla forma $P_j$; il vincolo $j \leq \lfloor(n-1)/2\rfloor$ viene dal fatto che $P_j$ e $P_{n-1-j}$ sono equivalenti (moltiplicando per $-1$ e rinominando $x_n \mapsto -x_n$).

\textbf{Passo 3 (unicità).} Le segnature di $F$ e di $F|_{\{x_0=0\}}$ distinguono le $E_i$ dalle $P_j$ e tra loro.
\end{proof}

\subsection{Le quadriche in dimensione 3}

Applichiamo la classificazione al caso $n = 3$, che è il più rilevante per la geometria intuitiva.

\begin{defbox}
\begin{definizione}[Quadrica non degenere]
\label{def:quadrica}
Una \emph{quadrica non degenere} in uno spazio affine euclideo $S$ di dimensione 3 è l'insieme $V(f)$ dove $f \colon S \to \R$ è polinomiale di grado 2 non degenere. Nelle coordinate $(x,y,z)$ di un opportuno riferimento ON, si presenta in una delle sei forme canoniche $Q_1, \ldots, Q_6$ (Tabella~\ref{tab:quadriche}).
\end{definizione}
\end{defbox}

\begin{table}[h!]
\centering
\caption{Quadriche non degeneri in $E^3(\R)$.}
\label{tab:quadriche}
\renewcommand{\arraystretch}{1.8}
\begin{tabular}{lll}
\toprule
\textbf{Equazione canonica} & \textbf{Nome} & \textbf{Osservazioni} \\
\midrule
$Q_1\colon \dfrac{x^2}{a^2}+\dfrac{y^2}{b^2}+\dfrac{z^2}{c^2}=1$ & Ellissoide & Compatto, convesso \\[4pt]
$Q_2\colon \dfrac{x^2}{a^2}+\dfrac{y^2}{b^2}-\dfrac{z^2}{c^2}=1$ & Iperboloide a una falda & Connesso; curvatura gaussiana $<0$ \\[4pt]
$Q_3\colon \dfrac{x^2}{a^2}-\dfrac{y^2}{b^2}-\dfrac{z^2}{c^2}=1$ & Iperboloide a due falde & Non connesso; curvatura gaussiana $>0$ \\[4pt]
$Q_4\colon -\dfrac{x^2}{a^2}-\dfrac{y^2}{b^2}-\dfrac{z^2}{c^2}=1$ & Insieme vuoto & $V(Q_4) = \emptyset$ \\[4pt]
$Q_5\colon z=\dfrac{x^2}{a^2}+\dfrac{y^2}{b^2}$ & Paraboloide ellittico & ``Coppa''; curvatura $>0$ \\[4pt]
$Q_6\colon z=\dfrac{x^2}{a^2}-\dfrac{y^2}{b^2}$ & Paraboloide iperbolico & ``Sella''; curvatura $<0$ \\
\bottomrule
\end{tabular}
\end{table}

\begin{osservazione}[Rette su $Q_2$ e $Q_6$]
L'iperboloide a una falda $Q_2$ e il paraboloide iperbolico $Q_6$ contengono ciascuno \emph{due schiere di rette} (famiglie di rette che ricoprono la superficie). Per $Q_2$, si fattorizza
\[
  \frac{x^2}{a^2} - \frac{z^2}{c^2} = \left(\frac{x}{a}+\frac{z}{c}\right)\left(\frac{x}{a}-\frac{z}{c}\right) = \left(1+\frac{y}{b}\right)\left(1-\frac{y}{b}\right).
\]
Le rette della prima schiera sono definite da $\lambda\left(\frac{x}{a}+\frac{z}{c}\right) = \mu\left(1+\frac{y}{b}\right)$, $\mu\left(\frac{x}{a}-\frac{z}{c}\right) = \lambda\left(1-\frac{y}{b}\right)$, al variare di $(\lambda:\mu) \in \mathbb{P}^1(\R)$. Analogamente per la seconda schiera. Questa proprietà ha importanti applicazioni in architettura (torri di raffreddamento) e in fisica (strutture reticolate).

Per $Q_6$: l'equazione $z = \left(\frac{x}{a}+\frac{y}{b}\right)\left(\frac{x}{a}-\frac{y}{b}\right)$ mostra direttamente le due schiere.
\end{osservazione}

\begin{exbox}
\begin{esempio}[Classificazione di una quadrica]
\label{es:quadrica}
Sia $Q \subseteq E^3(\R)$ la quadrica di equazione
\[
  f(x,y,z) \coloneqq xy - 6xz - z^2 + x + 2 = 0.
\]
L'omogeneizzazione è $F(w,x,y,z) = xy - 6xz - z^2 + wx + 2w^2$, con matrice simmetrica
\[
  A = \begin{pmatrix} 2 & \frac{1}{2} & 0 & 0 \\ \frac{1}{2} & 0 & \frac{1}{2} & -3 \\ 0 & \frac{1}{2} & 0 & 0 \\ 0 & -3 & 0 & -1 \end{pmatrix}.
\]
Si calcola $\det A = \frac{1}{2} \neq 0$, quindi $Q$ è non degenere.

Per determinare le segnature di $F$ usiamo i minori principali $A^{(k)}$ (sottomatrice $k\times k$ superiore sinistra):
\[
  \det A^{(1)} = 2 > 0, \quad \det A^{(2)} = -\tfrac{1}{4} < 0, \quad \det A^{(3)} = -\tfrac{1}{2} < 0, \quad \det A^{(4)} = \tfrac{1}{2} > 0.
\]
Per il criterio di Jacobi, $F$ ha segnatura $(2,2)$.

Tra $Q_1,\ldots,Q_6$, le quadriche con segnatura $(2,2)$ della forma quadratica associata sono $Q_2$ (iperboloide iperbolico) e $Q_6$ (paraboloide iperbolico). Per distinguerle, osserviamo la forma quadratica $F(0,x,y,z) = xy - 6xz - z^2$: il suo determinante è $-1 \neq 0$, quindi è non degenere. Dunque $f$ rientra nel caso (a) del Teorema~\ref{thm:classeq}, e $Q$ è un \textbf{iperboloide iperbolico} ($Q_2$).
\end{esempio}
\end{exbox}

% ============================================================
\section*{Esercizi}
\addcontentsline{toc}{section}{Esercizi del Capitolo 10}
\label{sec:es-cap10}
% ============================================================

\begin{esercizio}
La curva in $E^2(\R)$ di equazione
\[
  \mathcal{C} \colon 73x^2 + 72xy + 52y^2 + 30x - 40y - 75 = 0
\]
è un'ellisse.
\begin{enumerate}[label=(\arabic*)]
  \item Scrivete l'omogeneizzazione $F(x_0, x_1, x_2)$ del polinomio che definisce $\mathcal{C}$.
  \item Dimostrate che $\mathcal{C}$ è un'ellisse calcolando tre determinanti (applicando la Proposizione~\ref{prop:class-segn}).
  \item Determinate i fuochi di $\mathcal{C}$ e la somma costante delle distanze dei punti di $\mathcal{C}$ dai suoi fuochi.
\end{enumerate}
\end{esercizio}

\begin{esercizio}
Per $t \in \R$ sia $\mathcal{C}_t \subseteq E^2(\R)$ la curva di equazione
\[
  \mathcal{C}_t \colon x_1^2 + 6x_1 x_2 + t x_2^2 - 2x_1 + 4x_2 + 3 = 0.
\]
Determinate la natura di $\mathcal{C}_t$ al variare di $t$: per quali valori di $t$ è un'ellisse, un'iperbole, una parabola, o una conica degenere.

\emph{Suggerimento:} usate la Proposizione~\ref{prop:class-segn} calcolando i determinanti opportuni in funzione di $t$.
\end{esercizio}

\begin{esercizio}
Sia $\mathcal{C} \subseteq E^2(\R)$ la curva di equazione
\[
  \mathcal{C} \colon x_1^2 + 2x_1 x_2 + x_2^2 - 2x_1 - 6x_2 + 3 = 0.
\]
\begin{enumerate}[label=(\alph*)]
  \item Dimostrate che $\mathcal{C}$ è una parabola calcolando due determinanti.
  \item Determinate il fuoco e la direttrice di $\mathcal{C}$.
\end{enumerate}
\end{esercizio}

\begin{esercizio}
Siano $Q_1, Q_2, Q_3 \subseteq E^3(\R)$ le quadriche di equazioni
\begin{align*}
  Q_1 &\colon x^2 + xy + 2y^2 - z^2 + 3x - z + 25 = 0, \\
  Q_2 &\colon 2xy + 4xz + y^2 - 4z^2 + 2y + 4z + 3 = 0, \\
  Q_3 &\colon 2xy + 4xz + y^2 - 4z^2 + 2y + 5z + 3 = 0.
\end{align*}
Quale tra $Q_1, Q_2, Q_3$ è un paraboloide? Di quale tipo (ellittico o iperbolico)?
\end{esercizio}

\begin{esercizio}
Sia $\mathcal{C}$ l'ellisse di semiassi $a > b > 0$. Dimostrate che l'area dell'ellisse è $\pi ab$.

\emph{Suggerimento:} l'ellisse si ottiene dal cerchio di raggio $a$ tramite la trasformazione $(x,y) \mapsto (x, \frac{b}{a}y)$, che è un'applicazione lineare di determinante $b/a$. Usate la formula di cambio di variabile per gli integrali.
\end{esercizio}

\begin{esercizio}
Dimostrate che l'iperboloide a una falda $Q_2 \colon \frac{x^2}{a^2}+\frac{y^2}{b^2}-\frac{z^2}{c^2}=1$ è una superficie \emph{connessa} costruendo esplicitamente, per ogni coppia di punti $P, Q \in Q_2$, un cammino continuo su $Q_2$ che li connette.

\emph{Suggerimento:} mostrate prima che ogni punto di $Q_2$ si connette a un punto del ``cerchio equatoriale'' $\{z=0\} \cap Q_2$.
\end{esercizio}

\begin{esercizio}[Classificazione completa delle coniche tramite invarianti]
Sia $f(x,y) = ax^2 + bxy + cy^2 + dx + ey + f_0$ un polinomio di grado 2. Definite
\[
  \Delta \coloneqq \det \begin{pmatrix} a & b/2 \\ b/2 & c \end{pmatrix} = ac - \frac{b^2}{4}, \qquad
  \nabla \coloneqq \det \begin{pmatrix} a & b/2 & d/2 \\ b/2 & c & e/2 \\ d/2 & e/2 & f_0 \end{pmatrix}.
\]
Dimostrate che $\Delta$ e $\nabla$ sono invarianti per isometrie (cioè non cambiano se si applica una isometria a $f$), e usateli per riottenere la classificazione della Proposizione~\ref{prop:class-segn} nel caso $n=2$.
\end{esercizio}

\end{document}